\documentclass[11pt]{article}

\usepackage[margin=1.1in]{geometry}

\usepackage{amssymb,mathtools,natbib}

% Anchored formal environments for causalsmith papers.
% First mandatory argument = obj_id (machine anchor joining the paper to the
% Lean crosswalk; invisible in the rendered PDF). Optional argument = name.
% Each environment also defines \label{obj:<obj_id>} so prose can use
% \cref{obj:T-1}; cleveref derives the printed kind and number from the target.
\usepackage{amsmath}
\usepackage{mathrsfs}
\usepackage{amsthm}
\usepackage{booktabs}
% Long displays may break across pages; let TeX stretch a little harder before an
% overfull \hbox so wide formulas/tables are less likely to run off the page.
\allowdisplaybreaks
\setlength{\emergencystretch}{3em}
\newtheorem{theorem}{Theorem}
\newtheorem{lemma}{Lemma}
\newtheorem{assumption}{Assumption}
\newtheorem{definition}{Definition}
\newtheorem{proposition}{Proposition}
\newtheorem{remark}{Remark}
\newtheorem{citedresult}{Cited result}
% The amsthm counter is `algorithmbox` (not `algorithm`) so a later `\usepackage{algorithm}` cannot clash.
\newcounter{algorithmbox}
\renewcommand{\thealgorithmbox}{\arabic{algorithmbox}}
\NewDocumentEnvironment{theoremv}{m o s}{\begin{theorem}\label{obj:#1}{\normalfont[\texttt{\detokenize{#1}}]\IfValueT{#2}{\ (#2)}\IfBooleanT{#3}{\verificationfootnotemark}\ }}{\end{theorem}}
\NewDocumentEnvironment{lemmav}{m o s}{\begin{lemma}\label{obj:#1}{\normalfont[\texttt{\detokenize{#1}}]\IfValueT{#2}{\ (#2)}\IfBooleanT{#3}{\verificationfootnotemark}\ }}{\end{lemma}}
\NewDocumentEnvironment{assumptionv}{m o s}{\begin{assumption}\label{obj:#1}{\normalfont[\texttt{\detokenize{#1}}]\IfValueT{#2}{\ (#2)}\IfBooleanT{#3}{\verificationfootnotemark}\ }}{\end{assumption}}
\NewDocumentEnvironment{definitionv}{m o s}{\begin{definition}\label{obj:#1}{\normalfont[\texttt{\detokenize{#1}}]\IfValueT{#2}{\ (#2)}\IfBooleanT{#3}{\verificationfootnotemark}\ }}{\end{definition}}
% A `propositionv` is a proved result tiered as a Proposition (theorem-grade, machine-verified).
\NewDocumentEnvironment{propositionv}{m o s}{\begin{proposition}\label{obj:#1}{\normalfont[\texttt{\detokenize{#1}}]\IfValueT{#2}{\ (#2)}\IfBooleanT{#3}{\verificationfootnotemark}\ }}{\end{proposition}}
% A `remarkv` holds NON-load-bearing interpretive / open-question content (e.g. a causalsmith `oeq:`
% open question). It is neither proved nor assumed; no theorem may depend on it.
\NewDocumentEnvironment{remarkv}{m o s}{\begin{remark}\label{obj:#1}{\normalfont[\texttt{\detokenize{#1}}]\IfValueT{#2}{\ (#2)}\IfBooleanT{#3}{\verificationfootnotemark}\ }}{\end{remark}}
% An `algorithmv` is a DEFINITION-kind object presented as a boxed, numbered-step procedure (an
% estimator / selection rule built in stages). Same anchor contract as `definitionv`; its body is
% ordinary LaTeX whose steps are `\begin{enumerate}` items, so tex2html needs no extra machinery.
% Presented in the CONVENTIONAL algorithm style readers expect from `algorithm`+`algorithmic`:
% a rule above, an "Algorithm N (Title)" caption line, a rule under the caption, the numbered
% steps, and a closing rule. It is NOT a float — the anchor contract requires the object to stay
% where the outline places it, and floats drift. The obj-id tag is suppressed here (it is
% bookkeeping, and a caption line carrying it reads as debug output in a finished paper).
\NewDocumentEnvironment{algorithmv}{m o s}%
  {\par\addvspace{\medskipamount}\noindent\rule{\linewidth}{0.8pt}\par\nobreak\vspace{2pt}%
   \refstepcounter{algorithmbox}\label{obj:#1}%
   \noindent\textbf{Algorithm \thealgorithmbox}\IfValueT{#2}{ \textnormal{(#2)}}%
   \IfBooleanT{#3}{\verificationfootnotemark}\par\nobreak\vspace{2pt}%
   \noindent\rule{\linewidth}{0.4pt}\par\nobreak\vspace{2pt}\normalfont}%
  {\par\nobreak\vspace{2pt}\noindent\rule{\linewidth}{0.8pt}\par\addvspace{\medskipamount}}
% A `citedv` is an IMPORTED external result the paper relies on but does NOT prove
% (a causalsmith `gate_class:"cited"` node). It is machine-checked against the cited
% source at F2.5, never proved here; the fixed provenance note makes that explicit
% so the environment can never be read as a contribution of this paper. The \citep
% attribution is supplied by the surrounding prose (P2), keyed to references.bib.
\NewDocumentEnvironment{citedv}{m o s}{\begin{citedresult}\label{obj:#1}{\normalfont[\texttt{\detokenize{#1}}]\IfValueT{#2}{\ (#2)}\IfBooleanT{#3}{\verificationfootnotemark}\ \textit{(imported result; machine-checked against the cited source, not proved here.)}\ }}{\end{citedresult}}
% \leanref{obj_id}{display text}: an inline first-mention link to a from-note object's
% verified Lean code. In the PDF it renders the display text plainly (the Lean link is a
% web-only affordance); tex2html turns it into a drawer-opening span.
\newcommand{\leanref}[2]{#2}
% For a theorem/lemma whose Lean declaration accepts a source-matched published proposition as an
% input, the starred anchored environment puts the mark in its heading; the generated text remains
% outside the frozen mathematical environment. Keep the combined macro for older paper bundles.
\newcommand{\verificationfootnotemark}{\footnotemark}
\newcommand{\verificationfootnotetext}[1]{\footnotetext{#1}}
\newcommand{\verificationfootnote}[1]{\verificationfootnotemark\verificationfootnotetext{#1}}
% xcolor before hyperref (hyperref would otherwise pull in plain `color` for colorlinks, and a
% later xcolor load clashes). The page-1 identity stamp at the end of this file needs a grey.
\usepackage{xcolor}
\usepackage{graphicx}
\usepackage{eso-pic}
% hyperref follows natbib and the environment macros; cleveref must follow hyperref.
\usepackage[colorlinks=true,linkcolor=blue,citecolor=blue,urlcolor=blue]{hyperref}
\usepackage[capitalise,nameinlink,noabbrev]{cleveref}
% Pin the names explicitly: labels are placed inside the underlying amsthm environments, and these
% declarations make the PDF contract independent of cleveref's theorem-name inference. House style:
% reference names always print with a capital first letter ("Theorem 1", "Definition 2"), in any
% sentence position — hence `capitalise` above and capitalized \crefname forms below.
\crefname{theorem}{Theorem}{Theorems}
\Crefname{theorem}{Theorem}{Theorems}
\crefname{lemma}{Lemma}{Lemmas}
\Crefname{lemma}{Lemma}{Lemmas}
\crefname{assumption}{Assumption}{Assumptions}
\Crefname{assumption}{Assumption}{Assumptions}
\crefname{definition}{Definition}{Definitions}
\Crefname{definition}{Definition}{Definitions}
\crefname{proposition}{Proposition}{Propositions}
\Crefname{proposition}{Proposition}{Propositions}
\crefname{remark}{Remark}{Remarks}
\Crefname{remark}{Remark}{Remarks}
\crefname{citedresult}{Cited result}{Cited results}
\Crefname{citedresult}{Cited result}{Cited results}
\crefname{algorithmbox}{Algorithm}{Algorithms}
\Crefname{algorithmbox}{Algorithm}{Algorithms}
% ---------------------------------------------------------------------------
% Page-1 identity stamp (arXiv style) and the pinned title-block date.
%
% Split of responsibility: THIS file owns the FORMAT (placement, font, colour, punctuation),
% the generated `paper_stamp.tex` owns only the VALUES (number+version, area, revised date,
% DOI, link target). P4 refreshes this template on every emit, so a layout fix reaches every
% bundle from a recompile; and because `paper_stamp.tex` is not one of the P2 assembly
% digest's authored inputs, a DOI that arrives after publication needs only that one small
% file rewritten plus a recompile — never a redraft, never a reassembly.
%
% Defaults first, so a bundle WITHOUT a stamp file compiles exactly as it always did:
% `\cspaperdate` falls back to `\today` and no stamp is drawn.
\providecommand*{\cspaperdate}{\today}  % the \date{} target
\providecommand*{\csstampid}{}          % e.g. CSWP-2026-008v3 (empty ⇒ no stamp at all)
\providecommand*{\csstamparea}{}        % e.g. Stat
\providecommand*{\csstampdate}{}        % e.g. 27 Aug 2026
\providecommand*{\csstampdoi}{}         % e.g. 10.5281/zenodo.123 (empty ⇒ DOI part omitted)
\providecommand*{\csstamplink}{}        % href target (DOI url, else the paper's page; may be empty)
\IfFileExists{paper_stamp.tex}{% GENERATED by CausalSmith P4 — paper identity stamp values. Do not hand-edit.
%
% Field order on the stamp (owned by paper_macros.tex):
%   <number>v<version> · doi:<doi> · [<Area>] <date>   — the DOI is never the last token, so a
%   text extractor cannot run it into whatever follows. Without a DOI that part is omitted.
% Format lives in paper_macros.tex; only the values are here, and this file is NOT one of
% the P2 assembly digest's authored inputs. A DOI arriving after publication therefore needs
% this file rewritten plus a `latexmk` recompile — never a redraft or a reassembly.
\renewcommand*{\csstampid}{CSWP-2026-032v11}
\renewcommand*{\csstamparea}{Stat}
\renewcommand*{\csstampdate}{5 Oct 2026}
\renewcommand*{\csstampdoi}{}
\renewcommand*{\csstamplink}{https://causalsmith.org/papers/stat_noisydose_weakdesign_transition_v1}
\renewcommand*{\cspaperdate}{5 October 2026}
}{}
\makeatletter
% Field order is deliberate: the DOI is NEVER the last token on the line. A trailing identifier
% runs into whatever a text extractor emits next, which makes it unsafe to match mechanically
% (the Zenodo stamp matcher reads exactly this line); the date is a safe terminator.
%   <number>v<version> · doi:<doi> · [<Area>] <date>      — with a DOI
%   <number>v<version> · [<Area>] <date>                  — without
\newcommand{\cs@stampsep}{\,\textperiodcentered\,}
\newcommand{\cs@stamptext}{%
  \csstampid
  \ifx\csstampdoi\@empty\else\cs@stampsep doi:\csstampdoi\fi
  \ifx\csstamparea\@empty\else\cs@stampsep[\csstamparea]\fi
  \ifx\csstampdate\@empty\else\quad\csstampdate\fi
}
\ifx\csstampid\@empty\else
  \definecolor{csstampgrey}{gray}{0.45}
  % Background shipout material: it occupies no space, so the text block and the \thanks
  % footnote are byte-identical to an unstamped compile. The starred form applies to the
  % NEXT page shipped only — i.e. page 1. Placed in the left margin (the text block starts
  % at 1.1in), reading bottom-to-top, in a Times-like face at 8pt.
  \AddToShipoutPictureBG*{%
    \AtPageLowerLeft{%
      \put(\LenToUnit{0.42in},\LenToUnit{1.1in}){%
        \rotatebox{90}{%
          \fontfamily{ptm}\fontsize{8}{10}\selectfont
          \ifx\csstamplink\@empty
            \textcolor{csstampgrey}{\cs@stamptext}%
          \else
            \expandafter\href\expandafter{\csstamplink}{\textcolor{csstampgrey}{\cs@stamptext}}%
          \fi
        }%
      }%
    }%
  }
\fi
\makeatother


\title{Pointwise causal estimation with weak dose support and Gaussian measurement error}

\author{CausalSmith\thanks{Machine-generated by the CausalSmith research pipeline.}}

\date{\cspaperdate}

\begin{document}

\maketitle

\begin{abstract}
Under calibrated public density envelopes satisfying \(c_{\mathrm{lo}}\le(\kappa+1)2^\kappa\le c_{\mathrm{hi}}\), where \(c_{\mathrm{lo}}\) and \(c_{\mathrm{hi}}\) are the density-envelope constants and \(\kappa\) is the design-decay exponent, we characterize optimal estimation and honest interval length for the population causal mean at a prespecified latent dose when treatment is measured with Gaussian error and local dose support may vanish polynomially. The model has two observed strata, compact latent-dose support, bounded potential outcomes, schedule-level conditional exchangeability, and classical measurement error with known standard deviation \(\sigma\in[0,1/4]\). Conditional mean functions have Hölder exponent \(\beta\in(0,1]\), and conditional dose densities satisfy these public polynomial envelopes with \(\kappa\in[0,2]\). For this calibrated class and any fixed noncoverage probability \(\alpha\in(0,1/2)\), minimax expected absolute error and minimum worst-law expected length of honest connected \((1-\alpha)\)-coverage intervals have matching orders, uniformly over the error scale. For sample size \(n\), the characterization combines a direct weak-design regime with rate \(n^{-\beta/(2\beta+\kappa+1)}\), an intermediate Gaussian-inversion regime, and a compact-support regime. Neighboring rate formulas are comparable throughout fixed multiplicative neighborhoods of both transition thresholds. Every fixed error scale \(\sigma\in(0,1/4]\) yields order \((\log\log n/\log n)^\beta\) for this calibrated class. An observable estimator selects among Fourier, polynomial, and constant inverse weights using public error certificates, with the declared small-sample fallback; its associated interval has finite-sample coverage \(1-\alpha\) for every \(\alpha\in(0,1)\). Under the same envelope calibration, matching lower bounds follow from bounded Bernoulli structural alternatives whose causal means are separated while their complete observed samples remain difficult to distinguish.
\end{abstract}

\section{Introduction}

This paper characterizes how measurement precision and local treatment support jointly determine the accuracy of pointwise causal estimation under calibrated public density envelopes. We study the population mean response to setting a latent continuous treatment to the prespecified dose \(1/2\). Treatment is observed with independent Gaussian measurement error of known standard deviation \(\sigma\), the error scale. Within the bounded structural law class in \cref{obj:def:model-class}, assuming \(c_{\mathrm{lo}}\le(\kappa+1)2^\kappa\le c_{\mathrm{hi}}\), we establish matching orders for minimax expected absolute error and minimum worst-law expected length of honest connected \((1-\alpha)\)-coverage intervals for \(\alpha\in(0,1/2)\). This calibrated-class characterization is uniform over the maintained range \(\sigma\in[0,1/4]\), including error-free observation, for fixed public class constants and exponents in their declared ranges. An observable estimator and interval attain these orders, and structural alternatives evaluated through their induced observed distributions supply the converse under the same envelope calibration.

The model combines causal adjustment with explicit restrictions on dose availability and response regularity. Observations consist of a stratum indicator taking two values, a contaminated treatment measurement, and a realized outcome. The latent dose lies in \([0,1]\). Public constants \(K=(p_{\min},c_{\mathrm{lo}},c_{\mathrm{hi}})\) use \(p_{\min}\) as the minimum stratum probability and \(c_{\mathrm{lo}},c_{\mathrm{hi}}\) as the density-envelope constants. Thus stratum probabilities lie in \([p_{\min},1-p_{\min}]\), while the density of the centered latent-dose displacement \(t=A-1/2\) lies between \(c_{\mathrm{lo}}|t|^\kappa\) and \(c_{\mathrm{hi}}|t|^\kappa\), where \(0<p_{\min}\le1/2\) and \(0<c_{\mathrm{lo}}\le c_{\mathrm{hi}}\). Conditional potential-outcome means satisfy radius-one Hölder increment bounds with exponent \(\beta\in(0,1]\), the mean-smoothness exponent, and take values in \([0,1]\); potential outcomes are bounded in the same interval. We maintain \(\kappa\in[0,2]\). Conditional exchangeability applies to the entire measurable potential-outcome schedule, consistency links its random-dose evaluation to the realized outcome, and the Gaussian error is independent of the structural coordinates. Under these conditions, \cref{obj:prop:identification} identifies the population causal mean even when the dose-density envelopes vanish at the evaluation point.

The main result describes three ranges of measurement precision under the envelope calibration \(c_{\mathrm{lo}}\le(\kappa+1)2^\kappa\le c_{\mathrm{hi}}\). Let \(n\) denote sample size, let \(d=2\beta+\kappa+1\) be the effective smoothness-and-design exponent, and let \(L_n=\log(en)\) be the logarithmic sample-size scale. The localization scale \(b_{n,\sigma}\) is
\[
b_{n,\sigma}=
\begin{cases}
n^{-1/d},
&0\le\sigma\le n^{-1/d},\\[3pt]
\dfrac{\sigma}{\sqrt{\log(e+n\sigma^d)}},
&n^{-1/d}<\sigma\le L_n^{-1/2},\\[7pt]
\dfrac{\log(e+\sigma^2L_n)}{L_n},
&L_n^{-1/2}<\sigma\le1/4.
\end{cases}
\]
For this calibrated class and all sufficiently large samples, \cref{obj:thm:uniform-frontier} bounds both optimal criteria above and below by constant multiples of \(b_{n,\sigma}^{\beta}\), uniformly over the known error scale and the declared public potential-outcome spaces. The direct branch balances mean smoothness against the scarcity of nearby latent doses. Gaussian Fourier inversion governs the intermediate branch, while known compact support supports polynomial localization in the third branch. For every fixed multiplicative neighborhood of either threshold, the adjacent formulas are comparable once sample size is sufficiently large. Within the calibrated class, at zero error, both criteria have order \(n^{-\beta/(2\beta+\kappa+1)}\); at every fixed positive error scale in the stated range, both eventually have order \((\log\log n/\log n)^\beta\).

The upper bound comes from positive latent localization weights paired with observable inverses of the Gaussian channel. Weighted outcome and treatment moments estimate a localized mean within each stratum, and empirical stratum frequencies aggregate those estimates into the population target. A finite dictionary contains Fourier weights, polynomial inverse-heat weights, and a constant entry. The estimator selects a dictionary entry by minimizing a public expected-error certificate determined by sample size, the known error scale, the public class constant \(K\), and the smoothness and design-decay exponents \(\beta,\kappa\). Thus the construction in \cref{obj:def:total-estimator} uses one selection rule across the three regimes. \Cref{obj:thm:observable-upper} establishes its uniform large-sample risk bound. Dividing the same certificate by \(\alpha\) calibrates a connected interval with coverage at least \(1-\alpha\) for every sample size and every law in the model class, while its expected-length bound attains the frontier for sufficiently large samples, as stated in \cref{obj:thm:finite-honesty}.

Under the envelope calibration \(c_{\mathrm{lo}}\le(\kappa+1)2^\kappa\le c_{\mathrm{hi}}\), the matching lower bound uses Bernoulli threshold schedules within the bounded structural class. The alternatives share their latent treatment design and measurement channel, while opposite perturbations change the causal mean at the evaluation dose. In the inverse regimes, normalized filtered Jacobi packets preserve the perturbation's pointwise height and cancel low-order moments under the polynomial design weight. Their localization and regularity bounds control model membership, and Gaussian convolution attenuates the resulting change in the observed outcome-marked distribution. \Cref{obj:thm:observed-lower} supplies target separation at the frontier order together with a uniform total-variation bound for the complete observed sample under this calibration. This comparison yields lower bounds for both estimation and honest connected-interval length over the same calibrated class used for achievability.

The closest causal measurement-error comparison is the contaminated-treatment dose-response estimator of \citet{HuangZhang2023}. Under the conditions of their arXiv version 2, Theorem 4.4, its fixed-dose expansion has bias orders \(h^2+h_0^2\) and stochastic remainder \(O_p\{N^{-1/2}/\min(e(h),e(h_0))\}\), where \(e(b)=b^{1/2}\exp(-b^{-r}/\gamma)\) in the supersmooth case. Our contribution is a matched minimax characterization of absolute risk and honest interval length under explicit polynomial design envelopes satisfying \(c_{\mathrm{lo}}\le(\kappa+1)2^\kappa\le c_{\mathrm{hi}}\), uniform over the known noise scale. The direct endpoint shares its rate exponent with the degenerate-design Gaussian-response regression benchmark of \citet[Definition 2 and Theorem 1]{Gaiffas2005}, as made precise in \cref{obj:prop:error-free-reduction}. Gaussian attenuation connects the intermediate branch to errors-in-variables regression \citep[Sections 3--4]{FanTruong1993}, while compact-support density deconvolution provides a precedent for iterated-logarithm-over-logarithm localization \citep[Theorems 1--2]{Meister2007}. These comparisons concern their respective experiments, restrictions, and losses. Under the same calibrated envelopes, the present characterization places the three mechanisms within one causal model and couples their estimation orders to honest uncertainty about the same scalar target.



The paper proceeds through related work, setup and identification, and the optimal risk and interval-length characterization. It then presents the observable estimator and honest intervals, followed by the observed-law lower bounds. Discussion and limitations and future work follow; the appendices develop identification and error certificates, weighted packets and the observed-experiment converse, and proofs of the main results.


\section{Related work}

The closest causal comparison is the contaminated-treatment dose-response estimator of \citet{HuangZhang2023}. Under their positivity, smoothing, sieve, and supersmooth-error conditions, the fixed-dose expansion in arXiv version 2, Theorem 4.4 and equation (21) has bias orders
\[
h^2+h_0^2
\]
and stochastic remainder
\[
O_p\!\left\{\frac{N^{-1/2}}{\min(e(h),e(h_0))}\right\},
\qquad
e(b)=b^{1/2}\exp(-b^{-r}/\gamma),
\]
where \(h,h_0\) are bandwidths, \(N\) is sample size, \(r\) is the supersmooth exponent, and \(\gamma\) is the error-tail constant; the normal limit also uses the theorem's variance lower bound \citep[arXiv version 2, Theorem 4.4, equation (21)]{HuangZhang2023}. The present numerical bounds and schedule-level exchangeability specify a different structural class. Within that class, the delivered advance is a minimax characterization uniform over the known noise scale, allowing polynomially weak dose support and matching the absolute-risk order with honest connected-interval length.
% [Huang and Zhang, arXiv version 2, Theorem 4.4](https://arxiv.org/pdf/2211.04642v2)

The error-free endpoint connects to pointwise regression with a degenerate design. For sample size \leanref{sym:n}{\ensuremath{n}}, mean smoothness exponent \leanref{sym:beta(smoothness exponent in (0,1])}{\ensuremath{\beta}}, and design-decay exponent \leanref{sym:kappa}{\ensuremath{\kappa}}, our direct rate is \(n^{-\beta/(2\beta+\kappa+1)}\), as stated in \cref{obj:prop:error-free-reduction}. Definition 2 of \citet{Gaiffas2005} specifies bounded local polynomial-approximation classes, and Theorem 1 characterizes minimax pointwise \(L^p\) risk in direct random-design regression with independent Gaussian response errors. Under its locally symmetric, regularly varying design and regularly varying approximation modulus, the rate combines a power of sample size with a slowly varying factor. Pure-power smoothness and design behavior yield the same exponent as our direct endpoint \citep[Definition 2 and Theorem 1]{Gaiffas2005}. The connection is an exponent comparison across experiments: our target is a population causal mean with bounded outcomes and stratum-specific design envelopes.
% [Gaiffas, Definition 2 and Theorem 1](https://arxiv.org/pdf/math/0410354)

The intermediate inverse regime is related to the regression deconvolution analysis of \citet{FanTruong1993}. Their Sections 3--4 establish estimation rates and corresponding lower bounds for errors-in-variables regression, including pointwise squared-error risk, under positive-design, smoothness, moment, and characteristic-function conditions. In the supersmooth case, exponential attenuation of high frequencies leads to logarithmic estimation rates; Gaussian errors furnish the quadratic-exponential example \citep[Sections 3--4]{FanTruong1993}. Our intermediate branch uses this Fourier inversion mechanism within the weak-design causal class. The comparison therefore concerns how Gaussian attenuation controls localization, with each rate evaluated under its own design restrictions and loss.
% [Fan and Truong, Sections 3--4](https://archive.ymsc.tsinghua.edu.cn/pacm_download/468/10828-Jianqing_Fan_37.pdf)

Compact support supplies another relevant mechanism. \citet[Theorems 1--2]{Meister2007} studies density deconvolution under mean integrated squared error over Sobolev classes, exploiting support information and knowledge of the error characteristic function on a neighborhood of the origin. For source Sobolev smoothness \(s\), polynomial continuation yields the vanishing mean integrated squared error order \(\{(\log\log n)/(\log n)\}^{2s}\). The matching lower bound concerns the joint density-and-error classes of that experiment and carries additional restrictions on density smoothness, the Sobolev radius, and regularity of the local error transform. Our compact-support branch instead addresses pointwise causal evaluation under expected absolute loss, with the Gaussian error law known and the latent design constrained by polynomial envelopes.
% [Meister, Theorems 1--2](https://www.f08.uni-stuttgart.de/.content/media/downloads/Mathematik/mathematische_berichte/2005/2005-007.pdf)

A complementary compact-support result is the spline-assisted regression estimator of \citet{JiangMaCarroll2026}. Under their spline-order, dimension-growth, knot-spacing, bounded-density, compact-support, mean-smoothness, and approximation conditions, together with bounded score-derivative and second-moment operators and the locally unique nonsingular population root in (D5), Theorem 2 gives
\[
\sup_x |\widehat m(x)-m(x)|
=O_p\!\left((nh_b)^{-1/2}+h_b^q\right),
\qquad
E\{\widehat m(x)-m(x)\}^2
=O\!\left(h_b^{2q}+(nh_b)^{-1}\right).
\]
These algebraic spline bounds are conditional on the source experiment's operator and nonsingularity structure \citep[Theorem 2, conditions (C2)--(C4) and (D1)--(D5)]{JiangMaCarroll2026}. In particular, their increasing-dimensional spline system must retain the bounded score operators and locally unique nonsingular population root required by (D5). The measurable polynomial envelopes and Hölder mean restriction used here do not impose that increasing-dimensional operator condition uniformly over the model class. The present result instead gives a same-class minimax characterization of pointwise absolute error and honest interval length, uniform over the maintained Gaussian noise scale; its loss and uniformity quantifiers therefore differ from the source theorem.
% [Jiang, Ma and Carroll, Theorem 2 and its conditions](https://arxiv.org/pdf/1804.00793)

The interval construction connects to two strands of nonparametric inference. \citet[Sections 2--3]{ArmstrongKolesar2020} constructs bias-aware intervals for scalar functionals in direct regression, using a specified smoothness bound to calibrate coverage and assess interval length. Our construction likewise incorporates a uniform bias bound, together with an observable estimation-error certificate. In errors-in-variables regression, \citet[Assumption 3.1 and Theorem 3.2]{KatoSasaki2019} obtains asymptotically valid simultaneous bands through deconvolution and the multiplier bootstrap. Their result assumes ordinary-smooth measurement error, positive design and variance on the reporting interval, moment and regularity conditions, and bandwidth and auxiliary error-sample growth restrictions that include undersmoothing. Alongside these results, \cref{obj:thm:finite-honesty} supplies finite-sample coverage for a connected interval targeting one causal mean, while \cref{obj:thm:uniform-frontier} characterizes its optimal worst-law expected length.
% [Armstrong and Kolesar, Sections 2--3](https://www.princeton.edu/~mkolesar/papers/simple_cis.pdf)
% [Kato and Sasaki, Assumption 3.1 and Theorem 3.2](https://arxiv.org/pdf/1702.03377)

Continuous-treatment causal methods provide the broader identification and estimation context. Propensity-function adjustment identifies population treatment responses under ignorability and positivity \citep[Section 2]{Imai2004}, while doubly robust smoothing combines outcome regression and treatment-density estimation for dose-response estimation \citep[pp.~1229--1230]{Kennedy2017}. Vanishing-noise regression provides a separate precedent for studying rates jointly in sample size and noise scale: \citet[Sections 3--4]{DurotMukherjee2026} analyzes monotone-function estimation in shuffled and unlinked experiments under integrated absolute-error criteria and the respective sampling and error conditions. These connections motivate treating measurement precision as part of the statistical characterization.
% [Imai and van Dyk, Section 2](https://www.ma.imperial.ac.uk/~dvandyk/Research/04-jasa-propfunc.pdf)
% [Kennedy and coauthors, pp. 1229--1230](https://academic.oup.com/jrsssb/article/79/4/1229/7040663)
% [Durot and Mukherjee, Sections 3--4](https://arxiv.org/pdf/2404.09306)

Under the envelope calibration \(c_{\mathrm{lo}}\le(\kappa+1)2^\kappa\le c_{\mathrm{hi}}\), the contribution is a same-class characterization of minimax pointwise expected absolute error and minimum worst-law expected length of honest connected intervals, as defined in \cref{obj:def:minimax-risk,obj:def:honest-length}. For this calibrated class, fixed public \(K\), exponents in their declared ranges, and any noncoverage probability \(\alpha\in(0,1/2)\), \cref{obj:thm:uniform-frontier} gives matching orders uniformly over the maintained Gaussian error range \leanref{sym:sigma}{\ensuremath{\sigma}} in \([0,1/4]\). The observable construction and the finite-sample \((1-\alpha)\)-coverage guarantee in \cref{obj:thm:observable-upper,obj:thm:finite-honesty}, together with the observed-law converse in \cref{obj:thm:observed-lower} under this calibration, place direct estimation, Gaussian Fourier inversion, and compact-support estimation within one causal experiment and one set of model restrictions.

The Gaiffas result supplies a separate direct-regression benchmark for the error-free deduction in \cref{obj:prop:error-free-reduction}; it complements the frontier's own same-class zero-noise conclusion.


\section{Setup and assumptions}

Generic constants \(c,C>0\) may depend on the fixed public class parameters; dependence on sample size or error scale is displayed explicitly. Absolute value denotes the scalar norm, and \(\|f\|_\infty\) denotes the uniform norm of a function on its stated domain. The notation \(u\asymp v\) means that \(u/v\) is bounded above and below by positive constants uniformly over the stated index range. Stratum indices range over the two strata defined below, sample indices run from \(1\) to the sample size \(n\), and integer indices are nonnegative unless a different range is specified.

The \leanref{S-1}{structural coordinates} distinguish treatment assignment, measurement error, and potential outcomes. The set of observed strata is \leanref{sym:Xset}{\ensuremath{\mathcal X}}, with \(\mathcal X=\{0,1\}\), and \leanref{sym:X}{\ensuremath{X}} is the observed stratum coordinate. We write \leanref{sym:A}{\ensuremath{A}} for the latent dose, \leanref{sym:Z}{\ensuremath{Z}} for the standardized measurement error, \leanref{sym:Y}{\ensuremath{Y}} for the realized outcome, and \leanref{sym:U}{\ensuremath{U}} for an auxiliary rank coordinate used in the lower-bound subclass. Throughout, \leanref{S-2}{the smoothness exponent, design-decay exponent, and public error scale} are fixed at \(\beta\in(0,1]\), \(\kappa\in[0,2]\), and \(\sigma\in[0,1/4]\), respectively.

Potential outcomes are represented by a random schedule on a fixed public measurable space \leanref{sym:Yspace}{\ensuremath{(\mathcal S,\mathcal E)}}. Its evaluation map \leanref{sym:evalpath}{\ensuremath{e}} assigns an outcome to a schedule and a dose. The following definition makes both schedule-level independence and evaluation at a random dose measurable statements.

\begin{definitionv}{def:potential-path-space}[Potential-outcome space \((\mathcal S,\mathcal E)\)]
The public potential-outcome space \((\mathcal S,\mathcal E)\) is a measurable space of Borel functions from \([0,1]\) to \(\mathbb R\), satisfying the following conditions.
\begin{itemize}
\item \textbf{(Evaluation.)} The evaluation map
\[
e:\mathcal S\times[0,1]\longrightarrow\mathbb R,
\qquad e(f,a)=f(a),
\]
is jointly measurable for the product sigma algebra \(\mathcal E\otimes\mathcal B([0,1])\).
\item \textbf{(Extensionality.)} For \(f,g\in\mathcal S\),
\[
f=g
\quad\Longleftrightarrow\quad
e(f,a)=e(g,a)\ \text{for every }a\in[0,1].
\]
\item \textbf{(Threshold schedules.)} The space includes threshold schedules through the measurable map
\[
\iota:C([0,1],[0,1])\times[0,1]\longrightarrow\mathcal S,
\qquad
\iota(f,u)(a)=\mathbf1\{u\le f(a)\},
\]
where \(C([0,1],[0,1])\) carries its uniform-norm Borel sigma algebra.
\end{itemize}
The results apply to every fixed public space satisfying these conditions, with general schedules drawn from its Borel member functions. The potential-outcome schedule \(Y(\cdot)\) is a measurable random element of \(\mathcal S\), and its evaluation at dose \(a\) is
\[
Y(a)=e(Y(\cdot),a).
\]
\end{definitionv}

Write \leanref{sym:Yprocess}{\ensuremath{Y(\cdot)}} for the random potential-outcome schedule and \leanref{sym:P}{\ensuremath{P}} for the structural probability law of \((X,A,Z,Y(\cdot),Y,U)\) on the corresponding coordinate spaces. The stratum probability \leanref{sym:px}{\ensuremath{p_x}} is \(p_x=P(X=x)\). The contaminated dose \leanref{sym:W}{\ensuremath{W}} and the centered coordinates are defined by
\[
W=A+\sigma Z,
\qquad
a_0=\tfrac12,
\qquad
t=A-a_0,
\qquad
V=W-a_0=t+\sigma Z.
\]
Here \leanref{sym:a0}{\ensuremath{a_0}} is the prespecified evaluation dose, \leanref{sym:T}{\ensuremath{t}} is the centered latent dose, and \leanref{sym:V}{\ensuremath{V}} is the centered observed dose. For each stratum with positive probability, \leanref{sym:gx}{\ensuremath{g_x}} denotes a Lebesgue density of the conditional law of \(t\) given \(X=x\); its existence and bounds are imposed below.

The population causal target \leanref{sym:theta}{\ensuremath{\theta(P)}} is the mean potential outcome at the evaluation dose:
\[
\theta(P)=E_P[Y(a_0)].
\]
The boundedness condition below ensures that this expectation exists. To describe its stratum-specific components, define the conditional potential mean \leanref{sym:mux}{\ensuremath{\mu_x(a)}} before imposing restrictions on treatment assignment or mean smoothness.

\begin{definitionv}{def:conditional-potential-mean}[Conditional potential mean \(\mu_x(a)\)]
The conditional potential-outcome mean is
\[
\mu_x(a)
=
E_P[Y(a)\mid X=x]
=
\frac{E_P[\mathbf1\{X=x\}Y(a)]}{p_x},
\qquad a\in[0,1],\quad x\in\mathcal X,\quad p_x>0.
\]
Here \(P\) is the structural probability law, \(X\) is the observed stratum coordinate, \(\mathcal X\) is the set of two observed strata, and
\[
p_x=P(X=x).
\]
The potential outcome \(Y(a)\) is the schedule evaluation in \cref{obj:def:potential-path-space}. Under the bounded-potential-outcome condition, the expectation exists and belongs to \([0,1]\). The model's stratum-overlap condition supplies \(p_x>0\) for both strata.
\end{definitionv}

This definition separates the causal mean from the distribution of assigned doses. We first restrict the stratum probabilities and the latent-dose distribution so that both strata contribute and the amount of dose support near the evaluation point is controlled.

\begin{assumptionv}{ass:stratum-overlap}[Uniform stratum overlap]
For a public constant \(p_{\min}\in(0,1/2]\), the stratum probabilities satisfy
\[
p_{\min}\le p_x\le 1-p_{\min}
\qquad\text{for every }x\in\mathcal X.
\]
\end{assumptionv}

The finite-stratum overlap condition gives each observed stratum a probability bounded away from zero, uniformly over the structural laws under consideration \citep{HuangZhang2023}. The public constant \(p_{\min}\) fixes that common lower bound. Treatment assignment is then restricted to a common compact interval.

\begin{assumptionv}{ass:latent-support}[Latent dose support]
The latent dose \(A\) satisfies
\[
A\in[0,1]\qquad\text{almost surely}.
\]
\end{assumptionv}

Compact continuous-treatment support fixes the dose domain on which potential outcomes and their means are evaluated \citep{HuangZhang2023}. Centering this interval at \(a_0\) gives latent support \([-1/2,1/2]\). The next condition specifies how the conditional density behaves throughout that centered interval.

\begin{assumptionv}{ass:weak-design}[Local design density bounds]
For public constants \(0<c_{\mathrm{lo}}\le c_{\mathrm{hi}}\), each \(g_x\), \(x\in\mathcal X\), is a Lebesgue density satisfying
\[
c_{\mathrm{lo}}|t|^\kappa
\le g_x(t)\le
c_{\mathrm{hi}}|t|^\kappa
\qquad\text{for almost every }t\in[-1/2,1/2].
\]
\end{assumptionv}

Polynomially degenerate random design describes declining dose support near the evaluation point through a fixed density envelope \citep{Gaiffas2005}. For \(\kappa>0\), the upper envelope tends to zero as the latent dose approaches \(a_0\); for \(\kappa=0\), the density is bounded above and below by positive constants almost everywhere. The public constants \(c_{\mathrm{lo}}\) and \(c_{\mathrm{hi}}\) keep this comparison uniform across strata and structural laws.

Mean regularity supplies the corresponding control on changes in the causal response. For a real-valued function \(f\) on \([0,1]\), define the Hölder increment seminorm and its radius-one class by
\[
[f]_\beta
=
\sup_{\substack{a,a'\in[0,1]\\a\ne a'}}
\frac{|f(a)-f(a')|}{|a-a'|^\beta},
\qquad
\mathcal H^\beta([0,1])
=
\{f:[0,1]\to\mathbb R:[f]_\beta\le1\}.
\]
In particular, \([\mu_x]_\beta\) measures the regularity of the stratum-specific conditional potential mean.

\begin{assumptionv}{ass:holder-mean}[Hölder conditional mean increments]
For the potential-outcome schedules of \cref{obj:def:potential-path-space}, every stratum-specific conditional mean satisfies the radius-one Hölder increment condition
\[
\mu_x\in\mathcal H^\beta([0,1])
\quad\Longleftrightarrow\quad
|\mu_x(a)-\mu_x(a')|\le |a-a'|^\beta
\quad\text{for all }a,a'\in[0,1],
\qquad x\in\mathcal X.
\]
The class is defined by this increment condition; the unit-interval range is imposed separately.
\end{assumptionv}

The Hölder smoothness condition controls the approximation of a pointwise mean by means at nearby doses, as in pointwise estimation under a degenerate design \citep{Gaiffas2005}. Here the exponent and increment radius are fixed across the class. A separate restriction records the common range of these means.

\begin{assumptionv}{ass:mean-range}[Bounded conditional potential means]
For the potential-outcome schedules of \cref{obj:def:potential-path-space}, the conditional potential means satisfy
\[
0\le\mu_x(a)\le1
\qquad\text{for every }x\in\mathcal X
\text{ and }a\in[0,1].
\]
\end{assumptionv}

This explicit mean-range restriction is specific to this analysis and places all conditional causal means on the unit outcome scale. It also fixes the range used by the threshold subclass introduced below. The observation channel is governed by the following two restrictions.

\begin{assumptionv}{ass:gaussian-channel}[Standard Gaussian error law]
The standardized error coordinate \(Z\) satisfies
\[
Z\sim N(0,1).
\]
\end{assumptionv}

The Gaussian classical measurement error condition specifies the convolution law associated with the contaminated dose \citep{FanTruong1993}. Together with the known public scale \(\sigma\), it fixes the distribution of the additive error \(\sigma Z\), including the error-free case \(\sigma=0\). Independence specifies its relation to the structural coordinates.

\begin{assumptionv}{ass:error-independence}[Classical error independence]
The standardized error coordinate \(Z\) is independent of the observed stratum \(X\), latent dose \(A\), and potential-outcome schedule \(Y(\cdot)\):
\[
Z\perp(X,A,Y(\cdot)).
\]
\end{assumptionv}

Classical independent treatment measurement error makes the same additive noise law applicable across strata, latent doses, and potential-outcome schedules \citep{HuangZhang2023}. The causal interpretation of the latent response is supplied by schedule-level conditional exchangeability, bounded outcomes, and consistency.

\begin{assumptionv}{ass:conditional-unconfoundedness}[Conditional potential-outcome unconfoundedness]
The measurable random potential-outcome schedule \(Y(\cdot)\) and latent dose \(A\) are conditionally independent given the observed stratum \(X\):
\[
Y(\cdot)\perp A\mid X.
\]
\end{assumptionv}

Conditional exchangeability for individual potential outcomes is used in average dose-response estimation under measurement error \citep{HuangZhang2023}. In our model, functional conditional exchangeability requires the entire measurable potential-outcome schedule to be independent of dose assignment within each observed stratum. Its schedule-level formulation accommodates evaluation at the realized latent dose through the jointly measurable map in \cref{obj:def:potential-path-space}. Outcome boundedness provides a common scale for these evaluations.

\begin{assumptionv}{ass:bounded-potential-outcomes}[Bounded potential outcomes]
The potential outcomes in \cref{obj:def:potential-path-space} satisfy
\[
P\{Y(a)\in[0,1]\}=1
\qquad
\text{for every }a\in[0,1].
\]
\end{assumptionv}

The bounded potential outcomes condition ensures integrability at every fixed dose and fixes the unit scale of the causal target \citep{HuangZhang2023}. It also implies the conditional mean range stated in \cref{obj:ass:mean-range}. Consistency connects the potential-outcome schedule to the realized outcome.

\begin{assumptionv}{ass:consistency}[Realized outcome consistency]
The realized outcome \(Y\) equals the potential outcome from \cref{obj:def:potential-path-space} evaluated at the latent dose \(A\):
\[
Y=Y(A)\qquad P\text{-almost surely}.
\]
\end{assumptionv}

Continuous-treatment consistency assigns the observed outcome to the potential outcome at the actual latent dose \citep{HuangZhang2023}. Together, \cref{obj:ass:conditional-unconfoundedness,obj:ass:consistency} give the conditional potential mean its causal role in the latent-dose experiment.

For the observed-law lower bounds, we use a subclass with explicit rank-generated schedules. The auxiliary coordinate \(U\) provides a common rank across doses; the next restrictions specify its conditional distribution and relation to treatment assignment.

\begin{assumptionv}{ass:rank-uniform}[Conditional uniform rank law]
For every \(x\in\mathcal X\), the uniform-rank coordinate \(U\) satisfies
\[
U\mid X=x\sim\operatorname{Unif}[0,1].
\]
\end{assumptionv}

The conditional uniform rank law is specific to this analysis and supplies a unit-interval rank for constructing binary schedules with prescribed conditional means. Its relation to the latent dose is fixed separately.

\begin{assumptionv}{ass:rank-treatment-independence}[Conditional rank–dose independence]
The rank coordinate \(U\) and latent dose \(A\) are conditionally independent given the observed stratum \(X\):
\[
U\perp A\mid X.
\]
\end{assumptionv}

Conditional rank–dose independence is specific to this analysis and permits rank-generated schedules to satisfy the maintained conditional exchangeability restriction. The schedule construction uses the same rank at every dose.

\begin{assumptionv}{ass:structural-threshold}[Structural threshold schedules]
For the random potential-outcome schedule in \cref{obj:def:potential-path-space}, let \(U\) denote the uniform-rank coordinate and \(\mu_X(a)\) the conditional potential-outcome mean at the realized stratum. With \(P\)-probability one,
\[
Y(a)=\mathbf 1\{U\le \mu_X(a)\}
\qquad
\text{simultaneously for every }a\in[0,1].
\]
\end{assumptionv}

The structural threshold restriction is specific to this analysis and constructs an entire schedule from the conditional mean function and one rank. Its simultaneous statement ties all dose evaluations to a single measurable schedule. Under the conditional uniform rank law, each threshold has the prescribed conditional mean. The resulting subclass also satisfies the following outcome restriction.

\begin{assumptionv}{ass:binary-potential-outcomes}[Binary potential outcomes]
The potential outcomes in \cref{obj:def:potential-path-space} satisfy
\[
P\{Y(a)\in\{0,1\}\}=1
\qquad
\text{for every }a\in[0,1].
\]
\end{assumptionv}

Average dose-response estimation under measurement error provides a related causal setting \citep{HuangZhang2023}. Binary potential outcomes provide the outcome specification for the lower-bound witnesses within our bounded-outcome causal model. The threshold construction satisfies this restriction directly.

We now collect the defining restrictions into the bounded-outcome structural model class \leanref{sym:Pclass}{\ensuremath{\mathcal P_{K,\beta,\kappa,\sigma}}}. The public vector \(K\) records the common stratum-overlap and density-envelope constants.

\begin{definitionv}{def:model-class}[Bounded-outcome structural model class]
The class \(\mathcal P_{K,\beta,\kappa,\sigma}\) consists of structural laws \(P\) of \((X,A,Z,Y(\cdot),Y,U)\), where \(X\) takes values in the two observed strata \(\mathcal X\), \(A\) is the latent dose, \(Z\) is the standardized error coordinate, \(Y(\cdot)\) is the random potential-outcome schedule on the public path space of \cref{obj:def:potential-path-space}, \(Y\) is the realized outcome, and \(U\) is an auxiliary rank coordinate. The parameters \(\beta,\kappa,\sigma\in\mathbb R\) are fixed, and the public class constants satisfy
\[
K=(p_{\min},c_{\mathrm{lo}},c_{\mathrm{hi}}),
\qquad
0<p_{\min}\le \tfrac12,
\qquad
0<c_{\mathrm{lo}}\le c_{\mathrm{hi}}.
\]
Membership in the class is defined by the following conditions:
\begin{itemize}
\item \textbf{(Stratum overlap.)} \cref{obj:ass:stratum-overlap}.
\item \textbf{(Latent support.)} \cref{obj:ass:latent-support}.
\item \textbf{(Conditional design.)} \cref{obj:ass:weak-design}, with envelope constants \(c_{\mathrm{lo}},c_{\mathrm{hi}}\) and exponent \(\kappa\).
\item \textbf{(Mean smoothness.)} \cref{obj:ass:holder-mean}, with exponent \(\beta\).
\item \textbf{(Mean range.)} \cref{obj:ass:mean-range}.
\item \textbf{(Gaussian channel.)} \cref{obj:ass:gaussian-channel}.
\item \textbf{(Error independence.)} \cref{obj:ass:error-independence}.
\item \textbf{(Conditional unconfoundedness.)} \cref{obj:ass:conditional-unconfoundedness}.
\item \textbf{(Bounded potential outcomes.)} \cref{obj:ass:bounded-potential-outcomes}.
\item \textbf{(Consistency.)} \cref{obj:ass:consistency}.
\end{itemize}
The contaminated dose \(W\) is defined by
\[
W=A+\sigma Z.
\]
The conditional potential-outcome mean functions \(\mu_x\) are defined in \cref{obj:def:conditional-potential-mean}. The lower-bound subclass additionally satisfies \cref{obj:ass:binary-potential-outcomes,obj:ass:rank-uniform,obj:ass:rank-treatment-independence,obj:ass:structural-threshold}.
\end{definitionv}

This class fixes the structural restrictions under which causal estimation is evaluated. Its lower-bound subclass supplies explicit binary threshold laws satisfying those same restrictions.

An observed record \leanref{sym:O}{\ensuremath{O}} consists of \(O=(X,W,Y)\). For a sample of size \(n\), write \(O_i=(X_i,W_i,Y_i)\), \(i=0,\ldots,n-1\), for the corresponding records; this zero-based indexing is the one used by the modulo-three sample blocks of the estimator. Where a statement lists the sample as \(O_1,\ldots,O_n\), it denotes the same \(n\) records. The sampling restriction determines how their common observed law generates the statistical experiment.

\begin{assumptionv}{ass:iid}[Independent and identically distributed observations]
The observed records \(O_1,\ldots,O_n\), each consisting of the observed stratum, contaminated dose, and realized outcome, are independent and identically distributed under the structural law \(P\).
\end{assumptionv}

Independent and identically distributed sampling specifies repeated observations from the same contaminated-dose experiment \citep{HuangZhang2023}. The induced sample law \leanref{sym:Qn}{\ensuremath{Q_P^n}} is the experiment available to estimators and confidence intervals.

\begin{definitionv}{def:observed-experiment}[Observed experiment \(Q_P^n\)]
The sample observed-law experiment is
\[
Q_P^n=\mathcal L_P(O_1,\ldots,O_n),
\]
where \(O_1,\ldots,O_n\) are independent and identically distributed observed records under the structural law \(P\), each consisting of the observed stratum, contaminated dose, and realized outcome.
\end{definitionv}

Causal identification must account for both measurement error and the declining density near the target dose. \Cref{obj:prop:identification} establishes that, for common public class parameters and path space, equality of the observed-record laws determines the conditional potential means throughout the dose interval and hence the population causal target.

\begin{propositionv}{prop:identification}[Identification of the causal mean]
Fix the public potential-path space \((\mathcal S,\mathcal E)\) and evaluation map of \cref{obj:def:potential-path-space}, and public class constants \(K\) as in \cref{obj:def:model-class}. Let \(\mathcal X=\{0,1\}\) denote the two stratum labels, \(X\) the stratum coordinate, \(A\) the latent dose, \(Z\) the standardized error, and \(Y\) the realized outcome. Consider structural probability laws \(P\) and \(P'\) satisfying:
\begin{itemize}
\item \textbf{(Parameter ranges.)} \(\beta\in(0,1]\), \(\kappa\in[0,2]\), and \(\sigma\in[0,1/4]\).
\item \textbf{(Common model class.)} Both laws belong to the model class of \cref{obj:def:model-class} with the same \(K,\beta,\kappa,\sigma\) and public potential-path space.
\item \textbf{(Observed-law equality.)} Writing \(W=A+\sigma Z\) for the contaminated dose, the laws of the observed record coincide:
\[
\mathcal L_P(X,W,Y)=\mathcal L_{P'}(X,W,Y).
\]
\end{itemize}
For either structural law, write \(p_x(P)=P(X=x)\), and let \(\mu_x(a;P)\) denote its conditional potential-outcome mean from \cref{obj:def:conditional-potential-mean}. Set \(a_0=1/2\), the target dose, and \(t=A-a_0\), the centered latent dose. Define the population causal mean by
\[
\theta(P)=\int Y(a_0)\,dP,
\]
and analogously for \(P'\).

Then, for every \(x\in\mathcal X\), the stratum masses and conditional centered-dose measures agree:
\[
p_x(P)=p_x(P'),
\qquad
\mathcal L_P(A-a_0\mid X=x)
=
\mathcal L_{P'}(A-a_0\mid X=x).
\]
Their conditional latent mean measures also agree: for every Borel set \(B\subseteq\mathbb R\),
\[
\int_B \mu_x(t+a_0;P)\,
        \mathcal L_P(A-a_0\mid X=x)(dt)
=
\int_B \mu_x(t+a_0;P')\,
        \mathcal L_{P'}(A-a_0\mid X=x)(dt).
\]
Moreover,
\[
\mu_x(a;P)=\mu_x(a;P')
\qquad\text{for every }a\in[0,1].
\]
For every integer \(j\ge0\), the shrinking neighborhood has positive conditional latent mass:
\[
P\!\left(
-\frac1{j+3}<A-a_0<\frac1{j+3}
\,\middle|\,X=x
\right)>0,
\]
and its mean-measure to dose-measure ratio satisfies
\[
\lim_{j\to\infty}
\frac{
\displaystyle
\int_{(-1/(j+3),\,1/(j+3))}
\mu_x(t+a_0;P)\,
\mathcal L_P(A-a_0\mid X=x)(dt)
}{
\displaystyle
P\!\left(
-\frac1{j+3}<A-a_0<\frac1{j+3}
\,\middle|\,X=x
\right)
}
=
\mu_x(a_0;P).
\]
Finally,
\[
\theta(P)
=
\sum_{x\in\mathcal X}p_x(P)\mu_x(a_0;P),
\qquad
\theta(P)=\theta(P').
\]
\end{propositionv}

The shrinking-neighborhood conclusion explains how the target remains identified under weak dose support. Every displayed neighborhood has positive conditional latent mass, and its mean-measure to dose-measure ratio converges to the stratum-specific causal mean at \(a_0\). Thus the pointwise target has a population interpretation through local averages even when the density envelope declines to zero near that dose. The proposition also identifies the entire conditional mean function, while its final decomposition combines the two stratum-specific target values using the identified stratum probabilities. The realized-dose identity and Gaussian convolution argument supporting these conclusions are presented in \cref{obj:lem:realized-dose-mean,obj:lem:compact-gaussian-convolution-injective}; the positive neighborhood masses and the convergence of the local ratio are proved directly in steps 3 and 4 of the proof of \cref{obj:prop:identification}.


\section{Optimal estimation and interval length}

Identification in \cref{obj:prop:identification} makes the population causal mean a well-defined target of the observed experiment. We now characterize the precision attainable from that experiment as sample size and measurement precision vary. The first criterion is the minimax expected absolute error \leanref{sym:Rn}{\ensuremath{R_n(\sigma)}}, evaluated over the structural class in \cref{obj:def:model-class}.

\begin{definitionv}{def:minimax-risk}[Minimax absolute-error risk $R_{K,n}(\sigma)$]
The minimax expected absolute error is
\[
R_{K,n}(\sigma)
=
\inf_T
\sup_{P\in\mathcal P_{K,\beta,\kappa,\sigma}}
E_{Q_P^n}\!\left[|T-\theta(P)|\right].
\]
Here \(\theta(P)\) is the population causal mean at the prespecified latent dose, and \(Q_P^n\) is the observed-sample law induced by the structural law \(P\), with potential-outcome schedules on the public path space of \cref{obj:def:potential-path-space}. The infimum ranges over measurable, possibly randomized estimators \(T\).
\end{definitionv}

For uncertainty quantification, fix a noncoverage probability \leanref{sym:alpha}{\ensuremath{\alpha\in(0,1/2)}}, corresponding to coverage of at least \(1-\alpha\) under every law in the structural class. The second criterion, \leanref{sym:Hn}{\ensuremath{H_n(\sigma)}}, measures the smallest worst-law expected length among connected intervals satisfying this uniform coverage requirement. The separate finite-sample upper guarantee in \cref{obj:thm:finite-honesty} holds for the wider range \(\alpha\in(0,1)\); the restriction \(\alpha<1/2\) is used by our two-point argument for the honest-length lower bound, which requires \(1-2\alpha\) to be positive; we do not claim that the frontier fails when \(\alpha\ge1/2\).

\begin{definitionv}{def:honest-length}[Honest interval length $H_{K,\alpha,n}(\sigma)$]
The minimum worst-law expected length at fixed noncoverage probability \(\alpha\) is
\[
H_{K,\alpha,n}(\sigma)
=
\inf_{\substack{
C:\ \inf_{P\in\mathcal P_{K,\beta,\kappa,\sigma}}
Q_P^n\{\theta(P)\in C\}\ge 1-\alpha
}}
\sup_{P\in\mathcal P_{K,\beta,\kappa,\sigma}}
E_{Q_P^n}\!\left[\operatorname{length}(C)\right].
\]
The infimum ranges over measurable connected intervals \(C\). Here \(\theta(P)\) is the population causal mean at the prespecified latent dose, and \(Q_P^n\) is the observed-sample law induced by the structural law \(P\), with potential-outcome schedules on the public path space of \cref{obj:def:potential-path-space}.
\end{definitionv}

Both criteria are governed by a common rate \leanref{sym:rho}{\ensuremath{\rho_{n,\sigma}}}, obtained by raising the localization scale \leanref{sym:bns}{\ensuremath{b_{n,\sigma}}} to the mean-smoothness exponent. The effective exponent \leanref{sym:d}{\ensuremath{d}} combines smoothness and local design decay, while the logarithmic sample-size quantity \leanref{sym:Ln}{\ensuremath{L_n}} determines the second measurement-error threshold. Their definitions distinguish three ranges of measurement precision.

\begin{definitionv}{def:frontier-rate}[Frontier rate \(\rho_{n,\sigma}\)]
The frontier rate and its localization scale are
\[
\rho_{n,\sigma}=b_{n,\sigma}^{\beta},
\qquad
b_{n,\sigma}=
\begin{cases}
n^{-1/d},&0\le\sigma\le n^{-1/d},\\[2pt]
\dfrac{\sigma}{\sqrt{\log(e+n\sigma^d)}},
&n^{-1/d}<\sigma\le L_n^{-1/2},\\[6pt]
\dfrac{\log(e+\sigma^2L_n)}{L_n},
&L_n^{-1/2}<\sigma\le1/4.
\end{cases}
\]
Here the effective exponent \(d\) and logarithmic sample-size quantity \(L_n\) are
\[
d=2\beta+\kappa+1,
\qquad
L_n=\log(en).
\]
\end{definitionv}

The main characterization, \cref{obj:thm:uniform-frontier}, gives matching bounds for estimation error and honest interval length throughout the stated error-scale range. It also compares the neighboring formulas across multiplicative neighborhoods of both thresholds, making the rate characterization uniform through the transitions.

\begin{theoremv}{thm:uniform-frontier}[Uniform risk and length frontier]*
Fix public class constants \(K=(p_{\min},c_{\mathrm{lo}},c_{\mathrm{hi}})\), smoothness exponent \(\beta\), design-decay exponent \(\kappa\), and noncoverage probability \(\alpha\), satisfying:
\begin{itemize}
\item \textbf{(Class constants.)} \(0<p_{\min}\le 1/2\) and \(0<c_{\mathrm{lo}}\le c_{\mathrm{hi}}\).
\item \textbf{(Parameter ranges.)} \(0<\beta\le1\), \(0\le\kappa\le2\), and \(0<\alpha<1/2\).
\item \textbf{(Envelope calibration.)} \(c_{\mathrm{lo}}\le(\kappa+1)2^\kappa\le c_{\mathrm{hi}}\).
\end{itemize}
Write \(R_n(\sigma)\) for the minimax absolute risk in \cref{obj:def:minimax-risk} and \(H_n(\sigma)\) for the minimum honest expected connected-interval length in \cref{obj:def:honest-length}, at these fixed parameters. Define the effective exponent \(d\), logarithmic sample scale \(L_n\), transition thresholds \(a_n,b_n\), and comparison scales \(D_n,F_n(\sigma),P_n(\sigma)\) by
\[
d=2\beta+\kappa+1,\qquad L_n=\log(en),\qquad
a_n=D_n=n^{-1/d},\qquad b_n=L_n^{-1/2},
\]
\[
F_n(\sigma)=\frac{\sigma}{\sqrt{\log(e+n\sigma^d)}},
\qquad
P_n(\sigma)=\frac{\log(e+\sigma^2L_n)}{L_n}.
\]
Define the localization scale \(b_{n,\sigma}\) and absolute-risk rate \(\rho_{n,\sigma}\) by
\[
b_{n,\sigma}=
\begin{cases}
D_n,&\sigma\le a_n,\\
F_n(\sigma),&a_n<\sigma\le b_n,\\
P_n(\sigma),&\sigma>a_n\ \text{and}\ \sigma>b_n,
\end{cases}
\qquad
\rho_{n,\sigma}=b_{n,\sigma}^{\beta}.
\]

There exist constants \(0<c<C\) and an integer \(n_0\), depending on \(K,\beta,\kappa,\alpha\), such that, for every public potential-path space as in \cref{obj:def:potential-path-space}, every integer \(n\ge n_0\), and every \(\sigma\in[0,1/4]\),
\[
c\rho_{n,\sigma}\le R_n(\sigma)\le C\rho_{n,\sigma},
\qquad
c\rho_{n,\sigma}\le H_n(\sigma)\le C\rho_{n,\sigma}.
\]
For the same \(C\) and every \(n\ge n_0\), the transition comparisons satisfy
\[
D_n\le C F_n(a_n),\qquad F_n(a_n)\le C D_n,
\]
\[
F_n(b_n)\le C P_n(b_n),\qquad P_n(b_n)\le C F_n(b_n).
\]

Moreover, for every real \(W\ge1\), there exist a constant \(C_K\ge1\) and an integer \(n_K\) such that, for every integer \(n\ge n_K\),
\[
0<a_n/W,\qquad Wa_n\le1/4,\qquad
0<b_n/W,\qquad Wb_n\le1/4,
\]
and the two transition windows satisfy
\[
C_K^{-1}D_n\le F_n(\sigma)\le C_KD_n
\qquad
\text{for every }\sigma\in[a_n/W,Wa_n],
\]
\[
C_K^{-1}P_n(\sigma)\le F_n(\sigma)\le C_KP_n(\sigma)
\qquad
\text{for every }\sigma\in[b_n/W,Wb_n].
\]
Here \(C_K\) and \(n_K\) may depend on \(W\) and the fixed parameters.

Finally, for every fixed \(\sigma\in(0,1/4]\), there exist constants \(0<c<C\) and an integer \(n_0\), which may depend on \(\sigma\) and the fixed parameters, such that, for every public potential-path space as in \cref{obj:def:potential-path-space} and every integer \(n\ge n_0\),
\[
c\left(\frac{\log\log n}{\log n}\right)^\beta
\le R_n(\sigma)\le
C\left(\frac{\log\log n}{\log n}\right)^\beta,
\]
\[
c\left(\frac{\log\log n}{\log n}\right)^\beta
\le H_n(\sigma)\le
C\left(\frac{\log\log n}{\log n}\right)^\beta.
\]
\end{theoremv}
% CAUSALSMITH-CITED-SCOPE-BEGIN thm:uniform-frontier
\verificationfootnotetext{\textbf{Formalization scope.} This result uses the cited conclusion \citep{PetrushevXu2005} (Theorem 2.4, equation (2.14)), \citep{https://people.math.sc.edu/pencho/Publications/kpx-06-28-06-web.pdf} (Theorem 2.2, equation (2.4)), \citep{https://dlmf.nist.gov/18.3} (Table 18.3.1, Jacobi row and its caption; orthogonality and squared norm.) and \citep{https://dlmf.nist.gov/5.11.E12} (Equation 5.11.12, restricted to the positive real axis.). The cited source proof is not formalized here: Lean verifies its use and all remaining steps. If that cited conclusion is false, the portions of this result that depend on it are not certified.}
% CAUSALSMITH-CITED-SCOPE-END thm:uniform-frontier

The envelope calibration in \cref{obj:thm:uniform-frontier} has a concrete role: the density \((\kappa+1)2^\kappa|t|^\kappa\) integrates to one on \([-1/2,1/2]\), and bracketing its coefficient places the explicit lower-bound design inside the class. Thus the condition certifies a nonempty class containing the witnesses used for the converse.

Under the envelope calibration \(c_{\mathrm{lo}}\le(\kappa+1)2^\kappa\le c_{\mathrm{hi}}\), the three branches express the interaction between local dose support and measurement precision. In the direct regime, the measurement-error scale is at most the localization scale available from an error-free dose record. Local design decay enters through \(d=2\beta+\kappa+1\), and the resulting rate is \(n^{-\beta/d}\). In the intermediate Gaussian regime, recovering a localized mean from the contaminated dose requires Gaussian inversion, with attainable localization scale \(F_n(\sigma)\). For larger error scales, known compact latent support enables polynomial localization: inverse polynomial degree is the localization width, Gaussian inversion increases the inverse weight's second moment with degree, and the largest usable degree has order \(L_n/\log(e+\sigma^2L_n)\). Balancing these effects gives the polynomial localization scale \(P_n(\sigma)\) and explains the transition from the Fourier branch. The detailed public certificate appears in \cref{obj:lem:dictionary-score-rate}. For the calibrated class, raising each localization scale to the smoothness exponent gives both the estimation rate and the honest-length rate.

\Cref{tab:precision-regimes} displays these branches and their thresholds. The scales \(D_n\), \(F_n(\sigma)\), and \(P_n(\sigma)\) are those defined in \cref{obj:thm:uniform-frontier}; each row gives the rate up to the uniform constants in that result.

\begin{table}[htbp]
\centering
\caption{Measurement precision and optimal pointwise uncertainty}
\label{tab:precision-regimes}
\begin{tabular}{lll}
\hline
Regime & Measurement-error range & Risk and honest-length rate \\
\hline
Direct
& \(0\le \sigma\le a_n=n^{-1/d}\)
& \(D_n^\beta=n^{-\beta/d}\) \\[4pt]
Intermediate Gaussian
& \(a_n<\sigma\le b_n=L_n^{-1/2}\)
& \(\displaystyle F_n(\sigma)^\beta
=\left\{\frac{\sigma}{\sqrt{\log(e+n\sigma^d)}}\right\}^{\beta}\) \\[8pt]
Compact support
& \(b_n<\sigma\le1/4\)
& \(\displaystyle P_n(\sigma)^\beta
=\left\{\frac{\log(e+\sigma^2L_n)}{L_n}\right\}^{\beta}\) \\[6pt]
\hline
\end{tabular}
\end{table}

The transition comparisons give the thresholds a precise interpretation. For each fixed multiplicative factor \(W\ge1\), the direct and intermediate scales are comparable throughout \([a_n/W,Wa_n]\) once the sample size is sufficiently large. Likewise, the intermediate and compact-support scales are comparable throughout \([b_n/W,Wb_n]\). Thus neighboring descriptions yield the same order of uncertainty throughout each transition window. The comparison constants may depend on \(W\) and the fixed class parameters, while remaining uniform over the public potential-outcome path spaces.

For every fixed positive measurement-error scale in the stated range, increasing sample size eventually places the experiment in the compact-support regime. Both minimax error and minimum honest expected length then have order \((\log\log n/\log n)^\beta\), with comparison constants that may depend on the fixed error scale. The full uniform characterization also describes sequences of improving measurement precision: attaining the direct-regime order requires an error scale at most of order \(n^{-1/d}\), including its fixed multiplicative transition neighborhood.

The error-free endpoint follows directly from this characterization. The next result states that endpoint for both criteria and places its estimation rate alongside a direct regression benchmark with the same local design-decay exponent.

\begin{propositionv}{prop:error-free-reduction}[Error-free rates and benchmark comparison]*
Fix the public class constants \(K=(p_{\min},c_{\mathrm{lo}},c_{\mathrm{hi}})\) and parameters \(\beta,\kappa,\alpha\) satisfying:
\begin{itemize}
\item \textbf{(Class constants.)} \(0<p_{\min}\le 1/2\) and \(0<c_{\mathrm{lo}}\le c_{\mathrm{hi}}\).
\item \textbf{(Parameter ranges.)} \(0<\beta\le1\), \(0\le\kappa\le2\), and \(0<\alpha<1/2\), where \(\alpha\) is the noncoverage probability.
\item \textbf{(Envelope calibration.)}
\[
c_{\mathrm{lo}}\le(\kappa+1)2^\kappa\le c_{\mathrm{hi}}.
\]
\end{itemize}
There exist constants \(0<c<C\) and \(n_0\in\mathbb N\) such that, for every public potential-path space \((\mathcal S,\mathcal E)\) satisfying \cref{obj:def:potential-path-space} and every integer \(n\ge n_0\), the minimax absolute risk and minimum honest expected connected-interval length of \cref{obj:def:minimax-risk,obj:def:honest-length} satisfy
\[
c\,n^{-\beta/(2\beta+\kappa+1)}
\le R_n(0)\le
C\,n^{-\beta/(2\beta+\kappa+1)},
\qquad
c\,n^{-\beta/(2\beta+\kappa+1)}
\le H_n(0)\le
C\,n^{-\beta/(2\beta+\kappa+1)}.
\]

Moreover, fix any separate direct Gaussian-response benchmark with:
\begin{itemize}
\item \textbf{(Benchmark constants.)} \(c_{\mathrm G},r_{\mathrm G},M_{\mathrm G},\tau_{\mathrm G}>0\).
\item \textbf{(Evaluation neighborhood.)} \(x_*\in(0,1)\) and \(0<\delta\le\min\{x_*,1-x_*\}\).
\item \textbf{(Design.)} A fixed probability distribution \(\nu\) on \([0,1]\) with a nonnegative density \(g\) satisfying
\[
g(u)=c_{\mathrm G}|u-x_*|^\kappa
\quad\text{for }u\in[0,1]\text{ with }|u-x_*|\le\delta.
\]
\end{itemize}
The benchmark observations are \(B_i=f(D_i)+\xi_i\), \(i=1,\ldots,n\), where the \(D_i\) are independent draws from \(\nu\), and the \(\xi_i\) are independent \(N(0,\tau_{\mathrm G}^2)\) errors independent of the design. In the supremum below, \(f\) ranges over Borel functions on \([0,1]\) satisfying
\[
\|f\|_\infty\le M_{\mathrm G},
\qquad
\inf_{\substack{p\text{ a real polynomial}\\ \deg p\le\lfloor\beta\rfloor}}
\ \sup_{\substack{u\in[0,1]\\ |u-x_*|\le a}}
|f(u)-p(u-x_*)|
\le r_{\mathrm G}a^\beta
\]
for every
\[
0<a\le
\left\{\frac{\tau_{\mathrm G}^2(\kappa+1)}
{2c_{\mathrm G}r_{\mathrm G}^2n}\right\}^{1/(2\beta+\kappa+1)}.
\]
The infimum below ranges over measurable estimators from the benchmark observations, and \(E_f\) denotes expectation under their observed law. There exist constants \(0<c<C\) and \(n_0\in\mathbb N\), now also depending on the fixed benchmark, such that for every public potential-path space satisfying \cref{obj:def:potential-path-space} and every integer \(n\ge n_0\),
\[
c\,\inf_{\widehat f}\sup_f E_f|\widehat f(x_*)-f(x_*)|
\le R_n(0)\le
C\,\inf_{\widehat f}\sup_f E_f|\widehat f(x_*)-f(x_*)|.
\]
\end{propositionv}
% CAUSALSMITH-CITED-SCOPE-BEGIN prop:error-free-reduction
\verificationfootnotetext{\textbf{Formalization scope.} This result uses the cited conclusion \citep{PetrushevXu2005} (Theorem 2.4, equation (2.14)), \citep{https://people.math.sc.edu/pencho/Publications/kpx-06-28-06-web.pdf} (Theorem 2.2, equation (2.4)), \citep{https://dlmf.nist.gov/18.3} (Table 18.3.1, Jacobi row and its caption; orthogonality and squared norm.), \citep{https://dlmf.nist.gov/5.11.E12} (Equation 5.11.12, restricted to the positive real axis.) and \citep{Gaiffas2005} (Definition 2 and Theorem 1, equations (2.1), (2.4), and (2.5), PDF pages 3--4; exact power-law specialization with loss parameter one.). The cited source proof is not formalized here: Lean verifies its use and all remaining steps. If that cited conclusion is false, the portions of this result that depend on it are not certified.}
% CAUSALSMITH-CITED-SCOPE-END prop:error-free-reduction

\Cref{obj:prop:error-free-reduction} identifies the precision available when the latent dose is observed exactly. Its exponent reflects both mean smoothness and the declining frequency of observations near the evaluation dose. The comparison with Gaiffas's benchmark is a comparison of minimax absolute-error rates for the displayed function class and fixed benchmark parameters \citep[Definition 2 and Theorem 1, equations (2.1), (2.4), and (2.5), pp.~3--4]{Gaiffas2005}. At \(\beta=1\), that benchmark retains approximation by polynomials of degree at most one. The honest-length conclusion for the causal experiment is supplied by the first part of \cref{obj:prop:error-free-reduction}. Together, these conclusions provide an error-free reference for interpreting the measurement-error branches.

The upper and lower bounds have complementary statistical roles. Observable localization produces estimates of the stratum-specific causal means from contaminated-dose and outcome records; the construction and its honest intervals are given in \cref{obj:thm:observable-upper,obj:thm:finite-honesty}. The converse compares structural laws whose causal means are separated while their induced observed laws remain difficult to distinguish, as stated in \cref{obj:thm:observed-lower}. Working with the observed experiment makes the lower bound apply to every estimator and honest connected interval in the criteria above.

The dependency roadmap is short. Identification in \cref{obj:prop:identification} fixes the observed-law target. The public certificate in \cref{obj:lem:public-error-certificate} turns each inverse weight into a risk and interval-radius bound; the dictionary comparisons in \cref{obj:lem:dictionary-score-rate} select entries at the three frontier scales. Finally, the observed-law witnesses in \cref{obj:thm:observed-lower} convert packet cancellation into matching testing lower bounds. Detailed algebra and witness coordinates are deferred to the two technical appendices.

The weighted-packet construction supporting that converse is developed in \cref{obj:thm:weighted-packet-construction}. Its mathematical ingredients include filtered Jacobi localization \citep[Theorem 2.4, equation (2.14)]{PetrushevXu2005}, the corresponding local difference estimate \citep[Theorem 2.2, equation (2.4)]{https://people.math.sc.edu/pencho/Publications/kpx-06-28-06-web.pdf}, Jacobi orthogonality and squared norms \citep[Table 18.3.1, Jacobi row and caption]{https://dlmf.nist.gov/18.3}, and the positive-real-axis gamma-ratio asymptotic \citep[equation (5.11.12)]{https://dlmf.nist.gov/5.11.E12}. These ingredients supply localization and moment cancellation for the observed-law comparison; their detailed construction belongs with the auxiliary results in the appendix.


\section{Observable estimation and honest intervals}

The \leanref{S-3}{observable estimation layer} uses the public class constants, smoothness and design-decay indices, and Gaussian error scale to choose a localization weight before evaluating the sample. Exact Gaussian inversion turns that latent weight into an observable function of the contaminated dose. A finite dictionary then supplies a single procedure whose expected absolute error attains the rate in \cref{obj:def:frontier-rate}, uniformly over the stated error scales.

\subsection{Inverse weights and public scores}

Let \leanref{sym:q}{\ensuremath{q}} be a public nonnegative weight on the centered latent-dose interval, and let \leanref{sym:ellq}{\ensuremath{\ell_q}} be its observable Gaussian inverse. The inverse identity requires that averaging the observable inverse over the measurement error recover the latent weight at every dose in the support. The following definition groups this identity with the integrability conditions and the public quantities used to evaluate a candidate.

\begin{definitionv}{synth_4}[Public expected-error score \(A_q(K)\)]
The public expected-error score of a weight and inverse-weight pair \((q,\ell_q)\) is
\[
A_q(K)
=
B_q(K)
+\frac{12}{b_q(K)}\sqrt{\frac{V_q(K)}{n}}
+2\sqrt{\frac3n}.
\]
Fix the public class constants \(K=(p_{\min},c_{\mathrm{lo}},c_{\mathrm{hi}})\), the mean-smoothness index \(\beta\), the design-decay index \(\kappa\), the Gaussian scale \(\sigma\), and the sample size \(n\). The pair satisfies the following conditions.
\begin{itemize}
\item \textbf{(Public constants.)}
\[
0<p_{\min}\le\frac12,
\qquad
0<c_{\mathrm{lo}}\le c_{\mathrm{hi}},
\qquad n\ge1.
\]
The score is defined for every \(n\ge1\). The estimator and interval use it only when all three modulo-three sample blocks are nonempty, that is, when \(n\ge3\); for \(n\le2\) they return their fallback outputs, which do not depend on the scores.
\item \textbf{(Inverse identity.)}
The weight \(q:[-1/2,1/2]\to[0,\infty)\) and its inverse \(\ell_q:\mathbb R\to\mathbb R\) are public Borel functions, and
\[
E_Z[\ell_q(t+\sigma Z)]=q(t),
\qquad -1/2\le t\le1/2,
\qquad Z\sim N(0,1),
\]
with every expectation absolutely finite.
\item \textbf{(Weighted integrability.)}
With weighted moments
\[
I_s(q)=\int_{-1/2}^{1/2}|t|^s q(t)\,dt,
\]
one has
\[
0<I_\kappa(q)<\infty,
\qquad
I_{\kappa+\beta}(q)<\infty.
\]
Moreover, \(\ell_q(V)\) is absolutely integrable under every law in \(\mathcal P_{K,\beta,\kappa,\sigma}\).
\item \textbf{(Second moment.)}
The public inverse-weight second-moment quantity is
\[
V_q(K)
=
c_{\mathrm{hi}}\int_{-1/2}^{1/2}
|t|^\kappa E_Z[\ell_q(t+\sigma Z)^2]\,dt
<\infty.
\]
The integral is a nonnegative extended integral.
\end{itemize}
The public denominator and localization-bias quantities are
\[
b_q(K)=p_{\min}c_{\mathrm{lo}}I_\kappa(q),
\qquad
B_q(K)=\frac{c_{\mathrm{hi}}}{c_{\mathrm{lo}}}
\frac{I_{\kappa+\beta}(q)}{I_\kappa(q)}.
\]
For a second-moment quantity expressed with factor \(16\), use the distinct notation
\[
V_q^{(16)}
=
16\int_{-1/2}^{1/2}|t|^\kappa
E_Z[\ell_q(t+\sigma Z)^2]\,dt
=
\frac{16}{c_{\mathrm{hi}}}V_q(K).
\]
Thus the second-moment term in the score has the equivalent expression
\[
\frac{12}{b_q(K)}\sqrt{\frac{V_q(K)}n}
=
\frac{12}{b_q(K)}
\sqrt{\frac{c_{\mathrm{hi}}V_q^{(16)}}{16n}}.
\]
The two second-moment quantities differ by the factor \(c_{\mathrm{hi}}/16\) and are not interchangeable: \cref{obj:lem:public-error-certificate,obj:thm:observable-upper} write \(V_q\), without the argument \(K\), for the factor-\(16\) quantity \(V_q^{(16)}\), whereas the score uses \(V_q(K)\).
Dependence on the fixed quantities \(\beta,\kappa,n,\sigma\) is suppressed in the notation.
\end{definitionv}

The design-weighted moment \leanref{sym:Iq}{\ensuremath{I_\kappa(q)}} determines the public denominator bound \leanref{sym:bq}{\ensuremath{b_q(K)}}. Localization is measured by the bias quantity \leanref{sym:Bq}{\ensuremath{B_q(K)}}, while the inverse-weight second-moment quantity \(V_q(K)\), equal to \(c_{\mathrm{hi}}/16\) times the factor-\(16\) reference moment \leanref{sym:Vq}{\ensuremath{V_q^{(16)}}}, measures the sampling cost of Gaussian inversion. Their combination, the public expected-error score \leanref{sym:Aq}{\ensuremath{A_q(K)}}, also accounts for estimation of the stratum probabilities. All these quantities depend on the candidate and public parameters. The factor-\(16\) convention is related to the class-dependent second moment by the explicit conversion in \cref{obj:synth_4}.

For interpretation, define the localized population ratio \leanref{sym:mxq}{\ensuremath{m_{x,q}}} by
\[
m_{x,q}
=
\frac{\displaystyle\int_{-1/2}^{1/2}
\mu_x(a_0+t)q(t)g_x(t)\,dt}
{\displaystyle\int_{-1/2}^{1/2}q(t)g_x(t)\,dt},
\qquad x\in\mathcal X.
\]
This ratio averages the conditional causal mean over doses emphasized by the weight. The positive weighted moment and lower design envelope give a positive denominator. Exact inversion supplies observable weighted moments for estimating this ratio; \cref{obj:lem:public-error-certificate} connects the resulting estimator to the public score.

Two explicit families furnish the dictionary comparisons. The polynomial localization weight \leanref{sym:qM}{\ensuremath{q_m^{\mathrm M}}}, indexed by a positive odd integer \(m\), has the inverse-heat companion \leanref{sym:ellM}{\ensuremath{\ell_m^{\mathrm M}}}. Its polynomial form makes the inverse a finite derivative sum.

\begin{definitionv}{synth_2}[Polynomial localization weights and inverse-heat weights]
For a positive odd integer \(m\), define on \([-1/2,1/2]\)
\[
q_m^{\mathrm M}(t)=
\begin{cases}
\displaystyle
\left[\frac{\sin(m\arcsin(2t))}{2mt}\right]^6,
&t\ne0,\\[6pt]
1,&t=0.
\end{cases}
\]
For odd \(m\), this expression is a polynomial of degree \(6(m-1)\); use its polynomial extension to define \(q_m^{\mathrm M}\) on \(\mathbb R\). For known Gaussian error scale \(\sigma\in[0,1/4]\), define its companion observable inverse weight by
\[
\ell_m^{\mathrm M}(v)
=\sum_{j=0}^{3(m-1)}
\frac{(-\sigma^2/2)^j}{j!}
(q_m^{\mathrm M})^{(2j)}(v),
\qquad v\in\mathbb R,
\]
where the superscript \((2j)\) denotes the derivative of order \(2j\). For \(Z\sim N(0,1)\), this pair satisfies the exact Gaussian inversion identity
\[
E_Z\!\left[\ell_m^{\mathrm M}(t+\sigma Z)\right]
=q_m^{\mathrm M}(t),\qquad t\in\mathbb R.
\]
At sample size \(n\), the polynomial entries of the public dictionary use the odd indices
\[
m=2^j+1\le n^2,\qquad j\in\mathbb Z,\quad j\ge1.
\]
\end{definitionv}

The Fourier family instead uses the localization function \leanref{sym:Qkernel}{\ensuremath{Q}} and a bandwidth \leanref{sym:h}{\ensuremath{h}}. Its latent weight \leanref{sym:qF}{\ensuremath{q_h^{\mathrm F}}} is paired with the Gaussian inverse \leanref{sym:ellF}{\ensuremath{\ell_h^{\mathrm F}}} through a Fourier multiplier. Both families satisfy the same exact inversion requirement.

\begin{definitionv}{synth_1}[Fourier localization weights and Gaussian inverses]
For \(0<h\le1/4\) and known Gaussian error scale \(\sigma\in[0,1/4]\), define
\[
Q(u)=
\begin{cases}
(\sin u/u)^6,&u\ne0,\\
1,&u=0,
\end{cases}
\qquad
q_h^{\mathrm F}(t)=Q(t/h),\quad t\in\mathbb R.
\]
Use the Fourier-transform convention
\[
\widehat f(\omega)=\int_{\mathbb R}e^{-i\omega t}f(t)\,dt,
\qquad
\mathcal F^{-1}[g](v)=\frac1{2\pi}
\int_{\mathbb R}e^{i\omega v}g(\omega)\,d\omega.
\]
The companion observable inverse weight is
\[
\ell_h^{\mathrm F}(v)
=\mathcal F^{-1}\!\left[
\widehat q_h^{\mathrm F}(\omega)e^{\sigma^2\omega^2/2}
\right](v),\qquad v\in\mathbb R.
\]
For \(Z\sim N(0,1)\), this pair satisfies the exact Gaussian inversion identity
\[
E_Z\!\left[\ell_h^{\mathrm F}(t+\sigma Z)\right]
=q_h^{\mathrm F}(t),\qquad t\in\mathbb R.
\]
At sample size \(n\), the Fourier entries of the public dictionary use the dyadic bandwidths
\[
h=2^{-j},\qquad j\in\mathbb Z,\quad j\ge2,
\qquad n^{-2}\le2^{-j}\le1/4.
\]
\end{definitionv}

\subsection{Public selection and the observable procedure}

The finite public dictionary \leanref{sym:Ddict}{\ensuremath{\mathcal D_{n,\sigma}}} combines these two families with a constant entry. Each entry carries its inverse and its score, so selection compares the same public criterion across localization methods.

\begin{definitionv}{def:weight-dictionary}[Inverse-weight dictionary \(\mathcal D_{n,\sigma}\)]
The finite public dictionary \(\mathcal D_{n,\sigma}\) contains the following weights, each paired with its exact inverse and assigned the public expected-error score \(A_q\):
\begin{itemize}
\item \textbf{(Fourier weights.)} The weights \(q_h^{\mathrm F}\), paired with \(\ell_h^{\mathrm F}\), for dyadic bandwidths \(h\in[n^{-2},1/4]\).
\item \textbf{(Polynomial weights.)} The weights \(q_m^{\mathrm M}\), paired with \(\ell_m^{\mathrm M}\), for odd indices \(m=2^j+1\le n^2\).
\item \textbf{(Constant weight.)} The constant weight \(1\), paired with the inverse \(1\).
\end{itemize}
\end{definitionv}

The constant entry ensures that the dictionary is nonempty. A fixed order resolves ties and defines the selected entry \leanref{sym:qhat}{\ensuremath{\widehat q}} uniquely, including when differently tagged entries have the same latent weight.

\begin{definitionv}{synth_6}[Public dictionary order]
The fixed dictionary order on \(\mathcal D_{n,\sigma}\) in \cref{obj:def:weight-dictionary} is the following order of tagged entries: the constant entry first, then the Fourier entries \(q_{2^{-j}}^{\mathrm F}\) in increasing integer index \(j\), then the polynomial entries \(q_{2^j+1}^{\mathrm M}\) in increasing integer index \(j\). The Fourier indices are those integers \(j\ge2\) satisfying \(n^{-2}\le2^{-j}\le1/4\); the polynomial indices are those integers \(j\ge1\) satisfying \(2^j+1\le n^2\). Each entry carries its paired inverse weight and public score \(A_q(K)\). Entries retain their tags even when their latent weight functions coincide. The first minimizer \(\widehat q\) is the earliest entry in this order attaining \(\min_{q\in\mathcal D_{n,\sigma}}A_q(K)\).
\end{definitionv}

Because the scores and order use public inputs, the selected entry is determined by the sample size and declared parameters. The sample supplies the weighted moments used after selection. To state those formulas, partition the records into deterministic modulo-three blocks \leanref{sym:Iblocks}{\ensuremath{\mathcal I_r}}. The three blocks separately estimate the stratum frequency \leanref{sym:phatx}{\ensuremath{\widehat p_x}}, the weighted denominator \leanref{sym:Dhatx}{\ensuremath{\widehat D_x(q)}}, and the weighted outcome numerator \leanref{sym:Mhatx}{\ensuremath{\widehat M_x(q)}}.

The stratum-specific ratio estimate \leanref{sym:muhatx}{\ensuremath{\widehat\mu_x(q)}} uses a public denominator floor and projection onto the outcome range. Combining these ratios with estimated stratum frequencies gives the weight-specific causal-mean estimate \leanref{sym:thetahatq}{\ensuremath{\widehat\theta_q}}. The following definition supplies the intermediate formulas and the empty-block branch.

\begin{definitionv}{synth_3}[Split-sample ratios]
Index the observed records \(O_i=(X_i,W_i,Y_i)\) by \(i=0,\ldots,n-1\), put \(V_i=W_i-a_0\) with target dose \(a_0=1/2\), and define
\[
\mathcal I_r=\{i\in\{0,\ldots,n-1\}:i\bmod3=r\},
\qquad r\in\{0,1,2\}.
\]
For a public latent weight \(q\) with observable inverse \(\ell_q\), let
\[
I_\kappa(q)=\int_{-1/2}^{1/2}|t|^\kappa q(t)\,dt,
\qquad
b_q(K)=p_{\min}c_{\mathrm{lo}}I_\kappa(q)>0,
\qquad K=(p_{\min},c_{\mathrm{lo}},c_{\mathrm{hi}}).
\]
When all three blocks are nonempty, define, for each \(x\in\mathcal X=\{0,1\}\),
\[
\widehat p_x
=\frac1{|\mathcal I_0|}\sum_{i\in\mathcal I_0}\mathbf1\{X_i=x\},
\]
\[
\widehat D_x(q)
=\frac1{|\mathcal I_1|}\sum_{i\in\mathcal I_1}\mathbf1\{X_i=x\}\ell_q(V_i),
\qquad
\widehat M_x(q)
=\frac1{|\mathcal I_2|}\sum_{i\in\mathcal I_2}\mathbf1\{X_i=x\}Y_i\ell_q(V_i).
\]
Floor the denominator at \(b_q(K)/2\) and project the ratio onto the public outcome range:
\[
\widehat\mu_x(q)
=\Pi_{[0,1]}\!\left(
\frac{\widehat M_x(q)}{\max\{\widehat D_x(q),b_q(K)/2\}}
\right),
\qquad
\Pi_{[0,1]}(u)=\min\{1,\max\{0,u\}\}.
\]
The corresponding total estimator is
\[
\widehat\theta_q
=
\begin{cases}
\displaystyle\sum_{x\in\mathcal X}\widehat p_x\widehat\mu_x(q),
&\min_{r\in\{0,1,2\}}|\mathcal I_r|>0,\\
1/2,&\min_{r\in\{0,1,2\}}|\mathcal I_r|=0.
\end{cases}
\]
This is the ratio construction used by \cref{obj:def:total-estimator}. The floor uses the public constants \(p_{\min}\) and \(c_{\mathrm{lo}}\) of \(K\), which are arbitrary subject to \(0<p_{\min}\le1/2\) and \(0<c_{\mathrm{lo}}\le c_{\mathrm{hi}}\).
\end{definitionv}

The floor makes every computed ratio well defined, and projection keeps each stratum estimate within \([0,1]\). The full selected estimator \leanref{sym:thetahat}{\ensuremath{\widehat\theta}} combines these operations with public score minimization. Its numbered procedure also declares the midpoint output whenever a sample block is empty.

\begin{algorithmv}{def:total-estimator}[Total estimator $\widehat\theta$]
The inputs are the observed records \(O_i=(X_i,W_i,Y_i)\), indexed by \(i=0,\ldots,n-1\), the public class constants \(K=(p_{\min},c_{\mathrm{lo}},c_{\mathrm{hi}})\), and the ordered dictionary \(\mathcal D_{n,\sigma}\) with scores \(A_q(K)\) from \cref{obj:synth_4}.
\begin{enumerate}
\item Form the sample blocks, their sizes, and the centered observed doses:
\[
\mathcal I_r
=\{i\in\{0,\ldots,n-1\}:i\bmod3=r\},
\qquad
N_r=|\mathcal I_r|,
\qquad r\in\{0,1,2\},
\]
\[
V_i=W_i-a_0,\qquad i=0,\ldots,n-1.
\]
\item If \(\min\{N_0,N_1,N_2\}=0\), output \(1/2\) and stop. The remaining steps apply when all three blocks are nonempty.
\item For each dictionary weight \(q\), use the public denominator bound
\[
I_\kappa(q)=\int_{-1/2}^{1/2}|t|^\kappa q(t)\,dt,
\qquad
b_q(K)=p_{\min}c_{\mathrm{lo}}I_\kappa(q).
\]
Compute, for each stratum \(x\in\mathcal X=\{0,1\}\),
\[
\widehat p_x
=\frac1{N_0}\sum_{i\in\mathcal I_0}\mathbf1\{X_i=x\},
\]
\[
\widehat D_x(q)
=\frac1{N_1}\sum_{i\in\mathcal I_1}
\mathbf1\{X_i=x\}\ell_q(V_i),
\qquad
\widehat M_x(q)
=\frac1{N_2}\sum_{i\in\mathcal I_2}
\mathbf1\{X_i=x\}Y_i\ell_q(V_i).
\]
\item Floor each denominator, project the ratio onto \([0,1]\) as in \cref{obj:synth_3}, and combine strata:
\[
\widehat\mu_x(q)
=\min\left\{1,\max\left\{0,
\frac{\widehat M_x(q)}
{\max\{\widehat D_x(q),b_q(K)/2\}}
\right\}\right\},
\qquad
\widehat\theta_q
=\sum_{x\in\mathcal X}\widehat p_x\widehat\mu_x(q).
\]
\item Select the first public-score minimizer in the fixed dictionary order:
\[
\widehat q
=\operatorname*{arg\,min}_{q\in\mathcal D_{n,\sigma}}A_q(K),
\]
with ties resolved by that order.
\item Output \(\widehat\theta=\widehat\theta_{\widehat q}\).
\end{enumerate}
\end{algorithmv}

For \(n\ge3\), every modulo-three block is nonempty, and the estimator uses the selected inverse-weight ratios. At smaller sample sizes, the midpoint branch supplies the declared output. The rate guarantee below concerns uniformly sufficiently large samples and identifies the dictionary family supplying the comparison in each error-scale range.

\begin{theoremv}{thm:observable-upper}[Observable estimator and comparison bounds]
Fix the following public parameters:
\begin{itemize}
\item \textbf{(Class constants.)} \(K\) specifies a minimum stratum mass in \((0,1/2]\) and positive lower and upper weak-design envelope constants, with the lower constant at most the upper constant.
\item \textbf{(Smoothness.)} The smoothness exponent satisfies \(\beta\in(0,1]\).
\item \textbf{(Design decay.)} The design-decay exponent satisfies \(\kappa\in[0,2]\).
\end{itemize}
There exist \(C>0\) and \(n_0\in\mathbb N\), depending on \(K,\beta,\kappa\), such that the following conclusions hold uniformly over every measurable schedule space and potential-path structure in \cref{obj:def:potential-path-space}, every \(n\ge n_0\), and every public error scale \(\sigma\in[0,1/4]\). Let \(L_n=\log(en)\) be the logarithmic sample scale and let \(\rho_{n,\sigma}\) be the absolute-risk frontier in \cref{obj:def:frontier-rate}.

For every structural probability law \(P\) in the model class of \cref{obj:def:model-class} with parameters \(K,\beta,\kappa,\sigma\), under the sampling condition in \cref{obj:ass:iid}, the observed experiment \(Q_P^n\) in \cref{obj:def:observed-experiment} and the selected estimator \(\widehat\theta\) in \cref{obj:def:total-estimator} satisfy
\[
\mathbb E_{Q_P^n}\!\left[
  \left|\widehat\theta-\theta(P)\right|
\right]
\le C\rho_{n,\sigma},
\qquad
\theta(P)=\mathbb E_P[Y(a_0)].
\]
Here \(\theta(P)\) is the population causal mean, \(a_0\) is the prespecified evaluation dose, and \(Y(a_0)\) is the potential outcome obtained by evaluating the random schedule there.

The public dictionary \(\mathcal D_{n,\sigma}\) in \cref{obj:def:weight-dictionary} contains comparison witnesses as follows:
\begin{itemize}
\item \textbf{(Fourier range.)} If \(\sigma\le L_n^{-1/2}\), there is a Fourier entry indexed by a nonnegative integer whose associated inverse pair \((q,\ell_q)\) has public expected-error score \(A_q\) satisfying
\[
A_q\le C\rho_{n,\sigma}.
\]
\item \textbf{(Polynomial range.)} If \(L_n^{-1/2}<\sigma\), there is a polynomial entry indexed by a nonnegative integer whose associated inverse pair \((q,\ell_q)\) has public expected-error score \(A_q\) satisfying
\[
A_q\le C\rho_{n,\sigma}.
\]
\end{itemize}
The pairs and scores are those associated with the entries in \cref{obj:def:weight-dictionary}. For each witness, under the usual measurability and integrability conditions, the latent weight \(q\) and its Gaussian inverse weight \(\ell_q\) satisfy
\[
q(t)\ge0
\quad\text{for every }t\in[-1/2,1/2],
\qquad
0<I_\kappa(q)
:=\int_{-1/2}^{1/2}|t|^\kappa q(t)\,dt<\infty,
\qquad
V_q<\infty,
\]
where \(V_q\) denotes the separately normalized factor-\(16\) second-moment certificate of \cref{obj:lem:public-error-certificate}. Their Gaussian inversion identity holds pointwise:
\[
\int_{\mathbb R}\ell_q(t+\sigma z)
       \frac{e^{-z^2/2}}{\sqrt{2\pi}}\,dz
=q(t)
\quad\text{for every }t\in[-1/2,1/2].
\]
\end{theoremv}

The admissibility properties listed for each witness, namely measurability of \(q\) and \(\ell_q\), nonnegativity of \(q\), integrability of \(|t|^\kappa q(t)\) on \([-1/2,1/2]\) with \(0<I_\kappa(q)\), Gaussian integrability with exact inversion, and \(V_q<\infty\), are conclusions of \cref{obj:thm:observable-upper}, not additional hypotheses. The phrase ``under the usual measurability and integrability conditions'' in its statement refers to these proved properties.

The score expresses the tradeoff between localization bias and the sampling cost of the inverse weight. \Cref{obj:thm:observable-upper} supplies a Fourier comparison when \(\sigma\le L_n^{-1/2}\) and a polynomial comparison when \(\sigma>L_n^{-1/2}\). Public minimization gives the selected entry a score at most that of the relevant comparison. Together with the certificate in \cref{obj:lem:public-error-certificate}, these comparisons explain how one observable procedure attains the stated rate as the known error scale varies from zero to a fixed positive value. The constants and sample-size threshold are common to all error scales and path spaces covered by the theorem.

\subsection{Score-based honest intervals}

Fix a noncoverage probability \(0<\alpha<1\). The connected interval \leanref{sym:Cn}{\ensuremath{C_{n,\sigma}}} uses the selected public score as its radius after division by \(\alpha\), and intersects the resulting interval with the causal target range. The empty-block branch assigns the full target range.

\begin{definitionv}{def:honest-interval}[Score-based interval \(C_{n,\sigma}\)]
The connected interval is
\[
C_{n,\sigma}
=
\begin{cases}
\bigl[\widehat\theta-A_{\widehat q}(K)/\alpha,\,
      \widehat\theta+A_{\widehat q}(K)/\alpha\bigr]\cap[0,1],
      &\min_{r\in\{0,1,2\}}|\mathcal I_r|>0,\\
[0,1],&\min_{r\in\{0,1,2\}}|\mathcal I_r|=0.
\end{cases}
\]
Here \(K\) is public, \(\alpha>0\) is the noncoverage level, and \(\widehat\theta\) is the selected-weight estimator using the blocks
\[
\mathcal I_r
=\{i\in\{0,\ldots,n-1\}:i\bmod 3=r\},
\qquad r\in\{0,1,2\},
\]
and the first minimizer \(\widehat q\) of \(A_q(K)\) in the fixed dictionary order.
\end{definitionv}

The expected-error certificate calibrates the radius through Markov's inequality. Since the causal mean lies in \([0,1]\), intersection with that range preserves coverage. The following result separates honesty at every sample size from the expected-length rate for uniformly sufficiently large samples.

\begin{theoremv}{thm:finite-honesty}[Finite-sample coverage and large-sample length]
Fix public constants \(K=(p_{\min},c_{\mathrm{lo}},c_{\mathrm{hi}})\) and parameters satisfying:
\begin{itemize}
\item \textbf{(Class constants.)} \(0<p_{\min}\le 1/2\) and \(0<c_{\mathrm{lo}}\le c_{\mathrm{hi}}\).
\item \textbf{(Smoothness.)} \(0<\beta\le 1\).
\item \textbf{(Design decay.)} \(0\le\kappa\le 2\).
\item \textbf{(Honesty level.)} \(0<\alpha<1\), where \(\alpha\) is the noncoverage probability.
\end{itemize}
For any public potential-outcome path space as in \cref{obj:def:potential-path-space}, let \(P\) range over the structural laws in \(\mathcal P_{K,\beta,\kappa,\sigma}\) of \cref{obj:def:model-class}. Write
\[
\theta(P)=\mathbb E_P[Y(a_0)],
\qquad
\rho_{n,\sigma}=b_{n,\sigma}^{\beta},
\]
where \(a_0\) is the prespecified latent evaluation dose and \(b_{n,\sigma}\) is the localization scale in \cref{obj:def:frontier-rate}. Let \(Q_P^n\) be the observed sample law in \cref{obj:def:observed-experiment}, and let \(C_{n,\sigma}\) be the interval in \cref{obj:def:honest-interval}, constructed with \(K,\beta,\kappa,\alpha\).

For every such path space, every sample size \(n\ge 0\), every \(\sigma\in[0,1/4]\), and every \(P\in\mathcal P_{K,\beta,\kappa,\sigma}\),
\[
Q_P^n\!\left\{\theta(P)\in C_{n,\sigma}\right\}\ge 1-\alpha.
\]
Moreover, there exist \(C=C(K,\beta,\kappa,\alpha)>0\) and an integer \(n_0=n_0(K,\beta,\kappa,\alpha)\ge 0\) such that, simultaneously for every such path space, every \(n\ge n_0\), every \(\sigma\in[0,1/4]\), and every \(P\in\mathcal P_{K,\beta,\kappa,\sigma}\),
\[
\mathbb E_{Q_P^n}\!\left[\operatorname{length}(C_{n,\sigma})\right]
\le C\,\rho_{n,\sigma}.
\]
\end{theoremv}

The interval converts an expected absolute-error bound into a coverage guarantee using the public score. Its length is bounded by twice the score divided by \(\alpha\), with intersection at the target range providing further truncation. Thus the same dictionary comparisons that control estimation error also control expected interval length. \Cref{obj:thm:finite-honesty} guarantees coverage at each positive sample size, and also at \(n=0\) through the full-range fallback. For \(n=0,1,2\), an empty block triggers that fallback; for \(n\ge3\), the score-based branch applies. The expected-length bound and the estimation bound in \cref{obj:thm:observable-upper} apply once their respective public sample-size thresholds are reached, uniformly over the specified structural laws, error scales, and potential-outcome path spaces. At moderate sample sizes the interval can be uninformative: every dictionary score satisfies \(A_q(K)\ge2\sqrt{3/n}\), because its bias and second-moment terms are nonnegative, so the radius \(A_{\widehat q}(K)/\alpha\) is at least \(1\) whenever \(n\le12/\alpha^2\). Since \(\widehat\theta\in[0,1]\), the interval is then the whole range \([0,1]\); at \(\alpha=0.05\) this happens for every \(n\le4800\).


\section{Observed-law lower bounds}

The estimation and coverage guarantees in \cref{obj:thm:observable-upper,obj:thm:finite-honesty} are complemented by an observed-experiment converse. \Cref{obj:thm:observed-lower} establishes the statistical difficulty of causal evaluation from contaminated-dose records under its stated conditions. Its comparisons concern the sample observed laws \(Q_P^n\) defined in \cref{obj:def:observed-experiment}, and therefore apply directly to procedures based on the available observations. Together with the upper bounds, they supply the lower-bound component of \cref{obj:thm:uniform-frontier}.

The comparison uses the reference and perturbed structural laws \leanref{sym:Palt}{\ensuremath{P_\star,P_+,P_-}} within a bounded Bernoulli witness subclass. These laws share the stratum probabilities, latent-dose design, and Gaussian measurement channel. Their conditional potential means differ through admissible perturbations at the evaluation dose, producing separation of the causal targets. The common design isolates how changes in the conditional mean affect the distribution of observed outcomes and contaminated doses. The structural realization of these witnesses, including the rank coordinate and threshold schedules, is given with the comparison in \cref{obj:thm:observed-lower}.

The statistical argument connects this target separation to closeness of the sample observed laws. When two admissible laws assign separated causal means while generating records that are difficult to distinguish, estimation must incur error on at least one of them. The corresponding interval comparison combines observed-law closeness with the requirement of coverage across the model class. Thus the witness subclass links the causal target to both the estimation criterion in \cref{obj:def:minimax-risk} and the honest-length criterion in \cref{obj:def:honest-length}.

Weighted moment cancellation is the mechanism behind the two inverse-regime comparisons. The packets supplied by \cref{obj:thm:weighted-packet-construction} retain a controlled value at the evaluation point while cancelling moments against the design weight. Gaussian convolution attenuates the resulting signed change in the observed distribution, even as the causal means remain separated. Weighting the cancellation by the latent design makes this mechanism compatible with weak dose support. The two comparisons balance localization and cancellation at the scales relevant to the Gaussian error and the compact dose interval. Their packet formulas, regime-specific choices, and quantitative observed-law bounds are developed in the appendix before use in \cref{obj:thm:observed-lower}.

% Appendix equation labels to assign at their defining displays:
% eq:lower-witness-means -- reference and perturbed conditional means.
% eq:lower-witness-target-separation -- causal-target separation.
% eq:lower-weighted-moment-cancellation -- design-weighted cancellation identity.
% eq:lower-observed-law-comparison -- observed-law divergence comparison.
% eq:lower-gaussian-regime-comparison -- Gaussian inverse-regime tuning and bound.
% eq:lower-compact-regime-comparison -- compact-support inverse-regime tuning and bound.
% Later reader-facing references to these displays must use cleveref.


\section{Discussion}

The rate characterization in \cref{obj:thm:uniform-frontier} gives measurement precision a concrete interpretation for estimating the population causal mean at a prespecified dose. Under the structural model in \cref{obj:def:model-class}, with \(c_{\mathrm{lo}}\le(\kappa+1)2^\kappa\le c_{\mathrm{hi}}\) and fixed smoothness and design-decay exponents in their stated ranges, it describes how sample size and the known Gaussian error scale jointly determine attainable precision. Its application therefore rests on substantive justification for the compact latent-dose support, the weak-design envelope, and the causal identification conditions. Within this model, improving measurement precision can move the experiment between three regimes, with comparable rates throughout the transition neighborhoods specified in \cref{obj:thm:uniform-frontier}.

The design-decay exponent \(\kappa\) measures the scarcity of latent doses near the evaluation point. The envelope in \cref{obj:ass:weak-design} implies that the probability of a sufficiently small dose neighborhood of radius \(h\), within either stratum, has order \(h^{\kappa+1}\). Larger values of \(\kappa\) thus represent weaker local support. Under \(c_{\mathrm{lo}}\le(\kappa+1)2^\kappa\le c_{\mathrm{hi}}\), in the direct regime, this scarcity combines with mean smoothness through the effective exponent \(d=2\beta+\kappa+1\), yielding absolute-error and honest-length orders \(n^{-\beta/d}\). For fixed smoothness, increasing design decay slows this order and raises the measurement-error threshold \(n^{-1/d}\) below which direct estimation is attainable. The wider direct-regime range accompanies a coarser attainable localization scale.

In the intermediate Gaussian regime, measurement precision enters more directly. Under \(c_{\mathrm{lo}}\le(\kappa+1)2^\kappa\le c_{\mathrm{hi}}\), the common uncertainty order is
\[
\left\{\frac{\sigma}{\sqrt{\log(e+n\sigma^d)}}\right\}^{\beta},
\qquad
n^{-1/d}<\sigma\le L_n^{-1/2},
\]
with the notation defined in \cref{obj:def:frontier-rate}. Here the design-decay exponent enters through the logarithmic term as well as the first transition threshold. At a common sample size and error scale lying in the intermediate regime for the exponents being compared, a larger \(\kappa\) reduces \(n\sigma^d\) and increases the displayed uncertainty scale. The displayed scale depends on measurement precision, and once the direct regime is reached, the error-free weak-design order characterized in \cref{obj:prop:error-free-reduction} is attained. Further reductions within the direct regime share that minimax order.

Under the envelope calibration \(c_{\mathrm{lo}}\le(\kappa+1)2^\kappa\le c_{\mathrm{hi}}\), known compact support governs the eventual order when the measurement-error scale remains fixed and positive. Under \cref{obj:ass:latent-support}, every fixed \(\sigma\in(0,1/4]\) eventually lies in the compact-support regime as sample size increases. For this calibrated class, both optimal criteria then have order
\[
\left(\frac{\log\log n}{\log n}\right)^\beta,
\]
as established in \cref{obj:thm:uniform-frontier}. The design-decay exponent determines the direct and intermediate regimes, while smoothness determines the exponent in this eventual fixed-error order. Within the calibrated class, two fixed positive error scales therefore share this asymptotic order, with comparison constants allowed to depend on the error scale. The uniform characterization gives a more detailed account of improving precision with sample size: the compact-support expression depends on \(\sigma\) through \(\log(e+\sigma^2L_n)\), and sufficiently rapid improvement leads through the intermediate regime to the direct order.

Under the envelope calibration \(c_{\mathrm{lo}}\le(\kappa+1)2^\kappa\le c_{\mathrm{hi}}\), the matching estimation and interval-length orders concern uncertainty about the same scalar causal mean, \(\theta(P)\). Expected absolute error measures point-estimation accuracy, while honest expected length measures the width required for connected intervals with coverage of at least \(1-\alpha\) uniformly over the structural class; these criteria are defined in \cref{obj:def:minimax-risk,obj:def:honest-length}. For the calibrated class, their common rate connects improved measurement precision to both forms of uncertainty. The observable estimator and interval construction in \cref{obj:thm:observable-upper,obj:thm:finite-honesty} attain the corresponding upper bounds, and the observed-law comparison in \cref{obj:thm:observed-lower} supplies the converse under this calibration. Together, these results characterize the order of precision available for causal evaluation under the stated support, design, smoothness, and measurement conditions and the envelope calibration.


\section{Limitations and future work}

The characterization in \cref{obj:thm:uniform-frontier} concerns a scalar causal mean at a prespecified latent dose, with fixed public class constants, fixed mean-smoothness and design-decay exponents, a public Gaussian measurement-error scale, and the structural restrictions in \cref{obj:def:model-class}. The proved range \(\sigma\in[0,1/4]\) is an explicit restriction; the analysis does not claim the same uniform constants beyond that cap. The requirement \(\kappa\le2\) is also substantive rather than a normalization: it is used in the dictionary score-rate bounds of \cref{obj:lem:dictionary-score-rate} behind the observable upper bound, and in the comparison with the direct-regression benchmark in \cref{obj:prop:error-free-reduction}. The observable construction and coverage guarantee in \cref{obj:thm:observable-upper,obj:thm:finite-honesty} operate within that specification. These results provide a benchmark for extensions in which model parameters are learned, covariate adjustment is richer, the measurement-error range is widened, or inference covers several doses.

Adaptation over mean smoothness and design decay would require selecting localization weights while accounting for uncertainty about both exponents. The public-score selection in \cref{obj:def:total-estimator} uses the specified class parameters; adaptation would instead require observable criteria that remain valid across candidate classes. Point estimation and honest inference raise distinct questions here. For estimation, the task is to determine the attainable accuracy across classes and any cost of adaptation. For intervals, the task is to establish whether coverage over a union of classes can coexist with shorter expected length on individual classes, and which additional restrictions would permit that improvement.

An unknown measurement-error scale introduces a further source of uncertainty into Gaussian inversion. Replicate measurements or an external calibration sample could supply information about the scale, subject to assumptions linking those measurements to the main sample. An extension would need to establish joint identification, propagate calibration uncertainty through the inverse weights, and compare the required calibration precision with the localization scale. Honest intervals would then have to account for uncertainty from both samples. The transition between measurement-error regimes makes the relative sizes of the calibration and outcome samples a substantive statistical question.

Richer covariates would extend the current two-stratum adjustment to settings with many strata or continuous confounders. Such an analysis would require conditions on covariate support, conditional dose densities, and the regularity of the conditional potential-outcome mean. The statistical problem would combine dose inversion with estimation and aggregation across covariate values. Relevant questions include how covariate dimension affects attainable accuracy, how to control sparsely represented covariate regions, and how uncertainty in covariate adjustment enters an honest interval for the population causal mean.

Simultaneous dose inference would replace evaluation at one prespecified dose with inference over a collection or interval of doses. This extension requires a model describing how weak dose support varies across evaluation points, including its behavior near support boundaries. Uniform approximation and stochastic-error bounds would then be needed for an entire estimated curve. The corresponding questions concern the attainable width of honest confidence bands, the cost of simultaneous coverage, and the relationship between local design decay and uncertainty across doses.

Empirical use also calls for validation of the support and density envelopes. Because the dose is latent, checking these restrictions from contaminated measurements requires an explicit measurement model or additional information. Validation measurements, substantive support restrictions, and sensitivity analysis could help assess plausible specifications. A statistical treatment would need to incorporate uncertainty in estimated support boundaries and envelope constants into estimation and coverage guarantees. Distinguishing assumptions supported by validation data from restrictions supplied by substantive knowledge would make those guarantees interpretable in an application.

The exact public certificates also call for a numerical study separate from the rate theory proved here. A reproducible implementation should evaluate the Fourier and polynomial inverse weights with documented quadrature and stability controls, then report the selected dictionary entry, score, and interval radius over grids of sample sizes and error scales for representative calibrated \(K,\beta,\kappa\). Such diagnostics would measure finite-sample informativeness of the explicit constants without changing the uniform asymptotic and coverage guarantees established in this paper.

Dietary-intake applications are a direction requiring application-specific validation and analysis. A study would need to define the latent intake exposure and causal contrast, assess the measurement-error distribution and its independence conditions, and justify the covariates used for causal adjustment. Replicate assessments or reference measurements could inform the measurement model, while the observed study population would determine the relevant dose-support specification. Connecting these ingredients to the present benchmark requires an analysis of the resulting observed experiment and of uncertainty contributed by exposure calibration and confounder adjustment.

The secondary Gaiffas direct-regression comparison in \cref{obj:prop:error-free-reduction} has a narrower evidentiary status than the paper's own frontier. Its source interface was attested, but an independent source-retrieval check was inconclusive. The comparison is therefore retained as a benchmark for the error-free exponent, with its precise assumptions and formalization scope stated locally at the proposition.


\appendix

\section*{Appendices}

\section{Identification and observable risk certificates}

The auxiliary results connect the structural causal model to population localization and observable estimation. The argument begins with conditional means at the realized dose, then establishes uniqueness under Gaussian convolution and positivity of localized population denominators. A public expected-error certificate subsequently links these population quantities to the split-sample estimator, while dictionary comparisons supply the uniform rate guarantees in \cref{obj:thm:observable-upper,obj:thm:finite-honesty}.

\subsection{Random-dose means and identification}

Conditional unconfoundedness in \cref{obj:ass:conditional-unconfoundedness} concerns the measurable potential-outcome schedule. The regression identity needed for identification and weighting evaluates that schedule at the random latent dose. The following result supplies this bridge, including its integrated form for bounded measurable multipliers.

\begin{lemmav}{lem:realized-dose-mean}[Realized dose mean identities]
Suppose the following conditions hold:
\begin{itemize}
\item \textbf{(Public model primitives.)} Fix a potential-outcome path space as in \cref{obj:def:potential-path-space} and public class constants \(K\) as in \cref{obj:def:model-class}.
\item \textbf{(Parameter ranges.)} Let \(\beta\in(0,1]\), \(\kappa\in[0,2]\), and \(\sigma\in[0,1/4]\).
\item \textbf{(Structural law.)} Let \(P\) be a structural probability law in the model class \(\mathcal P_{K,\beta,\kappa,\sigma}\) of \cref{obj:def:model-class}. Write \(\mathcal X=\{0,1\}\) for the stratum space, and let \(X\in\mathcal X\), \(A\in\mathbb R\), \(Z\in\mathbb R\), and \(Y\in\mathbb R\) denote its stratum, latent dose, classical error, and realized outcome coordinates, respectively.
\end{itemize}
With \(\mu_x(a)\) denoting the conditional potential-outcome mean defined in \cref{obj:def:conditional-potential-mean}, the following statements hold \(P\)-almost surely:
\[
0\le Y\le1,\qquad
E_P[Y\mid A,X,Z]=\mu_X(A),\qquad
E_P[Y\mid A,X]=\mu_X(A).
\]
Moreover, for every bounded measurable function
\(F:\mathbb R\times\mathcal X\times\mathbb R\to\mathbb R\),
\[
E_P\!\left[F(A,X,Z)Y\right]
=
E_P\!\left[F(A,X,Z)\mu_X(A)\right].
\]
\end{lemmav}

\begin{proof}[Proof of \cref{obj:lem:realized-dose-mean}]
\begin{enumerate}
\item Define the projection onto the dose interval and the normalized stratum laws by
\[
\pi_{[0,1]}:\mathbb R\longrightarrow[0,1],
\qquad
\pi_{[0,1]}(a)=\max\{0,\min\{1,a\}\},
\]
\[
p_x=P(X=x),
\qquad
P_x=\frac{1}{p_x}P|_{\{X=x\}},
\qquad
\gamma_x=\mathcal L_{P_x}(A),
\qquad
\lambda_x=\mathcal L_{P_x}(Y(\cdot)),
\qquad x\in\mathcal X.
\]
By \cref{obj:ass:stratum-overlap}, \(p_x\ge p_{\min}>0\), so \(P_x\), \(\gamma_x\), and \(\lambda_x\) are probability measures. The law \(P_x\) is concentrated on \(\{X=x\}\).

Fix \(x\in\mathcal X\). For measurable sets
\[
\mathcal C_{\mathcal S}\in\mathcal E,
\qquad
\mathcal C_A\in\mathcal B(\mathbb R),
\]
\cref{obj:ass:conditional-unconfoundedness} gives
\[
\begin{aligned}
P_x\{Y(\cdot)\in\mathcal C_{\mathcal S},\,A\in\mathcal C_A\}
&=\frac{P\{Y(\cdot)\in\mathcal C_{\mathcal S},\,A\in\mathcal C_A,\,X=x\}}{p_x}\\
&=\frac{P\{Y(\cdot)\in\mathcal C_{\mathcal S},\,X=x\}}{p_x}
  \frac{P\{A\in\mathcal C_A,\,X=x\}}{p_x}\\
&=\lambda_x(\mathcal C_{\mathcal S})\gamma_x(\mathcal C_A).
\end{aligned}
\]
Thus equality on measurable rectangles yields
\[
\mathcal L_{P_x}(A,Y(\cdot))=\gamma_x\otimes\lambda_x.
\]
For every \(a\in\mathbb R\), \cref{obj:ass:bounded-potential-outcomes} transfers to the normalized restriction and gives
\[
\lambda_x\bigl\{f\in\mathcal S:
e(f,\pi_{[0,1]}(a))\notin[0,1]\bigr\}=0.
\]
The evaluation map from \cref{obj:def:potential-path-space} is jointly measurable. Integrating these fixed-dose assertions under the product law therefore gives
\[
\int_{\mathbb R}\int_{\mathcal S}
\mathbf1\{e(f,\pi_{[0,1]}(a))\notin[0,1]\}
\,\lambda_x(df)\,\gamma_x(da)=0,
\]
and hence
\[
P_x\{e(Y(\cdot),\pi_{[0,1]}(A))\in[0,1]\}=1.
\]
By \cref{obj:ass:latent-support}, \(\pi_{[0,1]}(A)=A\) almost surely under \(P_x\). Consistency from \cref{obj:ass:consistency} consequently yields
\[
P_x\{Y\in[0,1]\}=1.
\]
Combining the two strata,
\[
P\{Y\notin[0,1]\}
=\sum_{x\in\mathcal X}p_xP_x\{Y\notin[0,1]\}=0.
\]
% lean: CausalSmith.Stat.NoisydoseWeakdesignTransition.realized_outcome_mem_Icc

\item Define a mean extension on the entire real dose axis by
\[
\widetilde\mu_x(a)=\mu_x(\pi_{[0,1]}(a)),
\qquad (a,x)\in\mathbb R\times\mathcal X.
\]
The parameter restriction gives \(\beta>0\), and
\cref{obj:ass:holder-mean} supplies
\[
|\mu_x(a)-\mu_x(a')|\le |a-a'|^\beta,
\qquad a,a'\in[0,1],\quad x\in\mathcal X.
\]
Thus each \(\mu_x\) is continuous on \([0,1]\). Composition with the continuous projection, together with the discrete stratum space, makes
\((a,x)\mapsto\widetilde\mu_x(a)\) measurable.
By \cref{obj:ass:mean-range},
\[
0\le\widetilde\mu_x(a)\le1,
\qquad (a,x)\in\mathbb R\times\mathcal X.
\]
In particular,
\[
|Y|\le1\quad P\text{-almost surely},
\qquad
|\widetilde\mu_X(A)|\le1.
\]
Both random variables are measurable and integrable under the probability law \(P\).
Finally, latent support gives
\[
\widetilde\mu_X(A)=\mu_X(A)
\qquad P\text{-almost surely}.
\]
% lean: CausalSmith.Stat.NoisydoseWeakdesignTransition.projectedMean_measurable
% lean: CausalSmith.Stat.NoisydoseWeakdesignTransition.realized_outcome_integrable
% lean: CausalSmith.Stat.NoisydoseWeakdesignTransition.projectedMean_random_integrable
% lean: CausalSmith.Stat.NoisydoseWeakdesignTransition.realized_dose_mean (hactual)

\item Fix a bounded measurable test function
\[
F:\mathbb R\times\mathcal X\times\mathbb R\longrightarrow\mathbb R,
\qquad
|F(a,x,z)|\le B_F,
\qquad B_F\in[0,\infty).
\]
The preceding bounds imply
\[
|F(A,X,Z)Y|
=|F(A,X,Z)|\,|Y|\le B_F
\quad P\text{-almost surely},
\]
\[
|F(A,X,Z)\widetilde\mu_X(A)|
=|F(A,X,Z)|\,|\widetilde\mu_X(A)|\le B_F.
\]
These measurable products are integrable. Splitting the first expectation over the two strata gives
\[
\begin{aligned}
E_P[F(A,X,Z)Y]
&=\sum_{x\in\mathcal X}
 E_P[\mathbf1\{X=x\}F(A,X,Z)Y]\\
&=\sum_{x\in\mathcal X}p_x
 \int F(A,x,Z)Y\,dP_x.
\end{aligned}
\]
Here the second equality follows from the definition of \(P_x\) and \(X=x\) almost surely under \(P_x\). The same calculation gives
\[
E_P[F(A,X,Z)\widetilde\mu_X(A)]
=\sum_{x\in\mathcal X}p_x
 \int F(A,x,Z)\widetilde\mu_x(A)\,dP_x.
\]
% lean: CausalSmith.Stat.NoisydoseWeakdesignTransition.integral_eq_stratum_sum
% lean: CausalSmith.Stat.NoisydoseWeakdesignTransition.bounded_test_mean

\item Fix a stratum \(x\). Define the error law, clipped evaluation, and product integrand by
\[
\zeta=\mathcal L_P(Z),
\]
\[
e_{\mathrm{bd}}:\mathcal S\times\mathbb R\longrightarrow[0,1],
\qquad
e_{\mathrm{bd}}(f,a)
=\max\{0,\min\{1,e(f,\pi_{[0,1]}(a))\}\},
\]
\[
G_{x,F}:\mathbb R\times\mathbb R\times\mathcal S\longrightarrow\mathbb R,
\qquad
G_{x,F}(z,a,f)=F(a,x,z)e_{\mathrm{bd}}(f,a).
\]
These functions are measurable, and
\[
0\le e_{\mathrm{bd}}(f,a)\le1,
\qquad
|G_{x,F}(z,a,f)|
\le |F(a,x,z)|\le B_F.
\]
Thus \(G_{x,F}\) is integrable under \(\zeta\otimes(\gamma_x\otimes\lambda_x)\).

Error independence from \cref{obj:ass:error-independence} survives restriction to the positive-probability stratum. Indeed, for
\[
\mathcal C_Z\in\mathcal B(\mathbb R),
\qquad
\mathcal C_{A,\mathcal S}\in\mathcal B(\mathbb R)\otimes\mathcal E,
\]
\[
\begin{aligned}
&P_x\{Z\in\mathcal C_Z,\,(A,Y(\cdot))\in\mathcal C_{A,\mathcal S}\}\\
&\quad=
\frac{P\{Z\in\mathcal C_Z,\,(A,Y(\cdot))\in\mathcal C_{A,\mathcal S},\,X=x\}}{p_x}\\
&\quad=
\zeta(\mathcal C_Z)
\frac{P\{(A,Y(\cdot))\in\mathcal C_{A,\mathcal S},\,X=x\}}{p_x}\\
&\quad=
\zeta(\mathcal C_Z)P_x\{(A,Y(\cdot))\in\mathcal C_{A,\mathcal S}\}.
\end{aligned}
\]
Equality on rectangles and the schedule--dose product law already established yield
\[
\mathcal L_{P_x}(Z,A,Y(\cdot))
=\zeta\otimes(\gamma_x\otimes\lambda_x).
\]
The random-evaluation bound, latent support, and consistency also give
\[
e_{\mathrm{bd}}(Y(\cdot),A)=Y
\qquad P_x\text{-almost surely}.
\]
Consequently, change of variables under this product law and Fubini's theorem give
\[
\int F(A,x,Z)Y\,dP_x
=
\int_{\mathbb R}\int_{\mathbb R}\int_{\mathcal S}
G_{x,F}(z,a,f)\,\lambda_x(df)\,\gamma_x(da)\,\zeta(dz).
\]

For every \(a\in\mathbb R\), the fixed-dose bound established above makes clipping leave the evaluation unchanged \(\lambda_x\)-almost everywhere. The normalized restriction formula and
\cref{obj:def:conditional-potential-mean} therefore give
\[
\begin{aligned}
\int_{\mathcal S}e_{\mathrm{bd}}(f,a)\,\lambda_x(df)
&=\int_{\mathcal S}e(f,\pi_{[0,1]}(a))\,\lambda_x(df)\\
&=\frac{E_P[\mathbf1\{X=x\}Y(\pi_{[0,1]}(a))]}{p_x}\\
&=\mu_x(\pi_{[0,1]}(a))
=\widetilde\mu_x(a).
\end{aligned}
\]
Integrating out the schedule thus gives
\[
\int F(A,x,Z)Y\,dP_x
=
\int_{\mathbb R}\int_{\mathbb R}
F(a,x,z)\widetilde\mu_x(a)\,\gamma_x(da)\,\zeta(dz).
\]
Moreover,
\[
|F(a,x,z)\widetilde\mu_x(a)|\le B_F,
\qquad a,z\in\mathbb R,
\]
so this measurable integrand is integrable. Marginalizing the product law, using
\(\lambda_x(\mathcal S)=1\), gives
\[
\mathcal L_{P_x}(Z,A)=\zeta\otimes\gamma_x.
\]
It follows that
\[
\int F(A,x,Z)Y\,dP_x
=
\int F(A,x,Z)\widetilde\mu_x(A)\,dP_x.
\]
% lean: CausalSmith.Stat.NoisydoseWeakdesignTransition.bounded_test_mean_stratum

Substituting this equality into the two stratum sums proves, for every bounded measurable \(F\),
\[
E_P[F(A,X,Z)Y]
=
E_P[F(A,X,Z)\widetilde\mu_X(A)].
\]
% lean: CausalSmith.Stat.NoisydoseWeakdesignTransition.bounded_test_mean

\item Apply this identity to indicator tests. For every
\[
\mathcal C_{AXZ}\in\mathcal B(\mathbb R\times\mathcal X\times\mathbb R),
\]
it gives
\[
\int_{\{(A,X,Z)\in\mathcal C_{AXZ}\}}Y\,dP
=
\int_{\{(A,X,Z)\in\mathcal C_{AXZ}\}}\widetilde\mu_X(A)\,dP.
\]
Every event in \(\sigma(A,X,Z)\) has this preimage form. The candidate
\(\widetilde\mu_X(A)\) is integrable and \(\sigma(A,X,Z)\)-measurable, so the defining set-integral characterization and uniqueness of conditional expectation yield
\[
E_P[Y\mid A,X,Z]=\widetilde\mu_X(A)
=\mu_X(A)
\qquad P\text{-almost surely},
\]
where the second equality uses latent support.
% lean: CausalSmith.Stat.NoisydoseWeakdesignTransition.condExp_eq_of_bounded_tests
% lean: CausalSmith.Stat.NoisydoseWeakdesignTransition.realized_dose_mean (hXZ)

\item For any bounded measurable function
\[
F_{AX}:\mathbb R\times\mathcal X\longrightarrow\mathbb R,
\]
define its extension by
\[
F_{AXZ}:\mathbb R\times\mathcal X\times\mathbb R\longrightarrow\mathbb R,
\qquad
F_{AXZ}(a,x,z)=F_{AX}(a,x).
\]
This extension is measurable and has the same bound. The bounded-test identity therefore supplies
\[
E_P[F_{AX}(A,X)Y]
=
E_P[F_{AX}(A,X)\widetilde\mu_X(A)].
\]
The same indicator-test characterization, now for \(\sigma(A,X)\), gives
\[
E_P[Y\mid A,X]=\widetilde\mu_X(A)
=\mu_X(A)
\qquad P\text{-almost surely}.
\]
Finally, replacing the projected mean by its almost-surely equal value in the original bounded-test identity gives
\[
E_P[F(A,X,Z)Y]
=
E_P[F(A,X,Z)\widetilde\mu_X(A)]
=
E_P[F(A,X,Z)\mu_X(A)].
\]
% lean: CausalSmith.Stat.NoisydoseWeakdesignTransition.condExp_eq_of_bounded_tests
% lean: CausalSmith.Stat.NoisydoseWeakdesignTransition.realized_dose_mean
\end{enumerate}
\end{proof}

The conditional identities make the potential-outcome mean applicable to the observed outcome at the realized dose, including conditioning on the error coordinate. The almost-sure outcome bound also supplies the boundedness used in the observable moment certificates below. Joint measurability in \cref{obj:def:potential-path-space} gives random-dose evaluation its measurable meaning.

For the convolution argument, let \(\Phi_\sigma\) denote the centered Gaussian probability measure with variance \(\sigma^2\). At zero error scale, set \(\Phi_0=\delta_0\), where \(\delta_0\) is the unit point mass at zero. For a finite Borel measure \(\lambda\), its convolution with this error measure is defined, for each Borel set \(B\subseteq\mathbb R\), by
\[
(\lambda*\Phi_\sigma)(B)
=
\int_{\mathbb R}\int_{\mathbb R}
\mathbf 1\{t+u\in B\}\,d\Phi_\sigma(u)\,d\lambda(t).
\]
The centered latent-dose support in \cref{obj:ass:latent-support} motivates the following uniqueness statement for finite measures supported on \([-1/2,1/2]\).

\begin{lemmav}{lem:compact-gaussian-convolution-injective}[Injectivity of compact Gaussian convolution]
Let \(\sigma\in\mathbb R\), and let \(\Phi_\sigma\) denote the centered Gaussian probability measure with variance \(\sigma^2\), with \(\Phi_0=\delta_0\), the unit point mass at zero. Let \(\mu_1,\mu_2,\nu_1,\nu_2\) be nonnegative Borel measures on \(\mathbb R\). Suppose that:
\begin{itemize}
\item \textbf{(Error scale.)} \(0\le \sigma\le 1/4\).
\item \textbf{(Finiteness and support.)} Each of \(\mu_1,\mu_2,\nu_1,\nu_2\) is finite and assigns zero mass to \(\mathbb R\setminus[-1/2,1/2]\).
\item \textbf{(Convolution equality.)}
\[
\mu_1*\Phi_\sigma+\nu_2*\Phi_\sigma
=
\nu_1*\Phi_\sigma+\mu_2*\Phi_\sigma.
\]
\end{itemize}
Then
\[
\mu_1+\nu_2=\nu_1+\mu_2.
\]
\end{lemmav}

\begin{proof}[Proof of \cref{obj:lem:compact-gaussian-convolution-injective}]
\begin{enumerate}
\item Define the finite nonnegative Borel measures
\[
\lambda_{\mathrm L}:=\mu_1+\nu_2,
\qquad
\lambda_{\mathrm R}:=\nu_1+\mu_2.
\]
Distributivity of convolution over addition and the assumed convolution equality give
\[
\lambda_{\mathrm L}*\Phi_\sigma
=\mu_1*\Phi_\sigma+\nu_2*\Phi_\sigma
=\nu_1*\Phi_\sigma+\mu_2*\Phi_\sigma
=\lambda_{\mathrm R}*\Phi_\sigma.
\]
% lean: Measure.add_conv

\item For any finite nonnegative Borel measure \(\lambda\) on \(\mathbb R\), define its characteristic function by
\[
\widehat{\lambda}(\xi)
:=\int_{\mathbb R}e^{i\xi\tau}\,d\lambda(\tau),
\qquad \xi\in\mathbb R.
\]
Taking characteristic functions in the preceding equality yields
\[
\widehat{\lambda_{\mathrm L}*\Phi_\sigma}(\xi)
=
\widehat{\lambda_{\mathrm R}*\Phi_\sigma}(\xi).
\]
The characteristic function of a convolution is the product of the characteristic functions, so
\[
\widehat{\lambda_{\mathrm L}}(\xi)\widehat{\Phi_\sigma}(\xi)
=
\widehat{\lambda_{\mathrm R}}(\xi)\widehat{\Phi_\sigma}(\xi)
\qquad(\xi\in\mathbb R).
\]
% lean: charFun_conv

\item The characteristic-function formula for a centered Gaussian gives
\[
\widehat{\Phi_\sigma}(\xi)
=\exp\!\left(-\frac{\sigma^2\xi^2}{2}\right)\ne0
\qquad(\xi\in\mathbb R),
\]
because the complex exponential never vanishes. This formula includes \(\sigma=0\), when the factor equals \(1\). Cancelling this nonzero factor therefore gives
\[
\widehat{\lambda_{\mathrm L}}(\xi)
=
\widehat{\lambda_{\mathrm R}}(\xi)
\qquad(\xi\in\mathbb R).
\]
% lean: charFun_gaussianReal
% lean: Complex.exp_ne_zero
% lean: mul_right_cancel₀

\item Finite Borel measures on \(\mathbb R\) are determined by their characteristic functions. Hence
\[
\lambda_{\mathrm L}=\lambda_{\mathrm R},
\qquad\text{that is,}\qquad
\mu_1+\nu_2=\nu_1+\mu_2.
\]
% lean: Measure.ext_of_charFun
\end{enumerate}
\end{proof}

The grouped equality accommodates differences of finite nonnegative measures while retaining a statement entirely in terms of such measures. Its mathematical role is uniqueness of latent measures from their Gaussian-convolved counterparts. The formulation includes the error-free case through the point-mass convention.

Population localization additionally requires a positive denominator. Here positivity belongs to the latent weight \(q\), whereas empirical moments use its observable inverse \(\ell_q\). Exact Gaussian inversion connects the two under the stated integrability conditions. The next result records the denominator and approximation bounds needed for this construction.

\begin{lemmav}{lem:positive-ratio}[Positive denominators and approximation]
Fix a potential-path space \((\mathcal S,\mathcal E)\) as in \cref{obj:def:potential-path-space}. Let \(\mathcal X\) consist of two strata, and let \(P\) be a structural probability law with coordinates \(X\in\mathcal X\), \(A,Z,Y,U\in\mathbb R\), and potential-outcome schedule \(Y(\cdot)\in\mathcal S\). Set
\[
a_0=\frac12,\qquad t=A-a_0,\qquad W=A+\sigma Z,\qquad V=W-a_0=t+\sigma Z.
\]
Let \(q:\mathbb R\to\mathbb R\) be a public latent weight and \(\ell_q:\mathbb R\to\mathbb R\) its observable inverse weight. Define the weighted moments by
\[
I_\gamma(q)=\int_{-1/2}^{1/2}|t|^\gamma q(t)\,dt.
\]
Suppose the following conditions hold:
\begin{itemize}
\item \textbf{(Public constants.)} The constants \(K=(p_{\min},c_{\mathrm{lo}},c_{\mathrm{hi}})\) satisfy
\[
0<p_{\min}\le\frac12,\qquad 0<c_{\mathrm{lo}}\le c_{\mathrm{hi}}.
\]
\item \textbf{(Parameter ranges.)}
\[
\beta\in(0,1],\qquad \kappa\in[0,2],\qquad \sigma\in[0,1/4].
\]
\item \textbf{(Structural model.)} The law \(P\) satisfies \cref{obj:def:model-class} with path space \((\mathcal S,\mathcal E)\), constants \(K\), and parameters \((\beta,\kappa,\sigma)\).
\item \textbf{(Admissible weight.)} The functions \(q\) and \(\ell_q\) are measurable, and
\[
q(t)\ge0\quad\text{for every }t\in[-1/2,1/2],
\qquad 0<I_\kappa(q)<\infty.
\]
\item \textbf{(Exact Gaussian inversion.)} For every \(t\in[-1/2,1/2]\), the function \(z\mapsto\ell_q(t+\sigma z)\) is absolutely integrable under the standard Gaussian law, and
\[
\int_{\mathbb R}\ell_q(t+\sigma z)\frac{e^{-z^2/2}}{\sqrt{2\pi}}\,dz=q(t).
\]
The inverse weight is also absolutely integrable under the structural law:
\[
E_P[|\ell_q(V)|]<\infty.
\]
\end{itemize}
For every \(x\in\mathcal X\), the denominator is strictly positive and satisfies
\[
0<E_P[\mathbf1\{X=x\}\ell_q(V)],
\qquad
E_P[\mathbf1\{X=x\}\ell_q(V)]
\ge p_{\min}c_{\mathrm{lo}}I_\kappa(q).
\]
Consequently, the localized population ratio
\[
m_{x,q}
=
\frac{E_P[\mathbf1\{X=x\}Y\ell_q(V)]}
     {E_P[\mathbf1\{X=x\}\ell_q(V)]}
\]
is well defined and, with \(\mu_x(a_0)\) as in \cref{obj:def:conditional-potential-mean}, obeys
\[
|m_{x,q}-\mu_x(a_0)|
\le
\frac{c_{\mathrm{hi}}}{c_{\mathrm{lo}}}
\frac{I_{\kappa+\beta}(q)}{I_\kappa(q)}.
\]
These conclusions hold throughout the stated parameter ranges, including \(\sigma=0\).
\end{lemmav}

\begin{proof}[Proof of \cref{obj:lem:positive-ratio}]
\begin{enumerate}
\item Fix \(x\in\mathcal X\). Define the centered latent dose and stratum probability by
\[
T:=A-a_0,\qquad p_x:=P(X=x).
\]
Choose a measurable conditional density \(g_x:\mathbb R\to\mathbb R\) supplied by \cref{obj:ass:weak-design}, so that
\[
\mathcal L_P(T\mid X=x)(dt)
=
g_x(t)\mathbf1_{[-1/2,1/2]}(t)\,dt,
\qquad
c_{\mathrm{lo}}|t|^\kappa\le g_x(t)\le c_{\mathrm{hi}}|t|^\kappa
\]
for Lebesgue-almost every \(t\in[-1/2,1/2]\). To express the mean on the whole real line, define
\[
\bar\mu_x:\mathbb R\to\mathbb R,
\qquad
\bar\mu_x(a):=\mu_x\!\left(\min\{1,\max\{0,a\}\}\right).
\]
This extension is measurable by the Hölder condition, and \cref{obj:ass:mean-range} gives
\[
0\le\bar\mu_x(a)\le1\qquad(a\in\mathbb R).
\]
Define the latent weighted integrals
\[
L_{x,q}:=\int_{-1/2}^{1/2}g_x(t)q(t)\,dt,
\qquad
N_{x,q}:=\int_{-1/2}^{1/2}
g_x(t)\bar\mu_x(a_0+t)q(t)\,dt.
\]
Multiplying the design envelopes by the nonnegative weight gives
\[
0\le c_{\mathrm{lo}}|t|^\kappa q(t)
\le g_x(t)q(t)
\le c_{\mathrm{hi}}|t|^\kappa q(t)
\]
almost everywhere on the integration interval. The upper bound and the finite base moment ensure integrability of \(g_xq\); boundedness of \(\bar\mu_x\) then ensures integrability of the numerator integrand. Integrating the lower bound yields
\[
L_{x,q}\ge c_{\mathrm{lo}} I_\kappa(q)>0.
\]
Also, \cref{obj:ass:stratum-overlap} supplies
\[
p_x\ge p_{\min}>0.
\]
% lean: CausalSmith.Stat.NoisydoseWeakdesignTransition.latent_weight_lower

\item Define the normalized stratum law and the observed denominator by
\[
P_x:=\frac{1}{p_x}P|_{\{X=x\}},
\qquad
D_{x,q}:=E_P[\mathbf1\{X=x\}\ell_q(V)].
\]
Here \(P_x\) is a probability law on the structural coordinate space. Restriction and normalization preserve the assumed absolute integrability of \(\ell_q(V)\).

The error independence in \cref{obj:ass:error-independence} gives, for Borel sets \(C_Z,C_T\subseteq\mathbb R\),
\[
\begin{aligned}
P_x(Z\in C_Z,T\in C_T)
&=\frac{P(Z\in C_Z,X=x,T\in C_T)}{p_x}\\
&=P(Z\in C_Z)\frac{P(X=x,T\in C_T)}{p_x}\\
&=P(Z\in C_Z)P_x(T\in C_T).
\end{aligned}
\]
Together with \cref{obj:ass:gaussian-channel}, this identifies the joint law as
\[
\mathcal L_{P_x}(Z,T)
=
N(0,1)\otimes\mathcal L_{P_x}(T).
\]
% lean: CausalSmith.Stat.NoisydoseWeakdesignTransition.error_centered_product_stratum

Since \(V=T+\sigma Z\), Fubini's theorem, justified by absolute integrability, and the pointwise Gaussian-inverse identity give
\[
\begin{aligned}
\frac{D_{x,q}}{p_x}
&=E_{P_x}[\ell_q(T+\sigma Z)]\\
&=\int_{-1/2}^{1/2}
g_x(t)E_Z[\ell_q(t+\sigma Z)]\,dt\\
&=\int_{-1/2}^{1/2}g_x(t)q(t)\,dt
=L_{x,q}.
\end{aligned}
\]
Consequently,
\[
D_{x,q}=p_xL_{x,q}.
\]
% lean: CausalSmith.Stat.NoisydoseWeakdesignTransition.populationD_eq_latent

\item Define the observed numerator by
\[
M_{x,q}:=E_P[\mathbf1\{X=x\}Y\ell_q(V)].
\]
By \cref{obj:lem:realized-dose-mean},
\[
0\le Y\le1\quad P\text{-almost surely},
\qquad
E_P[Y\mid A,X,Z]=\mu_X(A)\quad P\text{-almost surely}.
\]
Thus \(M_{x,q}\) is finite. The multiplier
\(\mathbf1\{X=x\}\ell_q(V)\) is measurable with respect to \((A,X,Z)\) and integrable. Pulling it through conditional expectation gives
\[
\begin{aligned}
M_{x,q}
&=E_P\!\left[
E_P[\mathbf1\{X=x\}Y\ell_q(V)\mid A,X,Z]
\right]\\
&=E_P\!\left[
\mathbf1\{X=x\}\ell_q(V)E_P[Y\mid A,X,Z]
\right]\\
&=E_P[\mathbf1\{X=x\}\mu_x(A)\ell_q(V)]\\
&=E_P[\mathbf1\{X=x\}\bar\mu_x(a_0+T)\ell_q(T+\sigma Z)].
\end{aligned}
\]
The last equality uses the latent support in \cref{obj:ass:latent-support}, on which \(a_0+T=A\in[0,1]\).
% lean: CausalSmith.Stat.NoisydoseWeakdesignTransition.integrable_test_mean

The product \(\bar\mu_x(a_0+T)\ell_q(T+\sigma Z)\) is absolutely integrable because \(0\le\bar\mu_x\le1\). Applying the product-law integration from the preceding step with this bounded multiplier yields
\[
\begin{aligned}
\frac{M_{x,q}}{p_x}
&=E_{P_x}[\bar\mu_x(a_0+T)\ell_q(T+\sigma Z)]\\
&=\int_{-1/2}^{1/2}
g_x(t)\bar\mu_x(a_0+t)
E_Z[\ell_q(t+\sigma Z)]\,dt\\
&=\int_{-1/2}^{1/2}
g_x(t)\bar\mu_x(a_0+t)q(t)\,dt
=N_{x,q}.
\end{aligned}
\]
Hence
\[
M_{x,q}=p_xN_{x,q}.
\]
% lean: CausalSmith.Stat.NoisydoseWeakdesignTransition.populationM_eq_latent

\item The positivity established above implies
\[
D_{x,q}=p_xL_{x,q}>0.
\]
Moreover, multiplying the lower bounds for \(p_x\) and \(L_{x,q}\) gives
\[
D_{x,q}
\ge p_{\min}c_{\mathrm{lo}}I_\kappa(q)
=b_q(K).
\]
% lean: CausalSmith.Stat.NoisydoseWeakdesignTransition.positive_ratio

\item We next bound the centered latent numerator. The parameter assumptions give \(\beta>0\) and \(\kappa\ge0\). On \([-1/2,1/2]\),
\[
0\le |t|^\beta\le1,
\qquad
|t|^{\kappa+\beta}q(t)
=
\bigl(|t|^\kappa q(t)\bigr)|t|^\beta.
\]
Thus the higher-moment integrand is integrable, as the product of the integrable base-moment integrand and a bounded measurable function.
% lean: CausalSmith.Stat.NoisydoseWeakdesignTransition.higher_weightMoment_integrable

The integrability already established permits subtraction inside the integral:
\[
N_{x,q}-\mu_x(a_0)L_{x,q}
=
\int_{-1/2}^{1/2}
g_x(t)q(t)\bigl(\bar\mu_x(a_0+t)-\mu_x(a_0)\bigr)\,dt.
\]
For every \(t\in[-1/2,1/2]\), both doses belong to \([0,1]\), and \cref{obj:ass:holder-mean} gives
\[
a_0+t\in[0,1],
\qquad
\bar\mu_x(a_0+t)=\mu_x(a_0+t),
\qquad
|\bar\mu_x(a_0+t)-\mu_x(a_0)|\le |t|^\beta.
\]
Using \(g_x(t)q(t)\ge0\) and the upper design envelope, we obtain almost everywhere
\[
\begin{aligned}
\left|g_x(t)q(t)
\bigl(\bar\mu_x(a_0+t)-\mu_x(a_0)\bigr)\right|
&=g_x(t)q(t)
|\bar\mu_x(a_0+t)-\mu_x(a_0)|\\
&\le c_{\mathrm{hi}}|t|^\kappa q(t)|t|^\beta\\
&=c_{\mathrm{hi}}|t|^{\kappa+\beta}q(t).
\end{aligned}
\]
The absolute value of an integral is bounded by the integral of the absolute value, so
\[
\begin{aligned}
|N_{x,q}-\mu_x(a_0)L_{x,q}|
&\le
\int_{-1/2}^{1/2}
\left|g_x(t)q(t)
\bigl(\bar\mu_x(a_0+t)-\mu_x(a_0)\bigr)\right|\,dt\\
&\le c_{\mathrm{hi}}I_{\kappa+\beta}(q).
\end{aligned}
\]
% lean: CausalSmith.Stat.NoisydoseWeakdesignTransition.latent_weight_centered_bias

\item Nonnegativity of the higher-moment integrand gives
\[
I_{\kappa+\beta}(q)\ge0.
\]
Since \(p_x>0\) and \(L_{x,q}>0\), the observed numerator and denominator identities imply
\[
m_{x,q}
=\frac{p_xN_{x,q}}{p_xL_{x,q}}
=\frac{N_{x,q}}{L_{x,q}},
\qquad
m_{x,q}-\mu_x(a_0)
=
\frac{N_{x,q}-\mu_x(a_0)L_{x,q}}{L_{x,q}}.
\]
Therefore,
\[
\begin{aligned}
|m_{x,q}-\mu_x(a_0)|
&=\frac{|N_{x,q}-\mu_x(a_0)L_{x,q}|}{L_{x,q}}\\
&\le\frac{c_{\mathrm{hi}}I_{\kappa+\beta}(q)}{L_{x,q}}\\
&\le\frac{c_{\mathrm{hi}}I_{\kappa+\beta}(q)}{c_{\mathrm{lo}}I_\kappa(q)}\\
&=\frac{c_{\mathrm{hi}}}{c_{\mathrm{lo}}}
\frac{I_{\kappa+\beta}(q)}{I_\kappa(q)}
=B_q(K).
\end{aligned}
\]
The argument applies to every \(x\in\mathcal X\). Gaussian inversion was used directly throughout the stated range of \(\sigma\), including \(\sigma=0\), so the same conclusions hold at that endpoint.
% lean: CausalSmith.Stat.NoisydoseWeakdesignTransition.positive_ratio
\end{enumerate}
\end{proof}

The lower bound \(b_q\) makes population normalization explicit, and \(B_q\) quantifies localization around the target dose. Both certificates use moments of the nonnegative latent weight. Observable inverse weights are governed instead by inversion and integrability; their empirical denominators receive the public floor specified in \cref{obj:synth_3}. Thus population positivity and empirical stabilization have distinct roles.

The realized-dose identity and convolution injectivity support the identification conclusion in \cref{obj:prop:identification}: the first connects outcome-weighted latent measures to conditional potential means, and the second supplies uniqueness under the known measurement channel. The shrinking-neighborhood ratio statement of that proposition is proved directly in steps 3 and 4 of its proof. \Cref{obj:lem:positive-ratio} serves the estimation analysis instead: its weighted-denominator lower bound and localization-bias bound are the population inputs to the error certificate in \cref{obj:lem:public-error-certificate}.

\subsection{Split-sample error certification}

The observable construction in \cref{obj:synth_3} assigns separate modulo-three blocks to stratum frequencies, inverse-weight denominators, and outcome numerators. Its public denominator floor and projection define the empirical ratios for every sample with nonempty blocks. The following certificate quantifies expected absolute error for a specified weight pair before dictionary selection.

\begin{lemmav}{lem:public-error-certificate}[Public expected-error certificate]
Fix the public potential-path space of \cref{obj:def:potential-path-space} and public class constants \(K\) satisfying the restrictions in \cref{obj:def:model-class}. Suppose that:
\begin{itemize}
\item \textbf{(Parameter ranges.)} \(\beta\in(0,1]\), \(\kappa\in[0,2]\), \(n\ge3\), and \(\sigma\in[0,1/4]\).
\item \textbf{(Structural law.)} The structural law \(P\) belongs to \(\mathcal P_{K,\beta,\kappa,\sigma}\), the model class with constants \(K\) defined in \cref{obj:def:model-class}.
\item \textbf{(Sampling.)} The observed sample satisfies \cref{obj:ass:iid}, with joint law \(Q_P^n\) defined in \cref{obj:def:observed-experiment}.
\item \textbf{(Inverse pair.)} Let \(q\) be a public latent weight and \(\ell_q\) its observable inverse weight. Under the usual measurability and integrability conditions, they satisfy
\[
q(t)\ge0\qquad(t\in[-1/2,1/2]),
\qquad
0<I_\kappa(q):=\int_{-1/2}^{1/2}|t|^\kappa q(t)\,dt<\infty,
\]
and the pointwise Gaussian inversion identity
\[
\int_{\mathbb R}\ell_q(t+\sigma z)
       \frac{e^{-z^2/2}}{\sqrt{2\pi}}\,dz
=q(t)
\qquad(t\in[-1/2,1/2]),
\]
where the Gaussian integrand is absolutely integrable for every such \(t\).
\item \textbf{(Law integrability.)} Let \(W\) denote the contaminated dose and \(a_0\) the prespecified evaluation dose, and define the observed-dose displacement by
\[
V:=W-a_0.
\]
Then
\[
\mathbb E_P[|\ell_q(V)|]<\infty.
\]
\item \textbf{(Public second moment.)} Define the public inverse-weight second moment, as an extended nonnegative integral, by
\[
V_q:=
16\int_{-1/2}^{1/2}|t|^\kappa
   \left(
   \int_{\mathbb R}\ell_q(t+\sigma z)^2
       \frac{e^{-z^2/2}}{\sqrt{2\pi}}\,dz
   \right)dt.
\]
Assume \(V_q<\infty\).
\end{itemize}
Let \(\widehat\theta_q\) be the weight-indexed estimator of \cref{obj:def:total-estimator}, and let \(A_q(K)\) be its public expected-error score from that construction, evaluated at the fixed parameters \(\beta,\kappa,n,\sigma\). Define the population causal target by
\[
\theta(P):=\mathbb E_P[Y(a_0)],
\]
where \(Y(a_0)\) is the potential outcome at the evaluation dose in \cref{obj:def:potential-path-space}. Then
\[
\mathbb E_{Q_P^n}
 \left[|\widehat\theta_q-\theta(P)|\right]
\le A_q(K).
\]
\end{lemmav}

\begin{proof}[Proof of \cref{obj:lem:public-error-certificate}]
\begin{enumerate}
\item
Index the observed records and form the three blocks as in
\cref{obj:synth_3,obj:def:total-estimator}. Thus
\[
O_i=(X_i,W_i,Y_i),\qquad V_i=W_i-a_0,\qquad a_0=\frac12,
\qquad i=0,\ldots,n-1,
\]
and write
\[
N_r:=|\mathcal I_r|,
\qquad
\mathcal I_r=\{i\in\{0,\ldots,n-1\}:i\bmod3=r\},
\qquad r\in\{0,1,2\}.
\]
The indices \(3i+r\), for \(0\le i<\lfloor n/3\rfloor\), belong to
\(\mathcal I_r\). Consequently,
\[
N_r\ge\left\lfloor\frac n3\right\rfloor\ge1,
\qquad
n\le3\left\lfloor\frac n3\right\rfloor+2
\le5\left\lfloor\frac n3\right\rfloor
\le5N_r.
\]
In particular, every block is nonempty. The positive weight moment and public constants give
\[
b_q(K)=p_{\min}c_{\mathrm{lo}}I_\kappa(q)>0.
\]
For each \(x\in\mathcal X\), define the population moments
\[
D_x(q):=E_P[\mathbf1\{X=x\}\ell_q(V)],
\qquad
M_x(q):=E_P[\mathbf1\{X=x\}Y\ell_q(V)].
\]
The empirical moments are
\[
\widehat p_x=\frac1{N_0}\sum_{i\in\mathcal I_0}\mathbf1\{X_i=x\},
\]
\[
\widehat D_x(q)=\frac1{N_1}\sum_{i\in\mathcal I_1}
\mathbf1\{X_i=x\}\ell_q(V_i),
\qquad
\widehat M_x(q)=\frac1{N_2}\sum_{i\in\mathcal I_2}
\mathbf1\{X_i=x\}Y_i\ell_q(V_i).
\]
On every observed sample in \((\mathcal X\times\mathbb R\times\mathbb R)^n\), define
\[
\Delta_{D,x}:=|\widehat D_x(q)-D_x(q)|,\qquad
\Delta_{M,x}:=|\widehat M_x(q)-M_x(q)|,\qquad
\Delta_{p,x}:=|\widehat p_x-p_x|.
\]
All sample expectations below are taken under \(Q_P^n\).
% lean: CausalSmith.Stat.NoisydoseWeakdesignTransition.blocksNonempty_of_three_le
% lean: CausalSmith.Stat.NoisydoseWeakdesignTransition.splitBlock_card_lower
% lean: CausalSmith.Stat.NoisydoseWeakdesignTransition.splitBlock_card_fifth

\item
We first establish the population moments of these averages. Choose the
conditional densities \(g_x\) supplied by \cref{obj:ass:weak-design}.
Define the standard Gaussian probability measure by
\[
\Phi_1:=\mathcal N(0,1).
\]
The Gaussian channel and error independence imply that, conditional on
\(X=x\), the joint law of \(Z\) and \(A-a_0\) is the product of
\(\Phi_1\) and the measure with density \(g_x\) on \([-1/2,1/2]\).
The public second-moment quantity in \cref{obj:synth_4} satisfies
\[
V_q(K)
=c_{\mathrm{hi}}\int_{-1/2}^{1/2}|t|^\kappa
       \left(\int_{\mathbb R}\ell_q(t+\sigma z)^2\,d\Phi_1(z)\right)\,dt
=\frac{c_{\mathrm{hi}}}{16}V_q<\infty.
\]
Tonelli's theorem and the upper design envelope therefore give, initially
as nonnegative extended integrals,
\[
\begin{aligned}
E_P[\mathbf1\{X=x\}\ell_q(V)^2]
&=p_x\int_{-1/2}^{1/2}g_x(t)
       \left(\int_{\mathbb R}\ell_q(t+\sigma z)^2\,d\Phi_1(z)\right)\,dt\\
&\le c_{\mathrm{hi}}\int_{-1/2}^{1/2}|t|^\kappa
       \left(\int_{\mathbb R}\ell_q(t+\sigma z)^2\,d\Phi_1(z)\right)\,dt\\
&=V_q(K)<\infty,
\end{aligned}
\]
where we used \(p_x\le1\).
By \cref{obj:lem:realized-dose-mean}, \(0\le Y\le1\) almost surely, so
\[
E_P[\mathbf1\{X=x\}Y^2\ell_q(V)^2]
\le E_P[\mathbf1\{X=x\}\ell_q(V)^2]
\le V_q(K).
\]
Both observable marks are thus square integrable. For the frequency mark,
\[
0\le\mathbf1\{X=x\}\le1,
\qquad
\operatorname{Var}_P(\mathbf1\{X=x\})\le\frac14,
\]
by the variance bound for a random variable in \([0,1]\).

Linearity of expectation gives
\[
E_{Q_P^n}\widehat D_x(q)=D_x(q),\qquad
E_{Q_P^n}\widehat M_x(q)=M_x(q),\qquad
E_{Q_P^n}\widehat p_x=p_x.
\]
Independence of the sample records makes the mixed covariance terms vanish.
Hence
\[
\begin{aligned}
\operatorname{Var}_{Q_P^n}(\widehat D_x(q))
&=\frac1{N_1^2}\sum_{i\in\mathcal I_1}
  \operatorname{Var}_P(\mathbf1\{X=x\}\ell_q(V))
\le\frac{V_q(K)}{N_1},\\
\operatorname{Var}_{Q_P^n}(\widehat M_x(q))
&=\frac1{N_2^2}\sum_{i\in\mathcal I_2}
  \operatorname{Var}_P(\mathbf1\{X=x\}Y\ell_q(V))
\le\frac{V_q(K)}{N_2},\\
\operatorname{Var}_{Q_P^n}(\widehat p_x)
&=\frac1{N_0^2}\sum_{i\in\mathcal I_0}
  \operatorname{Var}_P(\mathbf1\{X=x\})
\le\frac1{4N_0}.
\end{aligned}
\]
Here the first two inequalities also use
\(\operatorname{Var}_P(\cdot)\le E_P[(\cdot)^2]\).
These averages and their absolute centered errors are square integrable,
and therefore integrable.
% lean: CausalSmith.Stat.NoisydoseWeakdesignTransition.Dhat_population_moments
% lean: CausalSmith.Stat.NoisydoseWeakdesignTransition.Mhat_population_moments
% lean: CausalSmith.Stat.NoisydoseWeakdesignTransition.phat_population_moments

\item
Distinct modulo-three blocks are disjoint: an index in both
\(\mathcal I_0\) and \(\mathcal I_r\) would have remainder both \(0\) and
\(r\). Under the product experiment, the records indexed by disjoint
blocks are independent. Taking measurable functions of those records
preserves independence, so, separately for the denominator and numerator
blocks,
\[
\widehat p_x\perp\Delta_{D,x},
\qquad
\widehat p_x\perp\Delta_{M,x}
\qquad\text{under }Q_P^n.
\]
The preceding integrability and independence give integrable products and
the factorizations
\[
E_{Q_P^n}[\widehat p_x\Delta_{D,x}]
=p_xE_{Q_P^n}\Delta_{D,x},
\qquad
E_{Q_P^n}[\widehat p_x\Delta_{M,x}]
=p_xE_{Q_P^n}\Delta_{M,x}.
\]
% lean: CausalSmith.Stat.NoisydoseWeakdesignTransition.phat_block_error_factorization

\item
Nonnegativity of the variance of an absolute centered error bounds its
mean by the standard deviation of the original average. For example,
unbiasedness gives
\[
\begin{aligned}
0
&\le\operatorname{Var}_{Q_P^n}(\Delta_{D,x})\\
&=E_{Q_P^n}\Delta_{D,x}^2
  -(E_{Q_P^n}\Delta_{D,x})^2\\
&=\operatorname{Var}_{Q_P^n}(\widehat D_x(q))
  -(E_{Q_P^n}\Delta_{D,x})^2.
\end{aligned}
\]
Since the expected absolute error is nonnegative, taking square roots
yields the first of the following bounds; the same argument gives the
other two:
\[
\begin{aligned}
E_{Q_P^n}\Delta_{D,x}
&\le\sqrt{\operatorname{Var}_{Q_P^n}(\widehat D_x(q))}
\le\sqrt{\frac{V_q(K)}{N_1}},\\
E_{Q_P^n}\Delta_{M,x}
&\le\sqrt{\operatorname{Var}_{Q_P^n}(\widehat M_x(q))}
\le\sqrt{\frac{V_q(K)}{N_2}},\\
E_{Q_P^n}\Delta_{p,x}
&\le\sqrt{\operatorname{Var}_{Q_P^n}(\widehat p_x)}
\le\sqrt{\frac1{4N_0}}.
\end{aligned}
\]
Because \(n\le5N_r\), \(n>0\), and \(V_q(K)\ge0\),
\[
\frac{V_q(K)}{N_r}\le9\frac{V_q(K)}{n},
\qquad
\frac1{4N_r}\le\frac3n.
\]
Monotonicity of the square root consequently gives
\[
E_{Q_P^n}\Delta_{D,x}\le3\sqrt{\frac{V_q(K)}{n}},
\qquad
E_{Q_P^n}\Delta_{M,x}\le3\sqrt{\frac{V_q(K)}{n}},
\qquad
E_{Q_P^n}\Delta_{p,x}\le\sqrt{\frac3n}.
\]
These inequalities also hold when \(V_q(K)=0\).
% lean: CausalSmith.Stat.NoisydoseWeakdesignTransition.expected_abs_deviation_le_sqrt_variance
% lean: CausalSmith.Stat.NoisydoseWeakdesignTransition.inverse_moments_expected_abs_error
% lean: CausalSmith.Stat.NoisydoseWeakdesignTransition.split_standard_deviation_le

\item
Partitioning the bounded potential outcome \(Y(a_0)\) by the two strata
and using \cref{obj:def:conditional-potential-mean} gives
\[
\theta(P)
=\sum_{x\in\mathcal X}E_P[\mathbf1\{X=x\}Y(a_0)]
=\sum_{x\in\mathcal X}p_x\mu_x(a_0).
\]
The estimator therefore satisfies the exact decomposition
\[
\widehat\theta_q-\theta(P)
=\sum_{x\in\mathcal X}\widehat p_x
       \bigl(\widehat\mu_x(q)-\mu_x(a_0)\bigr)
 +\sum_{x\in\mathcal X}(\widehat p_x-p_x)\mu_x(a_0).
\]
The frequencies are nonnegative by their defining averages. Applying the
triangle inequality gives, for every observed sample,
\[
|\widehat\theta_q-\theta(P)|
\le
\sum_{x\in\mathcal X}\widehat p_x
       |\widehat\mu_x(q)-\mu_x(a_0)|
+\left|\sum_{x\in\mathcal X}(\widehat p_x-p_x)\mu_x(a_0)\right|.
\]
% lean: CausalSmith.Stat.NoisydoseWeakdesignTransition.weightEstimator_error_decomposition

\item
For every stratum, \cref{obj:lem:positive-ratio} supplies
\[
D_x(q)\ge b_q(K)>0,\qquad
m_{x,q}:=\frac{M_x(q)}{D_x(q)},\qquad
|m_{x,q}-\mu_x(a_0)|\le B_q(K).
\]
We also need the range of \(m_{x,q}\). The conditional-mean identity in
\cref{obj:lem:realized-dose-mean}, applied by the conditional-expectation
pull-out rule to the integrable inverse-weight mark, gives
\[
M_x(q)
=E_P[\mathbf1\{X=x\}\mu_X(A)\ell_q(V)].
\]
Indeed, law integrability and the bounded outcome and mean make both
products integrable. Conditioning on the positive-mass stratum, using
the Gaussian product law, and applying Fubini and pointwise inversion
now yield
\[
D_x(q)
=p_x\int_{-1/2}^{1/2}g_x(t)
       \left(\int_{\mathbb R}\ell_q(t+\sigma z)\,d\Phi_1(z)\right)\,dt
=p_x\int_{-1/2}^{1/2}g_x(t)q(t)\,dt,
\]
and
\[
\begin{aligned}
M_x(q)
&=p_x\int_{-1/2}^{1/2}g_x(t)\mu_x(a_0+t)
       \left(\int_{\mathbb R}\ell_q(t+\sigma z)\,d\Phi_1(z)\right)\,dt\\
&=p_x\int_{-1/2}^{1/2}g_x(t)q(t)\mu_x(a_0+t)\,dt.
\end{aligned}
\]
Absolute integrability under \(P\) persists under each normalized
stratum restriction, which justifies these uses of Fubini. The
conditional mean is measurable on the dose interval by its Hölder
condition.

The design envelopes and nonnegativity of \(q\) give, almost everywhere
on \([-1/2,1/2]\),
\[
0\le g_x(t)q(t)\le c_{\mathrm{hi}}|t|^\kappa q(t),
\qquad
\int_{-1/2}^{1/2}g_x(t)q(t)\,dt
\ge c_{\mathrm{lo}}I_\kappa(q)>0.
\]
Thus this weight is integrable and has positive mass. Since
\(a_0+t\in[0,1]\), \cref{obj:ass:mean-range} gives
\[
\begin{aligned}
0
&\le\int_{-1/2}^{1/2}g_x(t)q(t)\mu_x(a_0+t)\,dt\\
&\le\int_{-1/2}^{1/2}g_x(t)q(t)\,dt.
\end{aligned}
\]
Cancelling \(p_x>0\) in the population ratio and dividing by the positive
weight mass proves
\[
m_{x,q}
=\frac{\displaystyle\int_{-1/2}^{1/2}g_x(t)q(t)\mu_x(a_0+t)\,dt}
       {\displaystyle\int_{-1/2}^{1/2}g_x(t)q(t)\,dt}
\in[0,1].
\]

For every observed sample, define the floored denominator
\[
\widehat D_x^{\mathrm{floor}}(q)
:=\max\{\widehat D_x(q),b_q(K)/2\}.
\]
It satisfies
\[
\widehat D_x^{\mathrm{floor}}(q)\ge b_q(K)/2>0,
\qquad
|\widehat D_x^{\mathrm{floor}}(q)-D_x(q)|
\le|\widehat D_x(q)-D_x(q)|=\Delta_{D,x}.
\]
The latter follows because \(D_x(q)\ge b_q(K)\), and taking the maximum
with \(b_q(K)/2\) contracts distance to any number at least \(b_q(K)/2\).

The clipping construction in \cref{obj:def:total-estimator} is
\[
\widehat\mu_x(q)
=\Pi_{[0,1]}
 \left(\frac{\widehat M_x(q)}
 {\widehat D_x^{\mathrm{floor}}(q)}\right),
\qquad
\Pi_{[0,1]}(u)=\max\{0,\min\{1,u\}\}.
\]
Projection onto this interval fixes \(m_{x,q}\) and contracts distance
to it. Moreover,
\[
\frac{\widehat M_x(q)}{\widehat D_x^{\mathrm{floor}}(q)}-m_{x,q}
=
\frac{\widehat M_x(q)-M_x(q)
      +m_{x,q}\bigl(D_x(q)-\widehat D_x^{\mathrm{floor}}(q)\bigr)}
     {\widehat D_x^{\mathrm{floor}}(q)}.
\]
Using \(|m_{x,q}|\le1\), the floor bound, and the triangle inequality
therefore gives
\[
|\widehat\mu_x(q)-m_{x,q}|
\le\frac{\Delta_{M,x}+\Delta_{D,x}}
          {\widehat D_x^{\mathrm{floor}}(q)}
\le\frac2{b_q(K)}(\Delta_{M,x}+\Delta_{D,x}).
\]
Combining this with the localization bound yields
\[
|\widehat\mu_x(q)-\mu_x(a_0)|
\le B_q(K)+\frac2{b_q(K)}(\Delta_{M,x}+\Delta_{D,x}).
\]
Every denominator used here is positive for every observed sample.
% lean: CausalSmith.Stat.NoisydoseWeakdesignTransition.localizedRatio_mem_Icc
% lean: CausalSmith.Stat.NoisydoseWeakdesignTransition.muhat_error_le

\item
Each record belongs to exactly one of the two strata, so
\[
\widehat p_x\ge0,
\qquad
\sum_{x\in\mathcal X}\widehat p_x
=\frac1{N_0}\sum_{i\in\mathcal I_0}
       \sum_{x\in\mathcal X}\mathbf1\{X_i=x\}
=1.
\]
Multiplying the stratum error bounds by these frequencies and summing
gives
\[
\sum_{x\in\mathcal X}\widehat p_x
       |\widehat\mu_x(q)-\mu_x(a_0)|
\le
B_q(K)+\frac2{b_q(K)}\sum_{x\in\mathcal X}
       \bigl(\widehat p_x\Delta_{M,x}
             +\widehat p_x\Delta_{D,x}\bigr).
\]
Also, \cref{obj:ass:mean-range} implies
\(|\mu_x(a_0)|\le1\), whence
\[
\left|\sum_{x\in\mathcal X}(\widehat p_x-p_x)\mu_x(a_0)\right|
\le\sum_{x\in\mathcal X}|\widehat p_x-p_x|\,|\mu_x(a_0)|
\le\sum_{x\in\mathcal X}\Delta_{p,x}.
\]
Substitution into the error decomposition proves the pointwise bound
\[
|\widehat\theta_q-\theta(P)|
\le
B_q(K)+\frac2{b_q(K)}\sum_{x\in\mathcal X}
       \bigl(\widehat p_x\Delta_{M,x}
             +\widehat p_x\Delta_{D,x}\bigr)
+\sum_{x\in\mathcal X}\Delta_{p,x}.
\]
The right side is integrable by the moment bounds and the block
factorizations. Integrating this bound and splitting the finite sums
therefore gives
\[
\begin{aligned}
E_{Q_P^n}|\widehat\theta_q-\theta(P)|
\le B_q(K)
&+\frac2{b_q(K)}\sum_{x\in\mathcal X}
 E_{Q_P^n}\bigl[\widehat p_x\Delta_{M,x}
                  +\widehat p_x\Delta_{D,x}\bigr]\\
&+\sum_{x\in\mathcal X}E_{Q_P^n}\Delta_{p,x}.
\end{aligned}
\]
% lean: CausalSmith.Stat.NoisydoseWeakdesignTransition.phat_sum_eq_one
% lean: CausalSmith.Stat.NoisydoseWeakdesignTransition.public_error_certificate

\item
The block factorizations and deviation bounds imply, for each stratum,
\[
\begin{aligned}
E_{Q_P^n}\bigl[\widehat p_x\Delta_{M,x}
                 +\widehat p_x\Delta_{D,x}\bigr]
&=p_x\bigl(E_{Q_P^n}\Delta_{M,x}
           +E_{Q_P^n}\Delta_{D,x}\bigr)\\
&\le p_x\,6\sqrt{\frac{V_q(K)}{n}}.
\end{aligned}
\]
Unbiasedness of the frequencies and their samplewise normalization give
\[
\sum_{x\in\mathcal X}p_x
=\sum_{x\in\mathcal X}E_{Q_P^n}\widehat p_x
=E_{Q_P^n}\!\left[\sum_{x\in\mathcal X}\widehat p_x\right]
=1.
\]
Consequently,
\[
\sum_{x\in\mathcal X}
 E_{Q_P^n}\bigl[\widehat p_x\Delta_{M,x}
                  +\widehat p_x\Delta_{D,x}\bigr]
\le6\sqrt{\frac{V_q(K)}{n}}.
\]
There are two strata, and hence
\[
\sum_{x\in\mathcal X}E_{Q_P^n}\Delta_{p,x}
\le2\sqrt{\frac3n}.
\]
Finally, \(2/b_q(K)\ge0\), so the integrated pointwise bound yields
\[
\begin{aligned}
E_{Q_P^n}|\widehat\theta_q-\theta(P)|
&\le B_q(K)+\frac2{b_q(K)}\,6\sqrt{\frac{V_q(K)}{n}}
             +2\sqrt{\frac3n}\\
&=B_q(K)+\frac{12}{b_q(K)}\sqrt{\frac{V_q(K)}{n}}
             +2\sqrt{\frac3n}
=A_q(K).
\end{aligned}
\]
% lean: CausalSmith.Stat.NoisydoseWeakdesignTransition.public_error_certificate
\end{enumerate}
\end{proof}

In the inverse-pair clause, the ``usual measurability and integrability conditions'' are the following hypotheses: \(q\) and \(\ell_q\) are Borel measurable, and \(t\mapsto|t|^\kappa q(t)\) is Lebesgue integrable on \([-1/2,1/2]\). Together with the displayed nonnegativity, positivity of \(I_\kappa(q)\), and Gaussian inversion with absolute integrability, they form the admissibility conditions on the pair.

The certificate separates localization error, inverse-weight sampling variability, and stratum-frequency error. In particular, the variability term involves the observable inverse through \(V_q(K)\), while localization uses the latent weight through \(B_q(K)\). The denominator bound \(b_q(K)\) connects population normalization to the stability of the clipped empirical ratio. For a public weight pair, all three terms in \(A_q(K)\) are available for the deterministic comparison rule in \cref{obj:def:total-estimator}.

\subsection{Dictionary admissibility and rate comparisons}

The Fourier and polynomial pairs in \cref{obj:synth_1,obj:synth_2} supply two ways to localize while inverting the Gaussian channel. Their finite dictionary also contains the constant pair, so the public minimization rule has a candidate for every \(n\ge1\); the estimator and interval use the selected entry only for \(n\ge3\), when all three sample blocks are nonempty. The next statement verifies the conditions required by the error certificate for every dictionary entry and records the comparisons that control the selected score.

\begin{lemmav}{lem:dictionary-score-rate}[Dictionary admissibility and score rate]
Fix the following public objects and parameters:
\begin{itemize}
\item \textbf{(Path space.)} A potential-path space \((\mathcal S,\mathcal E)\) as in \cref{obj:def:potential-path-space}.
\item \textbf{(Class constants.)} \(K=(p_{\min},c_{\mathrm{lo}},c_{\mathrm{hi}})\), with
\[
0<p_{\min}\le \tfrac12,
\qquad
0<c_{\mathrm{lo}}\le c_{\mathrm{hi}}.
\]
\item \textbf{(Exponents.)} \(\beta\in(0,1]\) and \(\kappa\in[0,2]\).
\end{itemize}
Use the public dictionary \(\mathcal D_{n,\sigma}\) from \cref{obj:def:weight-dictionary} and the public expected-error scores \(A_q(K)\) from \cref{obj:synth_4}. For every integer \(n\ge1\), every \(\sigma\in[0,1/4]\), and every inverse pair \((q,\ell_q)\in\mathcal D_{n,\sigma}\), both weights are measurable, \(q\) is nonnegative on \([-1/2,1/2]\), and its design-weighted moment satisfies
\[
0<I_\kappa(q)
=
\int_{-1/2}^{1/2}|t|^\kappa q(t)\,dt
<\infty.
\]
For every \(t\in[-1/2,1/2]\), the exact Gaussian inversion identity holds with absolute integrability:
\[
\int_{\mathbb R}|\ell_q(t+\sigma z)|\,\mathcal N(0,1)(dz)<\infty,
\qquad
\int_{\mathbb R}\ell_q(t+\sigma z)\,\mathcal N(0,1)(dz)=q(t),
\]
where \(\mathcal N(0,1)\) denotes the standard Gaussian probability measure. The inverse-weight second-moment integral satisfies
\[
16\int_{-1/2}^{1/2}|t|^\kappa
\left(\int_{\mathbb R}\ell_q(t+\sigma z)^2\,\mathcal N(0,1)(dz)\right)\,dt
<\infty.
\]
Let \(V\) denote the observed-dose displacement from the prespecified evaluation dose \(a_0\), and let \(W\) denote the contaminated dose:
\[
V=W-a_0,
\qquad
W=A+\sigma Z,
\]
where \(A\) is the latent dose and \(Z\) is the standardized Gaussian error. For every structural law \(P\in\mathcal P_{K,\beta,\kappa,\sigma}\) from \cref{obj:def:model-class}, the inverse weight is absolutely integrable under the observed centered dose:
\[
\mathbb E_P[|\ell_q(V)|]<\infty.
\]

Let \(\widehat q\) be the first public-score minimizer in the fixed dictionary order from \cref{obj:def:total-estimator}, and let \(\rho_{n,\sigma}\) and \(L_n\) be the frontier rate and logarithmic sample-size quantity from \cref{obj:def:frontier-rate}. There exist \(C>0\) and \(n_0\in\mathbb N\), uniform in \(\sigma\in[0,1/4]\), such that for every \(n\ge n_0\) and every \(\sigma\in[0,1/4]\),
\[
A_{\widehat q}(K)\le C\rho_{n,\sigma}.
\]
For the same \(C\) and \(n_0\), whenever \(\sigma\le L_n^{-1/2}\), there exists \(k\in\mathbb N\) such that the Fourier inverse pair indexed by \(k\) in \cref{obj:def:weight-dictionary} belongs to \(\mathcal D_{n,\sigma}\) and its public score is at most \(C\rho_{n,\sigma}\). Whenever \(L_n^{-1/2}<\sigma\), there exists \(j\in\mathbb N\) such that the polynomial inverse pair with index
\[
m=2^j+1
\]
belongs to \(\mathcal D_{n,\sigma}\) and satisfies
\[
A_{q_m^{\mathrm M}}(K)\le C\rho_{n,\sigma}.
\]
\end{lemmav}

\begin{proof}[Proof of \cref{obj:lem:dictionary-score-rate}]
Throughout, zeroth powers of nonnegative numbers are interpreted as \(1\). For the moment calculations, write
\[
I_s(q):=\int_{-1/2}^{1/2}|t|^s q(t)\,dt,
\qquad 0\le s\le3.
\]
For a public inverse pair, use the factor-\(16\) second-moment quantity of \cref{obj:synth_4}:
\[
V_q^{(16)}:=16\int_{-1/2}^{1/2}|t|^\kappa
E_Z[\ell_q(t+\sigma Z)^2]\,dt,
\qquad Z\sim N(0,1).
\]
The integral is understood as an extended nonnegative integral. The reference-score constants below are public and may depend on the fixed exponents \(\beta,\kappa\); the final public-score constant also depends on \(K\).

\begin{enumerate}
\item \emph{Admissibility of the dictionary pairs.}
Fix \(n\ge1\) and \(\sigma\in[0,1/4]\). We consider the three families in \cref{obj:def:weight-dictionary}.

For the constant pair, the polynomial construction in \cref{obj:synth_2} at index \(m=1\) gives
\[
q_1^{\mathrm M}(t)=\ell_1^{\mathrm M}(t)=1.
\]
The integrand defining \(I_\kappa(1)\) is continuous, nonnegative, and positive away from zero. Since the latent interval has length one and \(|t|^\kappa\le1\) there,
\[
0<I_\kappa(1)\le1,
\qquad
E_Z[|1|]=E_Z[1]=1,
\qquad
V_1^{(16)}=16I_\kappa(1)<\infty.
\]
Thus the constant pair satisfies the required conditions.

For a Fourier entry, let \(h=2^{-k}\in[n^{-2},1/4]\). Use the function \(Q\) and the pair \(q_h^{\mathrm F},\ell_h^{\mathrm F}\) from \cref{obj:synth_1}. The elementary bounds for sine imply, for \(0\le s\le3\),
\[
0\le Q(u)\le1,
\qquad
Q(u)\le |u|^{-6}\quad(u\ne0),
\qquad
0\le |u|^sQ(u)\le\frac{8}{(1+|u|)^3}.
\]
For the last bound, split at \(|u|=1\): on the central interval the product is at most one, and outside it the product is at most \(|u|^{-3}\). In particular, define
\[
c_{\mathrm F,s}:=\int_{-1/2}^{1/2}|u|^sQ(u)\,du,
\qquad
C_{\mathrm F,s}:=\int_{\mathbb R}|u|^sQ(u)\,du.
\]
Continuity and \(Q(0)=1\) give \(c_{\mathrm F,s}>0\), while the cubic envelope gives \(C_{\mathrm F,s}<\infty\). Changing variables \(t=hu\), and using \(h\le1\), yields
\[
\begin{aligned}
I_s(q_h^{\mathrm F})
&=h^{s+1}
  \int_{-1/(2h)}^{1/(2h)}|u|^sQ(u)\,du,\\
c_{\mathrm F,s}h^{s+1}
&\le I_s(q_h^{\mathrm F})
 \le C_{\mathrm F,s}h^{s+1}.
\end{aligned}
\]
These inequalities give the positive finite weighted moment.

We also record the spectral support needed below. With the transform convention of \cref{obj:synth_1}, define
\[
\varpi_h(\omega):=\frac h2\mathbf1_{[-1/h,\,1/h]}(\omega),
\qquad \omega\in\mathbb R.
\]
Direct integration gives
\[
\int_{\mathbb R}\varpi_h(\omega)e^{it\omega}\,d\omega
=\frac{\sin(t/h)}{t/h},
\]
with value one at \(t=0\). Taking six products and applying Fubini therefore gives
\[
q_h^{\mathrm F}(t)
=\int_{\mathbb R}\varpi_h^{*6}(\omega)e^{it\omega}\,d\omega,
\qquad
\operatorname{supp}\widehat q_h^{\mathrm F}
\subseteq[-6/h,6/h].
\]
Here convolution adds the support intervals; Fourier inversion gives the asserted support of the transform. Define
\[
G_{h,\sigma}(\omega)
:=\widehat q_h^{\mathrm F}(\omega)e^{\sigma^2\omega^2/2},
\qquad
\mathfrak l_{h,\sigma}(v)
:=\frac1{2\pi}\int_{\mathbb R}G_{h,\sigma}(\omega)e^{iv\omega}\,d\omega.
\]
The multiplier is continuous and compactly supported, so
\[
\ell_h^{\mathrm F}(v)=\operatorname{Re}\mathfrak l_{h,\sigma}(v),
\qquad
|\ell_h^{\mathrm F}(v)|
\le\frac1{2\pi}\int_{\mathbb R}|G_{h,\sigma}(\omega)|\,d\omega<\infty.
\]
This uniform bound proves Gaussian absolute integrability at every displacement and gives
\[
V_{q_h^{\mathrm F}}^{(16)}
\le
16I_\kappa(1)
\left(\frac1{2\pi}\int_{\mathbb R}|G_{h,\sigma}(\omega)|\,d\omega\right)^2
<\infty.
\]
The exact expectation identity is supplied by \cref{obj:synth_1}. Continuity supplies measurability of both functions.

For a polynomial entry, dictionary membership supplies a positive odd index \(m=2^j+1\). By \cref{obj:synth_2}, \(q_m^{\mathrm M}\) and \(\ell_m^{\mathrm M}\) are polynomials, and
\[
q_m^{\mathrm M}(t)\ge0,
\qquad
q_m^{\mathrm M}(0)=1.
\]
Consequently the weighted integrand is continuous on the latent interval, nonnegative, and positive on an interval away from zero. Hence
\[
0<I_\kappa(q_m^{\mathrm M})<\infty.
\]
At \(\sigma=0\), the inverse-heat sum reduces to
\[
\ell_m^{\mathrm M}(t)=q_m^{\mathrm M}(t),
\qquad
E_Z[\ell_m^{\mathrm M}(t+0Z)]=q_m^{\mathrm M}(t).
\]
For \(\sigma>0\), \cref{obj:synth_2} supplies the exact Gaussian expectation identity. In both cases, finite polynomial expansion and the Gaussian absolute moments give
\[
E_Z\!\left[|\ell_m^{\mathrm M}(t+\sigma Z)|\right]<\infty,
\qquad
16\int_{-1/2}^{1/2}|t|^\kappa
E_Z\!\left[\ell_m^{\mathrm M}(t+\sigma Z)^2\right]dt<\infty.
\]
Polynomial continuity supplies measurability. This proves pair admissibility and second-moment finiteness for every dictionary member, including when either indexed family is empty.
% lean: CausalSmith.Stat.NoisydoseWeakdesignTransition.dictionary_pair_admissible

\item \emph{Integrability under structural laws.}
Fix a dictionary pair \((q,\ell_q)\) and a law \(P\in\mathcal P_{K,\beta,\kappa,\sigma}\). Write \(X\) for the stratum coordinate, \(A\) for the latent dose, and \(Z\) for the standardized error, and define
\[
\mathcal X:=\{0,1\},
\qquad
a_0:=\frac12,
\qquad
T_{\mathrm{lat}}:=A-a_0,
\qquad
V:=A+\sigma Z-a_0.
\]
The independence and Gaussian conditions in \cref{obj:ass:error-independence,obj:ass:gaussian-channel} imply
\[
\mathcal L_P(T_{\mathrm{lat}},Z)
=\mathcal L_P(T_{\mathrm{lat}})\otimes N(0,1).
\]
For each \(x\in\mathcal X\), write \(g_x\) for the conditional Lebesgue density of \(T_{\mathrm{lat}}\) given \(X=x\), so that
\[
\mathcal L_P(T_{\mathrm{lat}}\mid X=x)(dt)=g_x(t)\,dt.
\]
The positive stratum probabilities supplied by \cref{obj:ass:stratum-overlap} make these conditional laws well defined. Since \(a_0=1/2\), \cref{obj:ass:latent-support} gives
\[
P\{T_{\mathrm{lat}}\in[-1/2,1/2]\}=1.
\]
Thus, for any nonnegative measurable function \(f\), conditioning on the stratum gives the following equality with the integral restricted to the latent interval; the upper density envelope in \cref{obj:ass:weak-design}, with \(p_x:=P(X=x)\le1\), gives the inequality:
\[
\begin{aligned}
E_P[\mathbf1\{X=x\}f(T_{\mathrm{lat}})]
&=p_x\int_{-1/2}^{1/2}g_x(t)f(t)\,dt\\
&\le c_{\mathrm{hi}}\int_{-1/2}^{1/2}|t|^\kappa f(t)\,dt.
\end{aligned}
\]
Summing the two strata and applying this inequality to
\[
f_q(t):=E_Z[\ell_q(t+\sigma Z)^2],
\qquad t\in\mathbb R,
\]
proves, by the product-law identity and Tonelli,
\[
E_P[\ell_q(V)^2]
=E_P[f_q(T_{\mathrm{lat}})]
\le\frac{c_{\mathrm{hi}}}{16}
\bigl(V_q^{(16)}+V_q^{(16)}\bigr)<\infty.
\]
The compositions involved are measurable. The inclusion \(L^2(P)\subseteq L^1(P)\) on a probability space now gives
\[
E_P[|\ell_q(V)|]
\le\sqrt{E_P[\ell_q(V)^2]}
\le\sqrt{\frac{c_{\mathrm{hi}}}{8}V_q^{(16)}}<\infty.
\]
This establishes the law-integrability assertion.
% lean: CausalSmith.Stat.NoisydoseWeakdesignTransition.dictionary_pair_law_admissible

\item \emph{The Fourier score bound.}
For finite-moment pairs and \(n\ge1\), define the reference quantities used to tune the dictionary:
\[
b_q^{\mathrm{ref}}:=\frac1{16}I_\kappa(q),
\qquad
B_q^{\mathrm{ref}}:=64\frac{I_{\kappa+\beta}(q)}{I_\kappa(q)},
\qquad
A_q^{\mathrm{ref}}:=B_q^{\mathrm{ref}}
+\frac{12}{b_q^{\mathrm{ref}}}\sqrt{\frac{V_q^{(16)}}{n}}
+2\sqrt{\frac3n}.
\]
Since \(0\le\kappa+\beta\le3\), the Fourier moment bounds already proved give
\[
B_{q_h^{\mathrm F}}^{\mathrm{ref}}
\le C_{\mathrm F,bias}h^\beta,
\qquad
C_{\mathrm F,bias}:=
\frac{64C_{\mathrm F,\kappa+\beta}}{c_{\mathrm F,\kappa}}>0.
\]

To control the variance, define the spectral profile constants
\[
a_{\mathrm F}:=\int_{\mathbb R}Q(u)\,du,
\qquad
b_{\mathrm F}:=\int_{\mathbb R}|u|Q(u)\,du.
\]
Differentiation under the Fourier integral, justified by the absolute moments through order three, gives
\[
\widehat q_h^{\mathrm F}\in C^3(\mathbb R),
\qquad
|\widehat q_h^{\mathrm F}(\omega)|\le a_{\mathrm F}h,
\qquad
|(\widehat q_h^{\mathrm F})'(\omega)|\le b_{\mathrm F}h^2.
\]
Both \(G_{h,\sigma}\) and its derivatives have compact support. Twice integrating by parts gives, for \(v\ne0\),
\[
|\mathfrak l_{h,\sigma}(v)|
\le\frac1{2\pi}
\min\left\{
\int_{\mathbb R}|G_{h,\sigma}(\omega)|\,d\omega,\,
\frac1{v^2}\int_{\mathbb R}|G_{h,\sigma}''(\omega)|\,d\omega
\right\}.
\]
Together with the direct bound at \(v=0\), this proves absolute integrability of the inverse. The same argument applied to \(G_{h,\sigma}'\) proves absolute integrability of its inverse.

For \(0\le s\le2\), define the physical-space energies
\[
\mathfrak e_{s,h,\sigma}
:=\int_{\mathbb R}|v|^s|\mathfrak l_{h,\sigma}(v)|^2\,dv.
\]
Parseval's identity and the inverse-transform derivative identity give
\[
\mathfrak e_{0,h,\sigma}
=\frac1{2\pi}\int_{\mathbb R}|G_{h,\sigma}(\omega)|^2\,d\omega,
\qquad
\mathfrak e_{2,h,\sigma}
=\frac1{2\pi}\int_{\mathbb R}|G_{h,\sigma}'(\omega)|^2\,d\omega.
\]
On the spectral support, \(|\omega|\le6/h\). Thus
\[
|G_{h,\sigma}(\omega)|
\le a_{\mathrm F}h\,e^{18(\sigma/h)^2}.
\]
Also,
\[
G_{h,\sigma}'
=e^{\sigma^2\omega^2/2}
\left((\widehat q_h^{\mathrm F})'
      +\sigma^2\omega\widehat q_h^{\mathrm F}\right),
\]
so
\[
\begin{aligned}
|G_{h,\sigma}'(\omega)|
&\le h^2\left(b_{\mathrm F}
             +6a_{\mathrm F}(\sigma/h)^2\right)e^{18(\sigma/h)^2}\\
&\le h^2(b_{\mathrm F}+6a_{\mathrm F})
       (1+\sigma/h)^2e^{18(\sigma/h)^2}.
\end{aligned}
\]
Integrating these squared bounds over an interval of length \(12/h\) yields
\[
\begin{aligned}
\mathfrak e_{0,h,\sigma}
&\le\frac6\pi a_{\mathrm F}^2h\,e^{36(\sigma/h)^2},\\
\mathfrak e_{2,h,\sigma}
&\le\frac6\pi(b_{\mathrm F}+6a_{\mathrm F})^2
h^3(1+\sigma/h)^4e^{36(\sigma/h)^2}.
\end{aligned}
\]
For \(0\le\kappa\le2\), splitting at \(|v|=h\) gives
\[
|v|^\kappa\le h^\kappa+h^{\kappa-2}v^2.
\]
Consequently,
\[
\begin{aligned}
\mathfrak e_{\kappa,h,\sigma}
&\le h^\kappa\mathfrak e_{0,h,\sigma}
    +h^{\kappa-2}\mathfrak e_{2,h,\sigma}\\
&\le\frac6\pi
\left(a_{\mathrm F}^2+(b_{\mathrm F}+6a_{\mathrm F})^2\right)
h^{\kappa+1}(1+\sigma/h)^4e^{36(\sigma/h)^2}.
\end{aligned}
\]

The power inequality
\[
|v-\sigma z|^\kappa
\le2\left(|v|^\kappa+\sigma^\kappa|z|^\kappa\right)
\]
holds throughout the stated range, including \(\kappa=0\). Moreover,
\[
E_Z[|Z|^\kappa]\le E_Z[1+Z^2]=2.
\]
These inequalities also supply joint integrability after translation: the dominating terms are integrable products of a Gaussian moment and one of the finite energies above. Extending the latent integral to the real line, translating \(v=t+\sigma z\), and using
\(|\operatorname{Re}\mathfrak l|\le|\mathfrak l|\), we obtain
\[
\begin{aligned}
V_{q_h^{\mathrm F}}^{(16)}
&\le16E_Z\int_{\mathbb R}|v-\sigma Z|^\kappa
                      |\mathfrak l_{h,\sigma}(v)|^2\,dv\\
&\le32\left(\mathfrak e_{\kappa,h,\sigma}
                 +2\sigma^\kappa\mathfrak e_{0,h,\sigma}\right).
\end{aligned}
\]
Since
\[
\sigma^\kappa h
=h^{\kappa+1}(\sigma/h)^\kappa,
\qquad
(\sigma/h)^\kappa\le(1+\sigma/h)^2,
\]
enlarging the powers \(2\) and \(4\) to \(10\) gives
\[
V_{q_h^{\mathrm F}}^{(16)}
\le C_{\mathrm F,var}\,
h^{\kappa+1}(1+\sigma/h)^{10}e^{36(\sigma/h)^2},
\]
where
\[
C_{\mathrm F,var}
:=\frac{192}{\pi}
\left(3a_{\mathrm F}^2+(b_{\mathrm F}+6a_{\mathrm F})^2\right)+1>0.
\]
% lean: CausalSmith.Stat.NoisydoseWeakdesignTransition.ellF_variance_bound

The denominator moment bound implies
\[
\frac{12}{b_{q_h^{\mathrm F}}^{\mathrm{ref}}}
\le\frac{192}{c_{\mathrm F,\kappa}h^{\kappa+1}}.
\]
Define, for \(n\ge1\), \(h>0\), and \(\sigma\ge0\),
\[
\mathfrak s_{n,\sigma}(h)
:=\sqrt{\frac{h^{-\kappa-1}(1+\sigma/h)^{10}
                         e^{36(\sigma/h)^2}}{n}},
\]
and set
\[
C_{\mathrm F,score}
:=C_{\mathrm F,bias}
+\frac{192\sqrt{C_{\mathrm F,var}}}{c_{\mathrm F,\kappa}}
+2\sqrt3+1.
\]
Combining the bias, denominator, and variance bounds in the reference score gives
\[
A_{q_h^{\mathrm F}}^{\mathrm{ref}}
\le C_{\mathrm F,score}
\left(h^\beta+\mathfrak s_{n,\sigma}(h)+n^{-1/2}\right),
\qquad
0<h\le\frac14,\quad 0\le\sigma\le\frac14.
\]
% lean: CausalSmith.Stat.NoisydoseWeakdesignTransition.qF_certificate_bound

\item \emph{The direct Fourier witness.}
Define
\[
d:=2\beta+\kappa+1,
\qquad
D_n:=n^{-1/d},
\qquad n\ge1.
\]
The parameter ranges imply \(1<d\le5\). The direct scale satisfies
\[
D_n^d=\frac1n,
\qquad
\frac{D_n^{-\kappa-1}}n=D_n^{2\beta},
\qquad
n^{-1/2}\le D_n^\beta.
\]
The last inequality follows from \(D_n\le1\) and \(d\ge2\beta\).

For sufficiently large \(n\),
\[
0<D_n,
\qquad
n^{-2}\le D_n\le\frac18;
\]
the lower bound follows by comparing exponents, and the upper bound follows from \(D_n\to0\). Choose a dyadic bandwidth by
\[
2^{-(k+1)}<D_n\le2^{-k},
\qquad
h:=2^{-k}.
\]
Then
\[
D_n\le h<2D_n,
\qquad
n^{-2}\le h\le\frac14,
\]
so the pair belongs to the dictionary by \cref{obj:def:weight-dictionary}.

Suppose \(\sigma\le D_n\). For fixed nonnegative \(\sigma\), every factor inside \(\mathfrak s_{n,\sigma}(h)^2\) decreases as \(h\) increases. Hence
\[
\begin{aligned}
h^\beta&\le2^\beta D_n^\beta,\\
\mathfrak s_{n,\sigma}(h)
&\le\mathfrak s_{n,\sigma}(D_n)
 \le\sqrt{2^{10}e^{36}}\,D_n^\beta,
\end{aligned}
\]
where the second bound uses \(0\le\sigma/D_n\le1\) and the variance balance above. Therefore, with
\[
C_{\mathrm{direct}}
:=C_{\mathrm F,score}
\left(2^\beta+\sqrt{2^{10}e^{36}}+1\right)>0,
\]
the rounded Fourier entry satisfies
\[
A_{q_h^{\mathrm F}}^{\mathrm{ref}}
\le C_{\mathrm{direct}}D_n^\beta
=C_{\mathrm{direct}}\rho_{n,\sigma}.
\]
The equality is the direct branch of \cref{obj:def:frontier-rate}. Thus there is an integer \(n_{\mathrm{direct}}\), independent of \(\sigma\), beyond which this witness is available whenever \(\sigma\le D_n\).
% lean: CausalSmith.Stat.NoisydoseWeakdesignTransition.direct_fourier_dictionary_witness

\item \emph{The intermediate Fourier witness.}
For \(n\ge1\) and \(0<\sigma\le1/4\), define
\[
N_{\mathrm F}:=n\sigma^d,
\qquad
\Lambda_{\mathrm F}:=\log(e+N_{\mathrm F}),
\qquad
F_n(\sigma):=\frac{\sigma}{\sqrt{\Lambda_{\mathrm F}}}.
\]
When \(D_n<\sigma\), the identity \(nD_n^d=1\) implies \(N_{\mathrm F}\ge1\), and hence \(\Lambda_{\mathrm F}\ge1\).

We first establish the uniform logarithmic bound used for tuning. Define
\[
C_{\log}:=1024\cdot8!\left(\frac43\right)^8(e+1)>0.
\]
For \(N_{\mathrm F}\ge1\), the inequalities \(d\le5\) and \(\Lambda_{\mathrm F}\ge1\) give
\[
\Lambda_{\mathrm F}^{d/2}\le\Lambda_{\mathrm F}^3,
\qquad
\left(1+\frac{\sqrt{\Lambda_{\mathrm F}}}{12}\right)^{10}
\le1024\Lambda_{\mathrm F}^5.
\]
The eighth term of the exponential series gives
\[
\Lambda_{\mathrm F}^8
\le8!\left(\frac43\right)^8e^{3\Lambda_{\mathrm F}/4}.
\]
Using \(e^{\Lambda_{\mathrm F}}=e+N_{\mathrm F}\) and
\(e+N_{\mathrm F}\le(e+1)N_{\mathrm F}\), we conclude
\[
\frac{\Lambda_{\mathrm F}^{d/2}
\left(1+\sqrt{\Lambda_{\mathrm F}}/12\right)^{10}
e^{\Lambda_{\mathrm F}/4}}{N_{\mathrm F}}
\le C_{\log}.
\]
% lean: CausalSmith.Stat.NoisydoseWeakdesignTransition.fourier_log_tuning_bound

At the target bandwidth \(12F_n(\sigma)\),
\[
\frac{\sigma}{12F_n(\sigma)}
=\frac{\sqrt{\Lambda_{\mathrm F}}}{12},
\qquad
36\left(\frac{\sigma}{12F_n(\sigma)}\right)^2
=\frac{\Lambda_{\mathrm F}}4,
\qquad
\frac{F_n(\sigma)^{-d}}n
=\frac{\Lambda_{\mathrm F}^{d/2}}{N_{\mathrm F}}.
\]
Since \(12F_n(\sigma)\ge F_n(\sigma)\) and \(d>0\), these identities yield
\[
\frac{\mathfrak s_{n,\sigma}(12F_n(\sigma))^2}
     {(12F_n(\sigma))^{2\beta}}
\le
\frac{\Lambda_{\mathrm F}^{d/2}
\left(1+\sqrt{\Lambda_{\mathrm F}}/12\right)^{10}
e^{\Lambda_{\mathrm F}/4}}{N_{\mathrm F}}
\le C_{\log}.
\]
Thus
\[
\mathfrak s_{n,\sigma}(12F_n(\sigma))
\le\sqrt{C_{\log}}\,(12F_n(\sigma))^\beta.
\]
Dropping the two factors that are at least one from the logarithmic bound also gives
\[
\frac{F_n(\sigma)^{-d}}n\le C_{\log}.
\]
Because \(0<F_n(\sigma)\le\sigma\le1/4\) and \(d\ge2\beta\), it follows that
\[
\frac1n\le C_{\log}F_n(\sigma)^d
\le C_{\log}F_n(\sigma)^{2\beta},
\qquad
n^{-1/2}\le\sqrt{C_{\log}}F_n(\sigma)^\beta.
\]

We next place the target bandwidth in the dictionary window. If
\(F_n(\sigma)\le n^{-2}\), then the inverse-power bound would imply
\[
n^{2d-1}\le C_{\log}.
\]
For \(n\ge1\), \(d>1\) gives \(n\le n^{2d-1}\). Thus, once \(n>C_{\log}\),
\[
n^{-2}\le12F_n(\sigma).
\]
If also \(\sigma\le L_n^{-1/2}\), then
\[
12F_n(\sigma)\le12\sigma\le12L_n^{-1/2}\le\frac18
\]
for sufficiently large \(n\), since \(L_n\to\infty\). These thresholds are independent of \(\sigma\).

Choose \(k\) and \(h\) by
\[
2^{-(k+1)}<12F_n(\sigma)\le2^{-k},
\qquad
h:=2^{-k}.
\]
Then
\[
12F_n(\sigma)\le h<24F_n(\sigma),
\qquad
n^{-2}\le h\le\frac14,
\]
so the pair belongs to the dictionary. Monotonicity of the stochastic factor and the Fourier reference-score bound give
\[
A_{q_h^{\mathrm F}}^{\mathrm{ref}}
\le C_{\mathrm{intermediate}}F_n(\sigma)^\beta,
\]
where
\[
C_{\mathrm{intermediate}}
:=C_{\mathrm F,score}
\left(24^\beta+\sqrt{C_{\log}}\,12^\beta+\sqrt{C_{\log}}\right)>0.
\]
Under \(D_n<\sigma\le L_n^{-1/2}\), \cref{obj:def:frontier-rate} identifies \(F_n(\sigma)^\beta=\rho_{n,\sigma}\). This proves the intermediate witness for all \(n\ge n_{\mathrm{intermediate}}\), with a threshold independent of \(\sigma\).
% lean: CausalSmith.Stat.NoisydoseWeakdesignTransition.intermediate_fourier_dictionary_witness

\item \emph{Combining the Fourier regimes.}
Set
\[
C_{\mathrm F}:=\max\{C_{\mathrm{direct}},C_{\mathrm{intermediate}}\},
\qquad
n_{\mathrm F}:=\max\{n_{\mathrm{direct}},n_{\mathrm{intermediate}}\}.
\]
For \(n\ge n_{\mathrm F}\) and \(\sigma\le L_n^{-1/2}\), split according to \(\sigma\le D_n\) or \(D_n<\sigma\). The preceding two steps supply, respectively, a dictionary Fourier pair satisfying
\[
A_{q_{2^{-k}}^{\mathrm F}}^{\mathrm{ref}}
\le C_{\mathrm F}\rho_{n,\sigma}.
\]
The enlargement of the constants is valid because \(\rho_{n,\sigma}\ge0\). The direct case includes \(\sigma=0\).
% lean: CausalSmith.Stat.NoisydoseWeakdesignTransition.fourier_dictionary_witness

\item \emph{The polynomial certificate.}
We derive the inverse-degree score bound before choosing the dictionary index. For any positive odd \(m\), the inequality
\[
|\sin(mz)|\le m|\sin z|
\]
follows by induction from the sine addition formula. Combining it with
\(|\sin(mz)|\le1\) and the representation in \cref{obj:synth_2} gives
\[
0\le q_m^{\mathrm M}(t)\le1,
\qquad
q_m^{\mathrm M}(t)\le(2m|t|)^{-6}\quad(t\ne0),
\qquad |t|\le\frac12.
\]
Jordan's inequality and \(\sin(\arcsin(2t))=2t\) imply
\[
2|t|\le|\arcsin(2t)|\le\pi|t|.
\]
Therefore, if \(|t|\le1/(2m)\), then
\[
|m\arcsin(2t)|\le\frac\pi2,
\qquad
|\sin(m\arcsin(2t))|
\ge\frac2\pi m|\arcsin(2t)|
\ge\frac4\pi m|t|,
\]
and consequently
\[
q_m^{\mathrm M}(t)\ge\left(\frac2\pi\right)^6.
\]
At \(t=0\), this follows from \(q_m^{\mathrm M}(0)=1\).

For \(0\le s\le3\), define
\[
c_{\mathrm M,s}:=\left(\frac2\pi\right)^6I_s(1)>0,
\qquad
C_{\mathrm{tail}}:=\int_{\mathbb R}\frac8{(1+|u|)^3}\,du<\infty.
\]
Integrating the central lower bound and rescaling gives
\[
I_s(q_m^{\mathrm M})\ge c_{\mathrm M,s}m^{-(s+1)}.
\]
For the upper bound, put \(u=mt\). Splitting at \(|u|=1\), the preceding upper bounds give
\[
|t|^s q_m^{\mathrm M}(t)
\le m^{-s}\frac8{(1+|mt|)^3}.
\]
Extending this nonnegative envelope to the real line and changing variables yields
\[
c_{\mathrm M,s}m^{-(s+1)}
\le I_s(q_m^{\mathrm M})
\le C_{\mathrm{tail}}m^{-(s+1)}.
\]
In particular, define
\[
C_{\mathrm M,bias}:=\frac{64C_{\mathrm{tail}}}{c_{\mathrm M,\kappa}},
\qquad
C_{\mathrm M,ratio}:=\frac{192}{c_{\mathrm M,\kappa}}.
\]
Then
\[
B_{q_m^{\mathrm M}}^{\mathrm{ref}}\le C_{\mathrm M,bias}m^{-\beta},
\qquad
\frac{12}{b_{q_m^{\mathrm M}}^{\mathrm{ref}}}
\le C_{\mathrm M,ratio}m^{\kappa+1}.
\]

For the variance estimate, define the degree bound
\[
D_{\mathrm{pol},m}:=6(m-1).
\]
We first bound the coefficients of the declared polynomial. Let the second-kind Chebyshev polynomials be defined by
\[
\mathsf U^{\mathrm{Ch}}_0(s):=1,\qquad
\mathsf U^{\mathrm{Ch}}_1(s):=2s,\qquad
\mathsf U^{\mathrm{Ch}}_{i+2}(s)
:=2s\mathsf U^{\mathrm{Ch}}_{i+1}(s)-\mathsf U^{\mathrm{Ch}}_i(s).
\]
For a polynomial \(p\), define its coefficient mass by
\[
\|p\|_{\mathrm{coef}}
:=\sum_{i\ge0}|[t^i]p(t)|;
\]
the sum is finite. The recurrence, together with the triangle inequality and submultiplicativity of coefficient mass, proves inductively that
\[
\|\mathsf U^{\mathrm{Ch}}_i\|_{\mathrm{coef}}\le3^i.
\]
Indeed, the initial masses are \(1,2\), and
\[
2\cdot3^{i+1}+3^i\le3^{i+2}.
\]
Define the even coefficients by
\[
u^{\mathrm{Ch}}_{m,i}:=[s^{2i}]\mathsf U^{\mathrm{Ch}}_{m-1}(s),
\qquad 0\le i\le m-1.
\]
The identity
\(\mathsf U^{\mathrm{Ch}}_{m-1}(\cos z)=\sin(mz)/\sin z\),
obtained from the same recurrence, gives the polynomial representation
\[
q_m^{\mathrm M}(t)
=\left[
\frac1m\sum_{i=0}^{m-1}
u^{\mathrm{Ch}}_{m,i}(1-4t^2)^i
\right]^6.
\]
Since \(\|1-4t^2\|_{\mathrm{coef}}=5\), every summand inside the brackets has coefficient mass at most
\[
\frac{3^{m-1}5^i}{m}\le\frac{15^{m-1}}m.
\]
Summing the \(m\) terms and taking the sixth power gives
\[
\|q_m^{\mathrm M}\|_{\mathrm{coef}}
\le15^{D_{\mathrm{pol},m}},
\qquad
\deg q_m^{\mathrm M}\le D_{\mathrm{pol},m}.
\]
% lean: CausalSmith.Stat.NoisydoseWeakdesignTransition.qMPoly_coeffOneNorm_le

Define the probabilists' Hermite polynomials by
\[
\mathsf{He}_i(z)
:=i!\sum_{a=0}^{\lfloor i/2\rfloor}
\frac{(-1)^a z^{i-2a}}{2^a a!(i-2a)!},
\qquad i\ge0,\ z\in\mathbb R.
\]
The finite coefficient formula gives the recurrence
\[
\mathsf{He}_0(z)=1,\qquad
\mathsf{He}_1(z)=z,\qquad
\mathsf{He}_{i+1}(z)=z\mathsf{He}_i(z)-\mathsf{He}_i'(z),
\qquad i\ge0,\ z\in\mathbb R.
\]
Differentiating this recurrence proves by induction that
\[
\mathsf{He}_0'(z)=0,\qquad
\mathsf{He}_{i+1}'(z)=(i+1)\mathsf{He}_i(z),
\qquad i\ge0.
\]
Indeed, the initial identity is \(\mathsf{He}_1'=1\); for \(i\ge1\), the preceding derivative identity gives
\[
\begin{aligned}
\mathsf{He}_{i+1}'(z)
&=\mathsf{He}_i(z)+z\mathsf{He}_i'(z)-\mathsf{He}_i''(z)\\
&=\mathsf{He}_i(z)
+i\bigl(z\mathsf{He}_{i-1}(z)-\mathsf{He}_{i-1}'(z)\bigr)\\
&=(i+1)\mathsf{He}_i(z).
\end{aligned}
\]
% lean: Causalean.Mathlib.Probability.derivative_hermite_succ

For every real polynomial \(p\), Gaussian decay gives
\[
\lim_{z\to\pm\infty}p(z)e^{-z^2/2}=0.
\]
Integration by parts, followed by Gaussian normalization, therefore yields
\[
\begin{aligned}
\int_{\mathbb R}zp(z)e^{-z^2/2}\,dz
&=-\int_{\mathbb R}p(z)\frac{d}{dz}e^{-z^2/2}\,dz\\
&=\int_{\mathbb R}p'(z)e^{-z^2/2}\,dz,\\
E_Z[Zp(Z)]&=E_Z[p'(Z)].
\end{aligned}
\]
All polynomial factors have finite Gaussian absolute moments. Applying this identity to the Hermite recurrence gives
\[
E_Z[\mathsf{He}_0(Z)]=1,\qquad
E_Z[\mathsf{He}_{i+1}(Z)]
=E_Z[Z\mathsf{He}_i(Z)]-E_Z[\mathsf{He}_i'(Z)]=0,
\qquad i\ge0.
\]
% lean: Causalean.Mathlib.Probability.integral_hermite_succ

Define the Gaussian inner products by
\[
\mathfrak h_{i,a}:=E_Z[\mathsf{He}_i(Z)\mathsf{He}_a(Z)],
\qquad i,a\in\mathbb Z_{\ge0}.
\]
Applying integration by parts to the polynomial product
\(\mathsf{He}_i\mathsf{He}_{a+1}\), the recurrence and derivative identity give
\[
\begin{aligned}
\mathfrak h_{i+1,a+1}
&=E_Z[Z\mathsf{He}_i(Z)\mathsf{He}_{a+1}(Z)]
  -E_Z[\mathsf{He}_i'(Z)\mathsf{He}_{a+1}(Z)]\\
&=E_Z[\mathsf{He}_i'(Z)\mathsf{He}_{a+1}(Z)
       +\mathsf{He}_i(Z)\mathsf{He}_{a+1}'(Z)]
  -E_Z[\mathsf{He}_i'(Z)\mathsf{He}_{a+1}(Z)]\\
&=(a+1)\mathfrak h_{i,a}.
\end{aligned}
\]
% lean: Causalean.Mathlib.Probability.integral_hermite_succ_mul_succ
The mean identities supply the boundary values
\[
\mathfrak h_{0,0}=1,\qquad
\mathfrak h_{i+1,0}=\mathfrak h_{0,i+1}=0,
\qquad i\ge0.
\]
Induct on the first index, allowing every second index. The boundary values handle either index being zero. When both indices are positive, the recurrence reduces both by one. On the diagonal it gives
\[
\mathfrak h_{i+1,i+1}=(i+1)\mathfrak h_{i,i}
=(i+1)i!=(i+1)!;
\]
for \(i\ne a\), the induction hypothesis gives
\[
\mathfrak h_{i+1,a+1}=(a+1)\mathfrak h_{i,a}=0.
\]
Thus their Gaussian orthogonality is
\[
E_Z[\mathsf{He}_i(Z)\mathsf{He}_a(Z)]
=\mathbf1\{i=a\}\,i!.
\]
% lean: Causalean.Mathlib.Probability.integral_hermite_mul
Taylor expansion gives
\[
q_m^{\mathrm M}(t+\sigma z)
=\sum_{i=0}^{D_{\mathrm{pol},m}}
\frac{\sigma^i(q_m^{\mathrm M})^{(i)}(t)}{i!}z^i.
\]
Applying the finite inverse-heat sum from \cref{obj:synth_2} to this expansion replaces \(z^i\) by \(\mathsf{He}_i(z)\): this is the displayed finite formula for \(\mathsf{He}_i\), with the chain rule supplying the powers of \(\sigma\). Thus
\[
\ell_m^{\mathrm M}(t+\sigma Z)
=\sum_{i=0}^{D_{\mathrm{pol},m}}
\frac{\sigma^i(q_m^{\mathrm M})^{(i)}(t)}{i!}\mathsf{He}_i(Z).
\]
Squaring and using orthogonality yields
\[
E_Z[\ell_m^{\mathrm M}(t+\sigma Z)^2]
=\sum_{i=0}^{D_{\mathrm{pol},m}}
\frac{\sigma^{2i}\bigl((q_m^{\mathrm M})^{(i)}(t)\bigr)^2}{i!}.
\]
These finite identities include \(\sigma=0\).
% lean: CausalSmith.Stat.NoisydoseWeakdesignTransition.ellM_gaussian_second_moment

For \(0\le i\le D_{\mathrm{pol},m}\), write
\[
(D_{\mathrm{pol},m})_i
:=\frac{D_{\mathrm{pol},m}!}{(D_{\mathrm{pol},m}-i)!}.
\]
Differentiating the monomial expansion and using \(|t|\le1\) gives
\[
|(q_m^{\mathrm M})^{(i)}(t)|
\le\|q_m^{\mathrm M}\|_{\mathrm{coef}}(D_{\mathrm{pol},m})_i
\le15^{D_{\mathrm{pol},m}}(D_{\mathrm{pol},m})_i.
\]
Moreover,
\[
\frac{(D_{\mathrm{pol},m})_i^2}{i!}
=\binom{D_{\mathrm{pol},m}}i(D_{\mathrm{pol},m})_i
\le\binom{D_{\mathrm{pol},m}}iD_{\mathrm{pol},m}^i.
\]
The binomial theorem and \(I_\kappa(1)\le1\) therefore give
\[
\begin{aligned}
V_{q_m^{\mathrm M}}^{(16)}
&\le16I_\kappa(1)\,15^{2D_{\mathrm{pol},m}}
       (1+D_{\mathrm{pol},m}\sigma^2)^{D_{\mathrm{pol},m}}\\
&\le16\,15^{2D_{\mathrm{pol},m}}
       (1+D_{\mathrm{pol},m}\sigma^2)^{D_{\mathrm{pol},m}}.
\end{aligned}
\]
The case \(m=1\) has \(D_{\mathrm{pol},m}=0\), and the same formulas reduce to the constant pair.
% lean: CausalSmith.Stat.NoisydoseWeakdesignTransition.ellM_Vq_explicit_bound

For \(n\ge1\) and \(\sigma\in[0,1/4]\), define
\[
\Lambda_{\mathrm M,n,\sigma}:=\log(e+\sigma^2L_n).
\]
Here \(L_n\ge1\) and \(\Lambda_{\mathrm M,n,\sigma}\ge1\). Suppose
\[
m\le\frac{L_n}{4096\Lambda_{\mathrm M,n,\sigma}}.
\]
Then
\[
D_{\mathrm{pol},m}
\le\frac{6L_n}{4096\Lambda_{\mathrm M,n,\sigma}},
\qquad
1+D_{\mathrm{pol},m}\sigma^2
\le6(e+\sigma^2L_n),
\]
so
\[
\log(1+D_{\mathrm{pol},m}\sigma^2)
\le\Lambda_{\mathrm M,n,\sigma}+\log6.
\]
Using \(\log15\le15\), \(\log6\le6\), and
\(\Lambda_{\mathrm M,n,\sigma}\ge1\), we obtain the explicit calibration
\[
\begin{aligned}
&D_{\mathrm{pol},m}\log15
+\frac{D_{\mathrm{pol},m}}2
 \log(1+D_{\mathrm{pol},m}\sigma^2)\\
&\quad\le
\frac{6L_n}{4096\Lambda_{\mathrm M,n,\sigma}}
\left(15+\frac{\Lambda_{\mathrm M,n,\sigma}+6}{2}\right)
\le\frac{L_n}{16}.
\end{aligned}
\]
The last numerical inequality follows from
\(6(18+\Lambda_{\mathrm M,n,\sigma}/2)
\le256\Lambda_{\mathrm M,n,\sigma}\).
Exponentiating twice the bound gives
\[
15^{2D_{\mathrm{pol},m}}
(1+D_{\mathrm{pol},m}\sigma^2)^{D_{\mathrm{pol},m}}
\le e^{L_n/8}\le n^{1/4}
\]
whenever \(L_n=1+\log n\ge2\). Consequently,
\[
V_{q_m^{\mathrm M}}^{(16)}\le16n^{1/4},
\qquad
\sqrt{\frac{V_{q_m^{\mathrm M}}^{(16)}}{n}}\le4n^{-3/8}.
\]

Finally, the standard logarithmic-power comparison gives
\[
\frac{L_n^4}{n^{3/8}}\longrightarrow0,
\qquad
L_n^4\le n^{3/8}
\quad\text{for all sufficiently large }n.
\]
The degree condition implies \(1\le m\le L_n\). Since
\(\kappa+\beta+1\le4\), it follows that
\[
m^{\kappa+\beta+1}\le m^4\le L_n^4\le n^{3/8},
\qquad
n^{-3/8}m^{\kappa+1}\le m^{-\beta}.
\]
Also,
\[
n^{-1/2}\le n^{-3/8}
\le n^{-3/8}m^{\kappa+1}\le m^{-\beta}.
\]
Substitution into the reference score proves
\[
A_{q_m^{\mathrm M}}^{\mathrm{ref}}
\le C_{\mathrm{poly}}m^{-\beta},
\qquad
C_{\mathrm{poly}}
:=C_{\mathrm M,bias}+4C_{\mathrm M,ratio}+2\sqrt3>0,
\]
uniformly for all sufficiently large \(n\), all public \(\sigma\), and all positive odd \(m\) satisfying the degree condition.
% lean: CausalSmith.Stat.NoisydoseWeakdesignTransition.qM_tuned_certificate_bound

\item \emph{The polynomial dictionary witness.}
Define the real target index
\[
m_{\mathrm{tar},n,\sigma}
:=\frac{L_n}{4096\Lambda_{\mathrm M,n,\sigma}}.
\]
Uniformly over \(\sigma\in[0,1/4]\),
\[
1\le\Lambda_{\mathrm M,n,\sigma}\le\log(e+L_n).
\]
Because \(\log(e+L_n)/L_n\to0\), eventually
\[
\log(e+L_n)\le\frac{L_n}{16384}.
\]
For sufficiently large \(n\), the inequality \(\log x\le x-1\) also gives
\[
L_n\le en\le n^2.
\]
Therefore the target lies in the uniform window
\[
4\le m_{\mathrm{tar},n,\sigma}\le L_n\le n^2.
\]
Choose \(j\ge0\) such that
\[
2^j\le\frac{m_{\mathrm{tar},n,\sigma}}2<2^{j+1},
\qquad
m:=2^j+1.
\]
The lower bound on the target implies \(j\ge1\), so \(m\) is odd. The rounding inequalities give
\[
\frac{m_{\mathrm{tar},n,\sigma}}4\le m\le m_{\mathrm{tar},n,\sigma},
\]
or equivalently,
\[
\frac{L_n}{16384\Lambda_{\mathrm M,n,\sigma}}
\le m
\le\frac{L_n}{4096\Lambda_{\mathrm M,n,\sigma}}.
\]
Since \(m\le n^2\), this polynomial pair belongs to the dictionary by \cref{obj:def:weight-dictionary}.
% lean: CausalSmith.Stat.NoisydoseWeakdesignTransition.polynomial_tuned_dictionary_degree

The same eventual logarithmic-power bound used above places the direct cutoff below the Fourier cutoff. Indeed,
\[
L_n^{d/2}\le L_n^4\le n^{3/8}\le n.
\]
Taking the power \(-1/d<0\) yields
\[
D_n=n^{-1/d}\le L_n^{-1/2}.
\]
Hence, when \(L_n^{-1/2}<\sigma\), the third branch of \cref{obj:def:frontier-rate} applies:
\[
b_{n,\sigma}
=\frac{\Lambda_{\mathrm M,n,\sigma}}{L_n},
\qquad
\rho_{n,\sigma}
=\left(\frac{\Lambda_{\mathrm M,n,\sigma}}{L_n}\right)^\beta.
\]
The lower rounding bound implies
\[
m^{-1}\le16384\,\frac{\Lambda_{\mathrm M,n,\sigma}}{L_n},
\qquad
m^{-\beta}
\le16384^\beta
\left(\frac{\Lambda_{\mathrm M,n,\sigma}}{L_n}\right)^\beta.
\]
Applying the polynomial certificate therefore gives
\[
A_{q_m^{\mathrm M}}^{\mathrm{ref}}
\le C_{\mathrm M}\rho_{n,\sigma},
\qquad
C_{\mathrm M}:=C_{\mathrm{poly}}\,16384^\beta>0.
\]
Taking a common integer threshold \(n_{\mathrm M}\) for the certificate, the target window, and the cutoff comparison proves the polynomial witness uniformly in \(\sigma\).
% lean: CausalSmith.Stat.NoisydoseWeakdesignTransition.inverseheat_dictionary_witness

\item \emph{Selection and the common bound.}
We first convert the reference scores to the public scores of \cref{obj:synth_4}. For every dictionary pair, nonnegativity of the weight and \(|t|^\beta\le1\) on the latent interval give
\[
0\le I_{\kappa+\beta}(q)\le I_\kappa(q)<\infty.
\]
Together with \(I_\kappa(q)>0\), this makes every summand of \(A_q^{\mathrm{ref}}\) nonnegative. The public quantities satisfy
\[
\begin{aligned}
b_q(K)&=p_{\min}c_{\mathrm{lo}}I_\kappa(q),\\
B_q(K)&=\frac{c_{\mathrm{hi}}}{64c_{\mathrm{lo}}}B_q^{\mathrm{ref}},\\
V_q(K)&=\frac{c_{\mathrm{hi}}}{16}V_q^{(16)},\\
\frac{12}{b_q(K)}\sqrt{\frac{V_q(K)}n}
&=\frac{\sqrt{c_{\mathrm{hi}}/16}}{16p_{\min}c_{\mathrm{lo}}}
\frac{12}{b_q^{\mathrm{ref}}}\sqrt{\frac{V_q^{(16)}}n}.
\end{aligned}
\]
The second identity follows from the two bias definitions, and the last follows by factoring the nonnegative constant out of the square root and using \(b_q^{\mathrm{ref}}=I_\kappa(q)/16\). Define
\[
S(K):=\max\left\{1,\frac{c_{\mathrm{hi}}}{64c_{\mathrm{lo}}},
\frac{\sqrt{c_{\mathrm{hi}}/16}}{16p_{\min}c_{\mathrm{lo}}}\right\}\ge1.
\]
Bounding the coefficient of each nonnegative summand by \(S(K)\) yields
\[
\begin{aligned}
A_q(K)
&=\frac{c_{\mathrm{hi}}}{64c_{\mathrm{lo}}}B_q^{\mathrm{ref}}
+\frac{\sqrt{c_{\mathrm{hi}}/16}}{16p_{\min}c_{\mathrm{lo}}}
\frac{12}{b_q^{\mathrm{ref}}}\sqrt{\frac{V_q^{(16)}}n}
+2\sqrt{\frac3n}\\
&\le S(K)\left(B_q^{\mathrm{ref}}
+\frac{12}{b_q^{\mathrm{ref}}}\sqrt{\frac{V_q^{(16)}}n}
+2\sqrt{\frac3n}\right)
=S(K)A_q^{\mathrm{ref}}.
\end{aligned}
\]
% lean: CausalSmith.Stat.NoisydoseWeakdesignTransition.Aq_le_scoreScale_mul_unitAq
Set
\[
C:=S(K)\max\{C_{\mathrm F},C_{\mathrm M}\}>0,
\qquad
n_0:=\max\{3,n_{\mathrm F},n_{\mathrm M}\}.
\]
For \(n\ge n_0\), \(L_n\ge1\). Every branch of the localization scale in \cref{obj:def:frontier-rate} is nonnegative, so
\[
\rho_{n,\sigma}\ge0.
\]
For either family witness \(q\) constructed above, the reference-score bound therefore gives
\[
A_q(K)\le S(K)A_q^{\mathrm{ref}}
\le S(K)\max\{C_{\mathrm F},C_{\mathrm M}\}\rho_{n,\sigma}
=C\rho_{n,\sigma}.
\]
Thus both public witness bounds hold with the same \(C\) and \(n_0\).
% lean: CausalSmith.Stat.NoisydoseWeakdesignTransition.score_le_of_unitScore_le

For completeness, the ordered selection rule gives the required minimizing inequality. Enumerate the dictionary in its fixed order as
\[
N_{\mathrm{dict}}:=|\mathcal D_{n,\sigma}|,
\qquad
q_{(1)},\ldots,q_{(N_{\mathrm{dict}})},
\]
with the constant pair first. Define the successive retained entries by
\[
q_{\mathrm{best},0}:=1,
\qquad
q_{\mathrm{best},i}:=
\begin{cases}
q_{(i)},&A_{q_{(i)}}(K)<A_{q_{\mathrm{best},i-1}}(K),\\
q_{\mathrm{best},i-1},&A_{q_{\mathrm{best},i-1}}(K)\le A_{q_{(i)}}(K),
\end{cases}
\quad 1\le i\le N_{\mathrm{dict}}.
\]
Induction on \(i\) proves
\[
q_{\mathrm{best},i}\in\{1,q_{(1)},\ldots,q_{(i)}\},
\qquad
A_{q_{\mathrm{best},i}}(K)\le A_q(K)
\quad
\text{for every }q\in\{1,q_{(1)},\ldots,q_{(i)}\}.
\]
If the new score is strictly smaller, the new entry satisfies the bound against every preceding entry by transitivity; otherwise the retained entry also bounds the new score. Equal scores retain the earlier entry. At the final index this is the specified first minimizer, and therefore
\[
\widehat q\in\mathcal D_{n,\sigma},
\qquad
A_{\widehat q}(K)\le A_q(K)
\quad(q\in\mathcal D_{n,\sigma}).
\]
% lean: CausalSmith.Stat.NoisydoseWeakdesignTransition.selectedTag_minimizes

Now split at the stated cutoff. If \(\sigma\le L_n^{-1/2}\), the Fourier witness supplies \(k\ge0\) with its pair in the dictionary and
\[
A_{\widehat q}(K)\le A_{q_{2^{-k}}^{\mathrm F}}(K)
\le C\rho_{n,\sigma}.
\]
If \(L_n^{-1/2}<\sigma\), the polynomial witness supplies \(j\ge0\), with \(m=2^j+1\), and
\[
A_{\widehat q}(K)\le A_{q_m^{\mathrm M}}(K)
\le C\rho_{n,\sigma}.
\]
These two cases establish the selected-score bound and both asserted family witnesses with the same constant and threshold.
% lean: CausalSmith.Stat.NoisydoseWeakdesignTransition.dictionary_score_rate
\end{enumerate}
\end{proof}

Admissibility holds for every positive sample size, whereas the score comparison holds beyond a threshold independent of the error scale. The Fourier comparison covers \(\sigma\le L_n^{-1/2}\), and the polynomial comparison covers \(L_n^{-1/2}<\sigma\), within the stated large-sample guarantee. These comparisons concern candidates in the same dictionary, so the selected score reflects the applicable localization family through a single public ordering and minimization rule.

\subsection{Observable achievability and interval calibration}

The certificate and dictionary bounds provide the mathematical components of the observable upper bound in \cref{obj:thm:observable-upper}: the certificate controls the error associated with an admissible entry, and the comparison bounds control the selected public score. That result states the expected absolute-error rate uniformly over the public error-scale range and the declared measurable potential-path spaces, for samples beyond its common threshold. The balance is between latent localization and the sampling variability of observable Gaussian inversion.

For interval construction, the same score determines the radius \(A_{\widehat q}(K)/\alpha\) in \cref{obj:def:honest-interval}. The finite-sample certificate and Markov's inequality supply the error control used for calibration, while intersection with \([0,1]\) incorporates the target range. Empty modulo-three blocks trigger the full target-range interval. Accordingly, \cref{obj:thm:finite-honesty} gives coverage of at least \(1-\alpha\) for every \(\alpha\in(0,1)\), every integer \(n\ge0\), every \(\sigma\in[0,1/4]\), and every structural law in the specified model class. Its expected-length bound applies beyond a common sample-size threshold and has order \(\rho_{n,\sigma}\), with a constant depending on \(\alpha\). Coverage and the uniform large-sample precision guarantee therefore retain their respective quantifiers.

\paragraph{Verification note.}
The theorem statements and proofs are machine-checked in Lean 4. The definition blocks written for exposition (\cref{obj:synth_1,obj:synth_2,obj:synth_3,obj:synth_4,obj:synth_5,obj:synth_6,obj:synth_7,obj:synth_8}) restate in prose the Lean definitions and constructions that the theorems use; these prose blocks are not themselves machine-checked statements.


\section{Weighted packets and observed-experiment converse}

The observable risk certificates are complemented by structural alternatives whose causal means are separated while their observed sample laws remain close. The construction belongs to the model class of \cref{obj:def:model-class} and compares the experiments in \cref{obj:def:observed-experiment}. Its analytic ingredient is a localized perturbation with exact cancellation under the weak-design weight.

The witness subclass uses the rank restrictions in \cref{obj:ass:rank-uniform,obj:ass:rank-treatment-independence,obj:ass:structural-threshold,obj:ass:binary-potential-outcomes}. Specifically, the auxiliary rank coordinate \(U\) is conditionally uniform on \([0,1]\) and independent of the latent dose given the stratum. A common rank generates the entire threshold schedule, and every potential outcome is binary. These restrictions give a measurable schedule with a prescribed conditional mean and place the witnesses within the bounded-outcome class. The threshold embedding and evaluation map are those of \cref{obj:def:potential-path-space}.

\subsection{Filtered Jacobi packets}

The packet construction separates localization from support. Its degree index \leanref{sym:m}{\ensuremath{m}} controls concentration and cancellation, while a fixed taper gives regularity at the boundary of the support. The filtered Jacobi sums \leanref{sym:PacketHandle}{\ensuremath{K_m,\widetilde K_m}} use an endpoint construction for positive design decay and a symmetric interior construction when the design-decay exponent is zero. The normalized packet \leanref{sym:psim}{\ensuremath{\psi_m}} is defined as follows.

\begin{definitionv}{def:packet-handle}[Filtered Jacobi packet \(\psi_m\)]
For each integer \(m\ge4\), the zero-extended filtered Jacobi packet is
\[
\psi_m(u)=
\begin{cases}
(1-u^2)^r K_m(2u^2-1)/K_m(-1),
& |u|\le1,\ \kappa>0,\\
(1-u^2)^r\widetilde K_m(u)/\widetilde K_m(0),
& |u|\le1,\ \kappa=0,\\
0,& |u|>1.
\end{cases}
\]
The fixed taper parameter and smooth spectral filter are
\[
r=4,\qquad
\eta(s)=
\begin{cases}
\exp\{-1/((s-9/8)(15/8-s))\},
& 9/8<s<15/8,\\
0,& \text{otherwise}.
\end{cases}
\]
The filter is nonnegative, supported inside \((1,2)\), and positive on \([5/4,7/4]\). Write \(\mathsf J_j^{(a,b)}\) for the standard Jacobi polynomial of degree \(j\), and define its squared norm by
\[
H_j^{a,b}
=
\int_{-1}^{1}
\bigl(\mathsf J_j^{(a,b)}(z)\bigr)^2
(1-z)^a(1+z)^b\,dz.
\]
For \(\kappa>0\), the endpoint parameter and filtered sum are
\[
b=\frac{\kappa-1}{2},\qquad
K_m(z)=\sum_{j\ge0}\eta(j/m)
\frac{\mathsf J_j^{(r,b)}(-1)\mathsf J_j^{(r,b)}(z)}
{H_j^{r,b}}.
\]
For \(\kappa=0\), the symmetric interior sum is
\[
\widetilde K_m(u)=\sum_{j\ge0}\eta(j/m)
\frac{\mathsf J_j^{(r,r)}(0)\mathsf J_j^{(r,r)}(u)}
{H_j^{r,r}}.
\]
\end{definitionv}

The two branches accommodate the parameter range of the localization estimates. For \(\kappa>0\), the endpoint parameter satisfies \(b>-1/2\). At \(\kappa=0\), the symmetric construction uses the interior point \(0\) and parameters \((r,r)\), both strictly above \(-1/2\). The fixed filter selects degrees proportional to \(m\), giving a common index for spatial concentration and weighted cancellation.

For the auxiliary polynomial identity, write \(P_j^{(a,b)}=\mathsf J_j^{(a,b)}\); both notations denote the same standard Jacobi polynomial. The parameters \(a,b\) in this identity are local polynomial parameters. For a real number \(v\) and an integer \(i\ge0\), the generalized binomial coefficient is the falling factorial divided by \(i!\), with value \(1\) at \(i=0\). The following statement records that convention together with the identity.

\begin{lemmav}{lem:jacobi-binomial}[Jacobi polynomial binomial representation]
Let \(a,b\in\mathbb R\) satisfy \(a,b>-1/2\), let \(j\ge0\) be an integer, and let \(z\in\mathbb R\). Write \(P_j^{(a,b)}=\mathsf J_j^{(a,b)}\) for the standard Jacobi polynomial. Define the generalized binomial coefficient, for \(v\in\mathbb R\) and integers \(i\ge0\), by
\[
\binom{v}{i}=\frac{\prod_{k=0}^{i-1}(v-k)}{i!},
\qquad \binom{v}{0}=1,
\]
where the empty product equals \(1\). Then
\[
P_j^{(a,b)}(z)
=2^{-j}\sum_{i=0}^{j}
\binom{j+a}{i}\binom{j+b}{j-i}
(z-1)^{j-i}(z+1)^i.
\]
\end{lemmav}

\begin{proof}[Proof of \cref{obj:lem:jacobi-binomial}]
\begin{enumerate}
\item Set
\[
y:=\frac{z-1}{2}\in\mathbb R.
\]
For \(v\in\mathbb R\) and \(\nu\in\mathbb N_0\), define the rising factorial by
\[
v^{\overline{\nu}}:=\prod_{k=0}^{\nu-1}(v+k),
\qquad v^{\overline{0}}:=1.
\]
The normalized shifted expansion of the Jacobi polynomial gives
\[
P_j^{(a,b)}(z)
=
\frac{(a+1)^{\overline j}}{j!}
\sum_{\nu=0}^{j}
(-1)^\nu\binom j\nu
\frac{(j+a+b+1)^{\overline\nu}}
     {(a+1)^{\overline\nu}}
\left(\frac{1-z}{2}\right)^\nu.
\]
For \(0\le \nu\le j\), splitting the rising product and using the factorial identity for the ordinary binomial coefficient yields
\[
(a+1)^{\overline j}
=(a+1)^{\overline\nu}
 (a+1+\nu)^{\overline{j-\nu}},
\qquad
\binom j\nu\,\nu!\,(j-\nu)!=j!.
\]
Since \(a+1>0\), every factor in \((a+1)^{\overline\nu}\) is positive, including the empty-product case. Thus cancellation is valid, and
\[
\begin{aligned}
\frac{(a+1)^{\overline j}}{j!}\binom j\nu
\frac{(j+a+b+1)^{\overline\nu}}
     {(a+1)^{\overline\nu}}
&=
\frac{(a+1+\nu)^{\overline{j-\nu}}}{(j-\nu)!}
\frac{(j+a+b+1)^{\overline\nu}}{\nu!}\\
&=\binom{j+a}{j-\nu}\binom{j+a+b+\nu}{\nu}.
\end{aligned}
\]
The last equality follows by reversing the factors in each falling product defining the generalized binomial coefficient. Moreover,
\[
(-1)^\nu\left(\frac{1-z}{2}\right)^\nu=y^\nu.
\]
Consequently,
\[
P_j^{(a,b)}(z)
=\sum_{\nu=0}^{j}
\binom{j+a}{j-\nu}\binom{j+a+b+\nu}{\nu}y^\nu.
\]
% lean: CausalSmith.Stat.NoisydoseWeakdesignTransition.jacobi_coefficient_binomial

\item We establish the polynomial identity
\[
\sum_{i=0}^{j}
\binom{\gamma}{j-i}\binom{\gamma_{\mathrm{aux}}}{i}
\xi^i(\xi+1)^{j-i}
=
\sum_{\nu=0}^{j}
\binom{\gamma}{j-\nu}
\binom{\gamma+\gamma_{\mathrm{aux}}-(j-\nu)}{\nu}\xi^\nu,
\qquad
\gamma,\gamma_{\mathrm{aux}}\in\mathbb R,
\]
where
\[
\xi\ \text{is a polynomial indeterminate over }\mathbb R.
\]
For \(0\le i\le\nu\le j\), splitting the falling product gives
\[
\begin{aligned}
\binom{\gamma}{j-i}\binom{j-i}{j-\nu}
&=
\frac{\prod_{k=0}^{j-i-1}(\gamma-k)}
     {(j-\nu)!\,(\nu-i)!}\\
&=
\binom{\gamma}{j-\nu}
\binom{\gamma-(j-\nu)}{\nu-i}.
\end{aligned}
\]
Ordinary binomial symmetry,
\[
\binom{j-i}{j-\nu}=\binom{j-i}{\nu-i},
\]
therefore gives the same product identity with the latter lower index.
% lean: CausalSmith.Stat.NoisydoseWeakdesignTransition.genBinom_choose_product

The generalized Vandermonde convolution for falling-factorial coefficients states
\[
\binom{\gamma+\gamma_{\mathrm{aux}}}{\nu}
=
\sum_{i=0}^{\nu}\binom{\gamma}{i}\binom{\gamma_{\mathrm{aux}}}{\nu-i},
\qquad \nu\in\mathbb N_0.
\]
% lean: CausalSmith.Stat.NoisydoseWeakdesignTransition.genBinom_vandermonde

For \(0\le\nu\le j\), the terms with \(i>\nu\) contribute zero to the coefficient of \(\xi^\nu\) on the left. Expanding \((\xi+1)^{j-i}\), then applying the product identity and Vandermonde convolution, gives that coefficient as
\[
\begin{aligned}
\sum_{i=0}^{\nu}
\binom{\gamma}{j-i}\binom{\gamma_{\mathrm{aux}}}{i}\binom{j-i}{\nu-i}
&=
\binom{\gamma}{j-\nu}
\sum_{i=0}^{\nu}
\binom{\gamma_{\mathrm{aux}}}{i}\binom{\gamma-(j-\nu)}{\nu-i}\\
&=
\binom{\gamma}{j-\nu}
\binom{\gamma+\gamma_{\mathrm{aux}}-(j-\nu)}{\nu}.
\end{aligned}
\]
For \(\nu>j\), every \(0\le i\le j\) satisfies
\[
\nu-i>j-i,
\qquad
\binom{j-i}{\nu-i}=0.
\]
Hence the coefficient is zero on both sides. Equality of all coefficients proves the polynomial identity.
% lean: CausalSmith.Stat.NoisydoseWeakdesignTransition.genBinom_bernstein_poly

\item Evaluate this identity at \(y\), with parameters \(j+a\) and \(j+b\). For \(0\le\nu\le j\),
\[
(j+a)+(j+b)-(j-\nu)=j+a+b+\nu.
\]
Comparison with the expansion in the first step therefore yields
\[
P_j^{(a,b)}(z)
=
\sum_{i=0}^{j}
\binom{j+a}{j-i}\binom{j+b}{i}
y^i(y+1)^{j-i}.
\]
% lean: CausalSmith.Stat.NoisydoseWeakdesignTransition.stdJacobi_binomial_repr

\item Reverse the finite sum by the bijection \(i\mapsto j-i\) of \(\{0,\ldots,j\}\). This gives
\[
P_j^{(a,b)}(z)
=
\sum_{i=0}^{j}
\binom{j+a}{i}\binom{j+b}{j-i}
y^{j-i}(y+1)^i.
\]
Since
\[
y+1=\frac{z+1}{2},
\qquad
2^{j-i}2^i=2^j
\quad(0\le i\le j),
\]
each power product satisfies
\[
y^{j-i}(y+1)^i
=
\frac{(z-1)^{j-i}(z+1)^i}{2^{j-i}2^i}
=
2^{-j}(z-1)^{j-i}(z+1)^i.
\]
Substitution and extraction of the common factor prove
\[
P_j^{(a,b)}(z)
=
2^{-j}\sum_{i=0}^{j}
\binom{j+a}{i}\binom{j+b}{j-i}
(z-1)^{j-i}(z+1)^i.
\]
All product splittings and sum reversals also apply when \(j=0\), using empty products and zeroth powers equal to \(1\).
% lean: CausalSmith.Stat.NoisydoseWeakdesignTransition.stdJacobi_binomial_repr
\end{enumerate}
\end{proof}

\Cref{obj:lem:jacobi-binomial} fixes the polynomial normalization used in the filtered sums. Positive squared norms and orthogonality are supplied by the standard Jacobi normalization \citep[Table 18.3.1, Jacobi row and caption]{https://dlmf.nist.gov/18.3}; the corresponding gamma-ratio asymptotic controls the degree dependence of normalization factors \citep[equation (5.11.12), positive real axis]{https://dlmf.nist.gov/5.11.E12}. The resulting packet properties are collected in the next statement.

\begin{theoremv}{thm:weighted-packet-construction}[Weighted packet construction]*
For every fixed \(\kappa\in[0,2]\), the degeneracy exponent, there exist real constants \(C>0\) and \(c>0\) such that, for every integer \(m\ge4\), the filtered-Jacobi construction has positive normalizing value:
\[
\begin{cases}
K_m(-1)>0, & \kappa>0,\\
\widetilde K_m(0)>0, & \kappa=0.
\end{cases}
\]
Here \(K_m(-1)\) and \(\widetilde K_m(0)\) are respectively the endpoint and symmetric interior normalizing values of the construction. The resulting normalized, tapered, zero-extended packet \(\psi_m\) is even and continuously differentiable on \(\mathbb R\), and satisfies
\[
\psi_m(u)=0\quad\text{for }u\notin[-1,1],
\qquad
\psi_m(0)=1,
\]
\[
|\psi_m(u)|\le C,\qquad
|\psi_m'(u)|\le Cm
\quad\text{for every }u\in\mathbb R,
\]
\[
\int_{-1}^{1}|u|^\kappa|\psi_m(u)|\,du
\le Cm^{-\kappa-1},
\qquad
\int_{-1}^{1}|u|^\kappa\psi_m(u)^2\,du
\le Cm^{-\kappa-1},
\]
and
\[
\int_{-1}^{1}|u|^\kappa u^j\psi_m(u)\,du=0
\quad\text{for every integer }0\le j<cm.
\]
\end{theoremv}
% CAUSALSMITH-CITED-SCOPE-BEGIN thm:weighted-packet-construction
\verificationfootnotetext{\textbf{Formalization scope.} This result uses the cited conclusion \citep{PetrushevXu2005} (Theorem 2.4, equation (2.14)), \citep{https://people.math.sc.edu/pencho/Publications/kpx-06-28-06-web.pdf} (Theorem 2.2, equation (2.4)), \citep{https://dlmf.nist.gov/18.3} (Table 18.3.1, Jacobi row and its caption; orthogonality and squared norm.) and \citep{https://dlmf.nist.gov/5.11.E12} (Equation 5.11.12, restricted to the positive real axis.). The cited source proof is not formalized here: Lean verifies its use and all remaining steps. If that cited conclusion is false, the portions of this result that depend on it are not certified.}
% CAUSALSMITH-CITED-SCOPE-END thm:weighted-packet-construction

The packet retains unit height at the evaluation point while its weighted absolute mass and squared mass decrease with degree. Its derivative grows at most linearly in degree, and its weighted moments vanish through an order proportional to degree. These properties serve distinct statistical purposes: normalization preserves target separation, regularity permits Hölder-admissible mean perturbations, and cancellation reduces the discrepancy after Gaussian convolution.

The localization component uses the filtered Jacobi estimate of \citet[Theorem 2.4, equation (2.14)]{PetrushevXu2005}. The local difference estimate supplies derivative control for the filtered sums \citep[Theorem 2.2, equation (2.4)]{https://people.math.sc.edu/pencho/Publications/kpx-06-28-06-web.pdf}. The taper makes the zero extension compatible with continuous differentiability. Together, these ingredients support the pointwise and weighted-integral bounds in \cref{obj:thm:weighted-packet-construction}.

Weighted cancellation is the role of Jacobi orthogonality \citep[Table 18.3.1, Jacobi row and caption]{https://dlmf.nist.gov/18.3}. In the positive-decay branch, the quadratic coordinate \(2u^2-1\) aligns the even weighted moments with the endpoint Jacobi weight; symmetry accounts for the odd moments. At zero design decay, the symmetric interior sum provides cancellation directly under the unweighted dose measure. The separate interior construction therefore supplies the same packet guarantees at that endpoint of the parameter range.

\subsection{Structural witnesses and observed distances}

We now combine the packet with independent structural coordinates. The tuning uses a localization width \(h\), a support radius \leanref{sym:ellrad}{\ensuremath{\ell}}, and a cancellation cutoff \leanref{sym:Jmom}{\ensuremath{J}}. In the inverse regimes, the relations \(h=\ell/m\) and \(J=\lfloor c_pm\rfloor\) connect these quantities to packet degree, where \(c_p>0\) is the packet's cancellation constant: its weighted moments of every order \(j<c_pm\) vanish. The grouped definition below specifies the reference and perturbed means, their structural laws, and their observed projections, and it lists the smallness and sample-size conditions under which the profiles have the asserted range and smoothness.

\begin{definitionv}{synth_5}[Structural lower-bound witness laws]
Fix \(0<\beta\le1\), \(0\le\kappa\le2\), \(0\le\sigma\le1/4\), a positive integer sample size \(n\), and the public measurable schedule space \((\mathcal S,\mathcal E)\) with threshold embedding \(\iota\). Put \(a_0=1/2\) and \(d=2\beta+\kappa+1\). For public class constants \(K=(p_{\min},c_{\mathrm{lo}},c_{\mathrm{hi}})\), require
\[
0<p_{\min}\le1/2,\qquad
c_{\mathrm{lo}}\le(\kappa+1)2^\kappa\le c_{\mathrm{hi}}.
\]
The three witness laws share independent seed coordinates \(X,T,Z,U\), where \(P\{X=0\}=P\{X=1\}=1/2\), \(Z\sim N(0,1)\), \(U\sim\operatorname{Unif}[0,1]\), and
\[
g_\kappa(t)=(\kappa+1)2^\kappa|t|^\kappa\mathbf1\{|t|\le1/2\}
\]
is the density of \(T\). Set \(A=a_0+T\).

The construction uses an amplitude \(\epsilon\ge0\), a comparison constant \(C>0\), and packet constants \(C_p,c_p>0\) such that, for every integer \(m\ge4\), the normalized filtered-Jacobi packet \(\psi_m\) of \cref{obj:def:packet-handle} is continuously differentiable on \(\mathbb R\) and satisfies
\[
|\psi_m(u)|\le C_p,\qquad
|\psi_m'(u)|\le C_pm
\quad(u\in\mathbb R),
\qquad
\int_{-1}^{1}|u|^\kappa u^j\psi_m(u)\,du=0\quad(j\in\mathbb N,\ j<c_pm).
\]
For example, the packet constants \(C_p,c_p\) of \cref{obj:thm:observed-lower} satisfy these conditions. The tuning data consist of \(h,\ell\in(0,1/4]\) and nonnegative integers \(m,J\). In the direct regime \(\sigma\le n^{-1/d}\), set
\[
h_0=n^{-1/d},\qquad h=\ell=h_0,\qquad m=J=0,
\]
and \(\delta(t)=\epsilon h_0^\beta(1-(t/h_0)^2)^2\mathbf1\{|t|\le h_0\}\), and require
\[
0<h_0\le1/4,\qquad 4\epsilon\le1,\qquad \epsilon h_0^\beta\le3/8.
\]
The bandwidth condition holds exactly when \(n\ge4^d\), so the direct-regime construction applies only to such \(n\). In either inverse regime \(\sigma>n^{-1/d}\), use the packet \(\psi_m\) and require
\[
m\ge4,\qquad h=\ell/m,\qquad J=\lfloor c_pm\rfloor\ge1,
\qquad 2\epsilon C_p\le1,\qquad \epsilon h^\beta C_p\le3/8,
\qquad \delta(t)=\epsilon h^\beta\psi_m(t/\ell).
\]
Under these requirements, the perturbation is supported on \([-\ell,\ell]\), is integrable against \(g_\kappa\), and satisfies
\[
\int_{\mathbb R}t^jg_\kappa(t)\delta(t)\,dt=0\qquad(0\le j<J);
\]
in the inverse regimes this follows from the packet cancellation because \(j<\lfloor c_pm\rfloor\) implies \(j<c_pm\). Also under these requirements, the continuous profiles
\[
\mu_\star(a)=1/2,\qquad \mu_\pm(a)=1/2\pm\delta(a-a_0)
\]
take values in \([1/8,7/8]\) and obey \(|\mu_\pm(a)-\mu_\pm(a')|\le|a-a'|^\beta\).

For \(s\in\{+,-,\star\}\), define \(P_s\) as the six-coordinate structural law of
\[
\bigl(X,A,Z,\iota(\mu_s,U),\mathbf1\{U\le\mu_s(A)\},U\bigr).
\]
Thus \(Y_s(a)=\mathbf1\{U\le\mu_s(a)\}\) simultaneously for all \(a\), \(Y_s=Y_s(A)\), and the observed projection is \((X,W,Y_s)\) with \(W=A+\sigma Z\). The retained rank \(U\) gives the common-seed representation required by the structural laws. The three laws have the same \((X,A)\) marginal and satisfy
\[
\mu_++\mu_-=2\mu_\star,\qquad
\theta(P_\pm)=1/2\pm\delta(0),\qquad \theta(P_\star)=1/2.
\]

Let \(Q_s=\mathcal L_{P_s}(X,W,Y_s)\). For every Borel set \(B\subseteq\mathbb R\) and \(x\in\{0,1\}\),
\[
Q_s\{X=x,W-a_0\in B,Y_s=1\}
=\frac12\int_{-1/2}^{1/2}g_\kappa(t)\mu_s(a_0+t)\Phi_\sigma(B-t)\,dt,
\]
with the analogous formula containing \(1-\mu_s\) when \(Y_s=0\). Here \(\Phi_\sigma=N(0,\sigma^2)\), with \(\Phi_0=\delta_0\), and \(Q_{P_s}^n=Q_s^{\otimes n}\). For \(\sigma>0\), the construction additionally requires the comparison bound below. It is a condition on the constant \(C\), not a consequence of the preceding requirements; the witnesses of \cref{obj:thm:observed-lower} satisfy it with that theorem's constant \(C\):
\[
\chi^2(Q_s,Q_\star)
\le C\sigma^{-\kappa-1}h^{2\beta+2\kappa+2}
\sum_{j=J}^{\infty}\frac{(\ell^2/\sigma^2)^j}{j!},
\qquad s\in\{+,-\}.
\]
\end{definitionv}

The calibrated design coefficient places the common density inside the public envelopes, and equal stratum probabilities satisfy the overlap restriction. Opposite perturbations of the constant reference mean create target separation while preserving the latent-design marginal. The observed formulas retain both outcome values: they compare the joint distribution of the contaminated dose and the binary outcome within each stratum.

The reference law is used to measure the discrepancy of each perturbed observed law. We make the divergence convention explicit, including its value when absolute continuity fails.

\begin{definitionv}{synth_7}[Chi-squared divergence of observed laws]
For probability measures \(Q\) and \(R\) on the observed-record space \(\{0,1\}\times\mathbb R\times[0,1]\), define
\[
\chi^2(Q,R)=
\begin{cases}
\displaystyle \int\left(\frac{dQ}{dR}-1\right)^2\,dR,
& Q\ll R,\\[4pt]
+\infty,& Q\not\ll R.
\end{cases}
\]
The integral is allowed to equal \(+\infty\). For the witness laws \(P_\pm\) and \(P_\star\), the notation \(\chi^2(P_\pm,P_\star)\) denotes
\[
\chi^2\!\left(\mathcal L_{P_\pm}(X,W,Y),
              \mathcal L_{P_\star}(X,W,Y)\right),
\qquad W=A+\sigma Z,
\]
the divergence between their single-record observed laws.
\end{definitionv}

Thus the structural-law notation in \cref{obj:synth_7} always refers to the induced observed-record comparison. The testing statement uses total variation on the observed sample space, with the following convention.

\begin{definitionv}{synth_8}[Total variation distance]
For probability measures \(\nu_1\) and \(\nu_2\) on a common measurable space \((\Omega,\mathcal F)\), define their total variation distance by
\[
\operatorname{TV}(\nu_1,\nu_2)
=\sup_{B\in\mathcal F}\bigl|\nu_1(B)-\nu_2(B)\bigr|.
\]
\end{definitionv}

Total variation controls the difference between probabilities of every observable event, including testing and interval-coverage events. The next result supplies structural witnesses with a prescribed total-variation bound and causal separation at the rate in \cref{obj:def:frontier-rate}.

\begin{theoremv}{thm:observed-lower}[Explicit observed lower witnesses]*
Fix public class constants \(K=(p_{\min},c_{\mathrm{lo}},c_{\mathrm{hi}})\) and exponents \(\beta,\kappa\) satisfying
\begin{itemize}
\item \textbf{(Class constants.)} \(0<p_{\min}\le1/2\) and \(0<c_{\mathrm{lo}}\le c_{\mathrm{hi}}\).
\item \textbf{(Exponents.)} \(0<\beta\le1\) and \(0\le\kappa\le2\).
\item \textbf{(Design calibration.)}
\[
c_{\mathrm{lo}}\le(\kappa+1)2^\kappa\le c_{\mathrm{hi}}.
\]
\end{itemize}
For every total-variation level \(\tau>0\), there exist constants \(c,C,\epsilon,C_p,c_p>0\) and an integer \(n_0\), chosen uniformly over the public potential-path space, sample size, and error scale, with the following properties.

For every integer \(m\ge4\), the packets of \cref{obj:def:packet-handle} have positive normalization:
\[
K_m(-1)>0\quad(\kappa>0),
\qquad
\widetilde K_m(0)>0\quad(\kappa=0).
\]
Each \(\psi_m\) is even and continuously differentiable on \(\mathbb R\), vanishes outside \([-1,1]\), and satisfies
\[
\psi_m(0)=1,\qquad
|\psi_m(u)|\le C_p,\qquad
|\psi_m'(u)|\le C_pm
\quad(u\in\mathbb R),
\]
\[
\int_{-1}^{1}|u|^\kappa|\psi_m(u)|\,du\le C_pm^{-\kappa-1},
\qquad
\int_{-1}^{1}|u|^\kappa\psi_m(u)^2\,du\le C_pm^{-\kappa-1},
\]
\[
\int_{-1}^{1}|u|^\kappa u^j\psi_m(u)\,du=0
\quad(j\in\mathbb N,\ j<c_pm).
\]

For every public potential-path space \((\mathcal S,\mathcal E)\) of \cref{obj:def:potential-path-space}, every integer \(n\ge n_0\), and every \(\sigma\in[0,1/4]\), there exist structural probability laws \(P_+,P_-,P_\star\) in the model class of \cref{obj:def:model-class}, with public constants \(K\) and parameters \((\beta,\kappa,\sigma)\). All three satisfy \cref{obj:ass:binary-potential-outcomes,obj:ass:rank-uniform,obj:ass:rank-treatment-independence,obj:ass:structural-threshold,obj:ass:iid}.

These laws admit the following common construction. Let \(X\) be the stratum coordinate on \(\mathcal X=\{0,1\}\), \(A\) the latent dose, \(Z\) the standardized Gaussian error, and \(U\) the rank coordinate. Write \(a_0=1/2\) for the evaluation dose and \(t=A-a_0\) for the centered latent dose. Under each law, \(X,A,Z,U\) are independent, with
\[
P(X=0)=P(X=1)=\frac12,\qquad
Z\sim N(0,1),\qquad U\sim\operatorname{Unif}[0,1].
\]
The conditional density \(g_x\) of \(A-a_0\) in either stratum is
\[
g_x(t)=
\begin{cases}
(\kappa+1)2^\kappa|t|^\kappa,& |t|\le1/2,\\
0,& |t|>1/2,
\end{cases}
\qquad x\in\mathcal X.
\]
There are continuous profiles \(\mu_+,\mu_-:[0,1]\to[0,1]\) and a reference profile \(\mu_\star(a)=1/2\). Under \(P_+\), \(P_-\), and \(P_\star\), respectively, the random schedule and realized outcome are
\[
Y(\cdot)=\iota(\mu_+,U),\quad
Y(\cdot)=\iota(\mu_-,U),\quad
Y(\cdot)=\iota(\mu_\star,U),
\qquad
Y=e(Y(\cdot),A).
\]
Here \(e\) and \(\iota\) are the evaluation and threshold-schedule maps of \cref{obj:def:potential-path-space}. In particular, the latent-design marginals agree:
\[
\mathcal L_{P_+}(X,A)=
\mathcal L_{P_-}(X,A)=
\mathcal L_{P_\star}(X,A).
\]

Writing
\[
\theta(P)=\mathbb E_P[Y(a_0)]
\]
for the population causal mean, and using the frontier rate \(\rho_{n,\sigma}\) of \cref{obj:def:frontier-rate} and the observed product experiment \(Q_P^n\) of \cref{obj:def:observed-experiment}, the witnesses satisfy
\[
|\theta(P_+)-\theta(P_-)|\ge c\rho_{n,\sigma},
\qquad
\operatorname{TV}(Q_{P_+}^n,Q_{P_-}^n)\le\tau.
\]

The same profiles have common tuning data \(h,\ell\in\mathbb R\) and \(m,J\in\mathbb N\), where \(h\) is the localization width, \(\ell\) the support radius, \(m\) the packet degree, and \(J\) the cancellation cutoff. Define the direct-regime bandwidth
\[
h_0=n^{-1/(2\beta+\kappa+1)}.
\]
The tuning satisfies
\[
0<h\le1/4,\qquad 0<\ell\le1/4,
\]
\[
\sigma\le h_0\ \Longrightarrow\
h=h_0,\quad \ell=h,\quad J=0,
\]
\[
h_0<\sigma\ \Longrightarrow\
m\ge4,\quad h=\ell/m,\quad
J=\lfloor c_pm\rfloor,\quad J\ge1.
\]
Define the mean perturbation \(\delta\) by
\[
\delta(t)=
\begin{cases}
\epsilon h_0^\beta\bigl(1-(t/h_0)^2\bigr)^2,
  &\sigma\le h_0,\ |t|\le h_0,\\
0,&\sigma\le h_0,\ |t|>h_0,\\
\epsilon h^\beta\psi_m(t/\ell),&h_0<\sigma.
\end{cases}
\]
For every \(a\in[0,1]\),
\[
\mu_\pm(a)=\frac12\pm\delta(a-a_0).
\]
Moreover,
\[
\delta(t)=0\quad(|t|>\ell),
\qquad
\int_{\mathbb R}t^jg_x(t)\delta(t)\,dt=0
\quad(x\in\mathcal X,\ j\in\mathbb N,\ j<J).
\]

For \(\sigma>0\), let \(\chi^2(P_\pm,P_\star)\) denote the chi-squared divergence between the respective one-observation laws induced by \(P_\pm\) and \(P_\star\) in \cref{obj:def:observed-experiment}. Both signs satisfy
\[
\chi^2(P_\pm,P_\star)
\le
C\sigma^{-\kappa-1}h^{2\beta+2\kappa+2}
\sum_{j=J}^{\infty}\frac{(\ell^2/\sigma^2)^j}{j!}.
\]
\end{theoremv}
% CAUSALSMITH-CITED-SCOPE-BEGIN thm:observed-lower
\verificationfootnotetext{\textbf{Formalization scope.} This result uses the cited conclusion \citep{PetrushevXu2005} (Theorem 2.4, equation (2.14)), \citep{https://people.math.sc.edu/pencho/Publications/kpx-06-28-06-web.pdf} (Theorem 2.2, equation (2.4)), \citep{https://dlmf.nist.gov/18.3} (Table 18.3.1, Jacobi row and its caption; orthogonality and squared norm.) and \citep{https://dlmf.nist.gov/5.11.E12} (Equation 5.11.12, restricted to the positive real axis.). The cited source proof is not formalized here: Lean verifies its use and all remaining steps. If that cited conclusion is false, the portions of this result that depend on it are not certified.}
% CAUSALSMITH-CITED-SCOPE-END thm:observed-lower

The witnesses change the causal response while keeping the distribution of stratum and latent dose fixed. Their opposite mean perturbations preserve a common reference profile, and the threshold representation carries those profiles into legal potential-outcome schedules. The comparison concerns the complete observed record, including its outcome mark. Consequently, the separation and total-variation clauses address the same experiment used by the estimation and interval criteria.

The analytic inputs to the packet component of \cref{obj:thm:observed-lower} are the localization estimate \citep[Theorem 2.4, equation (2.14)]{PetrushevXu2005}, the local difference estimate \citep[Theorem 2.2, equation (2.4)]{https://people.math.sc.edu/pencho/Publications/kpx-06-28-06-web.pdf}, Jacobi orthogonality and squared norms \citep[Table 18.3.1, Jacobi row and caption]{https://dlmf.nist.gov/18.3}, and the gamma-ratio asymptotic \citep[equation (5.11.12), positive real axis]{https://dlmf.nist.gov/5.11.E12}. Their roles remain normalization, localization, derivative control, and weighted cancellation.

\subsection{Regime-specific observed-law comparison}

The direct and inverse constructions in \cref{obj:thm:observed-lower} balance target separation against observable discrimination. In the direct regime, \(\sigma\le h_0\), the compact bump has support radius \(h_0=n^{-1/d}\) and height \(\epsilon h_0^\beta\). The witness separation is therefore \(2\epsilon h_0^\beta\), yielding the direct rate \(n^{-\beta/d}\). The theorem's observed-sample total-variation bound covers this entire regime, including \(\sigma=0\).

For \(h_0<\sigma\), packet normalization gives separation \(2\epsilon h^\beta\). The support radius and localization width differ by the packet degree, and cancellation removes all weighted moments below \(J\). In the displayed chi-squared comparison, the surviving series begins at that cutoff; its argument \(\ell^2/\sigma^2\) records the relation between perturbation support and measurement error. The comparison therefore incorporates both the cancellation order and the permitted support radius.

Across the intermediate range
\[
n^{-1/d}<\sigma\le L_n^{-1/2},
\]
the separation guarantee has order
\[
\left\{\frac{\sigma}{\sqrt{\log(e+n\sigma^d)}}\right\}^{\beta}.
\]
Across the compact-support range
\[
L_n^{-1/2}<\sigma\le1/4,
\]
it has order
\[
\left\{\frac{\log(e+\sigma^2L_n)}{L_n}\right\}^{\beta}.
\]
These are the two inverse branches of \cref{obj:def:frontier-rate}. Both are realized by packets supported within the stated radius bound, with a cancellation cutoff proportional to degree. The common total-variation guarantee in \cref{obj:thm:observed-lower} completes the observed-law comparison uniformly over the error-scale range and the declared potential-path spaces, beyond its common sample-size threshold.

\subsection{Risk, honest length, and transition consequences}

Let \(\Delta=|\theta(P_+)-\theta(P_-)|\) denote the witness target separation. The two-point testing bound converts a separated pair of targets and close sample laws into an absolute-error lower bound \citep[Section 2.2, Theorem 2.2]{Tsybakov2009}. With any fixed \(\tau\in(0,1)\), the witnesses in \cref{obj:thm:observed-lower} give
\[
R_n(\sigma)\ge \frac{\Delta}{4}(1-\tau)
\ge c_R\rho_{n,\sigma},
\]
for a positive constant \(c_R\), every sufficiently large sample size, every \(\sigma\in[0,1/4]\), and every public potential-path space in \cref{obj:def:potential-path-space}. The fixed parameters satisfy the design calibration and exponent restrictions of \cref{obj:thm:observed-lower}. This lower bound applies to the possibly randomized estimators allowed by \cref{obj:def:minimax-risk}.

For honest connected intervals, fix \(\alpha\in(0,1/2)\) and choose \(\tau\in(0,1-2\alpha)\). Uniform coverage at the two witness laws and their total-variation bound give the corresponding expected-length consequence
\[
H_n(\sigma)\ge \Delta(1-2\alpha-\tau)
\ge c_H\rho_{n,\sigma},
\]
with \(c_H>0\), under the same uniform large-sample quantifiers. Connectedness makes an interval containing both witness targets have length at least their separation. The factor \(1-2\alpha-\tau\) records the coverage requirement and the observable discrepancy of the pair. These consequences supply the lower bounds in \cref{obj:thm:uniform-frontier}; the observable upper bounds are provided by \cref{obj:thm:observable-upper,obj:thm:finite-honesty}.

The transition comparisons in \cref{obj:thm:uniform-frontier} connect the three descriptions throughout fixed multiplicative neighborhoods of their thresholds. Using the transition thresholds \(a_n,b_n\) and comparison scales \(D_n,F_n(\sigma),P_n(\sigma)\) defined there, for every fixed multiplicative factor \(M\ge1\),
\[
F_n(\sigma)\asymp D_n
\qquad
\bigl(\sigma\in[a_n/M,Ma_n]\bigr),
\]
and
\[
F_n(\sigma)\asymp P_n(\sigma)
\qquad
\bigl(\sigma\in[b_n/M,Mb_n]\bigr),
\]
once the sample size exceeds the stated window threshold. The constants and threshold may depend on \(M\) and the fixed parameters. Raising these comparisons to the smoothness exponent shows that neighboring rate formulas give comparable risk and honest length throughout each window.

At the error-free endpoint, \cref{obj:prop:error-free-reduction} gives
\[
R_n(0)\asymp H_n(0)\asymp n^{-\beta/(2\beta+\kappa+1)}.
\]
For every fixed \(\sigma\in(0,1/4]\), the final conclusion of \cref{obj:thm:uniform-frontier} instead gives
\[
R_n(\sigma)\asymp H_n(\sigma)
\asymp
\left(\frac{\log\log n}{\log n}\right)^\beta.
\]
The first comparison uses constants determined by the calibrated class and fixed coverage level; the second permits dependence on the fixed positive error scale. Both retain uniformity over the declared public potential-path spaces. The observed-law witnesses thus support the direct endpoint, both inverse regimes, and the transition comparisons within one structural model.

\paragraph{Verification note.}
The theorem statements and proofs are machine-checked in Lean 4. The definition blocks written for exposition (\cref{obj:synth_1,obj:synth_2,obj:synth_3,obj:synth_4,obj:synth_5,obj:synth_6,obj:synth_7,obj:synth_8}) restate in prose the Lean definitions and constructions that the theorems use; these prose blocks are not themselves machine-checked statements. The model assumptions specify the structural class under which the conclusions hold. The theorem-local verification notes identify the cited analytic conclusions, whose published source proofs remain external, and record the scope of their use.


\section{Proofs of the main results}\label{sec:deferred-proofs}

\begin{proof}[Proof of \cref{obj:prop:identification}]
Fix \(x\in\mathcal X\). Write
\[
p'_x=P'(X=x),
\qquad
\mu'_x(a)=\frac{E_{P'}[\mathbf1\{X=x\}Y(a)]}{p'_x},
\qquad a\in[0,1].
\]
By \cref{obj:ass:stratum-overlap}, both \(p_x\) and \(p'_x\) are strictly positive.

\begin{enumerate}
\item Define the centered observed dose by
\[
V=W-a_0=A-a_0+\sigma Z.
\]
For every Borel set \(B\subseteq\mathbb R\), define
\[
\begin{aligned}
\nu_{0x}(B)&=P(A-a_0\in B\mid X=x),
&
\nu'_{0x}(B)&=P'(A-a_0\in B\mid X=x),\\
M_{0x}(B)&=\frac{P(X=x,V\in B)}{p_x},
&
M'_{0x}(B)&=\frac{P'(X=x,V\in B)}{p'_x}.
\end{aligned}
\]
Equality of the observed laws equates their normalized restrictions to each stratum, and therefore
\[
M_{0x}=M'_{0x}.
\]

Let the Gaussian noise measure be
\[
\Phi_\sigma=\mathcal L(\sigma Z),
\qquad
\Phi_0=\delta_0,
\qquad
\delta_0(B)=\mathbf1\{0\in B\}.
\]
By \cref{obj:ass:gaussian-channel}, this is the centered Gaussian probability measure with variance \(\sigma^2\) under either structural law. For Borel sets \(B_Z,B_t\subseteq\mathbb R\), \cref{obj:ass:error-independence} gives
\[
\begin{aligned}
P(Z\in B_Z,A-a_0\in B_t\mid X=x)
&=\frac{P(Z\in B_Z)\,P(X=x,A-a_0\in B_t)}{p_x}\\
&=P(Z\in B_Z)\,\nu_{0x}(B_t).
\end{aligned}
\]
Scaling the error coordinate consequently yields
\[
\mathcal L_P(A-a_0,\sigma Z\mid X=x)
=\nu_{0x}\otimes\Phi_\sigma.
\]
The same factorization holds under \(P'\). Pushing these product measures forward under addition gives
\[
\begin{aligned}
M_{0x}(B)
&=\int_{\mathbb R}\int_{\mathbb R}
\mathbf1\{t+z\in B\}\,\Phi_\sigma(dz)\,\nu_{0x}(dt)
=(\nu_{0x}*\Phi_\sigma)(B),\\
M'_{0x}(B)&=(\nu'_{0x}*\Phi_\sigma)(B).
\end{aligned}
\]

The density representation in \cref{obj:ass:weak-design}, together with the support condition in \cref{obj:ass:latent-support} and \(a_0=1/2\), gives
\[
\nu_{0x}(B)=\int_{B\cap[-1/2,1/2]}g_x(t)\,dt.
\]
Both latent laws are probability measures, and the same support conclusion holds under \(P'\):
\[
\nu_{0x}(\mathbb R)=\nu'_{0x}(\mathbb R)=1,
\qquad
\nu_{0x}\bigl(\mathbb R\setminus[-1/2,1/2]\bigr)
=\nu'_{0x}\bigl(\mathbb R\setminus[-1/2,1/2]\bigr)=0.
\]
Thus
\[
\nu_{0x}*\Phi_\sigma=\nu'_{0x}*\Phi_\sigma
\quad\Longrightarrow\quad
\nu_{0x}=\nu'_{0x},
\]
by \cref{obj:lem:compact-gaussian-convolution-injective}, taking its two auxiliary measures to be zero. Its finiteness, support, and error-scale hypotheses are satisfied, including at \(\sigma=0\).
% lean: CausalSmith.Stat.NoisydoseWeakdesignTransition.latentDoseLaw_eq_of_obsLaw_eq

\item Since \(\beta>0\), the Hölder means are continuous on \([0,1]\). Define, for every Borel \(B\subseteq\mathbb R\),
\[
\begin{aligned}
\nu_{1x}(B)
&=\int_{B\cap[-1/2,1/2]}
\max\{\mu_x(t+a_0),0\}\,\nu_{0x}(dt),\\
\nu'_{1x}(B)
&=\int_{B\cap[-1/2,1/2]}
\max\{\mu'_x(t+a_0),0\}\,\nu'_{0x}(dt).
\end{aligned}
\]
The support of the unmarked laws makes these the marked latent measures in the statement. On this support, \cref{obj:ass:mean-range} makes the positive parts equal to the means themselves. Continuity and boundedness ensure integrability, and hence
\[
\begin{aligned}
\nu_{1x}(\mathbb R)
&=\int_{[-1/2,1/2]}\mu_x(t+a_0)\,\nu_{0x}(dt)\le1,
&
\nu'_{1x}(\mathbb R)&\le1,\\
\nu_{1x}\bigl(\mathbb R\setminus[-1/2,1/2]\bigr)
&=0,
&
\nu'_{1x}\bigl(\mathbb R\setminus[-1/2,1/2]\bigr)&=0.
\end{aligned}
\]

By \cref{obj:lem:realized-dose-mean},
\[
0\le Y\le1
\qquad P\text{-almost surely and }P'\text{-almost surely}.
\]
We may therefore define the finite nonnegative observed marked measures by
\[
M_{1x}(B)=E_P[\mathbf1\{V\in B\}Y\mid X=x],
\qquad
M'_{1x}(B)=E_{P'}[\mathbf1\{V\in B\}Y\mid X=x].
\]
They are determined by the normalized observed laws, so
\[
M_{1x}=M'_{1x}.
\]

For a fixed Borel \(B\), use the bounded measurable test
\[
(a,x',z)\longmapsto
\mathbf1\{x'=x,\ a+\sigma z-a_0\in B\},
\qquad
(a,x',z)\in\mathbb R\times\mathcal X\times\mathbb R.
\]
The bounded-test identity in \cref{obj:lem:realized-dose-mean} gives
\[
\begin{aligned}
M_{1x}(B)
&=\frac{E_P[\mathbf1\{X=x,V\in B\}Y]}{p_x}\\
&=\frac{E_P[\mathbf1\{X=x,V\in B\}\mu_X(A)]}{p_x}\\
&=E_P[\mathbf1\{V\in B\}\mu_x(A)\mid X=x].
\end{aligned}
\]
Using the conditional product law established above, and the supported latent mean, we obtain
\[
\begin{aligned}
M_{1x}(B)
&=\int_{[-1/2,1/2]}\int_{\mathbb R}
\mathbf1\{t+z\in B\}\mu_x(t+a_0)\,
\Phi_\sigma(dz)\,\nu_{0x}(dt)\\
&=(\nu_{1x}*\Phi_\sigma)(B).
\end{aligned}
\]
The same calculation under \(P'\) yields
\[
M'_{1x}=\nu'_{1x}*\Phi_\sigma.
\]
Consequently,
\[
\nu_{1x}*\Phi_\sigma=\nu'_{1x}*\Phi_\sigma
\quad\Longrightarrow\quad
\nu_{1x}=\nu'_{1x},
\]
again by \cref{obj:lem:compact-gaussian-convolution-injective} with the two auxiliary measures zero. The displayed finiteness and support properties discharge its measure hypotheses.
% lean: CausalSmith.Stat.NoisydoseWeakdesignTransition.latentMeanMeasure_eq_of_obsLaw_eq

\item For every integer \(j\ge0\), put
\[
\varepsilon_j=\frac1{j+3},
\qquad
B_j=(-\varepsilon_j,\varepsilon_j).
\]
Then
\[
0<\varepsilon_j\le\frac12,
\qquad
B_j\subseteq[-1/2,1/2].
\]
The lower envelope in \cref{obj:ass:weak-design} implies
\[
g_x(t)\ge c_{\mathrm{lo}}|t|^\kappa>0
\qquad
\text{for Lebesgue-almost every }t\in B_j\setminus\{0\}.
\]
The singleton \(\{0\}\) has Lebesgue measure zero, whereas
\[
\int_{B_j}1\,dt=2\varepsilon_j>0.
\]
If the integral of the nonnegative density over \(B_j\) were zero, the density would vanish almost everywhere there, contradicting this positivity. Thus
\[
\nu_{0x}(B_j)=\int_{B_j}g_x(t)\,dt>0.
\]
This proves every asserted shrinking-neighborhood mass is strictly positive.
% lean: CausalSmith.Stat.NoisydoseWeakdesignTransition.latentDoseLaw_local_mass_pos

\item For \(t\in B_j\), both \(t+a_0\) and \(a_0\) belong to \([0,1]\). The radius-one Hölder bound in \cref{obj:ass:holder-mean} gives
\[
|\mu_x(t+a_0)-\mu_x(a_0)|
\le |t|^\beta
\le\varepsilon_j^\beta.
\]
Integrating this pointwise bound against the latent probability law yields
\[
\left|
\int_{B_j}\bigl(\mu_x(t+a_0)-\mu_x(a_0)\bigr)\,\nu_{0x}(dt)
\right|
\le\varepsilon_j^\beta\,\nu_{0x}(B_j).
\]
Since the marked mass equals the integral of the mean, and the denominator is positive,
\[
\begin{aligned}
\left|
\frac{\nu_{1x}(B_j)}{\nu_{0x}(B_j)}-\mu_x(a_0)
\right|
&=
\frac{
\left|\displaystyle\int_{B_j}
\bigl(\mu_x(t+a_0)-\mu_x(a_0)\bigr)\,\nu_{0x}(dt)\right|
}{\nu_{0x}(B_j)}\\
&\le\varepsilon_j^\beta.
\end{aligned}
\]
% lean: CausalSmith.Stat.NoisydoseWeakdesignTransition.latentLocalRatio_error
Finally,
\[
\varepsilon_j\longrightarrow0,
\qquad
\varepsilon_j^\beta\longrightarrow0,
\]
because \(\beta>0\). Squeezing the nonnegative absolute error therefore proves
\[
\lim_{j\to\infty}\frac{\nu_{1x}(B_j)}{\nu_{0x}(B_j)}
=\mu_x(a_0).
\]
% lean: CausalSmith.Stat.NoisydoseWeakdesignTransition.latentLocalRatio_tendsto

\item Projecting the common observed law onto its stratum coordinate gives
\[
p_x=P(X=x)=P'(X=x)=p'_x.
\]
% lean: CausalSmith.Stat.NoisydoseWeakdesignTransition.strataProb_eq_of_obsLaw_eq

\item The equalities of the unmarked and marked latent measures imply, by uniqueness of densities with respect to their common unmarked probability measure,
\[
\max\{\mu_x(t+a_0),0\}
=\max\{\mu'_x(t+a_0),0\}
\qquad
\text{for }\nu_{0x}\text{-almost every }t\in[-1/2,1/2].
\]
The mean-range bounds remove the positive parts, giving
\[
\mu_x(t+a_0)=\mu'_x(t+a_0)
\qquad
\text{for }\nu_{0x}\text{-almost every }t\in[-1/2,1/2].
\]

The lower design envelope makes \(g_x(t)>0\) for Lebesgue-almost every \(t\) in the full latent interval: the possible exceptional point \(t=0\) is null. Thus, for every Borel \(B\subseteq\mathbb R\),
\[
\nu_{0x}(B)=\int_{B\cap[-1/2,1/2]}g_x(t)\,dt=0
\quad\Longrightarrow\quad
\int_{B\cap[-1/2,1/2]}1\,dt=0.
\]
Indeed, a zero integral of the nonnegative density forces it to vanish almost everywhere on that set, and its almost-everywhere positivity forces the set to be Lebesgue-null. The mean equality consequently transfers to Lebesgue-almost every centered dose:
\[
\mu_x(t+a_0)=\mu'_x(t+a_0)
\qquad
\text{for Lebesgue-almost every }t\in[-1/2,1/2].
\]

For \(t,t'\in[-1/2,1/2]\), the two Hölder bounds give
\[
\begin{aligned}
|\mu_x(t+a_0)-\mu_x(t'+a_0)|&\le|t-t'|^\beta,\\
|\mu'_x(t+a_0)-\mu'_x(t'+a_0)|&\le|t-t'|^\beta.
\end{aligned}
\]
Since \(\beta>0\), both centered mean functions are continuous on the closed interval. If they differed at any point, continuity would preserve their difference on a relative neighborhood of that point. Every such neighborhood has positive Lebesgue measure, including at either endpoint, contradicting their almost-everywhere agreement. Hence
\[
\mu_x(t+a_0)=\mu'_x(t+a_0)
\qquad
\text{for every }t\in[-1/2,1/2].
\]
For every \(a\in[0,1]\),
\[
a-a_0\in[-1/2,1/2],
\qquad
(a-a_0)+a_0=a.
\]
Substitution therefore proves
\[
\mu_x(a)=\mu'_x(a)
\qquad
\text{for every }a\in[0,1].
\]
Since \(x\) was arbitrary, all the preceding conclusions hold in both strata.
% lean: CausalSmith.Stat.NoisydoseWeakdesignTransition.condPotMean_eqOn_of_latent_measures_eq

\item Bounded potential outcomes ensure integrability at \(a_0\). The two strata partition the structural space, so pointwise
\[
Y(a_0)=\sum_{x\in\mathcal X}\mathbf1\{X=x\}Y(a_0).
\]
Taking expectations and using \cref{obj:def:conditional-potential-mean}, with the positive stratum probabilities, gives
\[
\begin{aligned}
\theta(P)
&=E_P[Y(a_0)]\\
&=\sum_{x\in\mathcal X}E_P[\mathbf1\{X=x\}Y(a_0)]\\
&=\sum_{x\in\mathcal X}p_x\mu_x(a_0).
\end{aligned}
\]
The same decomposition applies under \(P'\). Since \(a_0=1/2\in[0,1]\), the identified stratum probabilities and mean functions yield
\[
\theta(P')
=\sum_{x\in\mathcal X}p'_x\mu'_x(a_0)
=\sum_{x\in\mathcal X}p_x\mu_x(a_0)
=\theta(P).
\]
% lean: CausalSmith.Stat.NoisydoseWeakdesignTransition.causalTarget_eq_stratum_sum
% lean: CausalSmith.Stat.NoisydoseWeakdesignTransition.identification
\end{enumerate}
\end{proof}

\begin{proof}[Proof of \cref{obj:thm:uniform-frontier}]
Fix public constants \(K=(p_{\min},c_{\mathrm{lo}},c_{\mathrm{hi}})\), \(\beta\in(0,1]\), \(\kappa\in[0,2]\), and \(\alpha\in(0,1/2)\), under the envelope calibration in the statement. Throughout, use the rates and comparison scales defined there. In particular,
\[
d=2\beta+\kappa+1>1>0.
\]

\begin{enumerate}
\item \textbf{Upper risk bound.}
By \cref{obj:thm:observable-upper}, there are \(C_R>0\) and \(n_R\in\mathbb N\), independent of the public path space and of \(\sigma\), such that
\[
\sup_{P\in\mathcal P_{K,\beta,\kappa,\sigma}}
E_{Q_P^n}|\widehat\theta-\theta(P)|
\le C_R\rho_{n,\sigma}
\qquad
(n\ge n_R,\ 0\le\sigma\le1/4).
\]
The measurable estimator in \cref{obj:def:total-estimator} is admissible in \cref{obj:def:minimax-risk}. Viewed as a randomized estimator, it ignores the independent uniform seed, so its risk under the product experiment equals its observed-sample risk. Consequently,
\[
R_n(\sigma)
\le
\sup_{P\in\mathcal P_{K,\beta,\kappa,\sigma}}
E_{Q_P^n}|\widehat\theta-\theta(P)|
\le C_R\rho_{n,\sigma}.
\]
% lean: CausalSmith.Stat.NoisydoseWeakdesignTransition.minimaxRisk_frontier_upper

\item \textbf{Upper length bound.}
By \cref{obj:thm:finite-honesty}, the measurable connected interval in \cref{obj:def:honest-interval} satisfies
\[
Q_P^n\{\theta(P)\in C_{n,\sigma}\}\ge1-\alpha
\]
for every model law, every \(n\), and every allowed \(\sigma\). The same result supplies \(C_H>0\) and \(n_H\in\mathbb N\), uniform over the path spaces and error scales, such that
\[
\sup_{P\in\mathcal P_{K,\beta,\kappa,\sigma}}
E_{Q_P^n}\operatorname{length}(C_{n,\sigma})
\le C_H\rho_{n,\sigma}
\qquad(n\ge n_H).
\]
Thus this procedure belongs to the collection defining \cref{obj:def:honest-length}, and
\[
H_n(\sigma)
\le
\sup_{P\in\mathcal P_{K,\beta,\kappa,\sigma}}
E_{Q_P^n}\operatorname{length}(C_{n,\sigma})
\le C_H\rho_{n,\sigma}.
\]
% lean: CausalSmith.Stat.NoisydoseWeakdesignTransition.honestLength_frontier_upper

\item \textbf{A common lower bound.}
Use the published Jacobi localization bound
\citep[Theorem 2.4, equation (2.14)]{PetrushevXu2005},
the local kernel Lipschitz bound
\citep[Theorem 2.2, equation (2.4)]{https://people.math.sc.edu/pencho/Publications/kpx-06-28-06-web.pdf},
Jacobi orthogonality and the positive squared-norm formula
\citep[Table 18.3.1, Jacobi row and caption]{https://dlmf.nist.gov/18.3},
and the positive-real-axis gamma-ratio asymptotic
\citep[equation (5.11.12)]{https://dlmf.nist.gov/5.11.E12}
as the analytic inputs to \cref{obj:thm:observed-lower}.
The packet parameters are admissible: for \(\kappa>0\), they are
\(r=4\) and \(b=(\kappa-1)/2>-1/2\), while the \(\kappa=0\) construction uses the symmetric parameters \((4,4)\).

Set
\[
\tau:=\min\left\{\frac15,\frac{1-2\alpha}{2}\right\}.
\]
Since \(\alpha<1/2\), this choice gives
\[
0<\tau,\qquad \tau\le\frac15,\qquad
\tau\le1-2\alpha-\tau.
\]
That result supplies \(c_{\mathrm w}>0\) and \(n_{\mathrm w}\in\mathbb N\), uniform over the path spaces, such that for every \(n\ge n_{\mathrm w}\) and every allowed \(\sigma\), there are model laws \(P_+,P_-\) with
\[
\theta_\pm:=\theta(P_\pm),\qquad
\delta_{n,\sigma}:=c_{\mathrm w}\rho_{n,\sigma},
\]
\[
|\theta_+-\theta_-|\ge\delta_{n,\sigma},
\qquad
\operatorname{TV}(Q_{P_+}^n,Q_{P_-}^n)\le\tau.
\]
For \(n\ge1\), \(L_n\ge1\); each branch defining \(b_{n,\sigma}\) is nonnegative when \(\sigma\ge0\). Hence
\[
\rho_{n,\sigma}\ge0,\qquad \delta_{n,\sigma}\ge0.
\]

To handle the permitted randomization, define the seed law and product laws by
\[
\nu_{\mathrm{seed}}:=\operatorname{Unif}[0,1],
\qquad
Q_\pm^{n,\mathrm{rand}}:=Q_{P_\pm}^n\otimes\nu_{\mathrm{seed}}.
\]
For any measurable event \(\mathcal V\) in the sample--seed product space, define its section probability by
\[
f_{\mathcal V}(z)
:=
\nu_{\mathrm{seed}}\{u_{\mathrm{seed}}\in[0,1]:
(z,u_{\mathrm{seed}})\in\mathcal V\}.
\]
This is measurable and lies in \([0,1]\). Product integration and the layer-cake identity therefore give
\[
\begin{aligned}
\bigl|Q_+^{n,\mathrm{rand}}(\mathcal V)
      -Q_-^{n,\mathrm{rand}}(\mathcal V)\bigr|
&=
\left|\int f_{\mathcal V}\,dQ_{P_+}^n
      -\int f_{\mathcal V}\,dQ_{P_-}^n\right|\\
&=
\left|\int_0^1
\left[
Q_{P_+}^n\{f_{\mathcal V}\ge v\}
-
Q_{P_-}^n\{f_{\mathcal V}\ge v\}
\right]\,dv\right|\\
&\le \operatorname{TV}(Q_{P_+}^n,Q_{P_-}^n).
\end{aligned}
\]
Taking the supremum over such events yields
\[
\operatorname{TV}(Q_+^{n,\mathrm{rand}},Q_-^{n,\mathrm{rand}})
\le\tau\le\frac15.
\]

For any measurable real-valued estimator \(T\) on this product experiment, define
\[
\mathcal A_\pm
:=
\left\{(z,u_{\mathrm{seed}}):
|T(z,u_{\mathrm{seed}})-\theta_\pm|
\ge\frac{\delta_{n,\sigma}}2\right\}.
\]
The separation and triangle inequality imply
\(\mathcal A_+^c\subseteq\mathcal A_-\). Thus the total-variation bound gives
\[
\begin{aligned}
Q_+^{n,\mathrm{rand}}(\mathcal A_+)
+Q_-^{n,\mathrm{rand}}(\mathcal A_-)
&\ge
Q_+^{n,\mathrm{rand}}(\mathcal A_+)
+Q_-^{n,\mathrm{rand}}(\mathcal A_+^c)\\
&\ge
1-\operatorname{TV}(Q_+^{n,\mathrm{rand}},Q_-^{n,\mathrm{rand}})
\ge\frac45,
\end{aligned}
\]
and consequently
\[
\max\left\{
Q_+^{n,\mathrm{rand}}(\mathcal A_+),
Q_-^{n,\mathrm{rand}}(\mathcal A_-)
\right\}\ge\frac25.
\]
Pointwise domination by the threshold times its indicator now implies
\[
\begin{aligned}
\max_{\pm}E_{Q_\pm^{n,\mathrm{rand}}}|T-\theta_\pm|
&\ge
\frac{\delta_{n,\sigma}}2
\max_{\pm}Q_\pm^{n,\mathrm{rand}}(\mathcal A_\pm)\\
&\ge\frac{\delta_{n,\sigma}}5.
\end{aligned}
\]
These are inequalities between nonnegative integrals, so they also apply to infinite risks. Since both witnesses belong to the model class, taking the infimum over estimators proves
\[
R_n(\sigma)\ge\frac{\delta_{n,\sigma}}5
=\frac{c_{\mathrm w}}5\rho_{n,\sigma}.
\]

Next, let \(C^\circ\) be any measurable connected-interval procedure honest over the model class, and define its coverage events by
\[
\mathcal G_\pm:=\{z:\theta_\pm\in C^\circ(z)\}.
\]
Honesty and total variation imply
\[
Q_{P_+}^n(\mathcal G_+)\ge1-\alpha,
\qquad
Q_{P_+}^n(\mathcal G_-)
\ge Q_{P_-}^n(\mathcal G_-)-\tau
\ge1-\alpha-\tau.
\]
The complementary-event union bound therefore yields
\[
Q_{P_+}^n(\mathcal G_+\cap\mathcal G_-)
\ge1-2\alpha-\tau.
\]
On this intersection, connectedness gives containment of the closed segment joining the targets:
\[
[\min\{\theta_+,\theta_-\},\max\{\theta_+,\theta_-\}]
\subseteq C^\circ(z).
\]
Its Lebesgue length is \(|\theta_+-\theta_-|\), so
\[
\operatorname{length}(C^\circ(z))
\ge
\delta_{n,\sigma}\mathbf1_{\mathcal G_+\cap\mathcal G_-}(z).
\]
Integrating gives
\[
\sup_{P\in\mathcal P_{K,\beta,\kappa,\sigma}}
E_{Q_P^n}\operatorname{length}(C^\circ)
\ge
E_{Q_{P_+}^n}\operatorname{length}(C^\circ)
\ge(1-2\alpha-\tau)\delta_{n,\sigma}.
\]
Taking the infimum over honest procedures proves the corresponding length bound. The argument includes \(\delta_{n,\sigma}=0\).

Define
\[
c_L:=c_{\mathrm w}\tau>0,
\qquad
n_L:=\max\{n_{\mathrm w},1\}.
\]
For \(n\ge n_L\), the inequalities \(\tau\le1/5\) and \(\tau\le1-2\alpha-\tau\), together with \(\delta_{n,\sigma}\ge0\), give
\[
c_L\rho_{n,\sigma}=\tau\delta_{n,\sigma}
\le\frac{\delta_{n,\sigma}}5\le R_n(\sigma),
\]
\[
c_L\rho_{n,\sigma}=\tau\delta_{n,\sigma}
\le(1-2\alpha-\tau)\delta_{n,\sigma}\le H_n(\sigma).
\]
We have obtained, uniformly over the path spaces and allowed error scales,
\[
c_L\rho_{n,\sigma}\le R_n(\sigma),
\qquad
c_L\rho_{n,\sigma}\le H_n(\sigma)
\qquad(n\ge n_L).
\]
% lean: CausalSmith.Stat.NoisydoseWeakdesignTransition.frontier_both_lower

\item \textbf{Comparison at the second threshold.}
Define the positive constant and logarithmic denominator by
\[
\lambda_*:=\log(e+1)>0,
\qquad
\Lambda_n:=\log(e+n b_n^d)\quad(n\ge1).
\]
Because \(L_n\to\infty\) and \(\log L_n/L_n\to0\), choose \(n_E\ge1\) such that, for every \(n\ge n_E\),
\[
L_n\ge4,
\qquad
\frac d2\frac{\log L_n}{L_n}<\frac12.
\]
Since \(b_n=L_n^{-1/2}\) and \(\log n=L_n-1\), logarithmic monotonicity gives
\[
\Lambda_n
\ge\log(n b_n^d)
=L_n-1-\frac d2\log L_n
>\frac{L_n}{2}-1
\ge\frac{L_n}{4}.
\]
Also \(0<b_n\le1\) and \(d\ge0\), whence
\[
e+n b_n^d\le e+n\le2en.
\]
Using \(\log2\le2\), we obtain
\[
\Lambda_n\le\log2+L_n\le4L_n.
\]
Taking square roots and multiplying by \(\sqrt{L_n}\) yields
\[
\frac{\sqrt{L_n}}2\le\sqrt{\Lambda_n}\le2\sqrt{L_n},
\qquad
\frac{L_n}{2}
\le\sqrt{L_n}\sqrt{\Lambda_n}
\le2L_n.
\]
The defining formulas consequently give
\[
F_n(b_n)=\frac1{\sqrt{L_n}\sqrt{\Lambda_n}},
\qquad
\frac1{2L_n}\le F_n(b_n)\le\frac2{L_n}.
\]
Furthermore, \(b_n^2L_n=1\), so
\[
P_n(b_n)=\frac{\lambda_*}{L_n}.
\]
Define
\[
C_E:=\max\left\{\frac2{\lambda_*},\,2\lambda_*\right\}>0.
\]
Then, for \(n\ge n_E\),
\[
F_n(b_n)
\le\frac2{L_n}
\le C_E\frac{\lambda_*}{L_n}
=C_EP_n(b_n),
\]
\[
P_n(b_n)
=\frac{\lambda_*}{L_n}
\le\frac{C_E}{2L_n}
\le C_EF_n(b_n).
\]
% lean: CausalSmith.Stat.NoisydoseWeakdesignTransition.secondElbow_comparison

\item \textbf{Uniform constants and both threshold comparisons.}
Set
\[
c:=c_L,
\qquad
C:=\max\left\{
C_R,C_H,\sqrt{\lambda_*},\frac1{\sqrt{\lambda_*}},
C_E,c_L+1
\right\},
\]
\[
n_0:=\max\{n_R,n_H,n_L,n_E\}.
\]
Thus \(0<c<C\) and \(n_0\ge1\). The second-threshold inequalities just proved remain valid with \(C\) in place of \(C_E\), because the comparison scales are nonnegative.

For \(n\ge n_0\), positivity of \(n\) and \(d\) gives the exact identities
\[
n a_n^d
=n(n^{-1/d})^d
=n n^{-1}=1,
\qquad
F_n(a_n)=\frac{a_n}{\sqrt{\lambda_*}}.
\]
The choices \(C\ge\sqrt{\lambda_*}\) and
\(C\ge1/\sqrt{\lambda_*}\) therefore imply
\[
D_n=a_n\le C F_n(a_n),
\qquad
F_n(a_n)\le C D_n.
\]
For every \(\sigma\in[0,1/4]\), nonnegativity of the branches in
\cref{obj:def:frontier-rate} gives \(\rho_{n,\sigma}\ge0\). Hence enlarging the upper-bound constants preserves their inequalities:
\[
c\rho_{n,\sigma}\le R_n(\sigma)
\le C_R\rho_{n,\sigma}\le C\rho_{n,\sigma},
\]
\[
c\rho_{n,\sigma}\le H_n(\sigma)
\le C_H\rho_{n,\sigma}\le C\rho_{n,\sigma}.
\]
Together with the second-threshold bounds, this establishes the first group of assertions with one \(C\) and one \(n_0\).
% lean: CausalSmith.Stat.NoisydoseWeakdesignTransition.uniform_frontier

\item \textbf{Comparison factors for transition windows.}
Fix a real \(\omega_{\mathrm{win}}\ge1\), and define
\[
M_{\mathrm F,\omega_{\mathrm{win}}}:=\omega_{\mathrm{win}}(1+d\log \omega_{\mathrm{win}}),
\qquad
M_{\mathrm P,\omega_{\mathrm{win}}}:=1+2\log \omega_{\mathrm{win}},
\]
\[
C_{\mathrm a}:=\max\left\{\sqrt{\lambda_*},
\frac1{\sqrt{\lambda_*}}\right\},
\qquad
C_K:=\max\{1,M_{\mathrm F,\omega_{\mathrm{win}}}C_{\mathrm a},
M_{\mathrm F,\omega_{\mathrm{win}}}C_EM_{\mathrm P,\omega_{\mathrm{win}}}\}.
\]
All these factors are positive, and \(C_K\ge1\).

We first derive the scale comparisons used throughout either window. For positive radii \(\varrho_1,\varrho_2\), a nonnegative exponent \(\gamma\), and a nonnegative weight \(\omega_{\log}\), suppose
\[
\varrho_1,\varrho_2>0,\qquad
\gamma,\omega_{\log}\ge0,\qquad
\varrho_1\le \omega_{\mathrm{win}}\varrho_2.
\]
Since \(\omega_{\mathrm{win}}^\gamma\ge1\) and
\(\varrho_1^\gamma\le \omega_{\mathrm{win}}^\gamma\varrho_2^\gamma\),
\[
e+\omega_{\log}\varrho_1^\gamma
\le \omega_{\mathrm{win}}^\gamma(e+\omega_{\log}\varrho_2^\gamma).
\]
Taking logarithms gives
\[
\log(e+\omega_{\log}\varrho_1^\gamma)
\le\gamma\log \omega_{\mathrm{win}}+
\log(e+\omega_{\log}\varrho_2^\gamma).
\]
The logarithm on the right is at least \(1\), while
\(\gamma\log \omega_{\mathrm{win}}\ge0\). Therefore,
\[
\log(e+\omega_{\log}\varrho_1^\gamma)
\le
(1+\gamma\log \omega_{\mathrm{win}})
\log(e+\omega_{\log}\varrho_2^\gamma).
\]
Since \(1+\gamma\log \omega_{\mathrm{win}}\ge1\), this also implies
\[
\log(e+\omega_{\log}\varrho_1^\gamma)
\le
(1+\gamma\log \omega_{\mathrm{win}})^2
\log(e+\omega_{\log}\varrho_2^\gamma),
\]
and taking nonnegative square roots yields
\[
\sqrt{\log(e+\omega_{\log}\varrho_1^\gamma)}
\le
(1+\gamma\log \omega_{\mathrm{win}})
\sqrt{\log(e+\omega_{\log}\varrho_2^\gamma)}.
\]

Now let \(n\ge1\), \(\varrho>0\), and
\(\sigma\in[\varrho/\omega_{\mathrm{win}},\omega_{\mathrm{win}}\varrho]\). Both
\(\sigma\le \omega_{\mathrm{win}}\varrho\) and \(\varrho\le \omega_{\mathrm{win}}\sigma\) hold.
Applying the square-root comparison in both directions with
\(\gamma=d\) and \(\omega_{\log}=n\) gives
\[
\begin{aligned}
F_n(\sigma)
&=
\frac{\sigma}{\varrho}
\frac{\sqrt{\log(e+n\varrho^d)}}
     {\sqrt{\log(e+n\sigma^d)}}F_n(\varrho)
\le M_{\mathrm F,\omega_{\mathrm{win}}}F_n(\varrho),\\
F_n(\varrho)
&=
\frac{\varrho}{\sigma}
\frac{\sqrt{\log(e+n\sigma^d)}}
     {\sqrt{\log(e+n\varrho^d)}}F_n(\sigma)
\le M_{\mathrm F,\omega_{\mathrm{win}}}F_n(\sigma).
\end{aligned}
\]
Applying the logarithmic comparison in both directions with
\(\gamma=2\) and \(\omega_{\log}=L_n>0\), and dividing by \(L_n\), gives
\[
P_n(\sigma)\le M_{\mathrm P,\omega_{\mathrm{win}}}P_n(\varrho),
\qquad
P_n(\varrho)\le M_{\mathrm P,\omega_{\mathrm{win}}}P_n(\sigma).
\]
Thus both scales vary by the stated fixed factors throughout any such radius window.
% lean: CausalSmith.Stat.NoisydoseWeakdesignTransition.fourierScale_window_comparison
% lean: CausalSmith.Stat.NoisydoseWeakdesignTransition.polynomialScale_window_comparison

\item \textbf{One constant and threshold for both windows.}
The positive exponent \(d\) implies
\[
a_n=n^{-1/d}\longrightarrow0,
\qquad
b_n=L_n^{-1/2}\longrightarrow0.
\]
Multiplication by the fixed \(\omega_{\mathrm{win}}\) preserves these limits. Choose
\(n_{\omega_{\mathrm{win}}}\in\mathbb N\) so that, for every \(n\ge n_{\omega_{\mathrm{win}}}\),
\[
\omega_{\mathrm{win}}a_n\le\frac14,\qquad \omega_{\mathrm{win}}b_n\le\frac14,
\]
and define
\[
n_K:=\max\{1,n_{\omega_{\mathrm{win}}},n_E\}.
\]
For \(n\ge n_K\), positivity of \(n,L_n,\omega_{\mathrm{win}}\) ensures
\[
0<\frac{a_n}{\omega_{\mathrm{win}}},\qquad
\omega_{\mathrm{win}}a_n\le\frac14,\qquad
0<\frac{b_n}{\omega_{\mathrm{win}}},\qquad
\omega_{\mathrm{win}}b_n\le\frac14.
\]

At the first threshold, the exact identity from the preceding argument gives
\[
D_n\le C_{\mathrm a}F_n(a_n),
\qquad
F_n(a_n)\le C_{\mathrm a}D_n.
\]
For \(\sigma\in[a_n/\omega_{\mathrm{win}},\omega_{\mathrm{win}}a_n]\), the window bounds therefore imply
\[
\begin{aligned}
D_n
&\le C_{\mathrm a}F_n(a_n)
\le C_{\mathrm a}M_{\mathrm F,\omega_{\mathrm{win}}}F_n(\sigma)
\le C_KF_n(\sigma),\\
F_n(\sigma)
&\le M_{\mathrm F,\omega_{\mathrm{win}}}F_n(a_n)
\le M_{\mathrm F,\omega_{\mathrm{win}}}C_{\mathrm a}D_n
\le C_KD_n.
\end{aligned}
\]
Since \(C_K>0\), these are precisely
\[
C_K^{-1}D_n\le F_n(\sigma)\le C_KD_n.
\]

For \(\sigma\in[b_n/\omega_{\mathrm{win}},\omega_{\mathrm{win}}b_n]\), combine both window comparisons with the second-threshold bounds:
\[
\begin{aligned}
P_n(\sigma)
&\le M_{\mathrm P,\omega_{\mathrm{win}}}P_n(b_n)
\le M_{\mathrm P,\omega_{\mathrm{win}}}C_EF_n(b_n)\\
&\le M_{\mathrm P,\omega_{\mathrm{win}}}C_EM_{\mathrm F,\omega_{\mathrm{win}}}F_n(\sigma)
\le C_KF_n(\sigma),\\
F_n(\sigma)
&\le M_{\mathrm F,\omega_{\mathrm{win}}}F_n(b_n)
\le M_{\mathrm F,\omega_{\mathrm{win}}}C_EP_n(b_n)\\
&\le M_{\mathrm F,\omega_{\mathrm{win}}}C_EM_{\mathrm P,\omega_{\mathrm{win}}}P_n(\sigma)
\le C_KP_n(\sigma).
\end{aligned}
\]
Dividing the first inequality by \(C_K\) establishes
\[
C_K^{-1}P_n(\sigma)\le F_n(\sigma)\le C_KP_n(\sigma).
\]
The choices of \(C_K,n_K\) precede the choice of \(\sigma\), and the same choices apply to both entire windows.
% lean: CausalSmith.Stat.NoisydoseWeakdesignTransition.frontier_transition_windows

\item \textbf{Fixed positive noise and the frontier rate.}
Fix \(\sigma\in(0,1/4]\). Since both thresholds tend to zero, eventually
\[
a_n<\sigma,\qquad b_n<\sigma.
\]
The two branch tests in \cref{obj:def:frontier-rate} then select
\[
b_{n,\sigma}=P_n(\sigma)
=\frac{\log(e+\sigma^2L_n)}{L_n}.
\]

For natural \(n\) with \(\log n\ge1\), define
\[
\zeta_n:=\frac{\log\log n}{\log n}.
\]
Both \(\log n\) and \(\log\log n\) tend to infinity. Hence eventually
\[
\log n\ge1,\qquad
\log\log n\ge0,\qquad
\log\log n\ge-2\log(\sigma^2),
\qquad
\log\log n\ge\log(e+2\sigma^2).
\]
Also \(L_n=1+\log n>0\). Monotonicity of the logarithm gives the lower numerator bound
\[
\begin{aligned}
\log(e+\sigma^2L_n)
&\ge\log(\sigma^2\log n)\\
&=\log(\sigma^2)+\log\log n
\ge\frac12\log\log n.
\end{aligned}
\]
For the upper bound, \(\log n\ge1\) implies
\[
e+\sigma^2L_n
=e+\sigma^2(1+\log n)
\le(e+2\sigma^2)\log n.
\]
Thus
\[
\log(e+\sigma^2L_n)
\le\log(e+2\sigma^2)+\log\log n
\le2\log\log n.
\]
We also have
\[
0\le\zeta_n\le\log\log n.
\]
Consequently,
\[
\frac14\zeta_nL_n
=\frac14\log\log n+\frac14\zeta_n
\le\frac12\log\log n
\le\log(e+\sigma^2L_n),
\]
and
\[
\log(e+\sigma^2L_n)
\le2\log\log n
\le2\log\log n+2\zeta_n
=2\zeta_nL_n.
\]
Dividing by \(L_n>0\) proves
\[
0\le\zeta_n,\qquad
\frac14\zeta_n\le P_n(\sigma)\le2\zeta_n.
\]
Since \(\beta\ge0\), taking the \(\beta\)-th power preserves these comparisons. Together with the selected branch, this gives
\[
0\le\zeta_n^\beta,\qquad
4^{-\beta}\zeta_n^\beta
\le\rho_{n,\sigma}
\le2^\beta\zeta_n^\beta.
\]
Choose \(n_F\in\mathbb N\) so that all these assertions hold for every \(n\ge n_F\).
% lean: CausalSmith.Stat.NoisydoseWeakdesignTransition.fixedNoise_frontierRate_bounds

\item \textbf{Fixed-noise risk and length constants.}
Define
\[
c_\sigma:=c_L4^{-\beta},
\qquad
C_\sigma:=
\max\left\{\max\{C_R,C_H\}2^\beta,c_\sigma\right\}+1,
\]
\[
n_\sigma:=\max\{n_F,n_L,n_R,n_H\}.
\]
Then \(0<c_\sigma<C_\sigma\). For every public path space and every \(n\ge n_\sigma\), the lower-rate comparison and the common lower bound yield
\[
c_\sigma\zeta_n^\beta
\le c_L\rho_{n,\sigma}\le R_n(\sigma),
\qquad
c_\sigma\zeta_n^\beta
\le c_L\rho_{n,\sigma}\le H_n(\sigma).
\]
The upper-rate comparison and the two upper bounds yield
\[
R_n(\sigma)
\le C_R\rho_{n,\sigma}
\le C_R2^\beta\zeta_n^\beta
\le C_\sigma\zeta_n^\beta,
\]
\[
H_n(\sigma)
\le C_H\rho_{n,\sigma}
\le C_H2^\beta\zeta_n^\beta
\le C_\sigma\zeta_n^\beta.
\]
Substituting the definition of \(\zeta_n\) proves the fixed-positive-noise assertions. Every constant and sample threshold was chosen independently of the public path space, as required.
% lean: CausalSmith.Stat.NoisydoseWeakdesignTransition.uniform_frontier
\end{enumerate}
\end{proof}

\begin{proof}[Proof of \cref{obj:prop:error-free-reduction}]
\begin{enumerate}
\item Apply \cref{obj:thm:uniform-frontier}, using the filtered Jacobi localization estimate \citep[Theorem 2.4, equation (2.14)]{PetrushevXu2005}, the local Lipschitz estimate \citep[Theorem 2.2, equation (2.4)]{https://people.math.sc.edu/pencho/Publications/kpx-06-28-06-web.pdf}, Jacobi orthogonality and the squared-norm formula \citep[Table 18.3.1, Jacobi row and caption]{https://dlmf.nist.gov/18.3}, and the gamma-ratio asymptotic \citep[equation (5.11.12)]{https://dlmf.nist.gov/5.11.E12}. It supplies constants and a threshold satisfying
\[
0<c_{\mathrm F}<C_{\mathrm F},
\qquad n_{\mathrm F}\in\mathbb N,
\]
such that, for every public potential-path space, every \(n\ge n_{\mathrm F}\), and every \(\sigma\in[0,1/4]\),
\[
c_{\mathrm F}\rho_{n,\sigma}
\le R_n(\sigma)\le C_{\mathrm F}\rho_{n,\sigma},
\qquad
c_{\mathrm F}\rho_{n,\sigma}
\le H_n(\sigma)\le C_{\mathrm F}\rho_{n,\sigma},
\]
with the prescribed noncoverage probability \(\alpha\in(0,1/2)\). These constants and the threshold are uniform over the public potential-path space.
% lean: CausalSmith.Stat.NoisydoseWeakdesignTransition.uniform_frontier

\item Use the effective exponent from \cref{obj:def:frontier-rate},
\[
d=2\beta+\kappa+1.
\]
Since \(0\le n^{-1/d}\), the first branch of the localization scale at zero error gives
\[
b_{n,0}=n^{-1/d},
\qquad
\rho_{n,0}
=(n^{-1/d})^\beta
=n^{-\beta/d}.
\]
The last identity follows by multiplying the exponents. Substituting \(\sigma=0\in[0,1/4]\) into the bounds above yields
\[
c_{\mathrm F}n^{-\beta/d}
\le R_n(0)\le C_{\mathrm F}n^{-\beta/d},
\qquad
c_{\mathrm F}n^{-\beta/d}
\le H_n(0)\le C_{\mathrm F}n^{-\beta/d}
\]
for every public potential-path space and every \(n\ge n_{\mathrm F}\). This proves the first assertion.
% lean: CausalSmith.Stat.NoisydoseWeakdesignTransition.frontierRate_zero

\item Fix the benchmark parameters, evaluation point, neighborhood, and design distribution in the statement. Define its minimax absolute risk by
\[
\mathfrak R_{\mathrm G,n}
=
\inf_{\widehat f}\sup_f
E_f\bigl|\widehat f(x_*)-f(x_*)\bigr|,
\]
where the functions and measurable estimators range over the benchmark class and procedures specified there. The exponent restrictions imply
\[
0<\beta\le1,
\qquad 0\le\kappa\le2.
\]
Thus Gaiffas's direct Gaussian-response benchmark applies with smoothness \(\beta\), design exponent \(\kappa\), and loss parameter one \citep[Definition 2 and Theorem 1, equations (2.1), (2.4), and (2.5), pp.~3--4]{Gaiffas2005}. It supplies constants and a threshold satisfying
\[
0<c_{\mathrm B}<C_{\mathrm B},
\qquad n_{\mathrm B}\in\mathbb N,
\]
and, for every \(n\ge n_{\mathrm B}\),
\[
c_{\mathrm B}n^{-\beta/d}
\le \mathfrak R_{\mathrm G,n}
\le C_{\mathrm B}n^{-\beta/d}.
\]
This invocation uses the stated local polynomial-approximation class, whose polynomial degree is at most \(\lfloor\beta\rfloor\), including degree at most one when \(\beta=1\). The benchmark constants and threshold depend on the fixed benchmark data.
% lean: CausalSmith.Stat.NoisydoseWeakdesignTransition.GaiffasDirectBenchmark

\item Define the comparison constants and threshold by
\[
c_{\mathrm{cmp}}=\frac{c_{\mathrm F}}{C_{\mathrm B}},
\qquad
C_{\mathrm{cmp}}
=\frac{C_{\mathrm F}}{c_{\mathrm B}}+c_{\mathrm{cmp}}+1,
\qquad
n_{\mathrm{cmp}}=\max\{n_{\mathrm F},n_{\mathrm B}\}.
\]
Positivity of the preceding constants gives
\[
c_{\mathrm{cmp}}>0,
\qquad
C_{\mathrm{cmp}}-c_{\mathrm{cmp}}
=\frac{C_{\mathrm F}}{c_{\mathrm B}}+1>0.
\]
Moreover,
\[
c_{\mathrm{cmp}}C_{\mathrm B}=c_{\mathrm F},
\qquad
C_{\mathrm{cmp}}c_{\mathrm B}
=C_{\mathrm F}+(c_{\mathrm{cmp}}+1)c_{\mathrm B}
\ge C_{\mathrm F}.
\]
For every \(n\ge n_{\mathrm{cmp}}\), both preceding risk bounds apply. Multiplying the benchmark upper bound by \(c_{\mathrm{cmp}}\) gives
\[
c_{\mathrm{cmp}}\mathfrak R_{\mathrm G,n}
\le c_{\mathrm{cmp}}C_{\mathrm B}n^{-\beta/d}
=c_{\mathrm F}n^{-\beta/d}
\le R_n(0).
\]
Conversely, since \(n^{-\beta/d}\ge0\),
\[
R_n(0)
\le C_{\mathrm F}n^{-\beta/d}
\le C_{\mathrm{cmp}}c_{\mathrm B}n^{-\beta/d}
\le C_{\mathrm{cmp}}\mathfrak R_{\mathrm G,n}.
\]
The comparison constants and threshold are uniform over the public potential-path space: the frontier constants have that uniformity, and the benchmark constants depend on the fixed benchmark data. This proves the second assertion.
% lean: CausalSmith.Stat.NoisydoseWeakdesignTransition.error_free_reduction
\end{enumerate}
\end{proof}

\begin{proof}[Proof of \cref{obj:thm:observable-upper}]
\begin{enumerate}
\item Fix public class constants \(K\), \(\beta\in(0,1]\), and \(\kappa\in[0,2]\). The Fourier and polynomial comparison-score conclusions of \cref{obj:lem:dictionary-score-rate} supply positive constants \(C_{\mathrm F},C_{\mathrm M}\) and integer thresholds \(n_{\mathrm F},n_{\mathrm M}\), depending on \(K,\beta,\kappa\), such that
\[
\begin{aligned}
n\ge n_{\mathrm F},\quad \sigma\le L_n^{-1/2}
&\ \Longrightarrow\
\exists k\ge0:\ 
\text{the Fourier pair with index }k\text{ belongs to }\mathcal D_{n,\sigma},
\quad A_q(K)\le C_{\mathrm F}\rho_{n,\sigma},\\
n\ge n_{\mathrm M},\quad L_n^{-1/2}<\sigma
&\ \Longrightarrow\
\exists j\ge0:\ 
\text{the polynomial pair with index }2^j+1\text{ belongs to }\mathcal D_{n,\sigma},
\quad A_q(K)\le C_{\mathrm M}\rho_{n,\sigma},
\end{aligned}
\]
for every \(\sigma\in[0,1/4]\). These inequalities concern public scores and rates determined by \(\beta,\kappa,n,\sigma\), so the constants and thresholds can be fixed uniformly over the potential-path spaces. Set
\[
C=\max\{C_{\mathrm F},C_{\mathrm M}\},
\qquad
n_0=\max\{3,n_{\mathrm F},n_{\mathrm M}\}.
\]
Then \(C>0\). % lean: CausalSmith.Stat.NoisydoseWeakdesignTransition.fourier_dictionary_witness; CausalSmith.Stat.NoisydoseWeakdesignTransition.inverseheat_dictionary_witness

\item Fix a potential-path space, an integer \(n\ge n_0\), and \(\sigma\in[0,1/4]\). In particular,
\[
n\ge3,\qquad n\ge n_{\mathrm F},\qquad n\ge n_{\mathrm M},
\qquad L_n=\log(en)\ge0.
\]
The three expressions defining \(b_{n,\sigma}\) in \cref{obj:def:frontier-rate} are nonnegative: the direct expression is a power of a positive sample size, the Fourier expression has nonnegative numerator and positive denominator, and the polynomial expression has numerator \(\log(e+\sigma^2L_n)\ge0\) and denominator \(L_n>0\). Consequently,
\[
b_{n,\sigma}\ge0,
\qquad
\rho_{n,\sigma}=b_{n,\sigma}^{\beta}\ge0.
\]
This permits enlargement of either comparison constant to \(C\). % lean: CausalSmith.Stat.NoisydoseWeakdesignTransition.observable_upper

\item Suppose \(\sigma\le L_n^{-1/2}\). Choose the Fourier comparison pair from the first witness bound. Its dictionary membership and \(n>0\) allow application of the admissibility conclusion of \cref{obj:lem:dictionary-score-rate}. Thus its functions are measurable and satisfy
\[
\begin{gathered}
q(t)\ge0\quad(t\in[-1/2,1/2]),
\qquad 0<I_\kappa(q)<\infty,\\
E_Z|\ell_q(t+\sigma Z)|<\infty,
\qquad E_Z[\ell_q(t+\sigma Z)]=q(t)
\quad(t\in[-1/2,1/2]),
\qquad V_q<\infty.
\end{gathered}
\]
Moreover,
\[
A_q(K)\le C_{\mathrm F}\rho_{n,\sigma}
\le C\rho_{n,\sigma},
\]
where the last inequality follows from \(C_{\mathrm F}\le C\) and \(\rho_{n,\sigma}\ge0\). This proves the Fourier comparison guarantee, including equality at the cutoff. % lean: CausalSmith.Stat.NoisydoseWeakdesignTransition.observable_upper.hFourier

\item Suppose \(L_n^{-1/2}<\sigma\). Choose the polynomial comparison pair from the second witness bound, with
\[
m=2^j+1,
\qquad
(q,\ell_q)=(q_m^{\mathrm M},\ell_m^{\mathrm M}).
\]
Its dictionary membership again gives, by \cref{obj:lem:dictionary-score-rate}, measurability and
\[
\begin{gathered}
q(t)\ge0\quad(t\in[-1/2,1/2]),
\qquad 0<I_\kappa(q)<\infty,\\
E_Z|\ell_q(t+\sigma Z)|<\infty,
\qquad E_Z[\ell_q(t+\sigma Z)]=q(t)
\quad(t\in[-1/2,1/2]),
\qquad V_q<\infty.
\end{gathered}
\]
Its score satisfies
\[
A_q(K)\le C_{\mathrm M}\rho_{n,\sigma}
\le C\rho_{n,\sigma}.
\]
This proves the polynomial comparison guarantee. % lean: CausalSmith.Stat.NoisydoseWeakdesignTransition.observable_upper.hPoly

\item Let \(\widehat q\) be the selected latent weight in \cref{obj:def:total-estimator}. The fixed-order minimization gives
\[
(\widehat q,\ell_{\widehat q})\in\mathcal D_{n,\sigma},
\qquad
A_{\widehat q}(K)\le A_q(K)
\quad\text{for every }(q,\ell_q)\in\mathcal D_{n,\sigma}.
\]
If \(\sigma\le L_n^{-1/2}\), compare with the Fourier witness; otherwise \(L_n^{-1/2}<\sigma\), and compare with the polynomial witness. In either case,
\[
A_{\widehat q}(K)\le A_q(K)\le C\rho_{n,\sigma}.
\]
% lean: CausalSmith.Stat.NoisydoseWeakdesignTransition.selectedTag_minimizes; CausalSmith.Stat.NoisydoseWeakdesignTransition.observable_upper.hscore

\item Fix \(P\in\mathcal P_{K,\beta,\kappa,\sigma}\) satisfying \cref{obj:ass:iid}. Applied to the selected dictionary pair, \cref{obj:lem:dictionary-score-rate} supplies all its weight-admissibility conditions, finite extended second moment, and the law-specific integrability
\[
E_P\!\left[|\ell_{\widehat q}(W-a_0)|\right]<\infty.
\]
Hence \cref{obj:lem:public-error-certificate} yields
\[
E_{Q_P^n}\!\left[
|\widehat\theta_{\widehat q}-\theta(P)|
\right]\le A_{\widehat q}(K).
\]
For each \(r\in\{0,1,2\}\), the sample index \(i=r\) exists because \(n\ge3\), and \(i\bmod3=r\). Therefore
\[
r\in\mathcal I_r,\qquad
\mathcal I_r\ne\varnothing\quad(r=0,1,2),
\qquad
\widehat\theta=\widehat\theta_{\widehat q}.
\]
Using the total-estimator formula in \cref{obj:def:total-estimator} and the selected-score bound,
\[
E_{Q_P^n}\!\left[|\widehat\theta-\theta(P)|\right]
=
E_{Q_P^n}\!\left[
|\widehat\theta_{\widehat q}-\theta(P)|
\right]
\le A_{\widehat q}(K)
\le C\rho_{n,\sigma}.
\]
The choices of \(C,n_0\) apply to every permitted space, law, and public error scale, as required. % lean: CausalSmith.Stat.NoisydoseWeakdesignTransition.public_error_certificate; CausalSmith.Stat.NoisydoseWeakdesignTransition.blocksNonempty_of_three; CausalSmith.Stat.NoisydoseWeakdesignTransition.observable_upper
\end{enumerate}
\end{proof}

\begin{proof}[Proof of \cref{obj:thm:finite-honesty}]
Fix the parameters, a public potential-outcome space, an allowed error scale \(\sigma\), and a structural law \(P\) satisfying the stated conditions.

\begin{enumerate}
\item The observed-record map is measurable, so its image law is a probability measure. Its \(n\)-fold product \(Q_P^n\), defined in \cref{obj:def:observed-experiment}, is therefore a probability measure, including when \(n=0\).

The finite partition by the two observed strata and \cref{obj:def:conditional-potential-mean} give
\[
\theta(P)
=\sum_{x\in\mathcal X}E_P[\mathbf1\{X=x\}Y(a_0)]
=\sum_{x\in\mathcal X}p_x\mu_x(a_0).
\]
The expectations exist by \cref{obj:ass:bounded-potential-outcomes}, and the conditional means are defined because \cref{obj:ass:stratum-overlap} supplies positive stratum probabilities. Since
\[
p_x\ge0,\qquad \sum_{x\in\mathcal X}p_x=1,
\]
the bounds in \cref{obj:ass:mean-range} yield
\[
0
=0\sum_{x\in\mathcal X}p_x
\le\theta(P)
\le1\sum_{x\in\mathcal X}p_x
=1.
\]
% lean: CausalSmith.Stat.NoisydoseWeakdesignTransition.experiment_isProbabilityMeasure
% lean: CausalSmith.Stat.NoisydoseWeakdesignTransition.causalTarget_mem_Icc

\item First suppose that all three sample blocks in \cref{obj:def:total-estimator} are nonempty. An index in \(\mathcal I_2\) has remainder two modulo three, so
\[
i\in\mathcal I_2
\quad\Longrightarrow\quad
2\le i<n.
\]
Consequently \(n\ge3\).

The selection rule in \cref{obj:def:total-estimator} gives
\[
\widehat q\in\mathcal D_{n,\sigma},
\qquad
A_{\widehat q}(K)\le A_q(K)
\quad\text{for every }q\in\mathcal D_{n,\sigma}.
\]
By the admissibility assertion of \cref{obj:lem:dictionary-score-rate}, the selected pair satisfies the measurability and pointwise inversion conditions, together with
\[
\widehat q(t)\ge0\quad(t\in[-1/2,1/2]),
\qquad
0<I_\kappa(\widehat q)<\infty,
\qquad
V_{\widehat q}^{(16)}<\infty,
\qquad
E_P|\ell_{\widehat q}(V)|<\infty.
\]
Its public denominator and bias bounds have the formulas
\[
b_{\widehat q}(K)
=p_{\min}c_{\mathrm{lo}}I_\kappa(\widehat q)>0,
\qquad
B_{\widehat q}(K)
=\frac{c_{\mathrm{hi}}}{c_{\mathrm{lo}}I_\kappa(\widehat q)}
  \int_{-1/2}^{1/2}|t|^{\kappa+\beta}\widehat q(t)\,dt
\ge0.
\]
Here the bias bound is nonnegative because its integrand is nonnegative. Thus the score formula in \cref{obj:synth_4} gives
\[
A_{\widehat q}(K)
=
B_{\widehat q}(K)
+\frac{12}{b_{\widehat q}(K)}\sqrt{\frac{V_{\widehat q}(K)}{n}}
+2\sqrt{\frac3n}
\ge2\sqrt{\frac3n}>0.
\]
The product experiment satisfies the sampling condition in \cref{obj:ass:iid}. Applying \cref{obj:lem:public-error-certificate} to the selected pair, and using
\(\widehat\theta=\widehat\theta_{\widehat q}\) on nonempty blocks, therefore gives
\[
E_{Q_P^n}\!\left[|\widehat\theta-\theta(P)|\right]
\le A_{\widehat q}(K).
\]
% lean: CausalSmith.Stat.NoisydoseWeakdesignTransition.three_le_of_blocksNonempty
% lean: CausalSmith.Stat.NoisydoseWeakdesignTransition.selectedTag_minimizes
% lean: CausalSmith.Stat.NoisydoseWeakdesignTransition.Aq_pos
% lean: CausalSmith.Stat.NoisydoseWeakdesignTransition.public_error_certificate

\item The stratum estimates are clipped to \([0,1]\). The empirical stratum frequencies are nonnegative and sum to one, since the two strata partition the nonempty frequency block. Hence
\[
\widehat p_x\ge0,\qquad
\sum_{x\in\mathcal X}\widehat p_x=1,\qquad
0\le\widehat\mu_x(\widehat q)\le1,
\]
and
\[
0
\le
\widehat\theta
=\sum_{x\in\mathcal X}\widehat p_x\widehat\mu_x(\widehat q)
\le1.
\]
Together with the target range, this implies
\[
0\le|\widehat\theta-\theta(P)|\le1.
\]
The error is measurable and integrable under the probability law \(Q_P^n\).

Define the measurable event
\[
\mathcal G_{\mathrm{cover}}
=
\left\{\text{data}:
|\widehat\theta-\theta(P)|\le A_{\widehat q}(K)/\alpha\right\}.
\]
Markov's inequality and the certificate give
\[
A_{\widehat q}(K)/\alpha\,
Q_P^n\!\left\{
A_{\widehat q}(K)/\alpha\le|\widehat\theta-\theta(P)|
\right\}
\le
E_{Q_P^n}|\widehat\theta-\theta(P)|
\le A_{\widehat q}(K).
\]
Because \(A_{\widehat q}(K)>0\) and \(\alpha>0\), and the complement of
\(\mathcal G_{\mathrm{cover}}\) is contained in the displayed upper-tail event,
\[
Q_P^n(\mathcal G_{\mathrm{cover}}^c)
\le
Q_P^n\!\left\{
A_{\widehat q}(K)/\alpha\le|\widehat\theta-\theta(P)|
\right\}
\le\alpha,
\qquad
Q_P^n(\mathcal G_{\mathrm{cover}})\ge1-\alpha.
\]
On \(\mathcal G_{\mathrm{cover}}\),
\[
\widehat\theta-A_{\widehat q}(K)/\alpha
\le\theta(P)
\le\widehat\theta+A_{\widehat q}(K)/\alpha.
\]
Combining these inequalities with the target range gives
\[
\max\!\left\{0,\widehat\theta-A_{\widehat q}(K)/\alpha\right\}
\le\theta(P)
\le
\min\!\left\{1,\widehat\theta+A_{\widehat q}(K)/\alpha\right\}.
\]
Thus \cref{obj:def:honest-interval} implies
\[
Q_P^n\{\theta(P)\in C_{n,\sigma}\}
\ge Q_P^n(\mathcal G_{\mathrm{cover}})
\ge1-\alpha.
\]
% lean: CausalSmith.Stat.NoisydoseWeakdesignTransition.totalEstimator_mem_Icc
% lean: mul_meas_ge_le_integral_of_nonneg
% lean: CausalSmith.Stat.NoisydoseWeakdesignTransition.mem_honestInterval_of_abs_le

\item If some sample block is empty, \cref{obj:def:honest-interval} returns the full target range. Therefore
\[
C_{n,\sigma}=[0,1],
\qquad
Q_P^n\{\theta(P)\in C_{n,\sigma}\}=1\ge1-\alpha.
\]
This case includes \(n=0\), and completes the coverage assertion for every integer \(n\ge0\).
% lean: CausalSmith.Stat.NoisydoseWeakdesignTransition.finite_honesty

\item For the length assertion, choose a constant \(C_{\mathrm{up}}>0\) and a threshold \(n_{\mathrm{up}}\in\mathbb N\) supplied by \cref{obj:thm:observable-upper}. Set
\[
C=\frac{2}{\alpha}C_{\mathrm{up}},
\qquad
n_0=\max\{3,n_{\mathrm{up}}\}.
\]
Fix \(n\ge n_0\). For each block index \(r\in\{0,1,2\}\),
\[
r<n,\qquad r\bmod3=r,\qquad r\in\mathcal I_r,
\]
so all blocks are nonempty.

Write
\[
L_n=\log(en),
\]
as in \cref{obj:def:frontier-rate}. If \(\sigma\le L_n^{-1/2}\), the Fourier comparison in \cref{obj:thm:observable-upper} supplies a dictionary pair whose score satisfies
\[
A_q(K)\le C_{\mathrm{up}}\rho_{n,\sigma},
\qquad q\in\mathcal D_{n,\sigma}.
\]
The minimizing property consequently gives
\[
A_{\widehat q}(K)\le A_q(K)\le C_{\mathrm{up}}\rho_{n,\sigma}.
\]
If \(L_n^{-1/2}<\sigma\), the polynomial comparison supplies a dictionary pair with
\[
A_q(K)\le C_{\mathrm{up}}\rho_{n,\sigma},
\qquad q\in\mathcal D_{n,\sigma},
\]
and again
\[
A_{\widehat q}(K)\le A_q(K)\le C_{\mathrm{up}}\rho_{n,\sigma}.
\]
These two cases exhaust the allowed error scales. Dictionary admissibility and the score formula also give
\[
A_{\widehat q}(K)\ge2\sqrt{\frac3n}>0.
\]
% lean: CausalSmith.Stat.NoisydoseWeakdesignTransition.observable_upper
% lean: CausalSmith.Stat.NoisydoseWeakdesignTransition.blocksNonempty_of_three
% lean: CausalSmith.Stat.NoisydoseWeakdesignTransition.selectedTag_minimizes
% lean: CausalSmith.Stat.NoisydoseWeakdesignTransition.Aq_pos

\item Define the interval endpoints by
\[
c_{\mathrm{left}}
=\max\!\left\{0,\widehat\theta-\frac{A_{\widehat q}(K)}{\alpha}\right\},
\qquad
c_{\mathrm{right}}
=\min\!\left\{1,\widehat\theta+\frac{A_{\widehat q}(K)}{\alpha}\right\}.
\]
They satisfy
\[
c_{\mathrm{left}}\ge\widehat\theta-\frac{A_{\widehat q}(K)}{\alpha},
\qquad
c_{\mathrm{right}}\le\widehat\theta+\frac{A_{\widehat q}(K)}{\alpha},
\qquad
c_{\mathrm{right}}-c_{\mathrm{left}}\le\frac{2A_{\widehat q}(K)}{\alpha}.
\]
The length formula for a closed interval, together with
\(A_{\widehat q}(K)\ge0\), therefore yields the pointwise bound
\[
\operatorname{length}(C_{n,\sigma})
=\max\{c_{\mathrm{right}}-c_{\mathrm{left}},0\}
\le\frac{2A_{\widehat q}(K)}{\alpha}.
\]
Since the selected score is deterministic and \(Q_P^n\) is a probability measure,
\[
\begin{aligned}
E_{Q_P^n}\!\left[\operatorname{length}(C_{n,\sigma})\right]
&\le E_{Q_P^n}\!\left[\frac{2A_{\widehat q}(K)}{\alpha}\right]\\
&=\frac{2A_{\widehat q}(K)}{\alpha}\\
&\le\frac{2C_{\mathrm{up}}}{\alpha}\rho_{n,\sigma}
=C\rho_{n,\sigma}.
\end{aligned}
\]
The chosen constant and threshold depend only on \(K,\beta,\kappa\), and \(\alpha\), and the argument applies simultaneously to every permitted schedule space, error scale, and model law.
% lean: CausalSmith.Stat.NoisydoseWeakdesignTransition.honestInterval_volume_le
% lean: CausalSmith.Stat.NoisydoseWeakdesignTransition.finite_honesty
\end{enumerate}
\end{proof}

\begin{proof}[Proof of \cref{obj:thm:weighted-packet-construction}]
Fix \(\kappa\in[0,2]\). We use the filter, taper parameter, kernels, and packet from \cref{obj:def:packet-handle}. All constants below are independent of the integer \(m\ge4\).

\begin{enumerate}
\item \emph{Filter properties and cited kernel estimates.}
The standard smooth function
\[
\vartheta(v)=
\begin{cases}
e^{-1/v},&v>0,\\
0,&v\le0
\end{cases}
\qquad(v\in\mathbb R)
\]
gives
\[
\eta(s)=\vartheta\bigl((s-9/8)(15/8-s)\bigr).
\]
Consequently, \(\eta\) is smooth, \(0\le\eta\le1\), and
\[
\eta(s)\ne0\ \Longleftrightarrow\ 9/8<s<15/8.
\]
In particular,
\[
\eta(j/m)\ne0\quad\Longrightarrow\quad m<j<2m.
\]
On the central spectral interval,
\[
(s-9/8)(15/8-s)\ge\frac1{64}
\quad(5/4\le s\le7/4),
\]
so
\[
\eta(s)\ge c_\eta,\qquad c_\eta:=e^{-64}>0.
\]

For Jacobi parameters \(a_{\mathrm J},b_{\mathrm J}>-1/2\), define
\[
\mathcal K_m^{a_{\mathrm J},b_{\mathrm J}}(z_1,z_2)
:=
\sum_{j=0}^{2m}\eta(j/m)
\frac{\mathsf J_j^{(a_{\mathrm J},b_{\mathrm J})}(z_1)
      \mathsf J_j^{(a_{\mathrm J},b_{\mathrm J})}(z_2)}
     {H_j^{a_{\mathrm J},b_{\mathrm J}}}
\qquad(z_1,z_2\in\mathbb R),
\]
and, for \(z,z_1,z_2\in[-1,1]\), define
\[
\mathcal W_{a_{\mathrm J},b_{\mathrm J}}(m;z)
:=
(1-z+m^{-2})^{a_{\mathrm J}+1/2}
(1+z+m^{-2})^{b_{\mathrm J}+1/2},
\qquad
\delta_{\mathrm{ang}}(z_1,z_2)
:=
|\arccos z_1-\arccos z_2|.
\]
The spectral support permits truncation at degree \(2m\). Petrushev--Xu use squared norms normalized by the total Jacobi weight mass
\citep[pp.~1--2, equation~(1.2)]{PetrushevXu2005}. In the present notation, their norms and filtered kernel are
\[
\mathsf h_{\mathrm{PX},j}^{a_{\mathrm J},b_{\mathrm J}}
:=\frac{H_j^{a_{\mathrm J},b_{\mathrm J}}}
        {H_0^{a_{\mathrm J},b_{\mathrm J}}}
\qquad(j\ge0),
\]
\[
\begin{split}
\mathcal K_{\mathrm{PX},m}^{a_{\mathrm J},b_{\mathrm J}}(z_1,z_2)
&:=\sum_{j=0}^{2m}\eta(j/m)
\frac{\mathsf J_j^{(a_{\mathrm J},b_{\mathrm J})}(z_1)
      \mathsf J_j^{(a_{\mathrm J},b_{\mathrm J})}(z_2)}
     {\mathsf h_{\mathrm{PX},j}^{a_{\mathrm J},b_{\mathrm J}}}\\
&=H_0^{a_{\mathrm J},b_{\mathrm J}}
  \mathcal K_m^{a_{\mathrm J},b_{\mathrm J}}(z_1,z_2)
\qquad(z_1,z_2\in\mathbb R).
\end{split}
\]
Since \(H_0^{a_{\mathrm J},b_{\mathrm J}}>0\) depends only on the fixed Jacobi parameters, dividing their localization estimate by this factor absorbs its reciprocal into the localization constant below. The localization estimate of
\citep[Theorem~2.4, equation~(2.14)]{PetrushevXu2005}
provides, for every \(\tau>0\), a positive constant
\(C_{\mathrm{loc},\tau}^{a_{\mathrm J},b_{\mathrm J}}\) such that
\[
|\mathcal K_m^{a_{\mathrm J},b_{\mathrm J}}(z_1,z_2)|
\le
\frac{C_{\mathrm{loc},\tau}^{a_{\mathrm J},b_{\mathrm J}}m}
{\sqrt{\mathcal W_{a_{\mathrm J},b_{\mathrm J}}(m;z_1)
       \mathcal W_{a_{\mathrm J},b_{\mathrm J}}(m;z_2)}
 (1+m\delta_{\mathrm{ang}}(z_1,z_2))^\tau}.
\]
The local angular estimate of
\citep[Theorem~2.2, equation~(2.4)]{https://people.math.sc.edu/pencho/Publications/kpx-06-28-06-web.pdf},
with its fixed neighborhood constant equal to \(1\), provides a positive
\(C_{\mathrm{lip},\tau}^{a_{\mathrm J},b_{\mathrm J}}\) such that
\[
\begin{split}
&|\mathcal K_m^{a_{\mathrm J},b_{\mathrm J}}(z_1,z_2)
 -\mathcal K_m^{a_{\mathrm J},b_{\mathrm J}}(\xi,z_2)|\\
&\qquad\le
\frac{C_{\mathrm{lip},\tau}^{a_{\mathrm J},b_{\mathrm J}}m^2
      \delta_{\mathrm{ang}}(z_1,\xi)}
{\sqrt{\mathcal W_{a_{\mathrm J},b_{\mathrm J}}(m;z_2)
       \mathcal W_{a_{\mathrm J},b_{\mathrm J}}(m;\zeta)}
 (1+m\delta_{\mathrm{ang}}(z_2,\zeta))^\tau},
\end{split}
\]
whenever \(z_1,z_2,\xi,\zeta\in[-1,1]\) and
\[
\delta_{\mathrm{ang}}(z_1,\xi)\le m^{-1},
\qquad
\delta_{\mathrm{ang}}(\zeta,\xi)\le m^{-1}.
\]
The parameter and filter hypotheses hold for both Jacobi systems used below.
% lean: CausalSmith.Stat.NoisydoseWeakdesignTransition.FilteredJacobiLipschitz

\item \emph{Norms and uniform gamma-ratio bounds.}
Jacobi orthogonality and the positive squared-norm formula are supplied by
\citep[Table~18.3.1, Jacobi row and caption]{https://dlmf.nist.gov/18.3}:
\[
H_j^{a_{\mathrm J},b_{\mathrm J}}
=
\frac{2^{a_{\mathrm J}+b_{\mathrm J}+1}}
     {2j+a_{\mathrm J}+b_{\mathrm J}+1}
\frac{\Gamma(j+a_{\mathrm J}+1)\Gamma(j+b_{\mathrm J}+1)}
     {\Gamma(j+1)\Gamma(j+a_{\mathrm J}+b_{\mathrm J}+1)}
>0.
\]
Orthogonality applies against every polynomial of degree less than \(j\).

For fixed positive shifts \(\lambda_1,\lambda_2\), the asymptotic in
\citep[Equation~5.11.12]{https://dlmf.nist.gov/5.11.E12}
gives an integer cutoff \(j_\Gamma\) such that
\[
\frac12
\le
\frac{\Gamma(j+\lambda_1)}
     {\Gamma(j+\lambda_2)j^{\lambda_1-\lambda_2}}
\le2
\qquad(j\ge j_\Gamma,\ j\ge1).
\]
All these normalized ratios are positive. Thus the constants
\[
\begin{split}
c_\Gamma(\lambda_1,\lambda_2)
&:=
\min\left(
\{1/2\}\cup
\left\{
\frac{\Gamma(j+\lambda_1)}
     {\Gamma(j+\lambda_2)j^{\lambda_1-\lambda_2}}
:1\le j<j_\Gamma
\right\}\right),\\
C_\Gamma(\lambda_1,\lambda_2)
&:=
\max\left(
\{2\}\cup
\left\{
\frac{\Gamma(j+\lambda_1)}
     {\Gamma(j+\lambda_2)j^{\lambda_1-\lambda_2}}
:1\le j<j_\Gamma
\right\}\right)
\end{split}
\]
are positive and satisfy
\[
c_\Gamma(\lambda_1,\lambda_2)j^{\lambda_1-\lambda_2}
\le
\frac{\Gamma(j+\lambda_1)}{\Gamma(j+\lambda_2)}
\le
C_\Gamma(\lambda_1,\lambda_2)j^{\lambda_1-\lambda_2}
\qquad(j\ge1).
\]
The singleton sets in these definitions also cover an empty initial segment.
% lean: CausalSmith.Stat.NoisydoseWeakdesignTransition.gamma_ratio_integer_power_bounds

\item \emph{Endpoint normalization for \(\kappa>0\).}
Suppose first that \(\kappa>0\), and set
\[
b:=\frac{\kappa-1}{2}>-\frac12.
\]
The standard endpoint evaluation, followed by the gamma recurrence, gives
\[
\mathsf J_j^{(4,b)}(-1)
=
\frac{(-1)^j}{j!}\prod_{i=0}^{j-1}(b+1+i)
=
(-1)^j\frac{\Gamma(j+b+1)}
{\Gamma(b+1)\Gamma(j+1)}
\qquad(j\ge0),
\]
where the empty product is \(1\). Every factor in the product is positive, so the endpoint value is nonzero. Combining this identity with the norm formula yields
\[
\frac{\mathsf J_j^{(4,b)}(-1)^2}{H_j^{4,b}}
=
\frac{2j+b+5}{2^{b+5}\Gamma(b+1)^2}
\frac{\Gamma(j+b+1)}{\Gamma(j+1)}
\frac{\Gamma(j+b+5)}{\Gamma(j+5)}.
\]
For \(j\ge1\),
\[
2j\le2j+b+5\le(b+7)j.
\]
Applying the preceding gamma bounds to the shift pairs
\((b+1,1)\) and \((b+5,5)\), define
\[
\begin{split}
c_{\mathrm{spec,end}}
&:=
\frac{2c_\Gamma(b+1,1)c_\Gamma(b+5,5)}
     {2^{b+5}\Gamma(b+1)^2},\\
C_{\mathrm{spec,end}}
&:=
\frac{(b+7)C_\Gamma(b+1,1)C_\Gamma(b+5,5)}
     {2^{b+5}\Gamma(b+1)^2}.
\end{split}
\]
These constants are positive, and \(2b+1=\kappa\) gives
\[
c_{\mathrm{spec,end}}j^\kappa
\le
\frac{\mathsf J_j^{(4,b)}(-1)^2}{H_j^{4,b}}
\le
C_{\mathrm{spec,end}}j^\kappa
\qquad(j\ge1).
\]

Every summand of \(K_m(-1)\) is nonnegative. Strict positivity follows already from the degree
\[
\nu_{\mathrm{pos}}(m):=m+\lfloor m/2\rfloor,
\qquad
\frac98<\frac{\nu_{\mathrm{pos}}(m)}m<\frac{15}8
\quad(m\ge4),
\]
whose filter value and squared endpoint coefficient are positive.

For the quantitative lower bound, use the distinct degrees
\[
\nu_{\mathrm{end}}(m,i)
:=
m+\lfloor m/4\rfloor+1+i,
\qquad
0\le i<\lfloor m/4\rfloor.
\]
They satisfy
\[
\frac54\le\frac{\nu_{\mathrm{end}}(m,i)}m\le\frac74,
\qquad
m\le\nu_{\mathrm{end}}(m,i)\le2m.
\]
Moreover,
\[
m\le4\lfloor m/4\rfloor+3
\le7\lfloor m/4\rfloor
\qquad(m\ge4).
\]
Each selected summand is therefore at least
\(c_\eta c_{\mathrm{spec,end}}m^\kappa\). For the upper bound, inactive summands vanish, whereas each active summand is at most
\(C_{\mathrm{spec,end}}2^\kappa m^\kappa\); there are at most \(2m+1\le3m\) summands. Hence, with
\[
c_{\mathrm{end}}
:=\frac{c_\eta c_{\mathrm{spec,end}}}{7}>0,
\qquad
C_{\mathrm{end}}
:=3\,2^\kappa C_{\mathrm{spec,end}}>0,
\]
we obtain
\[
0<c_{\mathrm{end}}m^{\kappa+1}
\le K_m(-1)\le C_{\mathrm{end}}m^{\kappa+1}
\qquad(m\ge4).
\]
% lean: CausalSmith.Stat.NoisydoseWeakdesignTransition.packetKernel_endpoint_power_bounds

\item \emph{Endpoint localization.}
Define the polynomial coordinate and taper base on the whole line by
\[
z_{\mathrm{end}}(u):=2u^2-1,
\qquad
\Delta(u):=1-u^2
\qquad(u\in\mathbb R).
\]
For \(|u|\le1\), define
\[
\mathcal B_{\mathrm{end}}(m,u)
:=
\sqrt{\mathcal W_{4,b}(m;-1)
      \mathcal W_{4,b}(m;z_{\mathrm{end}}(u))}.
\]
Here \(0\le\Delta(u)\le1\), and
\[
\begin{split}
\mathcal W_{4,b}(m;-1)
&=(2+m^{-2})^{9/2}m^{-\kappa}
\ge m^{-\kappa},\\
\mathcal W_{4,b}(m;z_{\mathrm{end}}(u))
&=(2\Delta(u)+m^{-2})^{9/2}
  (2u^2+m^{-2})^{\kappa/2}.
\end{split}
\]
The inequalities
\[
\Delta(u)^8
\le\Delta(u)^{9/2}
\le(2\Delta(u)+m^{-2})^{9/2},
\qquad
m^{-\kappa}\le(2u^2+m^{-2})^{\kappa/2}
\]
therefore imply
\[
\mathcal B_{\mathrm{end}}(m,u)
\ge\Delta(u)^4m^{-\kappa}.
\]
Also,
\[
\delta_{\mathrm{ang}}(-1,z_{\mathrm{end}}(u))
=2\arcsin|u|
\ge|u|.
\]
Indeed, the equality follows from
\(\arccos(2u^2-1)=\pi-2\arcsin|u|\), and the inequality follows from
\(\sin v\le v\) for \(v\ge0\).

Use localization with \(\tau=\kappa+2\). On \([-1,1]\), positivity of the normalizer gives
\[
\begin{split}
|\psi_m(u)|
&=
\frac{\Delta(u)^4|K_m(z_{\mathrm{end}}(u))|}{K_m(-1)}\\
&\le
\frac{\Delta(u)^4 C_{\mathrm{loc},\kappa+2}^{4,b}m}
{\mathcal B_{\mathrm{end}}(m,u)
 (1+m\delta_{\mathrm{ang}}(-1,z_{\mathrm{end}}(u)))^{\kappa+2}
 K_m(-1)}\\
&\le
C_{\mathrm{dec,end}}(1+m|u|)^{-(\kappa+2)},
\end{split}
\qquad
C_{\mathrm{dec,end}}
:=
\frac{C_{\mathrm{loc},\kappa+2}^{4,b}}{c_{\mathrm{end}}}>0.
\]
Outside this interval \(\psi_m=0\), so the same bound holds for every real \(u\).
% lean: CausalSmith.Stat.NoisydoseWeakdesignTransition.packetPsi_endpoint_localization

\item \emph{Endpoint derivative and continuous differentiability.}
Fix \(0<|u|<1\). In the local angular estimate take the fixed second argument to be \(-1\), the reference and auxiliary arguments both to be \(z_{\mathrm{end}}(u)\), and the varying argument to be \(z_{\mathrm{end}}(v)\). For \(v\) sufficiently close to \(u\), all arguments belong to \([-1,1]\) and the required angular distance is at most \(m^{-1}\). Symmetry of the finite kernel sum then gives
\[
\begin{split}
&|K_m(z_{\mathrm{end}}(v))-K_m(z_{\mathrm{end}}(u))|\\
&\quad\le
\frac{C_{\mathrm{lip},1}^{4,b}m^2}
{\mathcal B_{\mathrm{end}}(m,u)
 (1+m\delta_{\mathrm{ang}}(-1,z_{\mathrm{end}}(u)))}
|\arccos z_{\mathrm{end}}(v)-\arccos z_{\mathrm{end}}(u)|.
\end{split}
\]
Dividing by \(|v-u|\) and passing to the limit is valid because the kernel is a finite polynomial sum. The identities
\[
1-z_{\mathrm{end}}(u)^2=4u^2\Delta(u),
\qquad
\left|\frac{d}{du}\arccos z_{\mathrm{end}}(u)\right|
=\frac2{\sqrt{\Delta(u)}}
\]
yield
\[
|(K_m\circ z_{\mathrm{end}})'(u)|
\le
\frac{2C_{\mathrm{lip},1}^{4,b}m^2}
{\mathcal B_{\mathrm{end}}(m,u)
 (1+m\delta_{\mathrm{ang}}(-1,z_{\mathrm{end}}(u)))
 \sqrt{\Delta(u)}}
\le
\frac{2C_{\mathrm{lip},1}^{4,b}m^2}
{\mathcal B_{\mathrm{end}}(m,u)\sqrt{\Delta(u)}}.
\]
Localization with exponent \(1\) likewise gives
\[
|K_m(z_{\mathrm{end}}(u))|
\le
\frac{C_{\mathrm{loc},1}^{4,b}m}
{\mathcal B_{\mathrm{end}}(m,u)}.
\]

For the derivative assembly, use the cubed taper. Since
\[
\Delta(u)^6
\le\Delta(u)^{9/2}
\le(2\Delta(u)+m^{-2})^{9/2},
\]
the same weight comparison as before gives
\[
\frac{\Delta(u)^3}{\mathcal B_{\mathrm{end}}(m,u)}
\le m^\kappa.
\]
Furthermore,
\[
\frac{\Delta(u)^4}{\sqrt{\Delta(u)}}\le\Delta(u)^3,
\]
because \(0<\Delta(u)\le1\). The product rule gives
\[
\psi_m'(u)=
\frac{-8u\Delta(u)^3K_m(z_{\mathrm{end}}(u))
+\Delta(u)^4(K_m\circ z_{\mathrm{end}})'(u)}
{K_m(-1)}.
\]
Using \(|u|\le1\), the preceding inequalities, and the normalization lower bound,
\[
\begin{split}
|\psi_m'(u)|
&\le
\frac{
8C_{\mathrm{loc},1}^{4,b}m
 \Delta(u)^3/\mathcal B_{\mathrm{end}}(m,u)
+
2C_{\mathrm{lip},1}^{4,b}m^2
 \Delta(u)^4/(\mathcal B_{\mathrm{end}}(m,u)\sqrt{\Delta(u)})
}{K_m(-1)}\\
&\le
\frac{8C_{\mathrm{loc},1}^{4,b}m^{\kappa+1}
      +2C_{\mathrm{lip},1}^{4,b}m^{\kappa+2}}
     {c_{\mathrm{end}}m^{\kappa+1}}\\
&\le C_{\mathrm{der,end}}m,
\end{split}
\qquad
C_{\mathrm{der,end}}
:=
\frac{8C_{\mathrm{loc},1}^{4,b}
      +2C_{\mathrm{lip},1}^{4,b}}{c_{\mathrm{end}}}>0.
\]

To obtain continuous differentiability on the whole line, define
\[
\mathcal T(u):=
\begin{cases}
(1-u^2)^4,&|u|\le1,\\
0,&|u|>1.
\end{cases}
\]
For every \(u\in\mathbb R\),
\[
\mathcal T(u)
=
\frac{(1-u^2)^4+(1-u^2)|1-u^2|^3}{2}.
\]
The function \(v\mapsto|v|^3\) is continuously differentiable, with derivative \(3v|v|\). Thus \(\mathcal T\) is continuously differentiable, and
\[
\psi_m(u)
=
\mathcal T(u)\frac{K_m(z_{\mathrm{end}}(u))}{K_m(-1)}
\qquad(u\in\mathbb R)
\]
is a product of continuously differentiable functions. Outside \([-1,1]\) the packet is locally zero, so
\[
\psi_m'(u)=0\qquad(|u|>1).
\]
Continuity of \(\psi_m'\) extends the derivative bound from the complement of
\(\{-1,0,1\}\) to all of \(\mathbb R\).
% lean: CausalSmith.Stat.NoisydoseWeakdesignTransition.packetPsi_endpoint_derivative_bound

\item \emph{Endpoint mass bounds and common constant.}
Set
\[
C_{\mathrm{env,end}}
:=
\max\{C_{\mathrm{dec,end}},C_{\mathrm{der,end}}\},
\qquad
C_{\mathrm{mid,end}}:=2C_{\mathrm{env,end}}.
\]
Both constants are positive. Since the coordinate \(z_{\mathrm{end}}\), taper, and support are invariant under reflection,
\[
\psi_m(-u)=\psi_m(u)\qquad(u\in\mathbb R).
\]
The envelope and \(1+m|u|\ge1\) give
\[
|\psi_m(u)|\le C_{\mathrm{env,end}}\le C_{\mathrm{mid,end}},
\qquad
|\psi_m'(u)|\le C_{\mathrm{mid,end}}m.
\]

For \(0\le u\le1\), the inequality
\((mu)^\kappa\le(1+mu)^\kappa\) gives
\[
u^\kappa(1+mu)^{-(\kappa+2)}
\le m^{-\kappa}(1+mu)^{-2}.
\]
The quadratic envelope has the explicit integral
\[
\int_0^1(1+mu)^{-2}\,du
=
\left[-\frac1{m(1+mu)}\right]_0^1
=
\frac1m-\frac1{m(1+m)}
\le\frac1m.
\]
By evenness,
\[
\begin{split}
\int_{-1}^1|u|^\kappa|\psi_m(u)|\,du
&=
2\int_0^1u^\kappa|\psi_m(u)|\,du\\
&\le
2C_{\mathrm{env,end}}m^{-\kappa}
\int_0^1(1+mu)^{-2}\,du\\
&\le C_{\mathrm{mid,end}}m^{-\kappa-1}.
\end{split}
\]
Continuity ensures integrability of these functions on the compact interval. Pointwise,
\[
\psi_m(u)^2
=
|\psi_m(u)|^2
\le C_{\mathrm{mid,end}}|\psi_m(u)|,
\]
and hence
\[
\int_{-1}^1|u|^\kappa\psi_m(u)^2\,du
\le
C_{\mathrm{mid,end}}
\int_{-1}^1|u|^\kappa|\psi_m(u)|\,du
\le C_{\mathrm{mid,end}}^2m^{-\kappa-1}.
\]
Define
\[
C_{\mathrm{fin,end}}
:=
\max\{C_{\mathrm{mid,end}},C_{\mathrm{mid,end}}^2\}>0.
\]
The supremum, derivative, absolute-mass, and squared-mass bounds all hold with this common constant.
% lean: CausalSmith.Stat.NoisydoseWeakdesignTransition.packetendpoint_bounds

\item \emph{Endpoint value, support, and moments.}
The definition and positive normalizer give
\[
\psi_m(u)=0\quad(|u|>1),
\qquad
\psi_m(0)=\frac{K_m(-1)}{K_m(-1)}=1.
\]

Let \(s\) be an integer with \(0\le s\le m\). Every active spectral degree exceeds \(m\), while the polynomial
\(((1+z)/2)^s\) has degree \(s\). Termwise Jacobi orthogonality therefore gives
\[
\int_{-1}^1
(1-z)^4(1+z)^bK_m(z)
\left(\frac{1+z}{2}\right)^s\,dz=0.
\]
For \(0<u<1\), the substitution \(z=2u^2-1\) has Jacobian \(4u\), and
\[
\begin{split}
&(1-(2u^2-1))^4(1+(2u^2-1))^b
\left(\frac{1+(2u^2-1)}2\right)^s(4u)\\
&\qquad=
2^{6+b}u^\kappa u^{2s}(1-u^2)^4.
\end{split}
\]
Consequently,
\[
\begin{split}
\int_0^1u^\kappa u^{2s}\psi_m(u)\,du
&=
\frac1{2^{6+b}K_m(-1)}
\int_{-1}^1
(1-z)^4(1+z)^bK_m(z)
\left(\frac{1+z}{2}\right)^s\,dz\\
&=0.
\end{split}
\]
The integrand of the full even moment is even, so reflection gives
\[
\int_{-1}^1|u|^\kappa u^{2s}\psi_m(u)\,du
=
2\int_0^1u^\kappa u^{2s}\psi_m(u)\,du=0.
\]
For an odd integer \(j\), reflection changes the sign of the integrand and gives
\[
\int_{-1}^1|u|^\kappa u^j\psi_m(u)\,du
=
-\int_{-1}^1|u|^\kappa u^j\psi_m(u)\,du
=0.
\]
If \(0\le j<m\) is even, writing \(j=2s\) gives \(s\le m\), so its moment vanishes as well. Thus the positive-exponent branch satisfies all assertions with
\[
C:=C_{\mathrm{fin,end}},
\qquad c:=1.
\]
% lean: CausalSmith.Stat.NoisydoseWeakdesignTransition.packetendpoint_bounds

\item \emph{Interior spectral lower bound.}
In the remaining case, \(0\le\kappa\) and \(\kappa\le0\), hence \(\kappa=0\). We use the symmetric Jacobi system with parameters \((4,4)\).

First, the norm formula and gamma bounds for the shift pairs \((5,1)\) and \((5,9)\) give
\[
\begin{split}
H_j^{4,4}
&=
\frac{512}{2j+9}
\frac{\Gamma(j+5)}{\Gamma(j+1)}
\frac{\Gamma(j+5)}{\Gamma(j+9)}\\
&\le
\frac{512}{2j+9}
C_\Gamma(5,1)j^4 C_\Gamma(5,9)j^{-4}
\le\frac{C_{\mathrm{norm,int}}}{j}
\qquad(j\ge1),
\end{split}
\]
where
\[
C_{\mathrm{norm,int}}
:=512C_\Gamma(5,1)C_\Gamma(5,9)>0.
\]

The binomial representation in \cref{obj:lem:jacobi-binomial}, evaluated at zero, gives
\[
\mathsf J_{2s}^{(4,4)}(0)
=
4^{-s}
\sum_{i=0}^{2s}
\binom{2s+4}{i}\binom{2s+4}{2s-i}(-1)^{2s-i}.
\]
The sum is the coefficient of degree \(2s\) in
\((1+z)^{2s+4}(1-z)^{2s+4}=(1-z^2)^{2s+4}\). Therefore
\[
\mathsf J_{2s}^{(4,4)}(0)
=
(-1)^s4^{-s}\binom{2s+4}{s}
\qquad(s\ge0).
\]
In particular, these even-degree center values are nonzero.

Wallis' inequality, in its finite-product form, states
\[
\prod_{i=1}^s\frac{(2i)^2}{(2i-1)(2i+1)}
=
\frac{2^{4s}(s!)^4}{((2s)!)^2(2s+1)}
\le\frac\pi2
\qquad(s\ge0).
\]
The factorial identity
\[
\frac{2^{4s}(s!)^4}{((2s)!)^2(2s+1)}
\left(4^{-s}\binom{2s}{s}\right)^2(2s+1)=1
\]
then gives, for \(s\ge1\),
\[
\left(4^{-s}\binom{2s}{s}\right)^2
\ge\frac2{\pi(2s+1)}
\ge\frac2{3\pi s}.
\]
Since \(\binom{2s+4}{s}\ge\binom{2s}{s}\),
\[
\mathsf J_{2s}^{(4,4)}(0)^2\ge\frac2{3\pi s}.
\]
Together with
\[
0<H_{2s}^{4,4}
\le\frac{C_{\mathrm{norm,int}}}{2s}
\le\frac{C_{\mathrm{norm,int}}}{s},
\]
this proves
\[
\frac{\mathsf J_{2s}^{(4,4)}(0)^2}{H_{2s}^{4,4}}
\ge c_{\mathrm{spec,int}},
\qquad
c_{\mathrm{spec,int}}
:=\frac2{3\pi C_{\mathrm{norm,int}}}>0
\qquad(s\ge1).
\]
% lean: CausalSmith.Stat.NoisydoseWeakdesignTransition.jacobi_four_even_spectral_lower

\item \emph{Interior normalization, including the smallest degrees.}
Define
\[
N_{\mathrm{even}}(m):=\max\{1,\lfloor m/8\rfloor\},
\qquad
\nu_{\mathrm{int}}(m,i)
:=2\bigl(\lceil5m/8\rceil+i\bigr),
\qquad
0\le i<N_{\mathrm{even}}(m).
\]
These are distinct positive even degrees satisfying
\[
\frac54\le\frac{\nu_{\mathrm{int}}(m,i)}m\le\frac74,
\qquad
\nu_{\mathrm{int}}(m,i)\le2m.
\]
For \(m\ge8\), the upper interval bound follows from
\[
\lceil5m/8\rceil+i
\le\frac{5m+7}{8}+\lfloor m/8\rfloor-1
\le\frac{6m-1}{8}\le\frac{7m}{8}.
\]
For \(m=4,5,6,7\), the construction selects respectively the single degree
\(6,8,8,10\), each in the stated interval.

The count satisfies
\[
N_{\mathrm{even}}(m)\ge\frac m{15}\qquad(m\ge4).
\]
Indeed, for \(m\ge8\),
\[
m\le8\lfloor m/8\rfloor+7
\le15\lfloor m/8\rfloor,
\]
and for \(4\le m\le7\) the count is \(1\).

Every summand of \(\widetilde K_m(0)\) is nonnegative. Each selected even degree contributes at least
\(c_\eta c_{\mathrm{spec,int}}\). Hence
\[
\widetilde K_m(0)
\ge
N_{\mathrm{even}}(m)c_\eta c_{\mathrm{spec,int}}
\ge c_{\mathrm{int}}m>0,
\qquad
c_{\mathrm{int}}
:=\frac{c_\eta c_{\mathrm{spec,int}}}{15}>0.
\]
% lean: CausalSmith.Stat.NoisydoseWeakdesignTransition.packetKernel_interior_lower_of_spectral_lower

\item \emph{Interior localization and derivative.}
For \(|u|\le1\), define
\[
\mathcal B_{\mathrm{int}}(m,u)
:=
\sqrt{\mathcal W_{4,4}(m;0)\mathcal W_{4,4}(m;u)}.
\]
The regularized weights satisfy
\[
\mathcal W_{4,4}(m;0)=(1+m^{-2})^9\ge1
\]
and
\[
\begin{split}
\mathcal W_{4,4}(m;u)
&=\bigl((1-u+m^{-2})(1+u+m^{-2})\bigr)^{9/2}\\
&\ge\Delta(u)^{9/2}\ge\Delta(u)^6.
\end{split}
\]
Thus
\[
\mathcal B_{\mathrm{int}}(m,u)\ge\Delta(u)^3\ge\Delta(u)^4.
\]
The angular distance dominates the coordinate distance:
\[
|u|
=
|\cos(\arccos u)-\cos(\arccos0)|
\le|\arccos u-\arccos0|
=
\delta_{\mathrm{ang}}(0,u).
\]
Localization with exponent \(2\) therefore gives
\[
\begin{split}
|\Delta(u)^4\widetilde K_m(u)|
&\le
\frac{\Delta(u)^4C_{\mathrm{loc},2}^{4,4}m}
{\mathcal B_{\mathrm{int}}(m,u)
 (1+m\delta_{\mathrm{ang}}(0,u))^2}\\
&\le C_{\mathrm{loc},2}^{4,4}m(1+m|u|)^{-2}.
\end{split}
\]
Dividing by the normalization lower bound yields
\[
|\psi_m(u)|
\le
\frac{C_{\mathrm{loc},2}^{4,4}}{c_{\mathrm{int}}}
(1+m|u|)^{-2}.
\]
This also holds outside \([-1,1]\), where the packet vanishes.

For \(|u|<1\), symmetry of the kernel and the local angular estimate, with reference and auxiliary arguments both equal to \(u\) and fixed second argument \(0\), give, for \(v\) sufficiently close to \(u\),
\[
|\widetilde K_m(v)-\widetilde K_m(u)|
\le
\frac{C_{\mathrm{lip},1}^{4,4}m^2}
{\mathcal B_{\mathrm{int}}(m,u)
 (1+m\delta_{\mathrm{ang}}(0,u))}
|\arccos v-\arccos u|.
\]
Passing to difference-quotient limits and using
\[
\left|\frac{d}{du}\arccos u\right|
=\frac1{\sqrt{\Delta(u)}}
\]
gives
\[
|\widetilde K_m'(u)|
\le
\frac{C_{\mathrm{lip},1}^{4,4}m^2}
{\mathcal B_{\mathrm{int}}(m,u)\sqrt{\Delta(u)}}.
\]
Localization with exponent \(1\) also gives
\[
|\widetilde K_m(u)|
\le
\frac{C_{\mathrm{loc},1}^{4,4}m}
{\mathcal B_{\mathrm{int}}(m,u)}.
\]
The product rule is
\[
\psi_m'(u)
=
\frac{-8u\Delta(u)^3\widetilde K_m(u)
+\Delta(u)^4\widetilde K_m'(u)}
{\widetilde K_m(0)}.
\]
Using
\[
\frac{\Delta(u)^3}{\mathcal B_{\mathrm{int}}(m,u)}\le1,
\qquad
\frac{\Delta(u)^4}{\sqrt{\Delta(u)}}\le\Delta(u)^3,
\]
we obtain
\[
\begin{split}
|\psi_m'(u)|
&\le
\frac{8C_{\mathrm{loc},1}^{4,4}m
      +C_{\mathrm{lip},1}^{4,4}m^2}
     {\widetilde K_m(0)}\\
&\le
\frac{8C_{\mathrm{loc},1}^{4,4}
      +C_{\mathrm{lip},1}^{4,4}}{c_{\mathrm{int}}}\,m.
\end{split}
\]
Here \(m\ge1\) absorbs the term independent of \(m\) after division.

On the whole line,
\[
\psi_m(u)
=
\mathcal T(u)\frac{\widetilde K_m(u)}{\widetilde K_m(0)}.
\]
The previously established regularity of \(\mathcal T\) proves that this packet is continuously differentiable. Its derivative vanishes outside \([-1,1]\), and continuity extends the bound to the two endpoints.

Set
\[
C_{\mathrm{env,int}}
:=
\frac{C_{\mathrm{loc},2}^{4,4}
      +8C_{\mathrm{loc},1}^{4,4}
      +C_{\mathrm{lip},1}^{4,4}}{c_{\mathrm{int}}}>0.
\]
We have proved, for every real \(u\),
\[
|\psi_m(u)|\le C_{\mathrm{env,int}}(1+m|u|)^{-2},
\qquad
|\psi_m'(u)|\le C_{\mathrm{env,int}}m.
\]
% lean: CausalSmith.Stat.NoisydoseWeakdesignTransition.packetinterior_bounds_of_normalization

\item \emph{Interior symmetry, masses, and normalization at zero.}
Reindexing the binomial representation in \cref{obj:lem:jacobi-binomial} by \(i\mapsto j-i\) gives
\[
\mathsf J_j^{(4,4)}(-u)
=(-1)^j\mathsf J_j^{(4,4)}(u).
\]
At \(u=0\), every odd-degree center value is therefore zero. The surviving even-degree terms give
\[
\widetilde K_m(-u)=\widetilde K_m(u),
\qquad
\psi_m(-u)=\psi_m(u).
\]

For the exponent \(\kappa=0\), we use \(|u|^0=1\), including at \(u=0\). Define
\[
C_{\mathrm{mid,int}}:=2C_{\mathrm{env,int}}>0.
\]
The envelope bounds the supremum by \(C_{\mathrm{mid,int}}\), and the derivative is bounded by \(C_{\mathrm{mid,int}}m\). Evenness and the explicit quadratic-envelope integral give
\[
\begin{split}
\int_{-1}^1|\psi_m(u)|\,du
&=
2\int_0^1|\psi_m(u)|\,du\\
&\le
2C_{\mathrm{env,int}}\int_0^1(1+mu)^{-2}\,du
\le C_{\mathrm{mid,int}}m^{-1}.
\end{split}
\]
The pointwise inequality
\(\psi_m(u)^2\le C_{\mathrm{mid,int}}|\psi_m(u)|\) then gives
\[
\int_{-1}^1\psi_m(u)^2\,du
\le
C_{\mathrm{mid,int}}\int_{-1}^1|\psi_m(u)|\,du
\le C_{\mathrm{mid,int}}^2m^{-1}.
\]
Thus, with
\[
C_{\mathrm{fin,int}}
:=
\max\{C_{\mathrm{mid,int}},C_{\mathrm{mid,int}}^2\}>0,
\]
all four analytic bounds hold with one constant. The definition and positive normalizer also give
\[
\psi_m(u)=0\quad(|u|>1),
\qquad
\psi_m(0)=\frac{\widetilde K_m(0)}{\widetilde K_m(0)}=1.
\]
% lean: CausalSmith.Stat.NoisydoseWeakdesignTransition.packetinterior_bounds

\item \emph{Interior moments and conclusion.}
Let \(j\) be an integer with \(0\le j<m\); in particular \(j\le m\). On \([-1,1]\), the taper equals the symmetric Jacobi weight:
\[
(1-u^2)^4=(1-u)^4(1+u)^4.
\]
Expanding the finite kernel sum therefore gives
\[
\begin{split}
\int_{-1}^1u^j\psi_m(u)\,du
&=
\frac1{\widetilde K_m(0)}
\sum_{\nu=0}^{2m}
\frac{\eta(\nu/m)\mathsf J_\nu^{(4,4)}(0)}
     {H_\nu^{4,4}}\\
&\qquad\qquad{}\times
\int_{-1}^1
(1-u)^4(1+u)^4
\mathsf J_\nu^{(4,4)}(u)u^j\,du
=0.
\end{split}
\]
Indeed, an inactive coefficient is zero, while an active degree satisfies
\(\nu>m\ge j\), so its integral vanishes by Jacobi orthogonality. Finite summation and compact-interval integrability justify the expansion.

The interior branch consequently satisfies every required assertion with
\[
C:=C_{\mathrm{fin,int}},
\qquad c:=1.
\]
Together with the positive-exponent branch, this proves the result for every fixed \(\kappa\in[0,2]\).
% lean: CausalSmith.Stat.NoisydoseWeakdesignTransition.weighted_packet_construction
\end{enumerate}
\end{proof}

\begin{proof}[Proof of \cref{obj:thm:observed-lower}]
Fix public constants \(K=(p_{\min},c_{\mathrm{lo}},c_{\mathrm{hi}})\), \(\beta\in(0,1]\), \(\kappa\in[0,2]\), and \(\tau>0\), under the envelope calibration in the statement. All constants chosen below may depend on \(K,\beta,\kappa,\tau\) and the fixed packet construction.

\begin{enumerate}
\item \textbf{Packet constants.}
For \(\kappa=0\), use the symmetric interior construction in \cref{obj:def:packet-handle}; for \(\kappa>0\), use the endpoint construction. In the latter case, its Jacobi parameter satisfies
\[
b=\frac{\kappa-1}{2}>-\frac12,
\]
and the remaining Jacobi parameters equal \(r=4\). The localization estimate is supplied by \citep[Theorem 2.4, equation (2.14)]{PetrushevXu2005}, and the local increment estimate by \citep[Theorem 2.2, equation (2.4)]{https://people.math.sc.edu/pencho/Publications/kpx-06-28-06-web.pdf}. The orthogonality and positive squared-norm formula are supplied by \citep[Table 18.3.1, Jacobi row and caption]{https://dlmf.nist.gov/18.3}, and the required Gamma-ratio asymptotic by \citep[equation (5.11.12), restricted to the positive real axis]{https://dlmf.nist.gov/5.11.E12}.

Thus \cref{obj:thm:weighted-packet-construction} supplies constants
\[
C_{\mathrm{packet}}>0,\qquad c_p>0
\]
uniformly for \(m\ge4\), including positive normalizing values, evenness, continuous differentiability of the zero extension, support in \([-1,1]\), and
\[
\psi_m(0)=1,\qquad
\|\psi_m\|_\infty\le C_{\mathrm{packet}},\qquad
\|\psi_m'\|_\infty\le C_{\mathrm{packet}}m,
\]
\[
\int_{-1}^{1}|u|^\kappa|\psi_m(u)|\,du
\le C_{\mathrm{packet}}m^{-\kappa-1},
\qquad
\int_{-1}^{1}|u|^\kappa\psi_m(u)^2\,du
\le C_{\mathrm{packet}}m^{-\kappa-1},
\]
\[
\int_{-1}^{1}|u|^\kappa u^j\psi_m(u)\,du=0
\qquad(0\le j<c_pm).
\]
% lean: CausalSmith.Stat.NoisydoseWeakdesignTransition.lower_regime_witnesses

\item \textbf{Common calibration.}
Define the cancellation constant, degree threshold, intermediate multiplier, and compact radius by
\[
\gamma:=\frac{c_p}{2},\qquad
M_{\mathrm{deg}}:=\left\lceil\max\left\{4,\frac{2}{c_p}\right\}\right\rceil,
\qquad
C_{\mathrm I}:=M_{\mathrm{deg}},\qquad
\ell_{\mathrm M}:=\frac18.
\]
Then
\[
M_{\mathrm{deg}}\ge4,\qquad
c_pM_{\mathrm{deg}}\ge2,\qquad
C_{\mathrm I}\ge4,\qquad c_pC_{\mathrm I}\ge2.
\]
Choose the intermediate radius constant explicitly as
\[
a_{\mathrm I}:=
\min\left\{
1,\frac{1}{4\sqrt{C_{\mathrm I}+1}},
\sqrt{\min\left\{\frac{\gamma}{2},
\gamma\exp(-1/\gamma-1)\right\}}
\right\}.
\]
Every entry in this minimum is positive, and hence
\[
0<a_{\mathrm I}\le1,\qquad
a_{\mathrm I}\sqrt{C_{\mathrm I}+1}\le\frac14,\qquad
a_{\mathrm I}^2\le\frac{\gamma}{2},\qquad
\frac{\exp(1)a_{\mathrm I}^2}{\gamma}\le\exp(-1/\gamma).
\]
For the compact regime, set
\[
C_{\mathrm M}:=
\max\left\{
1,\frac8\gamma,\frac{2\exp(2)\ell_{\mathrm M}^2}{\gamma}
\right\}.
\]
Consequently,
\[
C_{\mathrm M}\ge1,\qquad
\frac{\gamma C_{\mathrm M}}{\exp(1)\ell_{\mathrm M}^2}
\ge2\exp(1),\qquad
\gamma C_{\mathrm M}\ge8.
\]

Define the design coefficient and three comparison constants by
\[
c_\kappa:=(\kappa+1)2^\kappa,\qquad
B_{\mathrm D}:=8c_\kappa,\qquad
B_{\mathrm I}:=4\exp(1/2)(c_\kappa C_{\mathrm{packet}})^2,\qquad
B_{\mathrm G}:=4\exp(1/2)(2c_\kappa)^2,
\]
and set
\[
B_{\mathrm{amp}}:=\max\{B_{\mathrm D},2B_{\mathrm I}\},\qquad
\varepsilon:=
\min\left\{
\frac18,\frac{1}{8C_{\mathrm{packet}}},
\sqrt{\frac{\log(1+\tau^2)}{B_{\mathrm{amp}}}}
\right\}.
\]
These constants are positive. The amplitude satisfies
\[
\varepsilon\le\frac18,\qquad
4\varepsilon\le1,\qquad
\varepsilon C_{\mathrm{packet}}\le\frac18,
\]
\[
B_{\mathrm D}\varepsilon^2\le\log(1+\tau^2),
\qquad
2B_{\mathrm I}\varepsilon^2\le\log(1+\tau^2).
\]
Finally, choose
\[
C:=\varepsilon^2\max\{B_{\mathrm G},B_{\mathrm I}\},
\qquad
c:=2\varepsilon\min\left\{
1,
\left(\frac{a_{\mathrm I}}{\sqrt{C_{\mathrm I}+1}}\right)^\beta,
\left(\frac{\ell_{\mathrm M}}{2C_{\mathrm M}}\right)^\beta
\right\}.
\]
Both are positive, and these definitions give the comparison constants required in each regime.
% lean: CausalSmith.Stat.NoisydoseWeakdesignTransition.lower_common_amplitude_constant

\item \textbf{Uniform sample cutoff.}
Use the effective exponent and sample logarithm from \cref{obj:def:frontier-rate}, and define
\[
d:=2\beta+\kappa+1>0,\qquad
L_n:=\log(\exp(1)n),\qquad h_0:=n^{-1/d},
\]
\[
n_{\mathrm D}:=\max\{2,\lceil4^d\rceil\},\qquad
n_{\log}:=\max\{2,\lceil\exp(M_{\mathrm{deg}}^2)\rceil\},
\qquad
n_0:=\max\{n_{\mathrm D},n_{\log}\}.
\]
For every \(n\ge n_0\),
\[
0<h_0\le(4^d)^{-1/d}=\frac14,\qquad
nh_0^d=1,\qquad
M_{\mathrm{deg}}^2\le L_n.
\]
The last inequality follows by taking logarithms of
\[
\exp(M_{\mathrm{deg}}^2)\le n\le\exp(1)n.
\]
For every \(\sigma\in[0,1/4]\), also
\[
1\le\log(\exp(1)+n\sigma^d)\le L_n.
\]
Indeed, \(0\le\sigma^d\le1\), while \(n\ge2\) and \(\exp(1)\ge2\) give
\[
\exp(1)\le\exp(1)+n\sigma^d
\le\exp(1)+n\le\exp(1)n.
\]
Fix henceforth an arbitrary public space from \cref{obj:def:potential-path-space}, \(n\ge n_0\), and \(\sigma\in[0,1/4]\).
% lean: CausalSmith.Stat.NoisydoseWeakdesignTransition.lower_log_sample_cutoff

\item \textbf{Witness laws and observed densities.}
The common centered-dose density is
\[
g(t):=c_\kappa|t|^\kappa\mathbf1\{|t|\le1/2\},
\qquad t\in\mathbb R.
\]
We use the convention \(|t|^0=1\), including at \(t=0\). Its normalization follows from
\[
\int_{-1/2}^{1/2}|t|^\kappa\,dt
=\frac{2^{-\kappa}}{\kappa+1},
\qquad
\int_{\mathbb R}g(t)\,dt=1.
\]
Take mutually independent seed coordinates with
\[
P(X=0)=P(X=1)=\frac12,\qquad
\mathcal L(A-a_0)(dt)=g(t)\,dt,\qquad
Z\sim N(0,1),\qquad U\sim\operatorname{Unif}[0,1].
\]
Here the prescribed center is \(a_0=1/2\). For any continuous profile
\(\mu:[0,1]\to[1/8,7/8]\), define its structural law by the measurable pushforward
\[
P_\mu:=
\mathcal L\bigl(X,A,Z,\iota(\mu,U),\iota(\mu,U)(A),U\bigr).
\]
Thus, under \(P_\mu\),
\[
Y(\cdot)=\iota(\mu,U),\qquad
Y(a)=\mathbf1\{U\le\mu(a)\},\qquad
Y=Y(A),\qquad W=A+\sigma Z.
\]
The measurability requirements follow from the embedding and joint evaluation in \cref{obj:def:potential-path-space}. Uniformity of the independent rank gives
\[
E_{P_\mu}[Y(a)\mid X=x]=\mu(a),
\qquad
\theta(P_\mu)=\sum_{x\in\{0,1\}}\frac12\mu(a_0)=\mu(a_0).
\]

For a continuous compactly supported perturbation, write
\[
\mu_\pm(a):=\frac12\pm\delta(a-a_0),\qquad
\mu_\star(a):=\frac12,\qquad
P_\pm:=P_{\mu_\pm},\qquad P_\star:=P_{\mu_\star},
\]
whenever the profiles have the displayed range. Define
\[
Q_\pm:=\mathcal L_{P_\pm}(X,W,Y),\qquad
Q_\star:=\mathcal L_{P_\star}(X,W,Y),\qquad
v_\delta(t):=g(t)\delta(t).
\]
For \(\sigma>0\), define
\[
\phi_\sigma(w):=
\frac{1}{\sqrt{2\pi}\sigma}
\exp\left(-\frac{w^2}{2\sigma^2}\right),
\qquad
p_V:=g*\phi_\sigma,\qquad f_\delta:=v_\delta*\phi_\sigma.
\]
For \(\sigma=0\), define instead
\[
p_V:=g,\qquad f_\delta:=v_\delta.
\]
Conditioning on the latent dose and integrating the uniform rank gives the densities
\[
q_\pm(x,w,y)
=
p_x\left\{\frac{p_V(w-a_0)}2
\pm(2y-1)f_\delta(w-a_0)\right\},
\qquad
q_\star(x,w,y)=\frac{p_xp_V(w-a_0)}2,
\]
relative to counting measure in \(x\in\{0,1\}\), \(y\in\{0,1\}\), and Lebesgue measure in \(w\in\mathbb R\), with \(p_x=1/2\).

The marked probabilities belong to \([0,1]\), so
\[
0\le Q_\pm\le2Q_\star,\qquad Q_\pm\ll Q_\star.
\]
For the likelihood ratios
\[
\Lambda_\pm:=\frac{dQ_\pm}{dQ_\star},
\]
this implies
\[
0\le\Lambda_\pm\le2\quad Q_\star\text{-almost surely},
\qquad
\int(\Lambda_\pm-1)^2\,dQ_\star\le1.
\]
We use
\[
\chi^2(Q_\pm,Q_\star):=
\int(\Lambda_\pm-1)^2\,dQ_\star,
\qquad
\chi^2(P_\pm,P_\star):=\chi^2(Q_\pm,Q_\star),
\]
in agreement with the theorem's observed-law convention. Summing the squared likelihood deviations over the two marks and strata, and translating \(w\) by \(a_0\), yields
\[
\chi^2(Q_\pm,Q_\star)
=
4\int_{\{p_V>0\}}\frac{f_\delta(w)^2}{p_V(w)}\,dw.
\]
At \(\sigma=0\), both marked densities vanish wherever \(g=0\), and this identity reduces to
\[
\chi^2(Q_\pm,Q_\star)
=4\int_{\mathbb R}g(t)\delta(t)^2\,dt.
\]
% lean: CausalSmith.Stat.NoisydoseWeakdesignTransition.observed_marked_chiSq_identity

\item \textbf{Gaussian moment comparison.}
Suppose \(\sigma>0\), and let the support radius and cancellation cutoff satisfy
\[
\ell>0,\qquad J\in\mathbb N_0,\qquad
\operatorname{supp}(v_\delta)\subseteq[-\ell,\ell],\qquad
\int_{\mathbb R}t^jv_\delta(t)\,dt=0\quad(0\le j<J).
\]
Define the absolute mass and moment tail by
\[
\|v_\delta\|_1:=\int_{\mathbb R}|v_\delta(t)|\,dt,\qquad
\mathcal T_{\sigma,\ell,J}:=
\sum_{j=J}^{\infty}\frac{(\ell^2/\sigma^2)^j}{j!}.
\]

First consider \(w\ge0\). Restricting the convolution defining \(p_V(w)\) to \(0\le t\le\sigma\), which lies in the latent support, gives
\[
\begin{aligned}
p_V(w)
&=\phi_\sigma(w)\int_{\mathbb R}
g(t)\exp\left(\frac{wt}{\sigma^2}-\frac{t^2}{2\sigma^2}\right)\,dt\\
&\ge
c_\kappa\exp(-1/2)\phi_\sigma(w)\int_0^\sigma t^\kappa\,dt\\
&=
2^\kappa\exp(-1/2)\sigma^{\kappa+1}\phi_\sigma(w)\\
&\ge\exp(-1/2)\sigma^{\kappa+1}\phi_\sigma(w)>0.
\end{aligned}
\]
For \(w<0\), evenness of \(g\) and \(\phi_\sigma\) gives the same bound by replacing \(w\) with \(-w\). Hence the denominator bound holds on the entire real line.

Completing the square gives, for \(t,s,w\in\mathbb R\),
\[
\frac{\phi_\sigma(w-t)\phi_\sigma(w-s)}{\phi_\sigma(w)}
=
\exp(ts/\sigma^2)\phi_\sigma(w-t-s).
\]
Its integral in \(w\) is \(\exp(ts/\sigma^2)\). Compact support bounds the absolute triple integral by
\[
\exp(\ell^2/\sigma^2)\|v_\delta\|_1^2<\infty.
\]
Absolute Fubini therefore yields
\[
\int_{\mathbb R}\frac{f_\delta(w)^2}{\phi_\sigma(w)}\,dw
=
\iint_{\mathbb R^2}
v_\delta(t)v_\delta(s)\exp(ts/\sigma^2)\,dt\,ds.
\]
For each \(j\in\mathbb N_0\), the absolute integral of the \(j\)-th exponential-series term is bounded by
\[
\iint_{\mathbb R^2}
\frac{|t^jv_\delta(t)s^jv_\delta(s)|}{\sigma^{2j}j!}\,dt\,ds
\le
\|v_\delta\|_1^2\frac{(\ell^2/\sigma^2)^j}{j!}.
\]
The sum of these bounds is finite. Termwise integration and the cancellations give
\[
\begin{aligned}
\int_{\mathbb R}\frac{f_\delta(w)^2}{\phi_\sigma(w)}\,dw
&=
\sum_{j=0}^{\infty}
\frac{\left(\int_{\mathbb R}t^jv_\delta(t)\,dt\right)^2}
{\sigma^{2j}j!}\\
&=
\sum_{j=J}^{\infty}
\frac{\left(\int_{\mathbb R}t^jv_\delta(t)\,dt\right)^2}
{\sigma^{2j}j!}\\
&\le \|v_\delta\|_1^2\mathcal T_{\sigma,\ell,J},
\end{aligned}
\]
where the last inequality uses
\[
\left|\int_{\mathbb R}t^jv_\delta(t)\,dt\right|
\le\ell^j\|v_\delta\|_1.
\]
Combining this with the observed-density identity and the denominator bound proves
\[
\chi^2(Q_\pm,Q_\star)
\le
4\exp(1/2)\sigma^{-\kappa-1}
\|v_\delta\|_1^2\mathcal T_{\sigma,\ell,J}.
\]
% lean: CausalSmith.Stat.NoisydoseWeakdesignTransition.observed_perturbation_chiSq_moment_bound

\item \textbf{Direct profiles.}
Suppose first that \(\sigma\le h_0\). Set
\[
h=\ell=h_0,\qquad m=0,\qquad J=0,
\]
and define, for all \(t\in\mathbb R\),
\[
\delta(t):=
\varepsilon h_0^\beta
\left(1-\left(\frac{t}{h_0}\right)^2\right)^2
\mathbf1\{|t|\le h_0\}.
\]
For the increment calculation, define the clipped coordinate
\[
\zeta_{\mathrm D}(t):=\min\{|t/h_0|,1\}.
\]
Then
\[
\delta(t)=\varepsilon h_0^\beta
(1-\zeta_{\mathrm D}(t)^2)^2,
\qquad
0\le\zeta_{\mathrm D}(t)\le1,
\qquad
|\zeta_{\mathrm D}(s)-\zeta_{\mathrm D}(t)|
\le\frac{|s-t|}{h_0}.
\]
Factoring differences of squares twice gives
\[
\begin{aligned}
\left|(1-\zeta_{\mathrm D}(s)^2)^2
-(1-\zeta_{\mathrm D}(t)^2)^2\right|
&\le2|\zeta_{\mathrm D}(s)^2-\zeta_{\mathrm D}(t)^2|\\
&\le4|\zeta_{\mathrm D}(s)-\zeta_{\mathrm D}(t)|\\
&\le4\frac{|s-t|}{h_0}.
\end{aligned}
\]
Thus
\[
0\le\delta(t)\le\varepsilon h_0^\beta,\qquad
|\delta(s)-\delta(t)|
\le4\varepsilon h_0^\beta\frac{|s-t|}{h_0}.
\]
If \(|s-t|\le h_0\), then
\[
\frac{|s-t|}{h_0}
\le\left(\frac{|s-t|}{h_0}\right)^\beta,
\]
so the increment is at most \(4\varepsilon|s-t|^\beta\). If \(|s-t|>h_0\), the range bound gives
\[
|\delta(s)-\delta(t)|
\le\varepsilon h_0^\beta
\le\varepsilon|s-t|^\beta.
\]
Consequently,
\[
|\delta(s)-\delta(t)|\le4\varepsilon|s-t|^\beta.
\]
The explicit formula is continuous across its support endpoints. Since \(h_0^\beta\le1\), the common calibration ensures
\[
\mu_\pm(a)\in[1/8,7/8],\qquad
|\mu_\pm(a)-\mu_\pm(a')|\le|a-a'|^\beta
\quad(a,a'\in[0,1]).
\]
At the target,
\[
|\theta(P_+)-\theta(P_-)|
=2\delta(0)=2\varepsilon h_0^\beta
=2\varepsilon\rho_{n,\sigma}
\ge c\rho_{n,\sigma},
\]
using the direct branch of \cref{obj:def:frontier-rate}.

On \([-h_0,h_0]\), the bounds on \(g\) and \(\delta\) give
\[
|v_\delta(t)|\le c_\kappa\varepsilon h_0^{\beta+\kappa}.
\]
Outside that interval \(v_\delta=0\). Integrating over its length \(2h_0\) yields
\[
\|v_\delta\|_1
\le2c_\kappa\varepsilon h_0^{\beta+\kappa+1}.
\]
Moreover,
\[
\int_{\mathbb R}g(t)\delta(t)^2\,dt
\le\varepsilon h_0^\beta\|v_\delta\|_1
\le2c_\kappa\varepsilon^2h_0^d.
\]
For positive noise, the Gaussian moment comparison with \(J=0\) therefore gives
\[
\chi^2(Q_\pm,Q_\star)
\le
B_{\mathrm G}\varepsilon^2
\sigma^{-\kappa-1}h_0^{2\beta+2\kappa+2}
\mathcal T_{\sigma,h_0,0}
\le
C\sigma^{-\kappa-1}h^{2\beta+2\kappa+2}
\mathcal T_{\sigma,\ell,J}.
\]
The moment conditions for \(J=0\) have an empty index set.
% lean: CausalSmith.Stat.NoisydoseWeakdesignTransition.direct_lower_profile_frontier_construction

\item \textbf{Direct observed testing bound.}
At \(\sigma=0\), the exact observed comparison already gives
\[
\chi^2(Q_\pm,Q_\star)
=4\int_{\mathbb R}g(t)\delta(t)^2\,dt.
\]
At \(\sigma>0\), Cauchy--Schwarz under the finite positive measure
\(g(t)\phi_\sigma(w-t)\,dt\) gives
\[
f_\delta(w)^2
\le
p_V(w)\int_{\mathbb R}
g(t)\delta(t)^2\phi_\sigma(w-t)\,dt.
\]
Dividing by the positive denominator, integrating, and using the unit Gaussian integral yields
\[
\chi^2(Q_\pm,Q_\star)
\le4\int_{\mathbb R}g(t)\delta(t)^2\,dt.
\]
Thus, for every noise value in the direct regime,
\[
n\chi^2(Q_\pm,Q_\star)
\le8c_\kappa\varepsilon^2nh_0^d
=B_{\mathrm D}\varepsilon^2
\le\log(1+\tau^2).
\]

Here is the product comparison that will also be used for the inverse regimes. Independence in \cref{obj:def:observed-experiment} and the square-integrability of the likelihood ratios imply
\[
\begin{aligned}
1+\chi^2(Q_{P_\pm}^n,Q_{P_\star}^n)
&=
E_{Q_\star^{\otimes n}}
\left[\prod_{i=1}^n\Lambda_\pm(O_i)^2\right]\\
&=
\left(E_{Q_\star}\Lambda_\pm^2\right)^n
=\{1+\chi^2(Q_\pm,Q_\star)\}^n\\
&\le\exp\{n\chi^2(Q_\pm,Q_\star)\}.
\end{aligned}
\]
Hence the sample budget implies
\[
\chi^2(Q_{P_\pm}^n,Q_{P_\star}^n)\le\tau^2.
\]
For an absolutely continuous probability law, the signed likelihood deviation has integral zero. Its positive and negative parts consequently have equal integrals, and Cauchy--Schwarz gives
\[
\operatorname{TV}(Q_{P_\pm}^n,Q_{P_\star}^n)
=
\frac12E_{Q_{P_\star}^n}
\left|\prod_{i=1}^n\Lambda_\pm(O_i)-1\right|
\le\frac12\sqrt{\chi^2(Q_{P_\pm}^n,Q_{P_\star}^n)}
\le\frac\tau2.
\]
The triangle inequality now proves
\[
\operatorname{TV}(Q_{P_+}^n,Q_{P_-}^n)\le\tau.
\]
% lean: CausalSmith.Stat.NoisydoseWeakdesignTransition.direct_profile_pair_observed_tv_bound

\item \textbf{Inverse profiles and tail estimates.}
For the remaining cases, \(\sigma>h_0>0\). Given
\[
m\in\mathbb N,\qquad m\ge M_{\mathrm{deg}},\qquad
0<\ell\le\frac14,
\]
define
\[
h:=\frac{\ell}{m},\qquad
J:=\lfloor c_pm\rfloor,\qquad
\delta(t):=\varepsilon h^\beta\psi_m(t/\ell)
\quad(t\in\mathbb R).
\]
Then \(0<h\le1/4\). Since \(c_pm\ge2\) and
\(\lfloor c_pm\rfloor>c_pm-1\),
\[
J\ge\gamma m,\qquad J\ge1.
\]
The packet bounds imply continuity and support in \([-\ell,\ell]\), together with
\[
|\delta(t)|\le\varepsilon C_{\mathrm{packet}}h^\beta.
\]
The derivative bound and the mean-value inequality give
\[
|\delta(s)-\delta(t)|
\le\varepsilon C_{\mathrm{packet}}h^\beta\frac{|s-t|}{h}.
\]
If \(|s-t|\le h\), replacing \(|s-t|/h\) by its \(\beta\)-th power proves
\[
|\delta(s)-\delta(t)|
\le\varepsilon C_{\mathrm{packet}}|s-t|^\beta.
\]
If \(|s-t|>h\), the supremum bound gives
\[
|\delta(s)-\delta(t)|
\le2\varepsilon C_{\mathrm{packet}}h^\beta
\le2\varepsilon C_{\mathrm{packet}}|s-t|^\beta.
\]
Therefore
\[
|\delta(s)-\delta(t)|
\le2\varepsilon C_{\mathrm{packet}}|s-t|^\beta.
\]
The common amplitude calibration implies
\[
\mu_\pm(a)\in[1/8,7/8],\qquad
|\mu_\pm(a)-\mu_\pm(a')|\le|a-a'|^\beta.
\]
These arguments include the Lipschitz endpoint \(\beta=1\).

For \(j<J\), we have \(j<c_pm\). Changing variables \(t=\ell u\), using \(\ell\le1/4\), and applying the packet cancellations gives
\[
\int_{\mathbb R}t^jv_\delta(t)\,dt
=
c_\kappa\varepsilon h^\beta\ell^{\kappa+j+1}
\int_{-1}^{1}|u|^\kappa u^j\psi_m(u)\,du
=0.
\]
The weighted absolute-mass bound similarly gives
\[
\begin{aligned}
\|v_\delta\|_1
&=
c_\kappa\varepsilon h^\beta\ell^{\kappa+1}
\int_{-1}^{1}|u|^\kappa|\psi_m(u)|\,du\\
&\le
c_\kappa\varepsilon C_{\mathrm{packet}}
h^\beta\ell^{\kappa+1}m^{-\kappa-1}\\
&=
c_\kappa\varepsilon C_{\mathrm{packet}}h^{\beta+\kappa+1}.
\end{aligned}
\]
Thus the Gaussian moment comparison proves
\[
\chi^2(Q_\pm,Q_\star)
\le
B_{\mathrm I}\varepsilon^2
\sigma^{-\kappa-1}h^{2\beta+2\kappa+2}
\mathcal T_{\sigma,\ell,J}
\le
C\sigma^{-\kappa-1}h^{2\beta+2\kappa+2}
\mathcal T_{\sigma,\ell,J}.
\]
Normalization at zero gives the exact separation
\[
|\theta(P_+)-\theta(P_-)|
=2\varepsilon h^\beta.
\]

To control the sample budget, define
\[
\lambda:=\frac{\ell^2}{\sigma^2}>0.
\]
Whenever \(J\ge1\) and \(J\ge2\lambda\), consecutive factorial-series terms have ratio at most \(1/2\). Induction therefore gives
\[
\frac{\lambda^{J+k}}{(J+k)!}
\le\frac{\lambda^J}{J!}\,2^{-k}
\quad(k\in\mathbb N_0),
\qquad
\mathcal T_{\sigma,\ell,J}
\le2\frac{\lambda^J}{J!}.
\]
The nonnegative exponential series gives
\[
\frac{J^J}{J!}
\le\sum_{k=0}^\infty\frac{J^k}{k!}
=\exp(J).
\]
Multiplication by \((\lambda/J)^J\) yields
\[
\frac{\lambda^J}{J!}
\le\left(\frac{\exp(1)\lambda}{J}\right)^J,
\qquad
\mathcal T_{\sigma,\ell,J}
\le2\left(\frac{\exp(1)\lambda}{J}\right)^J.
\]
Finally, whenever \(h\le\sigma\),
\[
\sigma^{-\kappa-1}h^{2\beta+2\kappa+2}
=
h^d(h/\sigma)^{\kappa+1}
\le h^d.
\]
% lean: CausalSmith.Stat.NoisydoseWeakdesignTransition.inverse_lower_profile_chiSq_construction

\item \textbf{Intermediate regime.}
Suppose
\[
h_0<\sigma\le L_n^{-1/2}.
\]
Define the effective sample quantity and its logarithm by
\[
N_{\mathrm I}:=n\sigma^d,\qquad
S_{\mathrm I}:=\log(\exp(1)+N_{\mathrm I}),
\]
and choose
\[
m:=\lceil C_{\mathrm I}S_{\mathrm I}\rceil,\qquad
\ell:=a_{\mathrm I}\sigma\sqrt m,\qquad
h:=\frac{\ell}{m}=\frac{a_{\mathrm I}\sigma}{\sqrt m},
\qquad
J:=\lfloor c_pm\rfloor.
\]
The logarithmic window gives \(1\le S_{\mathrm I}\le L_n\). The ceiling bounds imply
\[
C_{\mathrm I}S_{\mathrm I}\le m
<C_{\mathrm I}S_{\mathrm I}+1
\le(C_{\mathrm I}+1)S_{\mathrm I}.
\]
In particular \(m\ge M_{\mathrm{deg}}\), so \(J\ge\gamma m\) and \(J\ge1\). Moreover,
\[
\begin{aligned}
0<\ell
&\le a_{\mathrm I}\sigma
\sqrt{C_{\mathrm I}+1}\sqrt{L_n}\\
&\le a_{\mathrm I}\sqrt{C_{\mathrm I}+1}
\le\frac14,
\end{aligned}
\qquad
0<h\le\ell\le\frac14,\qquad h\le\sigma.
\]
Thus the inverse profiles just constructed are legal. Their bandwidth satisfies
\[
h=\frac{a_{\mathrm I}\sigma}{\sqrt m}
\ge
\frac{a_{\mathrm I}}{\sqrt{C_{\mathrm I}+1}}
\frac{\sigma}{\sqrt{S_{\mathrm I}}},
\]
and therefore
\[
|\theta(P_+)-\theta(P_-)|
\ge
2\varepsilon
\left(\frac{a_{\mathrm I}}{\sqrt{C_{\mathrm I}+1}}\right)^\beta
\left(\frac{\sigma}{\sqrt{S_{\mathrm I}}}\right)^\beta
\ge c\rho_{n,\sigma},
\]
using the intermediate branch of \cref{obj:def:frontier-rate}.

For the Gaussian tail parameter,
\[
\lambda=\frac{\ell^2}{\sigma^2}=a_{\mathrm I}^2m,
\qquad
2\lambda\le\gamma m\le J,
\]
and the radius calibration gives
\[
\frac{\exp(1)\lambda}{J}
\le\frac{\exp(1)a_{\mathrm I}^2}{\gamma}
\le\exp(-1/\gamma).
\]
Consequently,
\[
\mathcal T_{\sigma,\ell,J}
\le2\exp(-J/\gamma)\le2\exp(-m).
\]
Since \(m\ge S_{\mathrm I}\),
\[
N_{\mathrm I}\exp(-m)
\le
N_{\mathrm I}\exp(-S_{\mathrm I})
=
\frac{N_{\mathrm I}}{\exp(1)+N_{\mathrm I}}
\le1.
\]
Using \(h\le\sigma\) and \(d>0\), we obtain the complete sample-factor bound
\[
\begin{aligned}
n\sigma^{-\kappa-1}h^{2\beta+2\kappa+2}
\mathcal T_{\sigma,\ell,J}
&\le2n\sigma^d\exp(-m)\\
&=2N_{\mathrm I}\exp(-m)\le2.
\end{aligned}
\]
Hence
\[
n\chi^2(Q_\pm,Q_\star)
\le2B_{\mathrm I}\varepsilon^2
\le\log(1+\tau^2).
\]
The product comparison established above proves
\[
\operatorname{TV}(Q_{P_+}^n,Q_{P_-}^n)\le\tau.
\]
% lean: CausalSmith.Stat.NoisydoseWeakdesignTransition.intermediate_lower_profile_sample_budget

\item \textbf{Compact regime.}
The remaining case satisfies
\[
\sigma>h_0,\qquad \sigma>L_n^{-1/2}.
\]
Define
\[
\tau_{\mathrm M}:=\sigma^2L_n>1,\qquad
B_{\mathrm M}:=\log(\exp(1)+\tau_{\mathrm M}),
\]
and choose
\[
m:=\left\lceil\frac{C_{\mathrm M}L_n}{B_{\mathrm M}}\right\rceil,
\qquad
\ell:=\ell_{\mathrm M},\qquad
h:=\frac{\ell_{\mathrm M}}m,\qquad
J:=\lfloor c_pm\rfloor.
\]

We first establish the logarithmic bounds used for this choice. Since
\(\tau_{\mathrm M}\ge1\),
\[
\exp(1)+\tau_{\mathrm M}
\le(\exp(1)+1)\tau_{\mathrm M}.
\]
Also \(\exp(1)+1\le\exp(2)\), so
\[
B_{\mathrm M}
\le\log\tau_{\mathrm M}+\log(\exp(1)+1)
\le\log\tau_{\mathrm M}+2.
\]
Applying \(\log x\le x-1\) to \(x=\sqrt{\tau_{\mathrm M}}\) gives
\[
\log\tau_{\mathrm M}\le2\sqrt{\tau_{\mathrm M}}-2,
\qquad
B_{\mathrm M}\le2\sqrt{\tau_{\mathrm M}}.
\]
Because \(\sigma\le1/4\),
\[
0<B_{\mathrm M}
\le2\sigma\sqrt{L_n}
\le\frac{\sqrt{L_n}}2,
\qquad
\frac{L_n}{B_{\mathrm M}}\ge2\sqrt{L_n}\ge1.
\]
The ceiling bounds consequently yield
\[
\frac{C_{\mathrm M}L_n}{B_{\mathrm M}}
\le m\le\frac{2C_{\mathrm M}L_n}{B_{\mathrm M}}.
\]
Furthermore,
\[
m\ge\frac{L_n}{B_{\mathrm M}}
\ge2\sqrt{L_n}\ge M_{\mathrm{deg}},
\]
where \(M_{\mathrm{deg}}^2\le L_n\) was enforced by the sample cutoff. Thus
\[
J\ge\gamma m\ge
\frac{\gamma C_{\mathrm M}L_n}{B_{\mathrm M}},
\qquad J\ge1.
\]
The bandwidth is positive and at most \(\ell_{\mathrm M}\le1/4\). Since
\(\sigma\sqrt{L_n}\ge1\),
\[
\sigma m\ge2\sigma\sqrt{L_n}\ge2\ge\ell_{\mathrm M},
\qquad h\le\sigma.
\]
The upper degree bound gives
\[
h=\frac{\ell_{\mathrm M}}m
\ge
\frac{\ell_{\mathrm M}}{2C_{\mathrm M}}\frac{B_{\mathrm M}}{L_n}.
\]
The inverse-profile separation therefore satisfies
\[
|\theta(P_+)-\theta(P_-)|
\ge
2\varepsilon
\left(\frac{\ell_{\mathrm M}}{2C_{\mathrm M}}\right)^\beta
\left(\frac{B_{\mathrm M}}{L_n}\right)^\beta
\ge c\rho_{n,\sigma},
\]
using the compact branch of \cref{obj:def:frontier-rate}.

For the tail calculation, define
\[
\lambda:=\frac{\ell_{\mathrm M}^2}{\sigma^2}
=\frac{\ell_{\mathrm M}^2L_n}{\tau_{\mathrm M}},
\qquad
D_{\mathrm M}:=
\frac{\gamma C_{\mathrm M}}
{\exp(1)\ell_{\mathrm M}^2}\ge2\exp(1).
\]
The cancellation order and the bound \(B_{\mathrm M}\le2\sqrt{\tau_{\mathrm M}}\) imply
\[
\frac{J}{\exp(1)\lambda}
\ge
D_{\mathrm M}\frac{\tau_{\mathrm M}}{B_{\mathrm M}}
\ge\exp(1)\sqrt{\tau_{\mathrm M}}
\ge2.
\]
In particular \(J\ge2\lambda\). Taking logarithms, and using
\(\log(\exp(1)+1)\le2\), gives
\[
\log\frac{J}{\exp(1)\lambda}
\ge1+\frac12\log\tau_{\mathrm M}
\ge\frac{B_{\mathrm M}}2.
\]
Hence
\[
J\log\frac{J}{\exp(1)\lambda}
\ge
\frac{\gamma C_{\mathrm M}L_n}{B_{\mathrm M}}
\frac{B_{\mathrm M}}2
=
\frac{\gamma C_{\mathrm M}L_n}{2}
\ge4L_n.
\]
The factorial-tail estimate now gives
\[
\mathcal T_{\sigma,\ell_{\mathrm M},J}
\le
2\exp\left(-J\log\frac{J}{\exp(1)\lambda}\right)
\le2\exp(-4L_n).
\]
Using \(h\le\sigma\) for the first inequality and \(0<h\le1\) for the second, the Gaussian prefactor satisfies
\[
\sigma^{-\kappa-1}h^{2\beta+2\kappa+2}\le h^d\le1.
\]
Also \(L_n\ge0\), so
\[
n\exp(-4L_n)
\le n\exp(-L_n)
=\exp(-1)\le1.
\]
It follows that
\[
n\sigma^{-\kappa-1}h^{2\beta+2\kappa+2}
\mathcal T_{\sigma,\ell_{\mathrm M},J}\le2,
\]
and therefore
\[
n\chi^2(Q_\pm,Q_\star)
\le2B_{\mathrm I}\varepsilon^2
\le\log(1+\tau^2).
\]
The same product comparison proves
\[
\operatorname{TV}(Q_{P_+}^n,Q_{P_-}^n)\le\tau.
\]
% lean: CausalSmith.Stat.NoisydoseWeakdesignTransition.compact_lower_profile_sample_budget

\item \textbf{Structural membership and assembly.}
Each regime has supplied continuous profiles satisfying
\[
\mu_\pm(a)\in[1/8,7/8],\qquad
|\mu_\pm(a)-\mu_\pm(a')|\le|a-a'|^\beta.
\]
The reference profile satisfies the same conditions, since
\[
\mu_\star(a)=\frac12,\qquad
|\mu_\star(a)-\mu_\star(a')|=0\le|a-a'|^\beta.
\]
We now verify membership for all three pushforward laws.

The common seed gives
\[
p_0=p_1=\frac12,\qquad
A=a_0+t\in[0,1]\quad\text{almost surely}.
\]
The common conditional centered-dose density has coefficient
\(c_\kappa=(\kappa+1)2^\kappa\). The assumed calibration therefore gives
\[
c_{\mathrm{lo}}|t|^\kappa\le g(t)
\le c_{\mathrm{hi}}|t|^\kappa
\qquad(t\in[-1/2,1/2]).
\]
Because the class constants satisfy \(p_{\min}\le1/2\), the equal stratum probabilities obey \(p_{\min}\le p_x=1/2\le1-p_{\min}\) for both strata. These facts establish \cref{obj:ass:stratum-overlap,obj:ass:latent-support,obj:ass:weak-design}.

For every one of the three profiles, the conditional-mean identity derived from the uniform seed is
\[
E_{P_\mu}[Y(a)\mid X=x]=\mu(a).
\]
The established profile bounds therefore give
\cref{obj:ass:holder-mean,obj:ass:mean-range}. The chosen Gaussian marginal gives \cref{obj:ass:gaussian-channel}. Independence of the Gaussian seed from the other seed coordinates is preserved under the measurable schedule map, yielding
\[
Z\perp(X,A,Y(\cdot)),
\]
as required by \cref{obj:ass:error-independence}.

The product seed also gives
\[
\mathcal L(U\mid X=x)=\operatorname{Unif}[0,1],
\qquad U\perp A\mid X,
\]
establishing \cref{obj:ass:rank-uniform,obj:ass:rank-treatment-independence}. More explicitly, for measurable schedule and dose events,
\[
\begin{aligned}
&P_\mu\{Y(\cdot)\in\mathcal B_{\mathrm{sched}},
A\in\mathcal B_{\mathrm{dose}}\mid X=x\}\\
&\qquad=
P\{\iota(\mu,U)\in\mathcal B_{\mathrm{sched}}\}
P\{A\in\mathcal B_{\mathrm{dose}}\}.
\end{aligned}
\]
Here the events have domains
\[
\mathcal B_{\mathrm{sched}}\in\mathcal E,\qquad
\mathcal B_{\mathrm{dose}}\in\mathcal B([0,1]).
\]
This factorization proves \cref{obj:ass:conditional-unconfoundedness} for the entire random schedule.

The measurable pushforward construction carries a single full-measure event on which
\[
Y(a)=\mathbf1\{U\le\mu(a)\}
=\mathbf1\{U\le\mu_X(a)\}
\qquad\text{for every }a\in[0,1].
\]
Thus \cref{obj:ass:structural-threshold,obj:ass:binary-potential-outcomes} hold. The same identity gives \(Y(a)\in[0,1]\), and the construction gives \(Y=Y(A)\); these establish \cref{obj:ass:bounded-potential-outcomes,obj:ass:consistency}. All ten conditions in \cref{obj:def:model-class} are therefore satisfied by \(P_+,P_-,P_\star\).

For each law, take the observed product measure
\[
Q_P^n:=\bigotimes_{i=1}^n\mathcal L_P(X,W,Y).
\]
This is the experiment in \cref{obj:def:observed-experiment} and satisfies \cref{obj:ass:iid}. Finally, projection of the seed pushforward onto \((X,A)\) removes the profile-dependent coordinates, so
\[
\mathcal L_{P_+}(X,A)
=\mathcal L_{P_-}(X,A)
=\mathcal L_{P_\star}(X,A).
\]
The direct branch includes \(\sigma=0\) and \(\sigma=h_0\); among the remaining noise values, the intermediate branch includes \(\sigma=L_n^{-1/2}\), and the compact branch covers the strict upper side. Each branch has supplied the required separation, observed product-TV bound, and, for positive noise, the displayed Gaussian moment-series comparison with the same constant \(C\). The three laws use the same seed and the same regime-specific tuning data, as asserted.
% lean: CausalSmith.Stat.NoisydoseWeakdesignTransition.observed_lower
\end{enumerate}
\end{proof}

\bibliographystyle{plainnat}

\bibliography{references}

\end{document}
